\documentclass[11pt]{amsart}
\usepackage{amsmath,amssymb,amsthm}
\usepackage[margin=1in]{geometry}
\usepackage{hyperref}
\usepackage{enumitem}

\newtheorem{theorem}{Theorem}[section]
\newtheorem{proposition}[theorem]{Proposition}
\newtheorem{lemma}[theorem]{Lemma}
\newtheorem{corollary}[theorem]{Corollary}
\newtheorem{conjecture}[theorem]{Expected result}
\theoremstyle{definition}
\newtheorem{definition}[theorem]{Definition}
\theoremstyle{remark}
\newtheorem{remark}[theorem]{Remark}
\numberwithin{equation}{section}

\newcommand{\bbH}{\mathbb{H}}
\newcommand{\Z}{\mathbb{Z}}
\newcommand{\R}{\mathbb{R}}
\newcommand{\C}{\mathbb{C}}
\newcommand{\E}{\mathbb{E}}
\newcommand{\Pb}{\mathbb{P}}
\newcommand{\T}{\mathbb{T}}
\newcommand{\cB}{\mathcal{B}}
\newcommand{\cC}{\mathcal{C}}
\newcommand{\cP}{\mathcal{P}}
\newcommand{\cS}{\mathcal{S}}
\newcommand{\cV}{\mathcal{V}}
\newcommand{\cW}{W}
\newcommand{\cZ}{\mathcal{Z}}
\newcommand{\cL}{\mathcal{L}}
\newcommand{\cT}{\mathcal{T}}
\newcommand{\cN}{\mathcal{N}}
\newcommand{\avg}[1]{\langle #1\rangle}
\newcommand{\Loc}{\operatorname{Loc}}
\newcommand{\rng}{\operatorname{rng}}
\newcommand{\conn}{\leftrightarrow}
\newcommand{\ip}[2]{\langle #1,#2\rangle}
\newcommand{\avint}{\mathop{\,\rlap{--}\!\!\int}\nolimits}
\newcommand{\frd}{\texttt{frd-redecompose.tex}}
\newcommand{\hr}{\texttt{h02rg.tex}}

\title[The continuum limit of the 2D $\bbH^{0|2}$ model]{The two-dimensional $\bbH^{0|2}$ model and arboreal gas:\\
an outline of the continuum limit on the unit torus}
\date{}

\begin{document}
\sloppy
\maketitle

\begin{abstract}
We outline a renormalisation group construction of the continuum limit of the $\bbH^{0|2}$
model, equivalently of the arboreal gas, on the discretised unit torus $\T_\epsilon$ in $d=2$.

In two dimensions the quartic coupling $1/\beta$ is marginal and asymptotically free. A
nontrivial limit requires
$\beta_\epsilon=\frac{3}{2\pi}\log\frac1\epsilon+O(\log\log\frac1\epsilon)$. The connection
probability $\Pb[0\conn x]$ then vanishes like $(\log\frac1\epsilon)^{-2/3}$. After division by
this factor it converges, and at short distances it behaves like $(\log\frac1{|x|})^{2/3}$.

The construction follows \hr{}. There are two changes:
\begin{itemize}
\item the diverging field strength is absorbed with the dressed finite range decomposition of
  \frd{};
\item the flow of the effective potential is computed to second order.
\end{itemize}
The one-loop coefficients are computed here: $\beta_1=\frac3{2\pi}$, and the anomalous dimension
of the spin field is $\gamma=\frac1{2\pi\beta}$. Statements marked \emph{Expected} are targets.
The document is organised so that each step can be formalised as an extension of the Lean
development of \hr{}.
\end{abstract}

\tableofcontents

%======================================================================
\section{Introduction}
\label{sec:intro}

\subsection{The model}
Fix an integer $L\ge L_0$. For $N\ge1$ let $\epsilon=L^{-N}$, and identify the discretised unit
torus $\T_\epsilon=\epsilon\Z^2/\Z^2$ with $\Lambda_N=(\Z/L^N\Z)^2$. Lattice units are used
inside the renormalisation group, and physical units (torus side $1$) in the statements. For
$x\in\T=\R^2/\Z^2$, $x_\epsilon$ is a nearest point of $\T_\epsilon$, and $|x|$ is the Euclidean
distance on $\T$.

The $\bbH^{0|2}$ model $\avg{\cdot}_{\beta,h}$ on $\Lambda_N$ is defined as in \hr{},
Section~\texttt{sec:spinintro}. By \texttt{prop:h02-forest}, at $h=0$,
\begin{equation}\label{e:P}
  \Pb_\beta[0\conn x]=\avg{\xi_0\eta_x}_{\beta,0}=1-\avg{\xi_0\eta_0\xi_x\eta_x}_{\beta,0},
  \qquad \avg{\xi_0\eta_0}_{\beta,0}=1 .
\end{equation}
The last identity is the Ward identity $\avg{z_0}_{\beta,0}=0$. The object of interest is
\[
  P_\epsilon(x)=\Pb^{\T_\epsilon}_{\beta_\epsilon}[0\conn x_\epsilon]
\]
for a suitable choice of $\beta_\epsilon$.

\subsection{Expected results}
\label{sec:results}
Throughout, $L\ge L_0$ and $\beta_R\ge\beta_0(L)$ are large. Constants $C,c$ may depend on $L$
and $\beta_R$ but not on $N$ or $x$. The running coupling $\beta_j$ and the spin renormalisation
$\vartheta_j$ are defined by the renormalisation group in Sections~\ref{sec:flow} and
\ref{sec:obs-flow}. Here $j(\ell)=\lfloor N+\log_L\ell\rfloor$ is the scale of a physical
length $\ell$.

\begin{theorem}[Choice of $\beta_\epsilon$; proved in Lean, \texttt{Spin2.thm\_beta},
\texttt{Spin2.thm\_flow\_master\_all}]\label{thm:beta}
For every $N$ there is a unique $\beta_\epsilon=\beta_\epsilon(\beta_R)$ such that the flow with
$\beta=\beta_\epsilon$ and $h=0$ stays perturbative at all scales and has $\beta_N=\beta_R$.
The map $\beta_R\mapsto\beta_\epsilon$ is continuous and increasing, and
\[
  \beta_{j+1}=\beta_j-\beta_1\log L+O(\beta_j^{-1}),\qquad\beta_1=\frac3{2\pi},\qquad
  \beta_\epsilon=\frac3{2\pi}\log\frac1\epsilon+O\Bigl(\log\log\frac1\epsilon\Bigr).
\]
\end{theorem}
\paragraph{Lean.} \texttt{Spin2.thm\_beta} (\texttt{H02UV/Spin2/Multiscale/BetaThm.lean}): one set of constants (those of
Theorem~\ref{thm:flow}, $\epsilon\le1$) such that \texttt{FlowC2Concl} (Theorem~\ref{thm:flow}) and \texttt{BetaConcl} hold:
for every $N\ge2$, $N_h=N-2$, $m^2\in(0,L^{-2(N-1)}]$ and $\beta_R\ge2/\epsilon$, with $b=1/\beta_\epsilon=$
\texttt{bFin}$(m^2,1/\beta_R)$ and $\mu_0=a_0^c(b)$:
\begin{itemize}
\item $b$ is the unique element of $U(m^2)$ (Theorem~\ref{thm:flow}) with $\operatorname{Re}\check g_N=1/\beta_R$, and the flow
  from $(b,\mu_0)$ is admissible;
\item $|\beta_{N_h}-\beta_R|\le\frac{16}3C_b$, $\beta_k=1/\check g_k$ (from $|\check g_{N-2}-1/\beta_R|\le4C_b\beta_R^{-2}$,
  \texttt{shoot\_beta\_near});
\item $|\beta_{k+1}-\beta_k+\beta_1\log L|\le K_ba_k+C_{gu}\beta_k^{-1}$ for $k<N_h$, $a_0=1$, $a_k=L^{-2k}+m^2L^{2(k+1)}$
  ($\sum_ka_k\le4$);
\item $|\beta_\epsilon-\beta_R-N_h\beta_1\log L|\le E+a\log\bigl(2(\beta_R+N_h\beta_1\log L+E+a|\log2a|)\bigr)$ with $E,a$
  depending only on $L$ (\texttt{beta\_eps\_shoot}, from \texttt{BetaAsym}); hence, with $\epsilon=L^{-N}$,
  $|\beta_\epsilon-\frac3{2\pi}\log\frac1\epsilon|\le\beta_R+2\beta_1\log L+E+a\log(\dots)$;
\item $\beta_R\mapsto\beta_\epsilon$ is strictly increasing and continuous on $[2/\epsilon,\infty)$; $m^2\mapsto\beta_\epsilon$ is
  continuous (Theorem~\ref{thm:flow}(ii)).
\end{itemize}
The constants are those of the master construction \texttt{Spin2.thm\_flow\_master} (hypothesis package
\texttt{Step2.flow\_master\_hyp0}): the conclusions of \texttt{thm\_flow\_complete2X} and of \texttt{thm\_flow\_final\_spinX}
(\texttt{SpinConcl}) hold with the same constants, and \texttt{thm\_flow\_master\_all} adds \texttt{BetaConcl} and the
magnetisation statement of \texttt{prop\_gamma\_M2} (\texttt{MagConcl}). Audit: \texttt{H02UV/Spin2/Audit.lean}.
Compared with the statement above: $d=2$, $i_0=0$, $m^2>0$; ``for every $N$'' is $N\ge2$ and $\beta_R\ge\beta_0=2/\epsilon$;
uniqueness is within $U(m^2)$; ``perturbative at all scales'' is admissibility for $j\le N_h=N-2$ together with the bounds of
the last two steps of Theorem~\ref{thm:flow}(i); the final condition is $\operatorname{Re}\check g_N=1/\beta_R$ (the final
coupling $\check g_N=\kappa_{N,g}/Z_N^2$ is complex), so $\beta_{N-2}=\beta_R+O(1)$ only; the per-step error contains the
summable $O(a_k)$ in addition to $O(\beta_k^{-1})$; the $O(\log\log\frac1\epsilon)$ is explicit, with constants depending on $L$
and linearly on $\beta_R$.

\begin{theorem}[Multiscale formula; proved in Lean, \texttt{Spin2.thm\_multiscale\_final\_h0}, \texttt{\_allx}]\label{thm:multiscale}
Let $d=2$ and $i_0=0$. There is $L_0$ such that for every $L\ge L_0$ there are $g_0>0$, $K$ and $N_0$ such that the
following holds for every $N\ge N_0$, every auxiliary mass $m^2>0$ and every $\beta_R>0$ with
\[
  4m^2L^{2N}\le\frac1{\beta_R}\le g_0 .
\]
There is $\beta_\epsilon=1/b>0$ such that the flow at $h=0$, i.e.\ from $(\check g_0,\mu_0)=(b,b-m^2)$, stays in the
domain $D_\varepsilon$ of Theorem~\ref{thm:flow} up to scale $N-2$ and has $\operatorname{Re}\check g_N=1/\beta_R$. With
$\beta_j=1/\check g_j$ and $\vartheta_j=(\beta_j/\beta_\epsilon)^{1/3}$, for every lattice point $x\ne0$,
\[
  \bigl|P_\epsilon(x)-\vartheta_{k}^2\bigr|\le K\,\frac{\vartheta_{k}^2}{\beta_{k}},\qquad k=\min(j(|x|),N-2).
\]
(For coalescence scale $j_x\le N-2$ this is $k=j(|x|)$.)
The constants $g_0,K,N_0$ depend on $L$ only.
\end{theorem}
\paragraph{Lean.} \texttt{Spin2.thm\_multiscale\_final\_h0} (\texttt{H02UV/Spin2/Final/FinalTheorem.lean}; audit \texttt{H02UV/Spin2/Audit.lean},
axioms \texttt{propext}, \texttt{Classical.choice}, \texttt{Quot.sound}). The statement has no named hypotheses: for $L\ge L_0$ there
are one master construction $M$ (\texttt{thm\_master\_all}) and $g_0>0$, $K$, $N_0=4+\lceil4K_p\rceil$ such that for every $N\ge N_0$,
every $R$ with the master parameters ($P=P_c$, $i_0=0$, $\ell,A,\rho$ of $M$), every $m^2>0$ and $\beta_R>0$ with
$4m^2L^{2N}\le1/\beta_R\le g_0$ there is $b>0$ with $\operatorname{Re}\check g_N(b,b-m^2)=1/\beta_R$ (\texttt{gFinH0}), the $h=0$
trajectory in $D_\varepsilon$ for $j\le N-2$, and for every $x\ne0$ with $j_x=\texttt{Obs.coal}\,L\,0\,x\le N-2$ a scale $j$ with
$j=j(|x|)$ (\texttt{jPhys} of the physical distance $|x|/L^N$), $j\le j_x\le j+1$ and
$|P_{1/b}[0\conn x]-\vartheta_j^2|\le K\vartheta_j^2/\beta_j$ ($P_\epsilon(x)=\texttt{Pconn}\,L\,N\,\beta_\epsilon\,x$, $\vartheta_j^2=\texttt{thetaSq}$,
$\beta_j=\texttt{betaJ}$). The side conditions of the intermediate theorems are folded into $g_0$ and $N_0$: $100C_b/\beta_R\le1$ and
the mass windows of the last step ($8K_Am^2L^{2N}\le1$, $K_tm^2L^{2N}\le\frac12$, and
$2\cdot\frac{81}{16}c_J(c_\rho+4/r_0^2)c_S\ell^2\,m^2L^{2N}\le1$ of the $ab$ last step) follow from $m^2L^{2N}\le g_0/4$;
$m^2\le L^{-2(N-1)}$ follows from $m^2L^{2N}\le1$; and $K_pa_{N_h-1}\le\frac12$ ($a_{N-3}=L^{-2(N-3)}+m^2L^{2(N-2)}$) follows from
$N\ge N_0$ for the first term and from $m^2L^{2N}\le1/(4(K_p+1))$ for the second.

Chain of proof: the master construction (Theorems~\ref{thm:flow} and~\ref{thm:beta}, the observable flow and the pair layer);
the window ends at $h=0$ with $r=1-b\avg{\bar\psi_0\psi_0}\ge0$ at the upper end (\texttt{MasterFacts.fb\_h0EndsGS},
\texttt{H02UV/Spin2/Final/FinalBeta.lean}); the pair layer without named inputs (\texttt{fw\_pairAbove6\_atD}, \texttt{H02UV/Spin2/Final/FinalTheorem.lean},
with \texttt{GamTopBd} from \texttt{fw\_gamTopBdD}: \texttt{fw\_pairAbove6\_MFD'}); on one $h=0$ trajectory the two-point input
\texttt{TwoPtSq} (hence Expected result~\ref{thm:rginput} at $h=0$) and the bound (\texttt{MasterFacts.fw\_multiscale\_traj6}, via
\texttt{twoPtSqD} and \texttt{MasterCore.multiscale\_h0}); the per-$b$ conclusion \texttt{FwRes6} (\texttt{fw\_atG6}, \texttt{fb\_consts6});
shooting on the $h=0$ family between $b_1=b_{\rm Fin}(\frac1{2\beta_R})$ and $b_2=b_{\rm Fin}(\frac2{\beta_R})$ of Theorem~\ref{thm:beta}
(\texttt{fw\_bracket6}, \texttt{thm\_multiscale\_h0\_beta6}); the folding of the side conditions (\texttt{thm\_multiscale\_final\_h0}).

\emph{All solutions} (\texttt{Spin2.thm\_multiscale\_final\_h0\_all}, \texttt{H02UV/Spin2/Final/FinalAll.lean}): with the same prefix, the
bound holds at \emph{every} $b$ in the bracket $[b_{\rm Fin}(\frac1{2\beta_R}),b_{\rm Fin}(\frac2{\beta_R})]$. Every solution $b$ of the
$h=0$ shooting problem (\texttt{H0Sol}: $\operatorname{Re}\check g_N(b,b-m^2)=1/\beta_R$, trajectory in $D_\varepsilon$) with
$b\in U_{\rm fix}$ and $m^2L^{2N}\le\check g^c_{N-2}(b)\le g_0$ lies in that bracket, so the bound holds there as well. The
localisation is not proved without these two hypotheses: the comparison $\check g^0\approx\check g^c$ is only available in the
small-$\check g^c$ window. A per-$b$ form, \texttt{Spin2.thm\_multiscale\_h0\_perb}, takes $y=\operatorname{Re}\check g_N(b,b-m^2)$:
if $b\in[b_{\rm Fin}(y/2),b_{\rm Fin}(2y)]$ and $4m^2L^{2N}\le y\le g_0$, the bound holds at $b$; if $b\in U_{\rm fix}$ and
$4m^2L^{2N}\le\check g^c(b)\le g_0$, then $b$ lies in this bracket.

\emph{Top coalescence scales} (\texttt{Spin2.thm\_multiscale\_final\_h0\_top}, \texttt{H02UV/Spin2/Final/FinalTop.lean}): with the same prefix
and constants $g_0,K,K',N_0$ depending on $L$ only, the solution $b$ satisfies in addition, for every $x\ne0$ with $j_x\ge N-1$
($j_x\in\{N-1,N\}$),
$|P_{1/b}[0\conn x]-\vartheta_{N-2}^2|\le K'\vartheta_{N-2}^2/\beta_{N-2}$; \texttt{Spin2.thm\_multiscale\_final\_h0\_allx} states the
bound for \emph{all} $x\ne0$ at the scale $\min(j(|x|),N-2)$. Route (a pair chain of length zero): the exact square at the split
$k=N_h+1=N-1$ ($\Pi^{<N-1}$ rank one, \texttt{obsHist\_coal\_validL}, \texttt{zz\_square\_split\_at}), $b\langle\bar\psi_0\psi_x\rangle=AB-E$ with
$A=\gamma_{0,N-1}$, $B=\gamma_{x,N-1}$, $E=n_{ab}(\mathrm{lastO})$; the sector reduction at scale $N-1$ (\texttt{redTop}, covering
polymer \texttt{dxTop}, \texttt{pair\_red\_top}): $E=n_{ab}X'+n_{ab}\Omega'$ with $n_{ab}X'$ the last-step defect
(\texttt{nAB\_zCstR\_top}, $r'_{N-1}=0$) and $\Omega'$ the reduced zero mode, bounded through the swap $\sigma_a\sigma_b\mapsto\sigma_a$
(\texttt{pairLastAB\_xT}, \texttt{pairTop\_E\_le}: $|E|\le K_EM_{N-2}^2\check g_{N-2}$); the seed
$A^{(|\mathcal B(X)|-4)_+}\|K'_{N-1}(X)\|_{w_M}\le(2\iota_TA)^{L^2}$ from the unreduced $ab$ chain up to $N-1$
(\texttt{kab\_hist\_top}, \texttt{wobs\_hist\_fullWUL}, \texttt{topSeed\_traj}, \texttt{MasterCore.topSeedD}); the one-point parts
$|\gamma_{\bullet,N-1}-M_{N-2}|\le K_OM_{N-2}\check g_{N-2}$ and $|\gamma_{a,N}|\le K_4M_{N-2}$ for every $x$ (\texttt{MasterFacts.oneTopD}).
The additional smallness (the seed and the mass window $2\cdot\frac{81}{16}c_J(c_\rho+4/r_0^2)(8c_K+1)\ell^2m^2L^{2N}\le1$) is folded into
$g_0$ (\texttt{fwGT}). The sharp $\check g_N-\check g_{N-1}=O(\check g^2)$ needed to state the bound with $\vartheta_{j(|x|)}$ for
$j(|x|)\in\{N-1,N\}$ is not proved; the bound is at $\vartheta_{N-2}$.

\emph{Every $b$ in the bracket, all $x\ne0$} (\texttt{Spin2.thm\_multiscale\_final\_h0\_allB}, \texttt{H02UV/Spin2/Final/FinalAllX.lean}):
the chain of \texttt{thm\_multiscale\_final\_h0\_top} (result predicate \texttt{FwResT}$=$\texttt{FwRes6}$\wedge$\texttt{FwTopBd},
\texttt{fb\_constsT}) run as in \texttt{fa\_main}/\texttt{ct\_main}: constants $M,g_0,K,N_0$ depending on $L$ only; a shooting
solution exists in the bracket $[b_{\rm Fin}(\frac1{2\beta_R}),b_{\rm Fin}(\frac2{\beta_R})]$; at \emph{every} $b$ of the bracket the $h=0$
trajectory is in $D_\varepsilon$ up to $N-2$, $\check g_{N-2}\ge1/(5\beta_R)$, and the bound holds for all $x\ne0$ at the scale
$\min(j(|x|),N-2)$ (\texttt{FwAllBd}, from \texttt{fwAll\_of}: $K=\max(K_{\rm pair},K_{\rm top})$); every \texttt{H0Sol} with
$b\in U_{\rm fix}$ and $m^2L^{2N}\le\check g^c_{N-2}(b)\le g_0$ lies in the bracket.

Differences from the original formulation ($P_\epsilon(x)=\vartheta_{j(|x|)}^2(1+O(\beta_{j(|x|)}^{-1}))$ uniformly in
$\epsilon\le|x|\le\frac12$, with $\vartheta_j^2=(\beta_j/\beta_\epsilon)^{2/3}(1+O(\beta_j^{-1}))$):
\begin{itemize}
\item the range is all $x\ne0$ (\texttt{thm\_multiscale\_final\_h0\_allx}, \texttt{\_allB}), but with $\vartheta_k$,
  $k=\min(j(|x|),N-2)$: for $j_x\le N-2$ (i.e.\ $|x|_\infty<1/(2L)$ in physical units) this is $\vartheta_{j(|x|)}$; for the remaining
  $x$ ($j_x\ge N-1$) the bound is at the scale $N-2$, since $\check g_N-\check g_{N-1}=O(\check g^2)$ is not proved;
\item $m^2>0$ is a fixed auxiliary mass in the window $4m^2L^{2N}\le1/\beta_R$, and $h=0$ is imposed through the windowed final
  condition (route (a1), \S\ref{sec:ward}); there is no limit $m^2\downarrow0$ or $h\downarrow0$;
\item existence, not uniqueness, of $\beta_\epsilon$ at $h=0$; the final condition is $\operatorname{Re}\check g_N=1/\beta_R$ (as in
  Theorem~\ref{thm:beta}), and $\beta_R\ge1/g_0$ is large;
\item $\vartheta_j$ is defined as $(\beta_j/\beta_\epsilon)^{1/3}$, so the second display holds by definition;
\item the error is stated as $|P_\epsilon(x)-\vartheta_j^2|\le K\vartheta_j^2/\beta_j$, the $O(\beta_j^{-1})$ form;
\item the constants depend on $L$ (and $N\ge N_0(L)$), but not on $N$, $m^2$, $\beta_R$ or $x$.
\end{itemize}
\emph{Ingredients (Case~(2) vertex).} \texttt{Spin2.thm\_multiscale\_case2} (\texttt{H02UV/Spin2/Case2/Case2.lean}) removes (H$_\lambda$) for the Case~(2) vertex
$\hat\lambda_k=\mu_{a,k}/\mu_{a,0}$ (the $\sigma_a\bar\psi_a\psi_a$ observable flow from \texttt{obsInit}): \texttt{MultiConcl2}
holds conditional on \texttt{RGInput} and two named inputs on the $\mu^{(2)}$ flow, \texttt{Case2Rel} (the observable
remainders, $|\mathrm{err}_k|\le c_\delta\check g_k^2|\beta_k(1-\mathfrak n^{(2)}_k)|$) and \texttt{Case2Loop} (the one-loop
identification $|-\mathfrak n^{(2)}_k-\frac\beta3\check g_k+(q_{k+1}-q_k)|\le B(\check g_k^2+\check g_ka_k)$, complex $q$,
$|q_k|\le c_q\check g_k$, $|\mathfrak n^{(2)}_k|\le\min(\frac18,B'\check g_k)$); the phase of $\hat\lambda$ is controlled by
the imaginary part of $\sum_k\log(1+x_k)$ (\texttt{lamHat\_near}).
\texttt{Spin2.thm\_multiscale\_case2B} (\texttt{H02UV/Spin2/Case2/Case2Loop.lean}) replaces \texttt{Case2Loop} by the bulk input
\texttt{Case2Bulk} (independent of the vertex): with the exact $\mathfrak n^{(2)}_k=\nu'_kS_{k+1}-\nu^\flat_kS_k
-g_k(\mathfrak l_{w_{k+1}}-\mathfrak l_{w_k})-(\text{drift})$ (\texttt{nn2\_eq}; $S_j=\sum_{k<j}\sum_x(C_k(a,x)+C_k(x,a))$,
$w_j=\sum_{k<j}C_k$, $\nu'_k=\nu^\flat_k+4g_k\delta_k$), \texttt{Case2Loop} holds with $q_k=\nu^\flat_kS_k$
(\texttt{case2Loop\_of}) if $|g_k(\mathfrak l_{w_{k+1}}-\mathfrak l_{w_k})-\frac\beta3\check g_k|\le B_1(\check g_k^2+\check g_ka_k)$
(the one-loop bubble), $|(\nu^\flat_{k+1}-\nu'_k)S_{k+1}|\le B_2(\check g_k^2+\check g_ka_k)$ (second-order $\nu$ flow),
drift $\le B_3\check g_k^2$, $|\nu^\flat_kS_k|\le c_q\check g_k$, and $B'\check g_k\le\frac18$.
\texttt{Spin2.thm\_multiscale\_case2C} (\texttt{H02UV/Spin2/Case2/Case2Thm.lean}) discharges \texttt{Case2Bulk} and the
smallness under the master constants (\texttt{LoopConcl}, $N_h=N-2$): the history bubble equals the dressed bubble
$\frac{\check g_k}{1+n_k}(\mathfrak l(w'_k+C'_{k+1})-\mathfrak l(w'_k))$ of Proposition~\ref{prop:d2}
(\texttt{bubble\_eq}, Parseval \texttt{Obs2.lfunC\_kmat}), the drifts and $\nu^\flat S$ follow from gradient,
range and row-sum bounds of each $C_i$ at its own Gram scale (\texttt{H02UV/Obs2/Flow/Tau2Bd.lean}, \texttt{H02UV/Spin2/Case2/Case2Bd.lean}),
and the second-order $\nu$ flow from \texttt{rawStep\_coup} (\texttt{nu\_sharp}, \texttt{nuflow\_le}); the constants
$B_i$ are fixed before $\epsilon$. So \texttt{MultiConcl2C} holds conditional only on \texttt{RGInput} and \texttt{Case2Rel}.
\texttt{Spin2.thm\_multiscale\_case2D} (\texttt{H02UV/Spin2/Case2/Case2Rel.lean}) discharges \texttt{Case2Rel} under the
master constants for every point configuration and every vertex with $\gamma_a=0$, $\mu_a\ne0$
(\texttt{RelConcl2}, constant $c_\delta c_R^2$), via the Case~(2) relative recursion
(\texttt{PtCorr2.obs\_hist\_rel2}, \texttt{H02UV/Obs2/Flow/BlockFlowE.lean}; field scale $s_j^2$ of $m_2$) and the Case~(2)
product bound (\texttt{PtCorr2.prodBd2\_hist}); \texttt{MultiConcl2D} holds conditional only on \texttt{RGInput}.
Vertices with $\gamma_a\ne0$ (e.g.\ the pair vertex $1-b\bar\psi\psi$) reduce exactly to $\gamma_a=0$: the observable
constant $-\gamma_a\sigma_a$ is a scalar of the field algebra, $\mathbb E_C\theta$ fixes it, $F_C(G,\cdot)=0$ for constants, so it
passes entirely into $u^{\rm obs}_1$ and (with $K_0=0$) the factors $e^{-\kappa(B)}$ of $I_0(B)$, $I_1(B)$ cancel against
$e^{u^{\rm obs}(U)}$ in $K_1$ (\texttt{H02UV/Obs2/Map/ConstSplit.lean}); hence $V_j,\mathcal W_j,K_j$ ($j\ge1$), $\beta_k$, $\mathfrak n^{(2)}_k$, the errors
and the $K,\mathcal W$ part of $\rho_k$ coincide with those of $\gamma_a=0$ (\texttt{H02UV/PtCorr2/Obs/ConstSplitHist.lean}). So \texttt{Case2Rel}
holds for every vertex with $\mu_a\ne0$ (\texttt{RelConcl2G}, \texttt{H02UV/Spin2/Case2/Case2RelPairThm.lean}).

\begin{conjecture}[Continuum limit]\label{thm:limit}
With $\vartheta_\epsilon^2=\vartheta_N^2$, one has
$\vartheta_\epsilon^2(\log\frac1\epsilon)^{2/3}\to z_*\in(0,\infty)$. Moreover,
$G^{\rm ren}_{\beta_R}(x)=\lim_{\epsilon\to0}\vartheta_\epsilon^{-2}P_\epsilon(x)$ exists, uniformly
on compact subsets of $\T\setminus\{0\}$, and it is continuous and invariant under the
symmetries of $\T$.
\end{conjecture}

\paragraph{Status of Conjecture~\ref{thm:limit}.} The $z_*$ part is proved (\texttt{thm\_theta\_eps}). For fixed
$\beta_R$ with $1/\beta_R\le g_0(L)$, the existence of subsequential limits uniformly on compact subsets of
$\T\setminus\{0\}$, their continuity, nonnegativity, point-group invariance and short-distance law are proved
(Theorem~\ref{thm:limit-fixed}, Lean: \texttt{Spin2.thm\_limit\_fixed\_betaR}). Open: convergence of the full sequence
and independence of the choice of $b_N$, $m_N$.

\paragraph{Lean status (finite-volume symmetry and continuity).} At every finite $\epsilon$ the
symmetries are exact: $\P_\beta[y\leftrightarrow x]=P_\beta(x-y)$ and $P_\beta$ is invariant under
the point group of $\Z^2$ (\texttt{Spin2.PconnXY\_eq}, \texttt{Pconn\_sgnPerm}, \texttt{Pconn\_neg},
\texttt{Pconn\_swap}, \texttt{Pconn\_refl}), and $\beta\mapsto P_\beta(x)$ is continuous on
$[0,\infty)$ (\texttt{Pconn\_continuousOn}), all in \texttt{H02UV/Spin2/Multiscale/ConnSym.lean}.

\paragraph{Lean status ($\vartheta_\epsilon$ part; \texttt{H02UV/Spin2/Cont/ContTheta.lean}).}
With $\vartheta_\epsilon^2:=$ \texttt{thetaEpsSq}~$\beta_R\,b=(\beta_Rb)^{2/3}$ ($\beta_\epsilon=1/b$),
$\beta_1=$ \texttt{betaOne}~$=3/(2\pi)$ (\texttt{betaL L}~$=\beta_1\log L$) and
$z_*=$ \texttt{thetaStar}~$\beta_R=(2\pi\beta_R/3)^{2/3}$:
\texttt{Spin2.thm\_theta\_eps} gives $L_0$ and, for $L\ge L_0$, constants $M,g_0,K,N_0$ depending on $L$
only such that for $N\ge N_0$, master parameters, $m^2>0$, $4m^2L^{2N}\le1/\beta_R\le g_0$ and
\emph{every} $b$ in the bracket $[\texttt{bFin}(1/(2\beta_R)),\texttt{bFin}(2/\beta_R)]$ (this contains
the shooting solutions and every \texttt{H0Sol} under the a priori hypotheses of
\texttt{thm\_multiscale\_final\_h0\_all}, \texttt{ct\_main}):
$|\beta_\epsilon-N\beta_1\log L|\le K(1+\beta_R+\log N)$, and, if $2K(1+\beta_R+\log N)\le N\beta_1\log L$,
\[
\Bigl|\vartheta_\epsilon^2\bigl(\log L^N\bigr)^{2/3}-z_*\Bigr|\le z_*\,\frac{2K(1+\beta_R+\log N)}{N\beta_1\log L}.
\]
\texttt{Spin2.thm\_theta\_eps\_tendsto}: for fixed $\beta_R$ ($1/\beta_R\le g_0$), along any
$N_k\to\infty$ and bracket points $b_k$, $\vartheta_\epsilon^2(\log L^{N_k})^{2/3}\to z_*$. The route is
\texttt{MasterFacts.ct\_beta\_scale}: along the $h=0$ trajectory in $D_\epsilon$,
$|\beta_j-(N-2-j)\beta_1\log L|\le(T+E)+a\log(2((N-2-j)\beta_1\log L+T+E+a|\log2a|))$ for $j\le N-2$,
$\beta_{N-2}\le T=5\beta_R$ (from $\check g^0_{N-2}\ge\frac78\check g^c_{N-2}\ge 7/(32\beta_R)$ on the bracket),
by summing the per-step band of \texttt{SpinConclD} from $j$ to $N-2$, the $\check g$-sum by
\texttt{LogSumD}, inverted by \texttt{beta\_eps\_log}.

\paragraph{Lean status (asymptotic Arzel\`a--Ascoli; \texttt{Cont/AscoliStep.lean}).}
\texttt{Cont.ascoli\_asymp\_exists\_subseq\_punctured} (instance
\texttt{ascoli\_asymp\_exists\_subseq\_torus} for $\mathbb T^2\setminus\{0\}$): if $f_N:\mathbb T^2\to\R$
(not necessarily continuous) are bounded on each $\{|x|\ge r\}$ and asymptotically equicontinuous there
($\forall\eta\,\exists\delta,N_0$: $N\ge N_0$, $|x-y|<\delta\Rightarrow|f_N(x)-f_N(y)|\le\eta$), then a
subsequence converges uniformly on compact subsets of $\mathbb T^2\setminus\{0\}$ to a limit continuous
on $\mathbb T^2\setminus\{0\}$; limits inherit asymptotic symmetries and bounds
$|f_N(x)-h(x)|\le e(x)+o(1)$ (\texttt{ascoli\_asymp\_limit\_invariant},
\texttt{ascoli\_asymp\_limit\_abs\_sub\_le}). This is step~6 of the route to continuous limit points
of $\vartheta_\epsilon^{-2}P_\epsilon$; the equicontinuity input (B3, tasks T1--T6) is open.
The boundedness input is proved (\texttt{Spin2.thm\_bound\_eps}, \texttt{H02UV/Spin2/Cont/ContLimit1.lean}): for fixed $\beta_R$ and
$\delta>0$ there is $C$ with $\vartheta_\epsilon^{-2}P_\epsilon(x)\le C$ for all $N\ge N_0$, all $b$ in the bracket and all
$x\ne0$ with $|x|/L^N\ge\delta$ (from (C) of \texttt{thm\_short\_epsX}). A bound uniform over all $x\ne0$ is false
($\vartheta_\epsilon^{-2}P_\epsilon\approx(\beta_1N\log L/\beta_R)^{2/3}$ at $|x|=1$). Norms:
$|x|_\infty\le|x|\le\sqrt2|x|_\infty$ (\texttt{cx\_supNorm\_eucl}).

\paragraph{Lean status (B3 equicontinuity, tasks T1--T3).}
\emph{Shift} (\texttt{H02UV/Spin2/Cont/ContShift.lean}): \texttt{cs\_shift}: $P(x)=\P[t\leftrightarrow t+x]$,
$P(y)=\P[t-w\leftrightarrow t+x]$, $w=y-x$; for the block centre $t=$ \texttt{csCentre}~$m$ and
$4|w|_\infty\le L^m$ (\texttt{cs\_four\_le}: from $|w|_\infty\le L^{m-r}$, $4\le L^r$), $t$ and $t-w$ lie in
the same block at every scale $k\ge m$ (\texttt{cs\_blk\_pair}, \texttt{cs\_blk\_ball}, \texttt{cs\_blk\_ball'});
$j_c(t,t+x)=j_c(0,x)$, $j_c(t-w,t+x)=j_c(0,y)$ (\texttt{cs\_coal\_eq}) and $|x|_\infty\ge L^{N-c}$ gives
$j_c\ge N-c$ (\texttt{cs\_coal\_ge}, \texttt{cs\_coal\_gt}, \texttt{cs\_coal\_ge\_phys}).
\emph{Scale-$k$ one-point states at any vertex position} (\texttt{H02UV/Spin2/Cont/ContScaleM.lean}, from \texttt{MasterFacts}):
\texttt{csm\_hlam} ($\|\hat\lambda^P_k-\Theta_k\|\le K_\lambda\Theta_k(\check g_k+a_k)$, $\Theta_k=(\check g_k/b)^{1/3}$
independent of $P$), the grid artefact \texttt{csm\_artefact}/\texttt{csm\_artefact\_rel}
($\|\beta^P_k-\beta^{P'}_k\|\le4K_\lambda(\check g_k+a_k)\|\beta^{P'}_k\|$, any two positions) and the $K$ part
\texttt{csm\_K} ($\le16c_WB_k(c\check g_k)^2s_k^2\|\beta_k\|$).
\emph{Relocation} (\texttt{H02UV/Obs2/Loc/Reloc.lean}): \texttt{reloc\_pbp}:
$\|\bar\psi_a\psi_a-\bar\psi_{a'}\psi_{a'}\|_{T_j(\ell)}\le2d|a-a'|_\infty\ell^2L^{-j}$ (\texttt{reloc\_coupling} with a
coupling); \texttt{reloc\_exact}: the exact step is independent of the Loc point. The restart theorem
for the one-point flow (T4) and the pair/top part (T5, top coalescence) are in the next paragraphs; the assembly (T6) is open.

\paragraph{Lean status (B3: support lemma and restart of the one-point flow, T4).}
\emph{Support} (\texttt{H02UV/PtCorr2/Restart/RestartLoc.lean}): $X^\square$ is a polymer, so along the one-point history the
$\sigma_a$, $\sigma_a\sigma_b$ parts of $K_j(Y)$ vanish unless $Y^\square$ contains the whole $j$-block of $a$
(\texttt{rs\_K\_a\_zero\_of\_block}); for $a,a'$ in one $m$-block, \texttt{OLocHyp} holds at the relocated
point for every step $j\ge m$ (\texttt{rs\_oLocHyp\_hist}), hence (\texttt{reloc\_exact}) the steps $j\ge m$ of the
$a$-history may be computed with the Loc point $a'$ (\texttt{rs\_reloc\_hist}).
\emph{Restart} (\texttt{H02UV/Obs2/Flow/RestartRec}, \texttt{RestartChain}, \texttt{RestartFlow}, \texttt{H02UV/PtCorr2/Restart/RestartStep},
\texttt{RestartHist}, \texttt{RestartProd}, \texttt{RestartThm}, \texttt{RestartBulk}, \texttt{H02UV/Spin2/Restart/RestartMaster}):
the relative recursion is replaced by an envelope $M_j=M_0\prod_{k<j}|1-\mathfrak n^F_k|$ with a free
$M_0>0$ dominating the datum ($|\beta_m|+D_m\le M_0$, $\|\mathcal W_m\|\le c_Wr_m2s_m^2M_0$,
$\|K_m\|\le k_0$ with $8k_0\le(c_k-60B)r_m^2c_W2s_m^2M_0$), and the transients (initial $K$, tadpoles of the
initial $\mathcal W_m$, memory of the scales $<m$ in $\mathfrak n^F$) decay geometrically
(\texttt{RelRecK.close}). \texttt{Spin2.MasterCore.restartD}: at the master constants (no new requirement), along
every trajectory in $D_\epsilon$, for every datum $o'$ at a scale $1\le m\le N_h$ (observable, local, symmetric,
$V_m(B)_a=c(B)\bar\psi_a\psi_a$) with $c_t=8c_Nk_0/(s_m^2M_0)+x_h+\frac{8}{3L}\le\frac1{48}$,
$x_h=c_{bn}(1280K^2\frac8{7L}+1)2c_Wc\check g_m$:
\[
  |\beta_{m+j}|\le\bigl(|\beta_m|+(\tfrac18+10c_t)M_0\bigr)\prod_{k<j}|1-\mathfrak n^F_{m+k}|\qquad(j\le N_h-m),
\]
$\mathfrak n^F_i=\mathfrak n^{(2)}_i-\sum_{k<m}\tau_{k,i}$; the $\mathcal W$/$K$ sizes in the envelope are in
\texttt{PtCorr2.obs\_hist\_restartB}. Open for T6: the linearity of the observable step in the datum
($\sigma_a^2=0$; not formalised) and the re-representation of the relocated coupling
$\lambda\sigma(\bar\psi_a\psi_a-\bar\psi_{a'}\psi_{a'})$ as a polymer activity (the restart theorem needs
$V_m(B)_a\propto\bar\psi_{a'}\psi_{a'}$).

\paragraph{Lean status (B3, T5: the pair/top restart, top coalescence).}
(\texttt{H02UV/PtCorr2/Restart/RestartPairTop}, \texttt{H02UV/Spin2/Restart/RestartPairTop1}, \texttt{RestartPairSeed}, \texttt{RestartPairTop2}.)
For the pair $(0,x)$ with $j_x\ge N_h+1$ and an \emph{arbitrary} local observable state $o$ at scale $N_h+1$
(e.g.\ the restarted history of a small $a$-side datum), the two-point output is
$\mathrm{out}(o)=-\epsilon(u^{\rm obs}_a)\,\gamma_{x,N-1}-n_{ab}(\mathrm{last}(o))$ (\texttt{rpTopOut}); on the
reference state it is $b\langle\bar\psi_0\psi_x\rangle$ (\texttt{rpTopOut\_ref}). The bulk part of $F_N$ does not
see the observables (\texttt{rp\_bulk\_finO}), so the reduction at $N-1$ holds for $o$
(\texttt{rp\_pair\_red\_gen}); the zero mode is bounded with an absolute $a$-side, normalised by the reference
one-point data ($|c|\le4|\gamma_{a,N}||s_0|/|\beta|$, \texttt{rp\_zero\_abs}); the seed of the reduced state
is $O(1)$ in $wM(M_a,M_b)$ for any $o$ (\texttt{rp\_topSeed\_gen}). \texttt{Spin2.MasterFacts.rp\_top\_traj}: under
the hypotheses of \texttt{fw\_top\_traj}, if $o$ has sizes $v_T,w_T,k_{xT},k_2$ in $wM(2s_{N_h}^2\sigma,M_b)$ and
$|\epsilon(u^{\rm obs}_a)|\le\gamma_\delta$, then
\[
  |\mathrm{out}(o)|\le\tfrac32(1+K_O)\gamma_\delta\vartheta_{N_h}
  +\tfrac94K_E'\,\frac{\sigma}{|\beta_{a,N_h}|}\,\vartheta_{N_h}^2\check g_{N_h},
\]
linear in the $a$-side data $(\gamma_\delta,\sigma)$ ($\sigma\asymp|\delta\lambda|+\|\delta K\|/s^2$).
Open for T5: the sizes of the restarted state at $N_h+1$ from the datum at $m$ (the $b$-sector equals
the reference; the $ab$ chain from $m$; the step $N_h\to N_h+1$; the accumulated $u_a$), and
coalescence $N-C(\delta)\le j_x\le N_h$ (the pair chain above coalescence with a general seed).

\paragraph{Lean status (B3, T6 inputs: linearity in the datum, re-representation, restart identity).}
\emph{Linearity} (\texttt{H02UV/Obs2/Map/LinearMrg}, \texttt{H02UV/Obs2/Map/LinearStep}, \texttt{H02UV/PtCorr2/Obs/LinearHist}): in Case~(2) the
substitution $\sigma_b\mapsto\sigma_a$, $F\mapsto F_\emptyset+\sigma_a(F_a+F_b)$, is an algebra endomorphism
($\sigma_a^2=0$) commuting with every map of the step when $P=(a,a)$ (\texttt{mrg\_obsStep}); with the one-source
projections this gives, for observable data $o_1,o_2$ and the same bulk history, additivity of the
$\sigma_a$ parts of the step (\texttt{lin\_obsStep}), of the history from $m$ (\texttt{lin\_obsHistFrom},
\texttt{\_sub}), of the input $Z_k$ (\texttt{lin\_ZhOF\_sub}) and of the chain $\Pi'^{\ge k}$ (\texttt{lin\_prodHiF},
\texttt{\_sub}); the abstract form is \texttt{lin\_a\_of\_comm}.
\emph{Re-representation} (\texttt{H02UV/PtCorr2/Restart/RerepZ}, \texttt{RerepState}): for $I=I'+\delta$ with
$\delta(B)\delta(B')=0$ and $K'(U)=K^U+\sum_{B\subseteq U}\delta(B)K^{U\setminus B}$,
$\sum_XI^{\Lambda\setminus X}K^X=\sum_UI'^{\Lambda\setminus U}K'^U$ exactly (\texttt{rr\_gZO}; $K'$ is multiplicative
on non-touching polymers, \texttt{rr\_mul}). With $\delta(B)=-\sigma_a\beta_B(\bar\psi_a\psi_a-\bar\psi_{a'}\psi_{a'})I'(B)$
this moves the relocated coupling into the activity: \texttt{rrSt} has $V'(B)_a=\beta_B\bar\psi_{a'}\psi_{a'}$,
the same total coupling, the same integrand $Z_m$ (\texttt{rr\_ZinO}), locality at the point $a'$
(\texttt{rrSt\_obsLocInv\_core}), and correction terms of $T_m(\ell)$ norm
$\le|\beta_B|\,4|a-a'|_\infty\ell^2L^{-m}\|I'(B)_\emptyset\|\,\|K^{Y\setminus B}_\emptyset\|$ (\texttt{rr\_corr\_norm},
\texttt{rr\_delta\_K\_norm}).
\emph{Restart identity} (\texttt{rr\_restart\_identity'}): for data $o$ (Loc point $a$) and $o'$ (Loc point $a'$,
same $m$-block) along a trajectory,
\[
  \Pi'^{\ge m}_{a}(o)_a-\Pi'^{\ge m}_{a'}(o')_a=\Pi'^{\ge m}_{a'}(\mathrm{rrSt}\,o-o')_a ,
\]
and for the one-point states of two vertex positions the difference datum has total coupling
$\beta^P_m-\beta^{P'}_m$, $|\beta^\delta_m|\le4K_\lambda(\check g_m+a_m)|\beta^{P'}_m|$ (\texttt{H02UV/Spin2/Cont/ContDiff},
\texttt{cd\_delta\_beta\_le}). Open for T6: the $\sigma_a\sigma_b$ (pair) part is bilinear, so linearity of the
two-point output in the $a$-datum at fixed $b$-datum still needs a three-source version; the size of
$\|\delta K\|$ ($\lesssim|\beta_m|(\check g_m^2+L^{-r})s_m^2$) is larger than the restart's
$k_0\lesssim\check g_m^2s_m^2M_0$, so $M_0\gtrsim|\beta_m|(\check g_m+L^{-r}\check g_m^{-2})$.

\paragraph{Lean status (B3, T6 preparation: sizes, symmetry, history link).}
\emph{History link} (\texttt{H02UV/PtCorr2/Restart/HistLink}): restarting the history at $m$ from its own state reproduces it,
$\mathrm{obsHistFrom}_m(\mathrm{obsHist}_m)_k=\mathrm{obsHist}_k$ for $k\ge m$ (\texttt{obsHistFrom\_obsHist}); hence
the restarted size predicates of such a datum are those of the history (\texttt{kbHF\_obsHist\_iff},
\texttt{wbHF\_obsHist\_iff}).
\emph{$U(1)$ invariance} (\texttt{H02UV/PtCorr2/Restart/RerepSym}): $\bar\psi_a\psi_a$ is invariant, so the re-represented state is
(\texttt{rrSt\_obsSymInv\_core}); with locality this gives all structural hypotheses of the restart theorem for
$\delta=\mathrm{rrSt}\,o-o'$.
\emph{Weighted activity} (\texttt{H02UV/PtCorr2/Restart/RerepK}): with the large-set regulator $A^{|\mathcal B(W)|}\|K^\flat(W)\|\le1$
and $\|I'(B)_\emptyset\|\le\iota$,
$A^{(|\mathcal B(Y)|-4)_+}\|K^\delta(Y)\|_{w_A}\le k_1+k_2+\bar\beta\,4|a-a'|_\infty s_m^2L^{-m}\,\iota A$
(\texttt{rr\_delta\_K\_wt}, \texttt{rr\_delta\_KbHF}), $k_{1,2}$ the bounds of $K^{o},K^{o'}$, $\bar\beta=\sum_B|\beta_B|$.
\emph{Restart compatibility} (\texttt{H02UV/Spin2/Cont/ContDiff2}): with $k_{1,2}\le16c_WB(c\check g_m)^2s_m^2|\beta^{P,P'}_m|$
(\texttt{csm\_K}) the condition $8k_0\le8Bc_W(c\check g_m)^2s_m^2M_0$ of \texttt{restartD} holds for
$M_0=16(|\beta^P_m|+|\beta^{P'}_m|)+b_0+E/(Bc_W(c\check g_m)^2s_m^2)$ (\texttt{cd2\_compat}) and \emph{forces}
$M_0\ge16(|\beta^P_m|+|\beta^{P'}_m|)$ (\texttt{cd2\_forced}): this corrects the estimate above, and the restart theorem
with the master value $c_k=64B$ gives no smallness for the difference datum. With a scale-dependent
$K$-step constant $c_k\asymp1/(c\check g_m)$, which the step estimate allows if its smallness conditions are read per step
($c_k\,c_R g_j\lesssim\varepsilon_0$ instead of $c_k\,c_R\varepsilon$), the envelope becomes
$M_0\asymp\check g_m|\beta_m|+L^{-r}|\beta_m|/\check g_m$ (\texttt{cd2\_compat\_ck}); this variant of
\texttt{restartD} is open.

\paragraph{Lean status (B3, T4$'$/T6: restart with large transients; the small output of $\delta$).}
(\texttt{H02UV/Obs2/Flow/RestartRec2}, \texttt{RestartRec2Flow}, \texttt{H02UV/PtCorr2/Restart/Restart2Step}, \texttt{Restart2Hist},
\texttt{Restart2Thm}, \texttt{Restart2D}, \texttt{H02UV/Spin2/Restart/Restart2Master}, \texttt{Restart2Lin}, \texttt{Restart2DMaster},
\texttt{ContDiff3}.) Route (A) is carried out. The $K$-step bound is linear in the activity size $m$, so a transient
is handled by taking the size $m'=m_j+\tilde w_j/r_j+8\tilde k_j/(\bar\kappa r_j)$ and a per-step constant
$c=60B+8\tilde k_j/(r_j^2m')$ with $c\,r_j\le c_kr_j+\bar\kappa$ (smallness read at the actual coupling,
$(8c_x+c)r_j^3\le\varepsilon_0$, \texttt{kobs\_state\_step2}); $\bar\kappa=64B(c\varepsilon)+\eta$ needs no new master
requirement ($c\varepsilon\le1/72$). \texttt{RelRecT}: with $\|\mathcal W_0\|\le c_Wr_0 2s_0M_0+w_0$, $K_0\le k_0$,
$k_0\ge120B\bar r w_0$, the transients decay as $6^{-j}$, $(8/15)^j$ and the envelope only has to dominate
$|\beta_0|+D_0$, $c_Kk_0/s_0$, $x_hM_0$, $5c_P\bar rw_0/s_0$. Master form (\texttt{MasterCore.restartDT},
\texttt{restartDT\_lin}): if $c_{bn}(\cdot)2c_W c\check g_k\le1/1344$, $|\beta_k|+D_k\le M_0$,
$1344(8c_Nk+\Gamma_ww)\le s_k^2M_0$ then $|\beta_{k+j}|\le\frac43M_0\prod_{i<j}|1-\mathfrak n^F_{k+i}|$.
The one-point share deviation is $O(\check g)$: $D_j\le6c_\delta c^2(\check g_j/(\beta/4))|\beta_j|$
(\texttt{MasterCore.relDev}, from $\sum_{k<j}\check g_k^2\le(\check g_j-\check g_0)/(\beta/4)$). For the difference
datum $\delta$ (\texttt{cd3\_restart}, \texttt{MasterFacts}): $0<k<N_h$, $a,a'$ in one $k$-block,
$|a-a'|_\infty\le L^{k-r}$, $K_\lambda(\check g_k+a_k)\le\frac12$, the scale condition, and the bulk inputs
$\|I'(B)_\emptyset\|\le\iota$, $A^{|\mathcal B(W)|}\|K^\flat_k(W)\|\le1$ (still named):
\[
|\beta^\delta_{k+j}|\le\tfrac43M^\delta\prod_{i<j}|1-\mathfrak n^F_{k+i}|,\qquad
M^\delta\lesssim(\check g_k+a_k+\iota A L^{-r})(|\beta^P_k|+|\beta^{P'}_k|),
\]
explicitly \texttt{MasterFacts.cd3M}. Route (B) (bounding $K^P_m-K^{P'}_m$ directly) would need a separate grid-artefact
analysis of $K,\mathcal W$ over all scales $<m$ and was not pursued.

\paragraph{Lean status (B3, T6: bulk inputs discharged; the relative form).}
(\texttt{H02UV/PtCorr2/Restart/Restart3a}, \texttt{Restart3b}, \texttt{H02UV/Spin2/Cont/ContDiff4a--c}.) The two bulk inputs of
\texttt{cd3\_restart} come from \texttt{MasterFacts}: the bulk part of $I'(B)$ sees only the bulk parts
($\mathrm{bulk}$ is an algebra map, $\mathrm{bulk}\,e^F=e^{\mathrm{bulk}F}$), so
$\|I'(B)_\emptyset\|\le e^{3/2}(1+1)=:\iota_0$ by $\|V^\flat+\Delta\|\le\frac32$, $\|W^{\rm lin}+D\|\le1$
(\texttt{MasterFacts.cd4\_iota}); the regulator $A^{|\mathcal B(W)|}\|K^\flat_k(W)\|\le1$ follows from
$\|K_k\|\le c_x\check g_k^3$ in $D_\varepsilon$ under the scale condition $A^4c_x\check g_k^3\le1$
(\texttt{cd4\_reg}). Hence \texttt{cd3\_restartM}: the statement of \texttt{cd3\_restart} with $\iota=\iota_0$ and no
named hypotheses beyond the scale conditions. The memory of the restarted one-loop coefficient is summable,
$|\mathfrak n^F_{k+i}-\mathfrak n_{k+i}|\le\frac43q2^{-i}$, so with $|\mathfrak n|\le\frac18$
$\prod_{i<j}|1-\mathfrak n^F_{k+i}|\le e^{\frac87\cdot\frac83q}\prod_{i<j}|1-\mathfrak n_{k+i}|$, $q=2/L$
(\texttt{cd4\_prod\_cmp}); and $e^{-2c_\delta(4/\beta_L)\varepsilon}|\beta_k|\prod_{i<j}|1-\mathfrak n_{k+i}|\le|\beta_{k+j}|$
(\texttt{cd4\_ratio}, from \texttt{bt2\_flow}, \texttt{Case2Rel}, $\sum\check g^2\le(4/\beta_L)\check g$). Relative form
(\texttt{cd3\_restartRel}): $|\beta^\delta_{k+j}|\,|\beta^{P'}_k|\le\frac43M^\delta e^{\frac{128}{21L}}e^{8c_\delta\varepsilon/\beta_L}|\beta^{P'}_{k+j}|$.

\paragraph{Lean status (B3, T5 cont.: the restarted state at $N_h+1$).}
(\texttt{H02UV/PtCorr2/Restart/RestartPairB}, \texttt{RestartPairAB}, \texttt{RestartPairGam}, \texttt{RestartPairGamB},
\texttt{H02UV/Spin2/Restart/RestartPairTopF1}.) Along $o_k=\mathrm{obsHistFrom}_m(\delta)_k$: if the datum carries the reference
$b$-part, so do all $o_k$, componentwise (\texttt{rpb\_b\_eq}, via the $\sigma_b$ projection); the coalescence
invariant, $V_{ab}=u_{ab}=0$ up to $\min(j_{ab},N_h+1)$, locality and $U(1)$ invariance up to $N_h+1$
(\texttt{rpb\_coalInv}, \texttt{rpb\_coal\_validL}, \texttt{rpb\_inv\_top}). The $ab$ chain restarts from a datum with
$ab$ size $\le c_{ab}r_mm_m^2$ ($c_{ab}\ge10B_2$) and keeps $\le c_{ab}r_jm_j^2$ up to $N_h$, and
$\le\frac{c_{ab}}5r_{N_h}m_{N_h}^2$ at $N_h+1$ (\texttt{rpab\_hist}, \texttt{rpab\_top}). The restarted $\gamma$-flow is
exactly $\gamma^F_{k+1}=\gamma^F_k-\sum_B\epsilon(V_k(B)_a)-\beta^F_kc_k+\rho^{WK}_k$ with the \emph{reference}
bulk factor $c_k=C_k(a,a)-\epsilon F_{C_k}(\bar\psi_a\psi_a,V'_k(\Lambda)_\emptyset)$ (\texttt{rpg\_gam\_succ}), so
$|\gamma^F_n|\le|\gamma^F_m|+|\sum_B\epsilon(V_m(B)_a)|+\sum_{m\le k<n}(|\beta^F_k||c_k|+|\rho^{WK}_k|)$
(\texttt{rpg\_gam\_le}), $|\rho^{WK}_k|\le c_N8\kappa_k(1+|C_k(a,a)|s_k^{-2})+c_BQ_1w_k$ for $m\le k\le N_h$
(\texttt{rpg\_gamWKF\_le}) and $|c_k|\le(1+c_e(\check g_k+z_k))|\Gamma_{k+1}(0,0)|$ for $k<N_h$
(\texttt{rpg\_C\_le}). Open for T5: the one-point $a$-sizes of the restarted state at $N_h+1$ (one more
restart step under \texttt{ValidAtL}) and $c_{N_h}$; coalescence $N-C(\delta)\le j_x\le N_h$; the combined bound.

\paragraph{Lean status (B3, T5, item (5): coalescence $j_x\le N_h$ for a general state).}
(\texttt{H02UV/PtCorr2/Restart/RestartPairChain1}, \texttt{H02UV/Spin2/Restart/RestartPairChain1--3}.) For an arbitrary observable state $o$ at
$j=j_x\le N_h$, local at $j$, the output is $\mathrm{out}(o)=-\epsilon(u_a(o))\,\gamma_{x,j}-n_{ab}(\omega_{\Gamma_j}Z_j(o))$;
on the reference state it is $b\langle\bar\psi_0\psi_x\rangle$ (\texttt{rpcOut\_ref}). Exactly: $\mathrm{sE}$ does not
see the state, and with $o'=\mathrm{red}_j(o)$ (reduced at $j$, covering polymer $D_x$),
$n_{ab}(\omega_{\Gamma_j}Z_j(o))=n_{ab}X'+n_{ab}\Omega'$ (\texttt{rpc\_pair\_red\_gen}); cut at $N_h+1$,
$\mathrm{out}(o)=-\epsilon(u_a(o))\gamma_{x,j}-n_{ab}(\prod_{j\le i\le N_h}\chi'_i)+\mathrm{out}_{\rm top}(o'')$ with
$o''$ the reduced chain at $N_h+1$ (\texttt{rpc\_out\_split}). The seed: from sizes $v_T,w_T,k_{xT},k_2$ of $o$ in
$wM(M_a,M_b)$ at $j$, $A^{(|\mathcal B(X)|-4)_+}\|K'_j(X)\|_{wM}\le(2\iota A)^{(L+18)^2}$, $\iota=e^{3/2+v_T}(2+w_T)$
(\texttt{rpc\_seed\_gen}). The zero mode: from the $wM(2s_{N_h}^2\sigma,M_b)$ sizes $c_V,c_W\check g,c_K$ of $o''$,
$|n_{ab}\Omega'|\le8(\frac12+\frac{81}{32}c_{Jab}c_S)c_uK_4(\sigma/|\beta_{a,N_h}|)\|M_{N_h}\|^2\check g_{N_h}$
(\texttt{rpc\_zm\_le}), hence $|\mathrm{out}(o)|\le|\epsilon(u_a(o))||\gamma_{x,j}|+|n_{ab}X'|+$ that
(\texttt{rpc\_out\_le}); linear in the $a$-side. Open: the sizes of the reduced chain ($o''$ and $n_{ab}X'$) from a
general seed, i.e.\ the pair chain above coalescence (\texttt{PairJointSizesP3}/\texttt{pl6\_*}) with a general
seed --- in Lean that layer is stated for the reference reduced state \texttt{oRed} only.

\paragraph{Lean status (B3, T5, items (2), (4$'$): the restarted state at $N_h+1$ in the top bound).}
(\texttt{H02UV/PtCorr2/Restart/RestartPairTopF1}, \texttt{H02UV/Spin2/Restart/RestartPairTopF1--2}.) One more restart step at $N_h$ under
\texttt{ValidAtL}, for a general envelope: the $K$-step (per-step smallness $8c_br^3+(r/m)8k\le\epsilon_0$, no
uniform $\epsilon$) and the full-norm $\mathcal W$-step along $\mathrm{obsHistFrom}$ (\texttt{rptf\_kobs\_stepF},
\texttt{rptf\_wfull\_stepF}); $|\beta^F_{N_h+1}|\le|1-\mathfrak n^F_{N_h}||\beta^F_{N_h}|+E_{N_h}$ and
$\sum_B|m_2V^F_{N_h+1}(B)_a|\le|1-\mathfrak n^F_{N_h}|(|\beta^F_{N_h}|+D_{N_h})+3E_{N_h}$ (\texttt{rptf\_beta\_shA2});
$T(V^F_j(B)_a)\le s_j^2\sum_B|m_2V^F_j(B)_a|$ (\texttt{rptf\_vobs\_block}). The top bulk factor:
$|c_{N_h}|\le|\Gamma_{N_h+1}(0,0)|+K_b\check g_{N_h}/|Z_{N_h+1}|$ (\texttt{rptf\_C\_top\_le},
\texttt{MasterFacts.rptf\_C\_topD}); the accumulated constant $|\epsilon(u^F_{N_h+1,a})|\le\gamma_\delta$ with
$\gamma_\delta$ the exact sum of \texttt{rpg\_gam\_le} up to $N_h$, bounded by the restart sizes
(\texttt{rptf\_gam\_top}, \texttt{rptf\_gam\_sizes}). Combined (\texttt{MasterFacts.rptf\_top\_traj}): for the
restarted state $o=\mathrm{obsHistFrom}_{m}(\delta)_{N_h+1}$ (datum with the reference $b$-part), given its sizes
at $N_h$ in $w=wM(2s_{N_h}^2\sigma,M_b)$ ($\|V\|_w\le m$, $\|\mathcal W\|_w\le rm$, $\|K\|_{wM1}\le k_1\le c_1r^2m$),
the share data $|1-\mathfrak n^F_{N_h}|\le a_N$, $|\beta^F_{N_h}|+D_{N_h}\le b_D$, $E_{N_h}\le e$, and the $ab$ size $k_2$
at $N_h+1$ (\texttt{rpab\_top}), all state hypotheses of \texttt{rp\_top\_traj} hold with
$v_T=\frac{81}{128}\frac{a_Nb_D+3e}\sigma+\frac32(\frac{81}{64})^2$, $w_T=e^{L^2C_w(r+m^2)}-1$,
$k_{xT}=(\theta_0+C_1(2+16c_x+8c_1)r)8k_1+Br^2m$, $k_2$, $\gamma_\delta$; hence the bilinear top bound for $o$.
Open: the restart sizes at $N_h$ themselves from the variant restart (and $E_{N_h}$, an \texttt{E\_top}-type
bound of the restarted \texttt{RelRecK}).

\paragraph{Lean status (B3, T5, item (5)(iii): the reduced chain of a general state, crude route).}
(\texttt{H02UV/Spin2/Restart/RestartPairChain4--7}.) The $V$-step of the joint induction and its top version
(\texttt{MasterCore.vstepD}, \texttt{vstepTopD}) are state-general; applied along the reduced chain of a general
state $o$ at $j_x\le N_h$ in a \emph{constant} weight $wM(M_a,M_b)$ (\texttt{rpc\_chain\_sizesD}): there are
$C_{1w},C_{1t}\ge0$ such that from the seed $A^{(|\mathcal B(X)|-4)_+}\|K'_{j_x}(X)\|_{wM}\le k_0$, a target
$\bar m>0$ with $8k_0\le\mathrm{vck}(C)(4C+1)\bar m$ ($C=C_{1w},C_{1t}$) and per-block masses
$|m_2V_n(B)_{ab}|\le\rho_n$, $(M_aM_b)^{-1}\rho_ns_n^2\le\bar m/2$ ($j_x<n\le N_h+1$): $\|V_n(B)\|_{wM}\le\bar m$,
$\|\mathcal W_{N_h+1}\|_{wM}\le(4C_{1t}+1)c\check g_{N_h}\bar m$, $\|K_{N_h+1}\|_{wM}\le(4C_{1t}+1)\bar m$. No smallness of
$\bar m$ is needed (the reduced chain is linear in its seed). The reduced constants exactly
(\texttt{rpc\_cst\_le}): $n_{ab}X'=\sum_{j_x\le k\le N_h}(-C_k(0,0)r'_k+d'_k)-C_N(0,0)r'_N+d'_N$ (no vertex, no
$\gamma$-cross terms), $|d'_N|\le c_{Lab}c_S\,2s_{N_h}^2\sigma M_b$. The seed from the one-point sizes of $o$ at $j_x$
under the master constants, $k_0=(2\iota A)^{(L+18)^2}$ (\texttt{MasterFacts.rpc\_seed\_MFD}). Combined
(\texttt{MasterFacts.rpc\_out\_traj\_sz}, \texttt{rpc\_out\_trajC}, \texttt{rpc\_out\_trajS}): along the $h=0$
trajectory, for $1\le j_x\le N_h$, $0<M_b\le c_u\|M_{N_h}\|\check g_{N_h}$, $|\epsilon(u_a(o))|\le\gamma_\delta$,
\[
|\mathrm{out}(o)|\le\gamma_\delta|\gamma_{x,j_x}|+|n_{ab}X'|+K_Z(c_S)\,\frac{\sigma}{|\beta_{a,N_h}|}\,
\vartheta_{N_h}^2\check g_{N_h},\qquad c_S=\bar m+(4C_{1t}+1)(c+8)\bar m+1,
\]
$K_Z(c_S)=\frac94\cdot8(\frac12+\frac{81}{32}c_{Jab}c_S)c_u\cdot2(1+K_F+K_{CL})$; and $|\gamma_{x,j_x}|\le(1+K_M\check g)
\|M_{j_x}\|$ from \texttt{ZBelowU1} (\texttt{rpc\_uB\_le}). Open: the per-block masses $\rho_n$ and the
$r'_k$, $d'_k$ ($k\le N_h$) of the reduced chain (its $r'$ flow \texttt{rRX\_flow} and $E$-step sources),
for a general seed.

\paragraph{Lean status (B3, T5, item (5)(iii) closed: masses, constants, the combined bound).}
(\texttt{H02UV/PtCorr2/Restart/RestartPairChain3}, \texttt{H02UV/Spin2/Restart/RestartPairChain8--12}.) \emph{Masses, one step at a time:}
exactly, $m_2(V_{k+1}(B_+)_{ab})=\sum_{B\subset B_+}(m_2(e^{\mathcal L_C}V_k(B)_{ab})-m_2(e^{\mathcal L_C}Q_k(B)_{ab}))
-[B_+\cap\{a,b\}\neq\emptyset]\,m_2(S_k(B_+)_{ab})$ (\texttt{rpc3\_mass\_Vnext}), hence
$|m_2V_{k+1}(B_+)_{ab}|s_{k+1}^2\le M_aM_b\Lambda_1\bar m$ whenever $\|V_i\|_{wM}\le\bar m$ for $j_x\le i\le k$
(\texttt{MasterCore.rpc8\_massD}; $\Lambda_1=\Lambda_1(L)$). No $\beta$-relative bound and no Gronwall are needed; the
joint induction (\texttt{rpc9\_chainD}) runs with growing targets $\bar m_i=\max(1,2\Lambda_1)^i\bar m_0$ over the
$N_h+1-j_x$ steps. \emph{Defects:} the $\epsilon$-tadpole $\epsilon F_C(V'_{ab},V'_\emptyset)=\sum_{B_+}\epsilon S(B_+)_{ab}
-\sum_B\epsilon(e^{\mathcal L_C}\mathcal W(B))_{ab}$ (no $\sigma_a$ parts), so $|d'_k|\le M_aM_b\Lambda_2\bar m$
(\texttt{rpc8\_def\_le}); with $|C_k(0,0)|\le s_k^2$, $|n_{ab}X'|\le M_aM_b\bigl((N_h+1-j_x)(2\Lambda_1+\Lambda_2)\bar m_{\rm top}
+L^2\Lambda_1\bar m_{\rm top}+c_{Lab}c_S\bigr)$ (\texttt{rpc10\_cst\_bd}). \emph{Combined}
(\texttt{MasterFacts.rpc11\_T5}, \texttt{rpc11\_T5$\delta$}, \texttt{rpc12\_T5S}; \texttt{ZBelowU1} exported as
\texttt{rpc11\_zBelowD}): for $N-C_\delta\le j_x\le N_h$,
\[
|\mathrm{out}(o)|\le\tfrac32(1+K_M)\,\gamma_\delta\,\vartheta_{j_x}+K_T(L,C_\delta)\,\frac{\sigma}{|\beta_{a,N_h}|}\,
\vartheta_{N_h}^2\check g_{N_h},
\]
with the one-point (restart) sizes of $o$ at $j_x$ in $wM(2s_{N_h}^2\sigma,M_b)$ as the only state input
(and $M_aM_b\le c_u^2(\sigma/|\beta|)\|M_{N_h}\|^2\check g^2$ from $2s_{N_h}^2|\beta_{a,N_h}|\le c_u\|M_{N_h}\|\check g_{N_h}$).

\paragraph{Lean status (B3, T6 (c): subsequential limits, conditional on equicontinuity).}
(\texttt{Cont/LatEmb}, \texttt{H02UV/Spin2/Cont/ContLimit1}, \texttt{H02UV/Spin2/Cont/ContLimit}.) Put
$f_N(x)=\vartheta_\epsilon^{-2}P_\epsilon(\lfloor L^Nx\rceil)$, that is
\texttt{clF}~$L\,\beta_R\,b\,N\,x=$ \texttt{Pconn}~$L\,N\,(1/b_N)\,(\texttt{latPtN}\,L\,N\,x)/\texttt{thetaEpsSq}\,\beta_R\,b_N$.
Translation invariance is built in, since $P_\epsilon(x)=\P[0\leftrightarrow x]$.
\texttt{Spin2.thm\_limit\_of\_equi} gives $L_0$ and, for $L\ge L_0$, a master construction $M$ and constants $g_0,N_0$
depending only on $L$. Fix $\beta_R$ with $1/\beta_R\le g_0$. There are then $\delta_0\in(0,1)$ and $K'\ge0$ with the
following property. Take any admissible choice (\texttt{ClAdm}): for $N\ge N_0$, master parameters $R_N$, $m_N^2>0$ with
$4m_N^2L^{2N}\le1/\beta_R$, and $b_N>0$ in the bracket. If $(f_N)$ is \emph{asymptotically equicontinuous} on every
$\{|x|\ge r\}$ (\texttt{ClEqui}), then a subsequence converges uniformly on compact subsets of $\T\setminus\{0\}$ to
some $G$. This $G$ satisfies:
\begin{itemize}
\item $G$ is continuous on $\T\setminus\{0\}$ and $G\ge0$;
\item $G$ is invariant under the point group, $G\circ\texttt{clSgn}\,\tau\,s=G$;
\item $G\le C(\delta)$ on $\{|x|\ge\delta\}$;
\item $G$ obeys the short-distance law of \texttt{thm\_short\_eps\_logX} for $|x|<\delta_0$:
$|G(x)-(\beta_1\lambda/\beta_R)^{2/3}|\le K'(\beta_1\lambda/\beta_R)^{2/3}\log\lambda/\lambda$ with $\lambda=\log(1/|x|)$.
\end{itemize}
Unconditional ingredients:
\begin{itemize}
\item uniform boundedness on $\{|x|\ge r\}$ for \emph{all} $N$ (\texttt{t6\_bdd}): \texttt{thm\_bound\_eps} for large $N$,
and finiteness of $\Lambda_N$ for the small $N$;
\item $|\lfloor L^Nx\rceil|/L^N\to|x|$ (\texttt{t6\_eucl\_tendsto});
\item rounding is compared with the point group through the embedding $z\mapsto z/L^N$ (\texttt{latPt\_latEmb},
\texttt{dist\_latEmb\_latPt\_le}, \texttt{t6\_latEmb\_sgnPerm});
\item one master construction for both \texttt{thm\_bound\_eps} and \texttt{thm\_short\_eps\_logX} (\texttt{t6\_short\_bound}).
\end{itemize}
Comparison with Conjecture~\ref{thm:limit}:
\begin{itemize}
\item the limit is only subsequential;
\item it may depend on the choice of $b_N$ and $m_N$;
\item the $z_*$ part is proved separately (\texttt{thm\_theta\_eps});
\item convergence of the full sequence would need Conjecture~\ref{thm:uv}.
\end{itemize}
\emph{Open: the equicontinuity hypothesis \texttt{ClEqui}.} Besides the wiring of the restart sizes into
\texttt{rpc12\_T5S}/\texttt{rptf\_top\_traj}, one exact identity is missing: linearity of the
\emph{pair} ($\sigma_a\sigma_b$) output in the $a$-datum at a fixed $b$-datum. Only the $\sigma_a$ parts are covered
(\texttt{lin\_obsHistFrom\_sub}, \texttt{rr\_restart\_identity$'$}). A route without a third source is as follows.
\begin{itemize}
\item The triples $(F,G,H)$ with equal $\emptyset$ and $b$ parts and $F_a-G_a=H_a$, $F_{ab}-G_{ab}=H_{ab}$ form a
subalgebra of $(\mathrm{OE})^3$. It is closed under $\exp$ and under componentwise linear maps.
\item The truncated output $s_{ab}/s_\emptyset-s_as_b/s_\emptyset^2$ is additive on such triples.
\item Hence only the scale-$m$ integrand $Z_m$ (\texttt{ZinO}) and the chain invariance
(\texttt{hist\_omegaOF\_from}, \texttt{rr\_ZinO}) are needed, not a commutation lemma for every map.
\end{itemize}

\paragraph{Lean status (B3, T6: pair linearity in the $a$-datum).}
(\texttt{H02UV/Obs2/Pair/PairLin1}, \texttt{H02UV/PtCorr2/Pair/PairLin1}, \texttt{H02UV/Spin2/Pair/PairLin1}.) The route above is carried out.
\texttt{PlRel}~$F\,G\,H$ (equal $\emptyset$, $b$ parts; $F_a-G_a=H_a$, $F_{ab}-G_{ab}=H_{ab}$, with $H$ carrying the
common $\emptyset$, $b$ parts) is a subalgebra of $(\mathrm{OE})^3$ (\texttt{plRelSub}), closed under $\exp$
(\texttt{pl\_rel\_oexp}: nilpotency is decided by the common bulk part) and $\omega_\Gamma$, hence under
$I^{\Lambda\setminus X}K^X$ (\texttt{pl\_rel\_gZO}, \texttt{pl\_rel\_ZinO}). The connected output
$n_{ab}-n_an_b$ is additive on triples (\texttt{pl\_outW\_sub}). The chain-invariant element is
$W_k=e^{-u_k}\omega_{\Gamma_k}Z_k$: $W_m=(\prod_{m\le i<k}\mathrm{ch}_i)\,W_k$ (\texttt{pl\_W\_chain}, from
\texttt{hist\_omegaOF\_from}), and $W_{N_h}=\mathrm{ch}_{N_h}\,e^{-u_{N_h+1}}\mathrm{lastOF}$ (\texttt{pl\_W\_top}).
With $\delta=$\,\texttt{plDiffSt}$(\mathrm{rrSt}\,o_1,o_2)$ ($a$, $ab$ parts the difference, $\emptyset$/$b$ parts
those of $o_2$; \texttt{pl\_diff\_prj} is the hypothesis of \texttt{rpb\_b\_eq}), the triple at scale $m$
(\texttt{rr\_ZinO}) propagates to every $m\le k\le N_h$ and to the top (\texttt{pl\_W\_rel}). For the
\texttt{rpcOut}/\texttt{rpTopOut} form $-\epsilon(u_a)\gamma-n_{ab}(Y)$ the identity is \emph{exact up to two explicit
defects} (\texttt{pl\_pair\_diff}, \texttt{pl\_pair\_diff\_top}; \texttt{pl\_rpc\_diff}, \texttt{pl\_rpTop\_diff} at
$P'=(0,y)$ with the reference $\gamma=u_B$, $s=s_\emptyset(W_2)\neq0$):
\[
\mathrm{out}(P,o_1)-\mathrm{out}(P',o_2)-\mathrm{out}(P',\delta)=-D_a\,n_b(W_2)-D_{ab},\qquad
D_\bullet=\epsilon(u^{(1)}_{j,\bullet})-\epsilon(u^{(2)}_{j,\bullet})-\epsilon(u^{(\delta)}_{j,\bullet}).
\]
$D_a$ does not vanish in general ($u_{j,a}$ depends on the $a$-point), but it is multiplied by the normalised
one-point part $n_b(W_2)$ of the reference, which is chain invariant and should vanish at $h=0$; $D_{ab}=0$ below coalescence when $u_{ab}=0$
along the three histories (\texttt{rpb\_coal\_validL}). \emph{Open:} $n_b(W)=0$ and $s\ne0$ for the reference at $h=0$;
$u_{ab}=0$ for $\delta$ (\texttt{ObsCoalInv}); \texttt{ObsLocInv} of $\delta$; and the identification of $P(x)$ with the
output of the translated pair $(y-x,y)$ (a general-pair version of \texttt{rpcOut\_ref}).

\paragraph{Lean status (B3, T6: $P(x)-P(y)$ as a restarted output).}
(\texttt{H02UV/PtCorr2/Pair/PairLin2}, \texttt{H02UV/Spin2/Pair/PairLin2}.) Three of the four open items above are closed and the
fourth is reduced to one hypothesis. \emph{(1)} For every pair $Q$ at $h=0$ and $k\le N_h$ (and at the top), the
chain-invariant element of the pair observable history has $s_\emptyset(W_k)\ne0$, $n_a(W_k)=n_b(W_k)=0$ and
$n_{ab}(W_k)=\langle\tilde z_{Q_a}\tilde z_{Q_b}\rangle$ (\texttt{pl2\_W\_facts}, \texttt{pl2\_WTop\_facts}: chain to
scale $0$, \texttt{zz\_rg\_ratio}, $r=0$). \emph{(4)} By translation invariance (\texttt{avg\_transl}),
$P(x)=-\langle\tilde z_{y-x}\tilde z_y\rangle$ (\texttt{pl2\_conn\_eq}); so the pair $(y-x,y)$ is used directly and no
general-pair \texttt{rpcOut\_ref} is needed. \emph{(3)} $\delta=\texttt{plDiffSt}(\mathrm{rrSt}\,o_1,o_2)$ with
$\beta_B=m_2(V(B)_a)$ is local at $(0,y)$ when $y-x$ and $0$ share the $m_s$-block (\texttt{pl2\_delta\_inv}; the
couplings of $z_a$ carry a constant, so the pure-form hypothesis of \texttt{rrSt\_obsLocInv} is replaced by
$\beta_B=0$ off the neighbourhood of $a$). Working with $W$ instead of the defect form gives, at $h=0$,
\emph{exactly} (\texttt{pl2\_conn\_diff}, $m_s\le j_y\le N_h$; \texttt{pl2\_conn\_diff\_top}, $m_s\le N_h$)
\[
P(x)-P(y)=\mathrm{rpcOut}_y(\delta\text{ restarted from }m_s\text{ at }j_y)+\epsilon(u^{\delta}_{j_y,ab}),
\]
and the same with \texttt{rpTopOut} at $N_h+1$. \emph{(2)} The $ab$ constant vanishes if
\texttt{ObsCoalInv}$\,(0,y)\,m_s\,\delta$ holds (\texttt{pl2\_conn\_diff\_coal}, \texttt{pl2\_conn\_diff\_top\_coal} for
$N_h+1\le j_y$). This hypothesis is still open (not proved in Lean): the $V_{ab}$, $W_{ab}$, $u_{ab}$ clauses are differences
of those of the two states, but the far-support clauses $V_{\rm no}$, $K_{\rm no}$ at radius $L^{m_s}$ around $0$
do not follow from those around $y-x$ (the relocated $a$-parts sit in the $m_s$-neighbourhood of the common
block), and the correction $\delta(B)(K^\flat+K)^{Y\setminus B}$ of $\mathrm{rrSt}$ can carry $ab$ parts.

\paragraph{Lean status (B3, T6 (a): restart sizes wired into the T5 bound).}
(\texttt{H02UV/PtCorr2/Cont/ContWire1}, \texttt{H02UV/Spin2/Cont/ContWire1--7}.) \emph{Exported envelopes:} the restart recursion with
large transients also bounds the share deviation, so along the restarted history $\sum_B|m_2V_{m_r+j}(B)_a|\le
|\beta|+D\le2M_j$ and $s_{m_r}^2M_0\le(7/5)^js_{m_r+j}^2M_j$, $M_j=M_0\prod_{i<j}|1-\mathfrak n^F_{m_r+i}|>0$
(\texttt{cw2\_obs\_hist\_restartBT}); at the master level, with the second conjunct of \texttt{obs\_hist\_restartBT}
($\mathcal W$, $K$ envelopes), \texttt{MasterCore.cw2\_restartDT}, \texttt{cw2\_restartDT\_lin} (no new master field).
In its own envelope $M_a=2s_j^2M_{j-m_r}$ the $a$-sector of the restarted state is $O(1)$: $V_a\le M_a$,
$\mathcal W_a\le(c_W c\check g_j+w/(s_{m_r}^2M_0))M_a$, $A^{(\cdot)}K_a\le((k+120B(c\epsilon)w)/(s_{m_r}^2M_0)+8c_WB(c\check g_j)^2)M_a$
(\texttt{cw2\_aSizes}). With the reference $b$-part the $b$-sector is the mirrored one-point history
(\texttt{cw2\_bSizes}, via \texttt{cw2\_relBaseD} = \texttt{relBase} along $D_\epsilon$), and
$2s_j^2|\beta^S_j|\le e^{2c_\delta(4/\beta_L)\epsilon}(81/56)^{N_h-j}M_b$ (\texttt{cw2\_pMS\_ratio}).
\texttt{MasterFacts.cw2\_T5M}: \texttt{rpc12\_T5S} with $M_b=pM(\mathrm{ptsS}\,x)_{N_h}$ and $h_{MaM}$, $h_{MbM}$,
$M_b>0$ discharged. \texttt{MasterFacts.cw2\_T5R}: for a datum $\delta$ at $m_r\le j_x$ (observable, local, symmetric,
\texttt{ObsCoalInv}, pure $\bar\psi_a\psi_a$ form, reference $b$-part) with the data of \texttt{restartDT\_lin} at $m_r$,
and $\sigma=s_{j_x}^2M_{j_x-m_r}/s_{N_h}^2$,
\[
|\mathrm{out}(o_{j_x})|\le\tfrac32(1+K_M)\gamma_\delta\vartheta_{j_x}+K_T(L,C_\delta)\frac{\sigma}{|\beta_{a,N_h}|}
\vartheta_{N_h}^2\check g_{N_h},
\]
the one-point sizes $v_T\ge1+\Theta_b$, $w_T$, $k_{xT}$ given explicitly ($\Theta_b=e^{2c_\delta(4/\beta_L)\epsilon}(81/56)^{N_h-j_x}$).
\emph{Open inputs:} the $\sigma_a\sigma_b$ parts $\mathcal W(q)_{ab}$, $K(X)_{ab}$ at $j_x$ (\texttt{rpab\_hist} needs the
datum's own $K_{ab}$ at $m_r$, i.e.\ the pair-linear datum); $\gamma_\delta\ge|\epsilon(u_a(o_{j_x}))|$ (needs
$\epsilon(\delta.u_a)$, cf.\ $D_a$ above); the restart data of $\delta$ at $m_r$ (\texttt{cd3\_restartM} covers
\texttt{cd3Delta}, not the datum with reference $b$-part, and is stated for admissible trajectories); the top case
(\texttt{rptf\_top\_traj}: $E_{N_h}$ of the restarted flow); the scalar smallness ($A^4(c_x\check g^3+k_{xT}+k_2)\le1$:
$\check g_{j_x}$ small depending on $C_\delta$; the window $m L^{2N}\le c(C_\delta)$).

\paragraph{Lean status (B3, T6: relaxed coalescence invariant; restart data of $\delta$ along $D_\epsilon$).}
(\texttt{H02UV/Obs2/Pair/CoalR}, \texttt{H02UV/Obs2/Pair/CoalHistR}, \texttt{H02UV/PtCorr2/Restart/CoalRestartR}, \texttt{H02UV/Spin2/Pair/PairLin3},
\texttt{H02UV/Spin2/Cont/ContDiffD1--3}, \texttt{H02UV/Spin2/Cont/ContWire8}, \texttt{H02UV/Spin2/Pair/PairLin4}.) \emph{The extra term vanishes.} The
coalescence invariant is relaxed to the support radius $L^j+\rho$ ($\rho\le L^j$; \texttt{ObsCoalInvR}); the margins of
the step absorb $\rho$ (a small set below coalescence is at distance $\ge L^k+\rho$ from $a$ or from $b$ if
$2(4+2^d)\le L$, \texttt{small\_farR}), so it propagates exactly as before (\texttt{obsStep\_coalInvR},
\texttt{crr\_coal\_validL}). The relocated state $\mathrm{rrSt}\,o_1$ satisfies it around $(0,y)$ with
$\rho=|y-x|_\infty<L^{m_s}$ (\texttt{crr\_rrSt\_coalInvR}: the correction $\delta(B)(K^\flat+K)^{Y\setminus B}$ vanishes if
$B$ is far from $y-x$ and has no $\sigma_b$ part if $Y$ is far from $y$), and the invariant is closed under
\texttt{plDiffSt}. Hence for $m_s<j_x$, $m_s<j_y$: $V^\delta_{ab}=u^\delta_{ab}=0$ up to $\min(j_y,N_h+1)$ and, at $h=0$,
\emph{exactly} $P(x)-P(y)=\mathrm{rpcOut}_y(\delta_{j_y})$ ($j_y\le N_h$, \texttt{pl3\_conn\_diff}) and
$=\mathrm{rpTopOut}_y(\delta_{N_h+1})$ ($N_h+1\le j_y$, \texttt{pl3\_conn\_diff\_top}); no \texttt{ObsCoalInv} hypothesis.
\emph{Restart data along $D_\epsilon$.} From scale $1$ on the $\sigma_a$ parts of $V,\mathcal W,K$ of $\delta$ are those of the
\texttt{cd3} datum for $\tilde z^0$ ($\gamma=0$; \texttt{cdd\_delta\_a}); with the $D_\epsilon$ copies of the bulk inputs
and of the share deviation (\texttt{cdd\_iotaD}, \texttt{cdd\_regD}, \texttt{cdd\_relDevD}) and the grid artefact from
(H$_\lambda$) for the pair vertex ($|\beta^{(y-x,y)}_{m_s}-\beta^{(0,y)}_{m_s}|\le4K_p(\check g+a)|\beta^{(0,y)}_{m_s}|$,
\texttt{cdd\_art\_rel}), \texttt{MasterFacts.cdd\_restart\_data} gives the inputs of the restart at $m_s$
($|\beta|+D\le M_0$, $\mathcal W$, $K$, $1344(8c_Nk+\Gamma_ww)\le s^2M_0$) with
$M_0\le\bigl(4K_p(\check g+a)+24c_\delta c^2\check g/(\beta_L/4)+1344(8c_N(64c_WB(c\check g)^2+48e^{3/2}A L^{-r})
+8\Gamma_wc_Wc\check g)+\check g\bigr)|\beta^{(0,y)}_{m_s}|$ (\texttt{cdd\_restart\_rel}), for $0<m_s<N_h$,
$|y-x|_\infty\le L^{m_s-r}$, $K_p(\check g_{m_s}+a_{m_s})\le1/2$, $A^4c_x\check g_{m_s}^3\le1$.
\texttt{MasterFacts.pl4\_T6\_inputs} collects, for $\delta$, the exact identity and every datum input of
\texttt{cw3\_T5R} (\texttt{cw2\_T5R} with the $\tilde z$ reference $\sigma_b$ part and $V_{ab}=0$ at $j_x$ in place of
\texttt{ObsCoalInv}). \emph{Still open:} the $\sigma_a\sigma_b$ sizes $w_2,k_2$ at $j_y$; $\gamma_\delta$ (the initial
$\epsilon(\delta.u_a)$ and the chain); the comparison $\sigma/|\beta_{a,N_h}|\lesssim M_0/|\beta_{m_s}|$ along $D_\epsilon$; the
top case; the final assembly.

\paragraph{Lean status (B3, T6: intermediate coalescence $m_s<j_y\le N_h$ --- comparison, $\gamma$ chain, reduction).}
(\texttt{H02UV/Spin2/Cont/ContWire9}, \texttt{H02UV/Spin2/Pair/PairLin5}, \texttt{H02UV/Spin2/Pair/PairLin5b}.) \emph{Comparison along $D_\epsilon$.}
With the $D_\epsilon$ copy of the product comparison (\texttt{cw4\_prod\_cmpD}), \texttt{cw2\_ratioD} and
$|\beta_j|\le e^{c}(8/7)^{N_h-j}|\beta_{N_h}|$ (\texttt{cw4\_beta\_up}), for every observable datum
$\sigma/|\beta_{a,N_h}|\le C_{\rm cmp}(N_h-j_y)\,M_0/|\beta^{(0,y)}_{m_s}|$,
$C_{\rm cmp}(n)=(15/14)^{2n}e^{(8/7)(8/3)(2/L)}e^{2c}(8/7)^n$ (\texttt{cw4\_sig\_cmp}); for $\delta$ and $M_0=$\texttt{cddM},
$\sigma/|\beta_{a,N_h}|\le C_{\rm cmp}(N_h-j_y)\,(\check g_{m_s}+a_{m_s}+48e^{3/2}AL^{-r})$ up to the constants of
\texttt{cddRel} (\texttt{cw4\_sig\_rel}). \emph{The $\gamma$ chain.} For a pure datum $\sum_B\epsilon(V(B)_a)=0$ at $m_s$;
with $|\beta^F_k|\le\mathrm{shA2}\le2E_{k-m_s}$, $|\mathrm{gamWKF}_k|$ from \texttt{rpg\_gamWKF\_le} and the envelopes of
\texttt{cw2\_restartDT\_lin} (\texttt{pl5\_gamWKF}), and the bulk factor $|c_k|\le(1+c_e(\check g_k+a_k))|\Gamma_{k+1}(0,0)|$:
$|\epsilon(u_a(\delta_{j_y}))|\le|\epsilon(\delta.u_a)|+\sum_{i<j_y-m_s}\bigl(2E_i(1+c_e(\check g+a))|\Gamma_{m_s+i+1}(0,0)|
+G_i\bigr)$ (\texttt{pl5\_gam\_jy}). \emph{Reduction} (\texttt{MasterFacts.pl5\_T6\_mid}, exact, $h=0$, $D_\epsilon$): there
are transients $k,w$ with $1344(8c_Nk+\Gamma_ww)\le s_{m_s}^2M_0$ such that, under the scalar conditions of
\texttt{cw3\_T5R}, the inputs $w_2,k_2$ (the $\sigma_a\sigma_b$ parts of $\mathcal W,K$ at $j_y$) and
$|\epsilon(\delta.u_a)|\le\gamma_0$,
$|P(x)-P(y)|\le(\gamma_0+\Gamma_{\rm chain})\tfrac32(1+K_M)\vartheta_{j_y}
+\max(K,0)\,C_{\rm cmp}(N_h-j_y)\,\mathrm{cddRel}\,\vartheta_{N_h}^2\check g_{N_h}$.
\emph{Open, and why.} (i) $w_2,k_2$: the restart transients of $\delta$ at $m_s$ are $\|\delta\mathcal W\|\simeq c_W c\check g\,
2s^2(|\beta^P|+|\beta^{P'}|)$ and $\|\delta K\|\simeq\check g^2s^2|\beta|+L^{-r}s^2|\beta|$ (triangle inequality), i.e.\
$O(1)\cdot s^2M_0$, a factor $\check g^{-1}$ (resp.\ $\check g^{-2}$) above the stationary envelopes $c_Wc\check g\,2s^2M$,
$c\check g^2s^2M$; the full-norm steps used for the $ab$ parts (\texttt{WObsBoundsU}, \texttt{rpab\_hist}) require
$\|\mathcal W\|_w\le r\,m$, $k\le c_1r^2m$ with $r=c\check g$, so they give $w_2,k_2$ of order $\check g^{-1}$ unless the
transients have decayed ($6^{-n}$, $(8/15)^n$, i.e.\ $j_y-m_s\gtrsim\log(1/\check g_{m_s})$, which is $N$-dependent for fixed
physical $x,y$). Needed: a $\mathcal W$/$K$ artefact bound for the relocated difference at $m_s$, or an $ab$ step linear in
the one-point transients. In addition the large-set $K^\delta_{ab}$ at $m_s$ is $K^{o_1}_{ab}-K^{o_2}_{ab}$ plus the
\texttt{rrSt} correction, not only the correction. (ii) $\epsilon(\delta.u_a)=\epsilon(u^{(y-x,y)}_{m_s,a})-
\epsilon(u^{(0,y)}_{m_s,a})$ (\texttt{rrSt} keeps $u$): a sum over all scales $k<m_s$ of the differences of the per-scale
$\gamma$ increments, relatively $O(\check g_k)$ by \texttt{cdd\_art\_rel} for the $\beta c$ part, not yet bounded.

\paragraph{Lean status (B3, T6: top coalescence $N_h+1\le j_y$ --- $E_{N_h}$, sizes at $N_h$, $\gamma_\delta$, assembly).}
(\texttt{H02UV/PtCorr2/Cont/ContTop1--3}, \texttt{H02UV/Spin2/Cont/ContTop1--8}.) \emph{The error of the top step.} For the restart with large
transients, $E_n\le(c_\delta r_n^2a_n+2c_t(3/4)^n)M_n$ from \texttt{close} at horizon $n$ and the inputs at the index $n$
only (\texttt{RelRecT.ct\_E\_top}, analogue of \texttt{RelRec.E\_top}); along the restarted history at $n=N_h$ (product
parts and initial tadpoles with horizon $N_h+1$, bulk inputs at $N_h$ from \texttt{C2Ctx.tadInc2\_leU} and
$|\mathfrak n^{(2)}_{N_h}|\le1/8$ of \texttt{TopConclD}): $|\beta_{m_s+j}|+D_{m_s+j}\le2M_j$,
$E_{N_h}\le(c_\delta(c\check g_{N_h})^2a_{N_h}+\tfrac1{24}(3/4)^{N_h-m_s})M_{N_h-m_s}$, $|\mathfrak n^F_{N_h}|\le1/7$
(\texttt{MasterFacts.ct\_restartTop}). \emph{The top step.} \texttt{rptf\_top\_traj} with $V_{ab}=0$ at $N_h+1$ in place of
\texttt{ObsCoalInv} (\texttt{ct\_top\_trajR}); the restart sizes at $N_h$ from \texttt{cw2\_aSizes}, \texttt{cw3\_bSizes},
\texttt{cw2\_pMS\_ratio}, the share data $a_N=8/7$, $b_D=2\sigma$, $e=(c_\delta r^2\cdot8/7+1/24)\sigma$, so that $v_T$ is
independent of the datum (\texttt{ct\_T5top}, $\sigma=M_{N_h-m_s}$). \emph{$\gamma_\delta$.} \texttt{rptf\_gam\_sizes} with
$B_k=2M_{k-m_s}$ and the envelopes of \texttt{cw2\_restartDT\_lin} (\texttt{ct\_gam\_top}); $\sum_B\epsilon(\delta V(B)_a)=0$,
the initial $\|\epsilon(\delta.u_a)\|$ stays explicit. \emph{Assembly} (\texttt{MasterFacts.ct\_conn\_diff\_top'}, exact, $h=0$,
every $b$ in the bracket, $0<m_s<N_h$, $m_s<j_x$, $N_h+1\le j_y$): $P(x)-P(y)=$ the restarted top output exactly
(\texttt{pl3\_conn\_diff\_top}); there are restart transients $k,w$ with $1344(8c_Nk+\Gamma_ww)\le s_{m_s}^2M_0$ such that, for
size constants $m,k_1,c_1$ with the inequalities of \texttt{ct\_T5top}, the input $k_2$ ($\sigma_a\sigma_b$ part of $K$ at
$N_h+1$) and the smallness conditions,
$|P(x)-P(y)|\le\tfrac32(1+K_O)\Gamma_\delta\vartheta_{N_h}+\tfrac94\max(K_E,0)\,C_{\rm cmp}(0)\,\mathrm{cddRel}\,
\vartheta_{N_h}^2\check g_{N_h}$ with $\mathrm{cddRel}\asymp\check g_{m_s}+a_{m_s}+48e^{3/2}AL^{-r}$ (\texttt{ct\_sig\_cmp}:
$\sigma/|\beta_{a,N_h}|\le C_{\rm cmp}(0)M_0/|\beta^{(0,y)}_{m_s}|$). \emph{Open.} $k_2$ (\texttt{ct\_k2\_top} reduces it to
the inputs of \texttt{rpab\_top}, including the datum's own $K^\delta_{m_s,ab}$; same obstruction as $w_2,k_2$ in the
intermediate case); $\|\epsilon(\delta.u_a)\|$; the scalar conditions (with the undecayed transient ratios of
\texttt{cw2\_aSizes}, $m\gtrsim\check g_{N_h}^{-1}$, so $w_T$ and the condition on $mL^{2N}$ are not $O(1)$; the transients do
decay as $6^{-i}$, $(8/15)^i$).

\paragraph{Lean status (B3, T6: the $\sigma_a\sigma_b$ inputs and the initial $a$-constant --- route; decayed transients).}
(\texttt{H02UV/Spin2/Cont/ContArt1--4}.) \emph{Math check.} The paper argument controls the difference datum at $m_s$ by the triangle
inequality and the decay of the transients, not by artefact bounds. The datum's own large-set $K^\delta_{m_s,ab}$ does not
vanish (\texttt{ObsCoalInvR} gives $K_{ab}=0$ on small sets only). Two routes need a new estimate and are rejected:
artefact bounds for $K^{o_1}-K^{o_2}$ (a Lipschitz-in-position estimate of the RG history below $m_s$), and a bilinear
$ab$ step (insufficient without decay). \emph{Route taken:} \texttt{rpab\_hist}/\texttt{rpab\_top} in \emph{one fixed}
weight with \emph{scale-dependent} sizes $m_j$ ($m_j\le2m_{j+1}$). The restart transients decay relative to the envelope as
$(7/30)^i$ ($\mathcal W$) and $(56/75)^i$ ($K$), so $m_j$ may start at $O(\check g^{-2}C^{j_y-m_s})$; then $K^\delta_{m_s,ab}$
(bounded by \texttt{kab\_hist} for both pair histories plus the \texttt{rrSt} correction) is admissible. Moreover
$k_2=O(\check g)$ and $w_2=O(1)$ once $j_y-m_s\ge d_0\approx7\log(1/\check g_{m_s})+C$ (resp.\ $N_h-m_s\ge d_0$). This
condition involves only $n=N-m_s$ (since $\check g_{m_s}\asymp1/(n\beta_1\log L)$), so it is uniform in $N$ and compatible
with the modulus. For the $a$-constant, at $h=0$ $n_a(W(\delta))=0$ gives exactly
$\epsilon(u^\delta_{N_h+1,a})=n_a(\mathrm{lastOF}(\delta_{N_h+1}))$ (no forward chain from $m_s$); the missing estimate is a
bound of $n_a(\mathrm{lastOF})$ for a general state, linear in its $a$-sizes (one-point analogue of the
\texttt{rpTopE} bound). \emph{Decayed exports (proved):} \texttt{MasterCore.art\_aSizesD} (\texttt{cw2\_aSizes} with
$(7/30)^i w/(s^2M_0)$ and $(56/75)^i(k+120B c\epsilon w)/(s^2M_0)$); with it, \texttt{MasterFacts.art\_T5topD}
(\texttt{ct\_T5top}), \texttt{art\_T5RD} (\texttt{cw3\_T5R}), \texttt{art\_T6\_midD} (\texttt{pl5\_T6\_mid}); with the waiting
condition the size constants $m$, $w_T$, $k_{xT}$ are $O(1)$. \emph{Open:} the wiring of $m_j$, the base bound
$K^\delta_{m_s,ab}$, $w_2$ by one \texttt{WObsBoundsU} step, and the $n_a(\mathrm{lastOF})$ bound.

\paragraph{Lean status (B3, T6, route (c): $w_2$, $k_2$ discharged under waiting; \texttt{H02UV/Spin2/Cont/ContK2a--e}).}
One fixed weight $w=w_M(M_a,M_b)$, $M_a=2s_J^2M_{J-m_s}$ ($J=j_y$, resp.\ $J=N_h$), $M_b=p_M(0,y)_{N_h}$, and the
explicit per-scale size $m_j=(c_W+1)\bigl((7/5)^{J-j}+\Theta_b(j)\bigr)+(8/15)^{j-m_s}Q$ ($m_j\le2m_{j+1}$;
\texttt{k2Mj}). The backward envelope ratio $2s_j^2M_{j-m_s}\le(7/5)^{J-j}M_a$ uses $|\mathfrak n^F_j|\le1/7$
along the restarted history (\texttt{k2\_n7}, \texttt{k2\_envRatio}); with \texttt{art\_aSizesD}, \texttt{cw3\_bSizes},
\texttt{cw2\_pMS\_ratio} this gives $\|V_j\|_w\le m_j$, $\|\mathcal W_j\|_w\le r_jm_j$, $\|K_j\|\le8Br_j^2m_j$ for
$m_s\le j\le J$ (\texttt{k2\_sizes}), provided $(7/5)^{J-m_s}\tau_w\le r_{m_s}Q$, $(7/5)^{J-m_s}\tau_k\le8Br_{m_s}^2Q$.
Then \texttt{rpab\_hist} gives $K_{j_y,ab}\le c_{ab}r_{j_y}m_{j_y}^2$ and one $\mathcal W$ step
(\texttt{k2\_wfull\_step}) gives $\mathcal W_{j_y,ab}\le e^{L^2C_w(r+m_{j_y-1}^2)}-1$ (\texttt{k2\_mid}); at the top
\texttt{ct\_k2\_top} gives $k_2=(c_{ab}/5)r_{N_h}m_{N_h}^2$ (\texttt{k2\_top}). Under the waiting hypothesis
$(8/15)^{J-m_s}Q\le1$ (i.e.\ $J-m_s\gtrsim\log(C/\check g_{m_s}^2)/\log(75/56)$ plus the $Z_{ab}$ term),
$m_J\le\bar m=(c_W+1)(1+\Theta_b(J))+1$ (\texttt{k2\_wait}). Final variants: \texttt{k2\_T6\_mid}
(= \texttt{art\_T6\_midD} without $hW_2$, $hK_2$) and \texttt{k2\_conn\_diff\_top} (= \texttt{ct\_conn\_diff\_top'} with
decayed transients, $k_2$ discharged). \emph{Still explicit:} the base bound $Z_{ab}$ of $K^\delta_{m_s,ab}$ in $w$
(with $Z_{ab}\le c_{ab}r_{m_s}Q^2$) and $\varepsilon(\delta.u_a)$.

\paragraph{Lean status (B3, T6: the initial $a$-constant at the top via $n_a(W(\delta))=0$; \texttt{H02UV/Spin2/Cont/ContNa1--7}).}
At $h=0$ the generating elements $W_{\rm top}=e^{-u_{N_h+1}}\,{\rm last}_{N_h+1}$ of the three histories form a
relative triple, $n_a$ is additive on triples (\texttt{pl\_nA\_sub}) and vanishes on the two reference entries
(\texttt{pl2\_WTop\_facts}); hence, for every $y$ and $m_s\le N_h$,
$\varepsilon(u^\delta_{N_h+1,a})=n_a({\rm lastOF}_{N_h+1}(\delta_{N_h+1}))$, i.e.\
$\texttt{rpTopU}(\delta_{N_h+1})=-n_a({\rm lastOF})$ (\texttt{na\_eps\_ua\_top}, \texttt{na\_rpTopU\_eq}).
The one-point analogue of the \texttt{rpTopE} bound (\texttt{na\_lastA\_le}, any $x$): for an observable state $o$
at $N_h+1$ with sizes $c_V,c_W,c_K$ in $w_M(2s_{N_h}^2\sigma,M_b)$,
$|n_a({\rm lastOF}(o))|\le\bigl(c_{Lab}c_Sc_u+8(\tfrac12+\tfrac{81}{32}c_{Jab}c_S)K_4\bigr)(\sigma/|\beta_{a,N_h}|)|M_{N_h}|$,
$c_S=c_V+c_W+8c_K+1$; proof: $n_a({\rm lastOF})=-\varepsilon(u^{\rm obs}_N(\Lambda))_a+n_a(\omega_{cQ}F'_N)$
(\texttt{na\_nA\_lastOF\_split}), the constant of the last step $\sum_B\varepsilon(V'_B)_a-\sum_B\varepsilon(S(B))_a$
(\texttt{na\_eps\_uXO\_a\_le}), the $\sigma_a$ jets of $F'_N$ (\texttt{na\_finG\_jets\_a}, \texttt{LastJetBd} directly), and
the zero mode with the reference normaliser (\texttt{rp\_zero\_abs}); one extra window
$2\cdot\tfrac{81}{16}c_J(c_\rho+4/r_0^2)c_S\ell^2m^2L^{2N}\le1$. Wired into the top case
(\texttt{na\_top\_traj}, \texttt{na\_top\_trajR}, \texttt{na\_T5top}, \texttt{na\_conn\_diff\_top}): the term
$\Gamma_\delta=$\texttt{ctGamB} (initial $\|\varepsilon(\delta.u_a)\|$ plus the chain from $m_s$) of
\texttt{ct\_conn\_diff\_top'} is replaced, and \texttt{na\_conn\_diff\_top'} reads
$|P(x)-P(y)|\le\tfrac94\bigl((1+K_O)\max(K_A^{na},0)+\max(K_E,0)\check g_{N_h}\bigr)\,\mathrm{cw4Cmp}(0)\,
\mathrm{cddRel}\,\vartheta_{N_h}^2$ (both terms relative, $\mathrm{cddRel}\asymp\check g_{m_s}+a_{m_s}+L^{-r}$).
\emph{Not done:} the intermediate case $j_y\le N_h$ ($hu_0$ of \texttt{pl5\_T6\_mid}/\texttt{art\_T6\_midD}); it needs
the $a$-only sizes of $\delta_{N_h+1}$ (weight $\langle M_a^{-1},0,0\rangle$, steps \texttt{rptf\_wfull\_stepF},
\texttt{rptf\_kobs\_stepF} from \texttt{art\_aSizesD} at $N_h$), the reference zero mode for $j_y\le N_h$ (from
$n_a(W_{\rm top}({\rm ref}))=0$ and \texttt{na\_nA\_lastOF\_split}), and the backward chain $j_y\to N_h+1$
(\texttt{pl5\_gamWKF} plus the top step of \texttt{rptf\_gam\_sizes}).

\paragraph{Lean status (B3, T6: the initial $a$-constant in both cases without $\varepsilon(\delta.u_a)$; \texttt{H02UV/Spin2/Cont/ContNb1--9}).}
\emph{Top case, merged} (\texttt{nb\_T5topD}, \texttt{nb\_conn\_diff\_top}): the decayed size conditions of
\texttt{art\_T5topD} ($(7/30)^{N_h-m_s}$, $(56/75)^{N_h-m_s}$), the $k_2$ data of \texttt{k2\_conn\_diff\_top}
($Q$, $Z_{ab}$, waiting $(8/15)^{N_h-m_s}Q\le1$) and the identity $\texttt{rpTopU}(\delta_{N_h+1})=-n_a({\rm lastOF})$
together: $|P(x)-P(y)|\le\tfrac94\bigl((1+K_O)\max(K_A^{na},0)+\max(K_E,0)\check g_{N_h}\bigr)\mathrm{cw4Cmp}(0)\,
\mathrm{cddRel}\,\vartheta_{N_h}^2$, no $\varepsilon(\delta.u_a)$. \emph{Intermediate case} $m_s<j_y\le N_h$:
(1) \texttt{nb\_lastA\_le} is \texttt{na\_lastA\_le} in general weights with $w_V.a=w_K.a=M_a^{-1}$, used with the
$a$-only weight $\langle M_a^{-1},0,0\rangle$ (\texttt{nbWa}); (2) the $a$-sizes of $\delta_{N_h+1}$ from
\texttt{art\_aSizesD} at $N_h$, \texttt{rptf\_vobs\_block}/\texttt{rptf\_beta\_shA2}, \texttt{rptf\_wfull\_stepF},
\texttt{rptf\_kobs\_stepF} (\texttt{nb\_naStep}, \texttt{nb\_naMid}); (3) the reference zero mode for every $y$,
$s_\emptyset\Omega\neq0$, $n_a\Omega=-\gamma_{a,N}$, from $n_a(W_{\rm top}({\rm ref}))=0$, $W_{\rm top}=e^{-u_{N_h+1}}X\Omega$
and $\gamma_{0,N_h+1}+n_aX=\gamma_{a,N}$ (\texttt{nb\_ref\_zm}); (4) the backward chain
$|\gamma^F_{j_y}|\le|\gamma^F_{N_h+1}|+\sum_{j_y\le k\le N_h}(|\beta^F_k||c_k|+|\mathrm{gamWKF}_k|)$ (\texttt{nb\_gam\_back},
\texttt{nb\_gam\_back\_top}, per-step bounds of \texttt{ct\_gam\_top}). Hence
$|\varepsilon(u^\delta_{j_y,a})|\le K_A^{na}(\sigma/|\beta_{a,N_h}|)\tfrac32\vartheta_{N_h}+\mathrm{nbGB}$
(\texttt{nb\_eps\_jy}), and \texttt{nb\_T6\_mid} is \texttt{k2\_T6\_mid} without $hu_0$ (the factor
$\gamma_0+\texttt{pl5Gam}$ replaced by this bound; new inputs: $a$-size constants $mm,k_1,c_1$ at $N_h$ with decayed
transients, the window of \texttt{na\_lastA\_le}). $Z_{ab}$ ($hZ$, $hQ_z$) remains a hypothesis in both cases.

\paragraph{Lean status (B3, T6, route (c): the base bound $Z_{ab}$ discharged; \texttt{H02UV/PtCorr2/Cont/ContZab1}, \texttt{H02UV/Spin2/Cont/ContZab1--5}).}
The $ab$ activity of $\delta$ at $m_s$ is $\delta.K(X)_{ab}=K^{o_1}(X)_{ab}+\mathrm{corr}(X)_{ab}-K^{o_2}(X)_{ab}$
(\texttt{zab\_delta\_ab\_eq}; $o_1$, $o_2$ the pair histories of $(y-x,y)$ and $(0,y)$).
(a) The pair-history parts along $D_\epsilon$ (no admissibility): \texttt{kab\_hist} up to any $1\le n<j_{ab}(P)$, $n\le N_h$,
for any pair $P$ in its natural weights $w_M(p_M(P)_n,p_M(P.b,0)_n)$ (\texttt{zab\_traj}, \texttt{MasterFacts.zab\_sizes};
one-point sizes from \texttt{cw2\_relBaseD}, $\beta$ ratio from \texttt{bt2\_ratio}, coalescence invariant from
\texttt{rpb\_coalInv}): $A^{(|\mathcal B|-4)_+}w_{ab}\|K_n(X)_{ab}\|\le c_{ab}r_n(2c_W)^2$. (b) The correction: since
\texttt{rrD} is pure $\sigma_a$, $\mathrm{corr}_{ab}=-\sum_B\beta_B(\bar\psi_a\psi_a-\bar\psi_0\psi_0)\,(I'(B)(K^\flat+K^{o_1})^{X\setminus B})_b$
(\texttt{zab\_rrD\_mul\_ab}); in the $\sigma_b$ weight $w_B=\langle0,M_b^{-1},0\rangle$, $\|I'(B)\|_{w_B}\le
\iota_Z=e^{7/2}(2+2c_W)$ (\texttt{oTj\_inI\_le} for \texttt{rrSt}\,$o_1$), the large-set regulator in $w_B$ from
\texttt{kabCond}, and $\sum_B\|\beta_B(\cdots)\|\le p_M(P_1)_{m_s}\,4|y-x|_\infty L^{-m_s}$ give
$A^{(|\mathcal B|-4)_+}\|\mathrm{corr}(X)_{ab}\|\le M_b\,p_M(P_1)_{m_s}\,4|y-x|_\infty L^{-m_s}\iota_ZA$
(\texttt{zab\_corr\_ab\_le}). With the conversion $w_{ab}(M_a,M_b)=(M_a'M_b'/M_aM_b)\,w_{ab}(M_a',M_b')$ and the ratios
$p_M(P)_{m_s}/M_a\le(7/5)^{J-m_s}|\beta^P_{m_s}|/M_0$ (\texttt{k2\_envRatio}),
$p_M(0,y)_{m_s}/p_M(0,y)_{N_h}\le\Theta_b(m_s)$ (\texttt{cw2\_pMS\_ratio}) this is
$Z_{ab}=(7/5)^{J-m_s}\Theta_b(m_s)M_0^{-1}\bigl(|\beta^{(y-x,y)}_{m_s}|(c_{ab}r_{m_s}(2c_W)^2+4|y-x|_\infty L^{-m_s}\iota_ZA)
+|\beta^{(0,y)}_{m_s}|c_{ab}r_{m_s}(2c_W)^2\bigr)$ (\texttt{zabZr}, \texttt{MasterFacts.zab\_Zfix}). Final variants with
$Q=Q_0+\sqrt{Z_{ab}/(c_{ab}r_{m_s})}$ explicit (\texttt{zabQ}; $Q_0$ only for the transient conditions):
\texttt{k2\_T6\_midZ}, \texttt{k2\_conn\_diff\_topZ} (new input $c_{ab}>0$; $hu_0$ unchanged). Since
$M_0\asymp(\check g_{m_s}+a_{m_s}+L^{-r})|\beta_{m_s}|$ and $\Theta_b(m_s)=(81/56)^{J-m_s}\Theta_b(J)$,
$Z_{ab}/(c_{ab}r_{m_s})\lesssim(567/280)^{J-m_s}\Theta_b(J)(1+L^{-r}/\check g_{m_s})/(\check g_{m_s}+L^{-r})$; as
$(8/15)^2\cdot567/280<0.58$, the waiting condition $(8/15)^{J-m_s}Q\le1$ still only asks
$J-m_s\gtrsim\log\bigl(C\,\Theta_b(J)/\check g_{m_s}^{2}\bigr)$ ($\Theta_b(J)=O(1)$ for fixed physical points; $N$-independent).

\paragraph{Lean status (B3, T6, final assembly: merged bounds; equicontinuity \emph{not} proved; \texttt{H02UV/Spin2/Cont/ContFin1--2}).}
The two coalescence cases are now available with both $Z_{ab}$ and the initial $a$-constant discharged:
\texttt{MasterFacts.fin\_T6\_mid} ($m_s<j_y\le N_h$; \texttt{nb\_T6\_mid} with $Q=\texttt{zabQ}\dots Q_0$ and $hZ$ from
\texttt{zab\_Zfix}) and \texttt{MasterFacts.fin\_conn\_diff\_top} ($N_h+1\le j_y$; \texttt{nb\_conn\_diff\_top} likewise).
Their remaining inputs are scalar, plus $c_{ab}>0$. The latter is not derivable from the master constants:
$c_{ab}=10\,k_{B2}(\theta_0(C_{0u}),C_1,c_x,8B)$, and the existentials of \texttt{kobs\_bounds\_U} and \texttt{kabCond}
only give $C_{0u},C_1\ge0$. The cheapest fix is to strengthen \texttt{kobs\_bounds\_U} to $C_0>0$ and use it in the master
construction (this edits the existing library).
\emph{Obstruction to \texttt{ClEqui} for fixed $\beta_R$.} \texttt{thm\_limit\_of\_equi} fixes $\beta_R$ with
$1/\beta_R\le g_0$ ($g_0$ depending only on $L$) and admissible masses with $4m_NL^{2N}\le1/\beta_R$. \texttt{ClEqui} then asks
for every radius $r>0$. The scalar conditions of the intermediate case need smallness that depends on $r$, through
$C_\delta\approx\log_L(1/r)$ and $N_h-j_y\le C_\delta$:
(i) $A^4(c_x\check g^3+k_{xT}+k_2)\le1$ with $k_2\ge c_{ab}\,c\,\check g_{j_y}\bigl((c_W+1)(1+\Theta_b(j_y))+1\bigr)^2$,
where $\Theta_b(j)=e^{c}(81/56)^{N_h-j}$; so $\check g_{j_y}\lesssim r^{2\log_L(81/56)}$, while
$\check g_{j_y}\asymp1/(\beta_R+\beta_1\log(1/r))$.
(ii) \texttt{hmAB}: $m L^{2N}\cdot\max(1,2\Lambda_1)^{C_\delta}m_0\lesssim1$ with
$m_0\gtrsim(2\iota(v_T,w_T)A)^{(L+18)^2}$, $v_T\ge1+\Theta_b(j_y)$; so $mL^{2N}\lesssim r^{c(L)}$.
For fixed $\beta_R$ and $r\to0$ neither condition follows from the hypotheses. The bounds therefore give
equicontinuity only on $\{\|x\|\ge r\}$ under $1/\beta_R\le g(r)$ and $m_NL^{2N}\le\mu(r)$. This is a statement of a
different form, and it would yield limits only on $\{\|x\|>r(\beta_R)\}$, with $r(\beta_R)\to0$ as $\beta_R\to\infty$.
Both growths are removed in Theorem~\ref{thm:limit-fixed}, which proves \texttt{ClEqui} for fixed $\beta_R$.

\subsubsection*{What the bounds give: equicontinuity away from the origin with $r$-dependent couplings}
Notation for the next proposition. $\T=\R^2/\Z^2$ with the Euclidean torus norm $|x|$, and
$f_N(x)=\vartheta_\epsilon^{-2}P_\epsilon(\lfloor L^Nx\rceil)$ (\texttt{clF}), where
$\vartheta_\epsilon^2=(\beta_Rb_N)^{2/3}$ and $\beta_\epsilon=1/b_N$. A sequence $(R_N,m_N^2,b_N)_{N\ge N_1}$ is
\emph{admissible for $\beta_R$} (\texttt{ClAdm}) if $R_N$ are master parameters, $m_N^2>0$, $4m_N^2L^{2N}\le1/\beta_R$ and
$b_N>0$ lies in the bracket of Theorem~\ref{thm:multiscale}. Put
\[
\kappa_L=\log_L\tfrac{81}{56},\qquad C_\delta(r)=\lceil\log_L(1/r)\rceil+3,\qquad
\Theta(r)=\bigl(\tfrac{81}{56}\bigr)^{C_\delta(r)}\asymp r^{-\kappa_L}.
\]

\begin{proposition}[Equicontinuity on $\{|x|\ge r\}$ and limits there; short distances at fixed $\epsilon$]\label{prop:equi-r}
Let $d=2$, $i_0=0$ and $L\ge L_0$. There are $c_L,C_L>0$, depending on $L$ only, such that the following hold.

\smallskip\noindent\textup{(i) Equicontinuity.} For every $r\in(0,\frac14]$ put
\[
g(r)=c_L\,\Theta(r)^{-2},\qquad \mu(r)=c_L\exp\bigl(-C_L\,\Theta(r)^{2}\bigr).
\]
There are $C(r)<\infty$, $\eta_0(r)\in(0,1)$ and $N_0(r)$, depending on $r$ and $L$ only, with the following property.
Let $1/\beta_R\le g(r)$, let $(R_N,m_N^2,b_N)$ be admissible for $\beta_R$ with $m_N^2L^{2N}\le\mu(r)$ for
$N\ge N_1$. Then for all $N\ge\max(N_0(r),N_1)$ and all $x,y\in\T$ with $|x|,|y|\ge r$ and $|x-y|\le\eta\le\eta_0(r)$,
\[
|f_N(x)-f_N(y)|\le\omega_r\bigl(\max(\eta,L^{-N})\bigr),\qquad
\omega_r(\eta)=C(r)\Bigl(\frac1{\beta_R+\frac12\beta_1\log(1/\eta)}+\eta^{1/2}\Bigr),
\]
with $\beta_1=3/(2\pi)$. More precisely, $\omega_r(\eta)=C(r)\bigl(\check g_{m_s}+a_{m_s}+L^{-r'}\bigr)$, where
$a_j=L^{-2j}+m^2L^{2j+2}$, and $m_s$, $r'$ are the integers with
$L^{N-m_s}\approx\eta^{-1/2}$ and $L^{m_s-r'}\approx\eta L^N$. Along $D_\epsilon$ one has
$\check g_{m_s}\le C_L/(\beta_R+\frac12\beta_1\log(1/\eta))$ for $\eta\le\eta_0$, and
$a_{m_s}\le L^{-N}+L^4\mu(r)\eta$. Consequently $(f_N)$ is asymptotically equicontinuous on $\{|x|\ge r\}$, and
\begin{itemize}
\item every subsequence of $(f_N)$ has a further subsequence that converges uniformly on $\{|x|\ge r\}$;
\item the limit $G_r$ is continuous on $\{|x|\ge r\}$, nonnegative, bounded and invariant under the point group, and it
satisfies $|G_r(x)-G_r(y)|\le\omega_r(|x-y|)$;
\item if $m_N^2L^{2N}\to0$, then for fixed $\beta_R$ a diagonal subsequence converges uniformly on compact subsets of
$\{|x|>r_*(\beta_R)\}$, where $r_*(\beta_R)=\inf\{r:g(r)\ge1/\beta_R\}\asymp(c_L\beta_R)^{-1/(2\kappa_L)}\to0$ as
$\beta_R\to\infty$.
\end{itemize}

\smallskip\noindent\textup{(ii) Short distances at fixed $\epsilon$.} There are $g_0,K,K_c,N_0$, depending on $L$ only,
such that for $N\ge N_0$, $4m^2L^{2N}\le1/\beta_R\le g_0$, every $b$ in the bracket and every lattice point $x\ne0$,
with $\ell=|x|/L^N<1$, $\lambda=\log(1/\ell)>0$ and $k=\min(j(\ell),N-2)$:
\begin{enumerate}[label=(\alph*)]
\item $\bigl|\vartheta_\epsilon^{-2}P_\epsilon(x)-(\beta_k/\beta_R)^{2/3}\bigr|\le K(\beta_k/\beta_R)^{2/3}/\beta_k$ and
$|\beta_k-\beta_1\lambda|\le K_c(1+\beta_R+\log(1+\lambda))$;
\item if $2K_c(1+\beta_R+\log(1+\lambda))\le\beta_1\lambda$, then
\[
\Bigl|\vartheta_\epsilon^{-2}P_\epsilon(x)-\Bigl(\frac{\beta_1\lambda}{\beta_R}\Bigr)^{2/3}\Bigr|\le
\Bigl(\frac{\beta_1\lambda}{\beta_R}\Bigr)^{2/3}\frac{K_c(1+\beta_R+\log(1+\lambda))+2K}{\beta_1\lambda};
\]
\item $\vartheta_\epsilon^{-2}P_\epsilon(x)\le(1+|K|)\bigl((\beta_1\lambda+K_c(1+\beta_R+\log(1+\lambda)))/\beta_R\bigr)^{2/3}$;
in particular, for every $\delta>0$, $\sup\{\vartheta_\epsilon^{-2}P_\epsilon(x):N\ge N_0,\ \ell\ge\delta\}\le C(\delta,\beta_R)$;
\item for fixed $\beta_R$ there are $\delta_0\in(0,1)$ and $K'\ge0$ with
$\bigl|\vartheta_\epsilon^{-2}P_\epsilon(x)-(\beta_1\lambda/\beta_R)^{2/3}\bigr|\le K'(\beta_1\lambda/\beta_R)^{2/3}\log\lambda/\lambda$
for $\ell\le\delta_0$, and every limit of $\vartheta_\epsilon^{-2}P_\epsilon(x_N)$ along samples with $\ell_N\to\ell\in(0,\delta_0)$
satisfies the same bound at $\ell$;
\item $|\beta_\epsilon-N\beta_1\log L|\le K(1+\beta_R+\log N)$, and
$\vartheta_\epsilon^2(\log L^N)^{2/3}\to z_*=(2\pi\beta_R/3)^{2/3}$ with rate $C(\beta_R,L)\log N/N$.
\end{enumerate}
In particular every limit $G_r$ of (i) satisfies the short-distance law of (d) on $\{r\le|x|<\delta_0\}$:
$|G_r(x)-(\beta_1\log(1/|x|)/\beta_R)^{2/3}|\le K'(\beta_1\log(1/|x|)/\beta_R)^{2/3}\log\log(1/|x|)/\log(1/|x|)$.
\end{proposition}

\begin{proof}[Proof sketch of (i)]
Translation invariance gives $P_\epsilon(x)=\P[0\leftrightarrow x]$. Let $X=\lfloor L^Nx\rceil$ and $Y=\lfloor L^Ny\rceil$. Then
$|X|,|Y|\ge\frac r2L^N$ for $N\ge N_0$, and $|Y-X|_\infty\le\eta L^N+1$. So the coalescence scales of $X$ and $Y$ with
$0$ satisfy $j_X,j_Y\ge N-C_\delta(r)$.
\begin{enumerate}
\item \emph{One coordinate at a time.} The pair bound needs $Y-X$ in the scale-$m_s$ block of $0$. Pass from $X$ to $Y$
through $Z=(Y_1,X_2)$, which changes one coordinate per step. In each step exchange the roles of the two points, if
necessary, so that the difference has a nonnegative representative $<L^{m_s}$. All three points have norm
$\ge\frac r4L^N$.
\item \emph{Relocation and the difference datum.} At scale $m_s<\min(j_X,j_Y)$, relocate the $\sigma_X$ observable to
$Y$ (\texttt{rrSt}; $\operatorname{loc}$ of the relocated coupling is exact). This gives
$P(X)-P(Y)=\mathrm{out}(\delta)$, by pair linearity in the $a$-datum. The datum $\delta$ has restart envelope
$M_0\le\mathrm{cddRel}\cdot|\beta^{(0,Y)}_{m_s}|$, where
\[
\mathrm{cddRel}=4K_p(\check g_{m_s}+a_{m_s})+24c_\delta c^2\frac{\check g_{m_s}}{\beta_1\log L/4}
+1344\Bigl(8c_N\bigl(64c_WB(c\check g_{m_s})^2+24\iota_0AL^{-r'}\bigr)+8\Gamma_wc_Wc\check g_{m_s}\Bigr)+\check g_{m_s},
\]
with $|Y-X|_\infty\le L^{m_s-r'}$.
\item \emph{Restarted flow to the top.} In both coalescence cases ($m_s<j_Y\le N_h$ and $N_h+1\le j_Y$),
\[
|P(X)-P(Y)|\le C(r)\bigl(\mathrm{cddRel}\cdot\vartheta_{N_h}\vartheta_{j_Y}+\mathrm{GB}\cdot\vartheta_{j_Y}\bigr)
+C(r)\,\mathrm{cddRel}\,\vartheta_{N_h}^2\check g_{N_h}.
\]
Here $\sigma/|\beta_{a,N_h}|\le(15/14)^{2n}(8/7)^ne^{c}\,\mathrm{cddRel}$ with $n=N_h-j_Y\le C_\delta(r)$. The backward
$\gamma$ chain $\mathrm{GB}$ has $\le C_\delta(r)$ terms. Each is linear in the envelope, since
$M_i\le e^{c}(M_0/|\beta_{m_s}|)|\beta_{m_s+i}|$. So $\mathrm{GB}\le C(r)\,\mathrm{cddRel}\,\vartheta_{N_h}$.
\item \emph{The scalar conditions.} The waiting conditions, $(8/15)^{j_Y-m_s}Q\le1$ and the transients
$(7/30)^{j_Y-m_s}$, $(56/75)^{j_Y-m_s}$, hold for $j_Y-m_s\ge\frac12\log_L(1/\eta)-C_\delta(r)$ once
$\eta\le\eta_0(r)$. This uses $(8/15)^2\cdot\frac{567}{280}<1$. The other conditions depend on $r$:
\begin{itemize}
\item $A^4(c_x\check g^3+k_{xT}+k_2)\le1$ with $k_2\ge c_{ab}c\check g_{j_Y}\bigl((c_W+1)(1+\Theta_b)+1\bigr)^2$, where
$\Theta_b=e^c(81/56)^{N_h-j_Y}\le e^c\Theta(r)$. This holds if $\check g_{j_Y}\le C_L/\beta_R\le c\,\Theta(r)^{-2}$, i.e.\ if
$1/\beta_R\le g(r)$.
\item The mass windows $\mathtt{hmAB}$, $\mathtt{hmAa}$ read $C\,\bar m\,m^2L^{2N}\le1$, where
$\bar m=\max(1,2\Lambda_1)^{C_\delta}m_0$ and $m_0\ge(2\iota(v_T,w_T)A)^{(L+18)^2}$, with
$\iota(v,w)=e^{3/2+v}(2+w)$, $v_T\ge1+\Theta_b$ and $w_T\ge e^{L^2C_w(2\bar k)^2}-1$, $\bar k\asymp\Theta_b$. These hold
if $m^2L^{2N}\le\mu(r)$.
\end{itemize}
The remaining conditions are those of Theorem~\ref{thm:multiscale}.
\item \emph{Normalisation.} $\vartheta_j^2/\vartheta_\epsilon^2=(\beta_j/\beta_R)^{2/3}$ and
$\beta_j\le\beta_R+\beta_1\log(1/r)+C$ for $j\ge N-C_\delta(r)$, so these ratios are $\le C(r)$. By
$|\beta_{m_s}-(N-2-m_s)\beta_1\log L|\le C_L(1+\beta_R+\log(N-m_s))$ (\texttt{ct\_beta\_scale}) and
$\beta_{m_s}\ge\beta_{N-2}\ge\beta_R/C_L$, one gets $\check g_{m_s}\le C_L/(\beta_R+\frac12\beta_1\log(1/\eta))$ for
$\eta\le\eta_0$.
\item \emph{Limits.} The set $\{|x|\ge r\}$ is compact, and $(f_N)$ is uniformly bounded there by (ii)(c) together with
finiteness of $\Lambda_N$ for small $N$. Asymptotic Arzel\`a--Ascoli on this compact set gives the subsequence. The
point-group invariance passes to the limit as in \texttt{t6\_sym\_tendsto}. The diagonal statement follows by applying
(i) along $r_k\downarrow r_*(\beta_R)$; this needs $m_N^2L^{2N}\le\mu(r_k)$ eventually for every $k$.
\end{enumerate}
Part (ii) collects \texttt{thm\_short\_epsX}, \texttt{thm\_short\_eps\_logX}, \texttt{thm\_bound\_eps} and
\texttt{thm\_theta\_eps}.
\end{proof}

\begin{theorem}[Continuum limit at fixed $\beta_R$ on $\T\setminus\{0\}$; proved in Lean, \texttt{Spin2.thm\_limit\_fixed\_betaR}]\label{thm:limit-fixed}
Let $d=2$, $i_0=0$ and $L\ge L_0$. There are $g_0>0$ and $N_0$, depending on $L$ only, such that for every $\beta_R$
with $1/\beta_R\le g_0$ there are $\delta_0\in(0,1)$ and $K'\ge0$ with the following properties. Let
$(R_N,m_N^2,b_N)_{N\ge N_0}$ be admissible for $\beta_R$, i.e.\ $4m_N^2L^{2N}\le1/\beta_R$ and $b_N$ in the bracket. Then
\begin{enumerate}[label=(\alph*)]
\item $(f_N)$ is asymptotically equicontinuous on $\{|x|\ge r\}$ for every $r>0$: for all $r,\eta>0$ there are
$\delta>0$ and $N_1$ with $|f_N(x)-f_N(y)|\le\eta$ whenever $N\ge N_1$, $|x|,|y|\ge r$ and $|x-y|<\delta$;
\item a subsequence of $(f_N)$ converges uniformly on compact subsets of $\T\setminus\{0\}$ to a function $G$ that is
continuous on $\T\setminus\{0\}$, nonnegative, invariant under the point group, bounded on $\{|x|\ge\delta\}$ for every
$\delta>0$, and satisfies, for $0<|x|<\delta_0$ and $\lambda=\log(1/|x|)$,
\[
\Bigl|G(x)-\Bigl(\frac{\beta_1\lambda}{\beta_R}\Bigr)^{2/3}\Bigr|\le K'\Bigl(\frac{\beta_1\lambda}{\beta_R}\Bigr)^{2/3}\frac{\log\lambda}{\lambda}.
\]
\end{enumerate}
In particular Conjecture~\ref{thm:limit} holds for $1/\beta_R\le g_0$ up to the choice of a subsequence (and possibly
of $b_N$, $m_N$).
\end{theorem}

\begin{proof}[Proof sketch]
We keep the notation of the proof of Proposition~\ref{prop:equi-r}(i); only the scalar conditions of step~4 change.
Their dependence on $r$ is removed by two observations, after which $g(r)$ and $\mu(r)$ can be replaced by constants
$g,\mu$ depending on $L$ only.
\begin{enumerate}
\item \emph{Uniform $b$-envelope.} Along $D_\epsilon$ the field scale satisfies $s_k^2=\ell^2/|Z_k|$ exactly, and
$\bigl|\log|Z_{n+1}|+\frac13\log(\check g_n/b)\bigr|\le K_Z\check g_n$. Together with $(H_\lambda)$,
$|\beta_k/b-\theta_k|\le K_p\theta_k(\check g_k+a_k)$ for $k<N_h$, and the monotonicity of $\theta_k=(\check g_k/b)^{1/3}$,
this gives $2s_j^2|\beta_j|\asymp\ell^2b\,\theta_{j-1}\theta_j$, hence for $1\le j\le N_h$
\[
2s_j^2|\beta_j|\le\Theta_*\,p_{M,b}(N_h),\qquad\Theta_*=\tfrac{16}{7}e^{2K_Z}e^{c}(1+3K_p).
\]
The lower bound at $N_h$ uses one step $|\beta_{N_h}|\ge e^{-c}\frac78|\beta_{N_h-1}|$ and
$K_p(\check g_{N_h-1}+a_{N_h-1})\le\frac12$. So $\Theta_b(j_y)$ may be replaced by $\Theta_*$ everywhere: $v_T=1+\Theta_*$,
$\bar k=(c_W+1)(1+\Theta_*)+1$, $w_2=e^{L^2C_w(c+(2\bar k)^2)}-1$, $w_T$ and the seed mass
$m_0\ge(2\iota(v_T,w_T)A)^{(L+18)^2}$ become constants depending on $L$ only.
\item \emph{Homogeneity of the reduced pair chain.} The reduced chain of the difference datum is purely $ab$. For such
elements the norm in the weight $w(KM_a,M_b)$ is exactly $K^{-1}$ times the norm in $w(M_a,M_b)$. Apply the chain bound with
$\sigma'=K\sigma$, seed $k_0/K$ and $\bar m_0/K$, where $K=\max(1,2\Lambda_1)^{C_\delta}$. Then the mass window
\texttt{hmAB} is required at $\bar m_0$ itself, and the factor $K$ appears only in the constant of the conclusion.
\item \emph{Scalar conditions.} By 1 and 2, all smallness conditions read $\check g_{N_h},a_{m_s},a_{N_h-1}\le g_M$ and
$m^2L^{2N}\le\mu_M$ with $g_M,\mu_M>0$ depending on $L$ only. The radius enters only through $J=N-C_\delta(r)$, through
the waiting condition $(3/4)^{J-m_s}K_W(C_\delta(r))\le\check g_{m_s}^2$, and through the constant $C(r)$ in
$|P(X)-P(Y)|\le C(r)\,\mathrm{cddRel}\,\vartheta_{N_h}^2$.
\item \emph{Modulus.} Let $r,\eta>0$. Since $\check g_{m_s}\ge1/(\alpha d+\gamma)$ with $d=N_h-m_s$ and $\alpha,\gamma$
depending on $L,\beta_R$ (\texttt{ct\_beta\_scale}), the waiting condition holds for all $d\ge d_0(r,\beta_R)$; choose
$d$ also so large that $\check g_{m_s}\le\eta'$, and $r'$ with $L^{-r'}\le\eta'$. Then $d,r'$ and
$\delta=\frac12L^{-(d+r'+2)}$ do not depend on $N$, and for $N$ large
$|f_N(x)-f_N(y)|\le10C(r)K_c(\check g_{m_s}+a_{m_s}+L^{-r'})\le\eta$, where $\check g_{N_h}\le4/\beta_R\le4g$,
$a_{m_s},a_{N_h-1}\le L^{-2m_s}+m^2L^{2N}$. This proves (a) on $\{|x|\ge r\}$ whenever $1/\beta_R\le g$ and
$m_N^2L^{2N}\le\mu$.
\item \emph{The window.} Take $g_0\le\min(g,4\mu)$. Admissibility gives $4m_N^2L^{2N}\le1/\beta_R\le4\mu$, so the
mass condition of 4 holds for all $N\ge N_0$, and (a) holds for every $r>0$.
\item \emph{Limits.} Exhaust $\T\setminus\{0\}$ by the compact sets $\{|x|\ge r\}$. By (a), uniform boundedness on each
of them (Proposition~\ref{prop:equi-r}(ii)(c) and finiteness for small $N$) and asymptotic Arzel\`a--Ascoli, a diagonal
subsequence converges uniformly on compact subsets of $\T\setminus\{0\}$. Nonnegativity, the bound and point-group
invariance pass to the limit as in Proposition~\ref{prop:equi-r}(i); the short-distance law is
Proposition~\ref{prop:equi-r}(ii)(d).
\end{enumerate}
\end{proof}

\paragraph{Lean status (Theorem~\ref{thm:limit-fixed}).} Proved:
\texttt{Spin2.thm\_limit\_fixed\_betaR} (\texttt{H02UV/Spin2/Cont/ContFx5b}), on the master construction \texttt{eqs\_master}, with no
hypothesis beyond $1/\beta_R\le g_0$ and \texttt{ClAdm}; \texttt{MasterFacts.fx\_clEqui} gives \texttt{ClEqui} (part (a)).
\begin{itemize}
\item Step 1: \texttt{MasterFacts.fxTh} ($\Theta_*$), \texttt{fx\_pMS\_ratioU} (\texttt{H02UV/Spin2/Cont/ContFx1}).
Step 2: \texttt{PtCorr2.fx\_oTo\_wM\_ab\_scale}, \texttt{fx\_oTo\_wM\_ab\_le} (\texttt{H02UV/PtCorr2/Cont/ContFx1}), \texttt{fx\_T5S}
(\texttt{H02UV/Spin2/Cont/ContFx2a}).
\item The intermediate case with $\Theta_*$, constant $\bar k$ and the window at $\bar m_0$: \texttt{MasterFacts.fx\_T6\_mid}
(\texttt{H02UV/Spin2/Cont/ContFx2a--f}); package \texttt{FxMidPkg}, \texttt{fx\_eqc\_mid\_pkg}, \texttt{fx\_eqc\_all\_pairs}
(\texttt{H02UV/Spin2/Cont/ContFx3a--b}). The top case is unchanged.
\item Step 3: \texttt{fx\_mid\_pkg}, \texttt{fx\_pairs} (\texttt{H02UV/Spin2/Cont/ContFx4a--b}) with thresholds \texttt{fxGg}~$r_0$,
\texttt{fxGm}~$r_0\,K_t\,C_{1*}$ that do not depend on $C_\delta$; the waiting constant \texttt{eqsKWa}~$C_\delta$ is unchanged.
\item Steps 4--6: \texttt{fx\_equiOn} (\texttt{ClEquiOn} for every $r$ with $r$-independent $g,\mu$; \texttt{H02UV/Spin2/Cont/ContFx5a}),
\texttt{fx\_clEqui} ($g_0=\min(g,4\mu)$, \texttt{clEqui\_iff\_on}) and \texttt{thm\_limit\_fixed\_betaR}
(\texttt{eqs\_limitAnn} at $r_0=0$; \texttt{H02UV/Spin2/Cont/ContFx5b}).
\end{itemize}

\subsubsection{The observable renormalisation condition}\label{sec:obs}
Let $x_0=(\frac12,0)\in\T$ and let $x_{0,N}\in\Lambda_N$ be its lattice point (so $|x_{0,N}|/L^N\ge\frac14$
and $j(|x_{0,N}|/L^N)\ge N-2$). For $\beta>0$ and $N\ge1$ define
\begin{equation}\label{e:tdef}
  t(\beta,N)=\beta\,\bigl(\mathbb{P}^{\Lambda_N}_\beta[0\conn x_{0,N}]\bigr)^{3/2}.
\end{equation}
By the multiscale formula at the scale $N-2$, $\mathbb{P}_\beta[0\conn x_{0,N}]=(\tilde\beta_{N-2}/\beta)^{2/3}(1+O(1/\tilde\beta_{N-2}))$,
where $\tilde\beta_{N-2}$ is the running inverse temperature at the macroscopic scale; hence
$t(\beta,N)=\tilde\beta_{N-2}+O(1)$. Thus $t(\beta,N)$ is the renormalised inverse temperature, defined through the
model alone. For $\beta>0$, $t(\cdot,N)$ is continuous. The renormalisation condition $t(\beta_N,N)=t$ replaces the
condition on $\beta_R$ of Theorem~\ref{thm:limit-fixed}. Under it, the natural normalisation of the two-point
function is the ratio
\begin{equation}\label{e:fN-obs}
  f_N(x)=\frac{\mathbb{P}_{\beta_N}[0\conn x_N]}{\mathbb{P}_{\beta_N}[0\conn x_{0,N}]}=\Bigl(\frac{\beta_N}{t}\Bigr)^{2/3}\mathbb{P}_{\beta_N}[0\conn x_N],
  \qquad x\in\T,
\end{equation}
which satisfies $f_N(x_0)=1$.

\paragraph{Lean.} $t(\beta,N)$ is \texttt{Spin2.obsBeta L N $\beta$}, $x_0$ is \texttt{Spin2.obsX0}
(\texttt{torusProj ![1/2, 0]}), $x_{0,N}$ is \texttt{Spin2.obsX0N} (\texttt{latPtN L N obsX0}), and \eqref{e:fN-obs}
is \texttt{Spin2.obsF L $(\beta_N)_N$ N x} (all in \texttt{H02UV/Spin2/Final/ObsBeta1.lean}). The continuity of
$\beta\mapsto\mathbb{P}_\beta[0\conn x]$ is \texttt{Spin2.Pconn\_continuousOn} (\texttt{H02UV/Spin2/Multiscale/ConnSym.lean}).
The identity $\mathbb{P}_{\beta_N}[0\conn x_{0,N}]=(t/\beta_N)^{2/3}$ under $t(\beta_N,N)=t$ is \texttt{ob\_Pconn\_of\_obsBeta},
and $f_N(x_0)=1$ is \texttt{obsF\_obsX0}.

Here $t(\beta,N)$ is \eqref{e:tdef} and $f_N$ is \eqref{e:fN-obs}. For $x\in\Lambda_N$ we write
$\lambda=\log(L^N/|x|)$, $k=\min(j(|x|/L^N),N-2)$, and $\beta_1=\frac3{2\pi}$.

\begin{theorem}[Observable renormalisation condition; proved in Lean, \texttt{Spin2.thm\_arboreal\_obs}]\label{thm:obs}
There are $t_0\ge1$, $N_0\ge3$ and $K,K_c,K_s\in\R$, depending only on $L$, such that for every $N\ge N_0$ and every
$t\ge t_0$ there is $\beta_N>0$ with
\begin{equation}\label{e:obs-cond}
  \beta_N\,\mathbb{P}^{\Lambda_N}_{\beta_N}[0\conn x_{0,N}]^{3/2}=t
\end{equation}
and the following properties.
\begin{enumerate}[label=(\roman*)]
\item $\bigl|\beta_N-\frac3{2\pi}\log L^N\bigr|\le K(1+t+\log N)$.
\item There are $\tilde\beta_0=\beta_N,\tilde\beta_1,\dots,\tilde\beta_{N-2}$ with $\tilde\beta_j\ge t/32$ and
$|\tilde\beta_{N-2}-t|\le K$, such that for every $x\in\Lambda_N\setminus\{0\}$,
\begin{equation}
  |\tilde\beta_k-\beta_1\lambda|\le K_c(1+t+\log(1+\lambda)),\qquad
  \bigl|\mathbb{P}_{\beta_N}[0\conn x]-\vartheta_k^2\bigr|\le K\,\frac{\vartheta_k^2}{\tilde\beta_k},\qquad
  \vartheta_k^2=\Bigl(\frac{\tilde\beta_k}{\beta_N}\Bigr)^{2/3}.
\end{equation}
\item For every $x\in\Lambda_N\setminus\{0\}$, with $\delta(\lambda)=K_c(1+t+\log(1+\lambda))$: if
$2\delta(\lambda)\le\beta_1\lambda$ then
\begin{equation}
  \Bigl|\frac{\mathbb{P}_{\beta_N}[0\conn x]}{\mathbb{P}_{\beta_N}[0\conn x_{0,N}]}-\Bigl(\frac{\beta_1\lambda}{t}\Bigr)^{2/3}\Bigr|
  \le\Bigl(\frac{\beta_1\lambda}{t}\Bigr)^{2/3}\frac{\delta(\lambda)+2K_s}{\beta_1\lambda},
\end{equation}
and in any case
$\mathbb{P}_{\beta_N}[0\conn x]/\mathbb{P}_{\beta_N}[0\conn x_{0,N}]\le(1+|K_s|)\bigl((\beta_1\lambda+\delta(\lambda))/t\bigr)^{2/3}$.
\end{enumerate}
\end{theorem}

\begin{theorem}[Continuum limit under the observable renormalisation condition; proved in Lean, \texttt{Spin2.thm\_arboreal\_limit\_obs}]\label{thm:limit-obs}
There are $t_0\ge1$, $N_0\ge3$ and $K,K_c,K_s$ as in Theorem~\ref{thm:obs} such that for every $t\ge t_0$ there are
$\delta_0\in(0,1)$, $K'\ge0$, depending only on $L$ and $t$, and a sequence $(\beta_N)_N$ with
\eqref{e:obs-cond} and the conclusions of Theorem~\ref{thm:obs} for every $N\ge N_0$, such that:
there are a subsequence $(N_k)$ and a function $G$ on $\T$ with $f_{N_k}\to G$ uniformly on compact subsets of
$\T\setminus\{0\}$; $G$ is continuous on $\T\setminus\{0\}$, nonnegative, invariant under coordinate permutations
and reflections, bounded on $\{|x|\ge\delta\}$ for every $\delta>0$, $G(x_0)=1$, and for $0<|x|<\delta_0$,
\begin{equation}
    \Biggl|G(x)-\Biggl(\frac{3\log\frac1{|x|}}{2\pi t}\Biggr)^{2/3}\Biggr|
    \le K'\Biggl(\frac{3\log\frac1{|x|}}{2\pi t}\Biggr)^{2/3}\frac{\log\log\frac1{|x|}}{\log\frac1{|x|}}.
\end{equation}
\end{theorem}

\paragraph{Lean status (Theorems~\ref{thm:obs}, \ref{thm:limit-obs}).} Theorem~\ref{thm:obs} is \texttt{Spin2.thm\_arboreal\_obs} (\texttt{H02UV/Spin2/Final/ObsBeta2.lean}) and
Theorem~\ref{thm:limit-obs} is \texttt{Spin2.thm\_arboreal\_limit\_obs} (\texttt{H02UV/Spin2/Final/ObsBeta3.lean}). The condition
\eqref{e:obs-cond} is \texttt{obsBeta L N $\beta_N$ = t}; (i), (ii) are the predicate \texttt{ObsScale L N t K $K_c$ $K_\epsilon$ E a $\beta_N$ $\tilde\beta$}
(with the constant of (i) called $K_\epsilon$, and in addition the scale-by-scale bound
$|\tilde\beta_j-(N-2-j)\beta_1\log L|\le(40t+E)+a\log(2((N-2-j)\beta_1\log L+40t+E+a|\log2a|))$); (iii) is
\texttt{ArbShort L N t $K_s$ $K_c$ $\beta_N$}, the predicate of Theorem~\ref{thm:limit-fixed} with $t$ in place of $\beta_R$,
since $\mathbb{P}_{\beta_N}[0\conn x]/\mathbb{P}_{\beta_N}[0\conn x_{0,N}]=\mathbb{P}_{\beta_N}[0\conn x]/(t/\beta_N)^{2/3}$ under \eqref{e:obs-cond}.
In Theorem~\ref{thm:limit-obs}, $f_N$ is \texttt{obsF L $(\beta_N)_N$ N}, and the Lean statement also contains the
asymptotic equicontinuity of $(f_N)$ on every $\{\|x\|\ge r\}$ (\texttt{ClEqui}). Choices made in the proof:
\begin{itemize}
\item The auxiliary mass is $m_N^2=1/(32tL^{2N})$ (\texttt{obsMass}), so that $4m_N^2L^{2N}=1/(8t)\le1/\beta_R$ for the
  four values $\beta_R=t4^i/8$, $i=0,\dots,3$ (\texttt{obsBR}). The brackets $[b_{\rm Fin}(1/(2\beta_R)),b_{\rm Fin}(2/\beta_R)]$
  at these values share endpoints, and their union is an interval $[a_4,a_0]$ (\texttt{MasterFacts.obs\_union}).
\item On a bracket at $\beta_R$ the multiscale formula at $x_{0,N}$, scale $N-2$, gives
  $|t(1/b,N)-\tilde\beta_{N-2}(b)|\le3K$ (\texttt{ob\_obsBeta\_close}, from \texttt{arb\_perN}); with $\tilde\beta_{N-2}\ge\beta_R/4$ and
  $\tilde\beta_{N-2}\le5\beta_R$ this gives $t(1/a_4,N)>t>t(1/a_0,N)$, and $\beta_N=1/b_N$ is obtained from the intermediate value
  theorem (\texttt{MasterFacts.ob\_master}). $b_N$ is one solution; uniqueness is not claimed. The constants $t/32$, $40t$,
  $K_c$ (eight times the constant of \texttt{arb\_perN}) and $K$ (three times) come from $\beta_R\in[t/8,8t]$
  (\texttt{obsScale\_of}, \texttt{ArbShort.rescale}); $t_0$ includes $t\ge128K$, $t\ge32/\varepsilon$ and $8/t\le g_0$.
\item In Theorem~\ref{thm:limit-obs} the $b_N$ lie in brackets at $N$-dependent $\beta_R(N)\in\{t4^i/8\}$, whereas the
  equicontinuity, boundedness and short-distance inputs (\texttt{fx\_clEqui}, \texttt{t6\_bdd}, \texttt{arb\_eqs\_shortBound}) are for
  one fixed $\beta_R$ and an admissible sequence (\texttt{ClAdm}). For each of the four values an auxiliary sequence
  $b^{(i)}_N$ (equal to $b_N$ if $b_N$ lies in the bracket at $t4^i/8$, the left end of that bracket otherwise) is
  admissible at fixed $\beta_R=t4^i/8$, and $f_N=(\beta_R/t)^{2/3}\vartheta_\epsilon^{-2}P_\epsilon$ with
  $(\beta_R/t)^{2/3}\le8$ (\texttt{clF\_rescale}); the four families are combined by \texttt{ob\_equiOn\_of\_family},
  \texttt{ob\_bdd\_of\_family}, and the short-distance law uses $(\beta_R/t)^{2/3}(\beta_1\lambda/\beta_R)^{2/3}=(\beta_1\lambda/t)^{2/3}$
  (\texttt{ob\_csA\_rescale}); $\delta_0\le e^{-1}$ and $K'$ is the largest of the four constants.
  The limit assembly (Arzel\`a--Ascoli on annuli, symmetry, continuity of the short-distance law) is that of
  \texttt{thm\_limit\_fixed\_betaR}. $G(x_0)=1$ holds since $f_N(x_0)=1$ exactly ($x_{0,N}$ is by definition the lattice
  point of $x_0$ used in \eqref{e:tdef}).
\end{itemize}
Convergence is along a subsequence; in the boundedness the distance is the max-norm, in the short-distance bound the
Euclidean norm.

\paragraph{The summed condition.} The condition \eqref{e:tdef} refers to the point $x_0$. The preferred,
reference-point-free version uses the expected size of the tree $T_0$ of the origin: by linearity of the expectation,
$\E_\beta|T_0|=\sum_{x\in\Lambda_N}\mathbb{P}_\beta[0\conn x]$ (the term $x=0$ equals $1$). For $\beta>0$ and $N\ge1$ define
\begin{equation}\label{e:tsum-def}
  \rho(\beta,N)=\frac{\E_\beta|T_0|}{|\Lambda_N|}=\frac1{L^{2N}}\sum_{x\in\Lambda_N}\mathbb{P}^{\Lambda_N}_\beta[0\conn x],
  \qquad t_\Sigma(\beta,N)=\beta\,\rho(\beta,N)^{3/2}.
\end{equation}
By the multiscale formula, $\mathbb{P}_\beta[0\conn x]=\vartheta_{k}^2(1+O(1/\tilde\beta_{k}))$ with $k=\min(j(|x|/L^N),N-2)$,
and the running couplings are monotone, $\tilde\beta_k\ge\tilde\beta_{N-2}$ (equivalently $\check g_k\le\check g_{N-2}$),
with the per-step recursion $\tilde\beta_k\le\tilde\beta_{N-2}+(N-2-k)\beta_1\log L+O(1)$; since only a fraction
$O(L^{-2(1+i)})$ of the points of $\Lambda_N$ has $k=N-2-i$, this gives
$\rho(\beta,N)=\vartheta_{N-2}^2(1+O(1/\tilde\beta_{N-2}))$ and hence $t_\Sigma(\beta,N)=\tilde\beta_{N-2}+O(1)$, as for
\eqref{e:tdef}. For $\beta>0$, $t_\Sigma(\cdot,N)$ is continuous, and $0<\rho(\beta,N)\le1$. Under the renormalisation
condition $t_\Sigma(\beta_N,N)=t$, the normalised two-point function is
\begin{equation}\label{e:fN-sum}
  f^\Sigma_N(x)=\frac{\mathbb{P}_{\beta_N}[0\conn x_N]}{\rho(\beta_N,N)}=\Bigl(\frac{\beta_N}{t}\Bigr)^{2/3}\mathbb{P}_{\beta_N}[0\conn x_N],
  \qquad x\in\T,
\end{equation}
whose lattice average is $1$: $L^{-2N}\sum_{x\in\Lambda_N}\mathbb{P}_{\beta_N}[0\conn x]/\rho(\beta_N,N)=1$.

\begin{theorem}[Summed observable renormalisation condition; proved in Lean, \texttt{Spin2.thm\_arboreal\_sum}]\label{thm:obs-sum}
There are $t_0\ge1$, $N_0\ge3$ and $K,K_c,K_s\in\R$, depending only on $L$, such that for every $N\ge N_0$ and every
$t\ge t_0$ there is $\beta_N>0$ with
\begin{equation}\label{e:obs-sum-cond}
  \beta_N\,\rho(\beta_N,N)^{3/2}=\beta_N\Bigl(\frac{\E_{\beta_N}|T_0|}{|\Lambda_N|}\Bigr)^{3/2}=t
\end{equation}
and the following properties.
\begin{enumerate}[label=(\roman*)]
\item $\bigl|\beta_N-\frac3{2\pi}\log L^N\bigr|\le K(1+t+\log N)$.
\item There are $\tilde\beta_0=\beta_N,\tilde\beta_1,\dots,\tilde\beta_{N-2}$ with $\tilde\beta_j\ge t/32$ and
$|\tilde\beta_{N-2}-t|\le K$, such that for every $x\in\Lambda_N\setminus\{0\}$,
\begin{equation}
  |\tilde\beta_k-\beta_1\lambda|\le K_c(1+t+\log(1+\lambda)),\qquad
  \bigl|\mathbb{P}_{\beta_N}[0\conn x]-\vartheta_k^2\bigr|\le K\,\frac{\vartheta_k^2}{\tilde\beta_k},\qquad
  \vartheta_k^2=\Bigl(\frac{\tilde\beta_k}{\beta_N}\Bigr)^{2/3}.
\end{equation}
\item For every $x\in\Lambda_N\setminus\{0\}$, with $\delta(\lambda)=K_c(1+t+\log(1+\lambda))$: if
$2\delta(\lambda)\le\beta_1\lambda$ then
\begin{equation}
  \Bigl|\frac{\mathbb{P}_{\beta_N}[0\conn x]}{\rho(\beta_N,N)}-\Bigl(\frac{\beta_1\lambda}{t}\Bigr)^{2/3}\Bigr|
  \le\Bigl(\frac{\beta_1\lambda}{t}\Bigr)^{2/3}\frac{\delta(\lambda)+2K_s}{\beta_1\lambda},
\end{equation}
and in any case
$\mathbb{P}_{\beta_N}[0\conn x]/\rho(\beta_N,N)\le(1+|K_s|)\bigl((\beta_1\lambda+\delta(\lambda))/t\bigr)^{2/3}$.
\end{enumerate}
\end{theorem}

\begin{theorem}[Continuum limit under the summed observable renormalisation condition; proved in Lean, \texttt{Spin2.thm\_arboreal\_limit\_sum}]\label{thm:limit-obs-sum}
There are $t_0\ge1$, $N_0\ge3$ and $K,K_c,K_s$ as in Theorem~\ref{thm:obs-sum} such that for every $t\ge t_0$ there are
$\delta_0\in(0,1)$, $K'\ge0$, depending only on $L$ and $t$, and a sequence $(\beta_N)_N$ with
\eqref{e:obs-sum-cond} and the conclusions of Theorem~\ref{thm:obs-sum} for every $N\ge N_0$, such that:
$L^{-2N}\sum_{x\in\Lambda_N}f^\Sigma_N(x/L^N)=1$ for every $N\ge N_0$;
there are a subsequence $(N_k)$ and a function $G$ on $\T$ with $f^\Sigma_{N_k}\to G$ uniformly on compact subsets of
$\T\setminus\{0\}$; $G$ is continuous on $\T\setminus\{0\}$, nonnegative, invariant under coordinate permutations
and reflections, bounded on $\{|x|\ge\delta\}$ for every $\delta>0$, and for $0<|x|<\delta_0$,
\begin{equation}
    \Biggl|G(x)-\Biggl(\frac{3\log\frac1{|x|}}{2\pi t}\Biggr)^{2/3}\Biggr|
    \le K'\Biggl(\frac{3\log\frac1{|x|}}{2\pi t}\Biggr)^{2/3}\frac{\log\log\frac1{|x|}}{\log\frac1{|x|}}.
\end{equation}
\end{theorem}

The normalisation $\int_\T G=1$ (from the lattice averages, uniform convergence on $\{|x|\ge\delta\}$ and the smallness
of the contribution of $\{|x|<\delta\}$, uniformly in $N$, by the short-distance bound (iii)) is expected but not
included in the Lean statement; see below.

\paragraph{Lean status (Theorems~\ref{thm:obs-sum}, \ref{thm:limit-obs-sum}).}
$\rho(\beta,N)$ is \texttt{Spin2.obsDens L N $\beta$}, $t_\Sigma(\beta,N)$ is \texttt{Spin2.obsBetaSum L N $\beta$} and
\eqref{e:fN-sum} is \texttt{Spin2.obsFS L $(\beta_N)_N$ N x} (\texttt{H02UV/Spin2/Final/ObsSum1.lean}; $\mathbb{P}_\beta[0\conn0]=1$ is
\texttt{Pconn\_zero\_eq\_one}, $0<\rho\le1$ is \texttt{obsDens\_pos}, \texttt{obsDens\_le\_one}, continuity is
\texttt{obsDens\_continuousOn}, the identity $\rho(\beta_N,N)=(t/\beta_N)^{2/3}$ under \eqref{e:obs-sum-cond} is
\texttt{os\_dens\_of\_obsBetaSum}, $f^\Sigma_N=\texttt{clF L t b N}$ is \texttt{obsFS\_eq\_clF}, and the lattice average is
\texttt{obsFS\_avg}). Theorem~\ref{thm:obs-sum} is \texttt{Spin2.thm\_arboreal\_sum} (\texttt{H02UV/Spin2/Final/ObsSum2.lean}) and
Theorem~\ref{thm:limit-obs-sum} is \texttt{Spin2.thm\_arboreal\_limit\_sum} (\texttt{H02UV/Spin2/Final/ObsSum3.lean}); their
statements are those of \texttt{thm\_arboreal\_obs}, \texttt{thm\_arboreal\_limit\_obs} with \texttt{obsBeta} replaced by
\texttt{obsBetaSum}, \texttt{obsF} by \texttt{obsFS}, the same predicates \texttt{ObsScale}, \texttt{ArbShort L N t}
(reused unchanged), and the clause $G(x_0)=1$ replaced by the lattice average
$L^{-2N}\sum_{x\in\Lambda_N}\mathbb{P}_{\beta_N}[0\conn x]/\rho(\beta_N,N)=1$ (in the per-$N$ clause).
\begin{itemize}
\item \emph{The key estimate} (\texttt{os\_dens\_close}): for $b$ in a bracket at $\beta_R$, with
$\vartheta_k^2=(\tilde\beta_k/\beta_\epsilon)^{2/3}$ along the $h=0$ trajectory,
$|\rho(1/b,N)-\vartheta_{N-2}^2|\le(4C_1+2K)\,\vartheta_{N-2}^2/\tilde\beta_{N-2}$, where $K$ is the constant of the
multiscale formula (\texttt{arb\_perN}) and $C_1=\beta_1\log L+C_{gu}+4K_b$ (\texttt{MasterFacts.osC1}). Inputs: the
multiscale formula for every $x\ne0$ (\texttt{ArbScale}), $\mathbb{P}[0\conn0]=1=\vartheta_0^2$; the monotonicity
$\check g_k\le\check g_{k+1}$ of the running coupling along every trajectory in $D_\epsilon$ (\texttt{gmonoD}), giving the
lower bound $\vartheta_k^2\ge\vartheta_{N-2}^2$; the per-step recursion
$|\tilde\beta_{k+1}-\tilde\beta_k+\beta_1\log L|\le K_b a_k+C_{gu}\check g_k$ of \texttt{spinD} (the flow-level clause
behind the per-step statement of \texttt{thm\_beta}; it holds along every trajectory in $D_\epsilon$), summed over the
steps $k,\dots,N-3$ with $\sum_ka_k\le4$ and $\check g_k\le\epsilon\le1$:
$\tilde\beta_k\le\tilde\beta_{N-2}+C_1(1+(N-2-k))$ (\texttt{MasterFacts.os\_beta\_rec}); this is the sharp
version (error $O(1)$, not $O(\log\tilde\beta_{N-2})$). Hence
$\vartheta_k^2/\vartheta_{N-2}^2\le\tilde\beta_k/\tilde\beta_{N-2}\le1+C_1(1+i)/\tilde\beta_{N-2}$, $i=N-2-k$ (the
exponent $2/3$ is replaced by $1$). Counting: $k_x\le N-2-i$ with $i\ge1$ forces $|x|_\infty<L^{N-1-i}$
(\texttt{os\_supNorm\_lt}, from $j(|x|/L^N)=\lfloor\log_L|x|\rfloor$), and
$\#\{|x|_\infty\le L^{m+1}-1\}\le4L^{2(m+1)}$ (\texttt{os\_card\_ball}, from \texttt{FRD.card\_supNorm\_le}), so
$\sum_x(1+i_x)\le2L^{2N}$ (\texttt{os\_geom}); the point $x=0$ is counted with $k=0$.
\item From the key estimate, $|t_\Sigma(1/b,N)-\tilde\beta_{N-2}|\le3C_L$, $C_L=4C_1+2K$, once $\beta_R\ge16C_L$
(\texttt{os\_obsBetaSum\_close}, \texttt{MasterFacts.os\_perM}); the constants of \texttt{ObsScale} are then
$3C_L$ (in place of $3K$), $8K_c$, $8K_\epsilon$, and $t_0$ contains $128C_L$. The rest is a copy of the proofs of
Theorems~\ref{thm:obs}, \ref{thm:limit-obs}: the same auxiliary mass $m_N^2=1/(32tL^{2N})$, the same four brackets and
their union (\texttt{obs\_union}), continuity of $b\mapsto t_\Sigma(1/b,N)$ (finite sum of
\texttt{Pconn\_continuousOn}), the intermediate value theorem (\texttt{MasterFacts.os\_master}), the renormalisation
$\beta_R\to t$ (\texttt{obsScale\_of}, \texttt{ArbShort.rescale}), $f^\Sigma_N=\texttt{clF L t b N}$ exactly, the
four-family combination and Arzel\`a--Ascoli.
\item \emph{Not formalised}: $\int_\T G=1$. It would follow from the lattice averages, the uniform convergence on
$\{|x|\ge\delta\}$ and a bound $L^{-2N}\sum_{|x_N|<\delta L^N}f^\Sigma_N\le C(\delta)\to0$ uniformly in $N$ (shell
counting with the bound (iii)), together with Riemann-sum convergence on $\T$; this was not attempted.
\end{itemize}

\paragraph{Lean status (Proposition~\ref{prop:equi-r}).} Part (ii) is proved: \texttt{Spin2.thm\_short\_epsX},
\texttt{thm\_short\_eps\_logX}, \texttt{thm\_bound\_eps}, \texttt{thm\_theta\_eps(\_tendsto)}. For part (i), the pair
bounds of step 3 are proved with all non-scalar inputs discharged (\texttt{MasterFacts.fin\_T6\_mid},
\texttt{fin\_conn\_diff\_top}).
\begin{itemize}
\item \emph{$c_{ab}>0$} (\texttt{H02UV/Spin2/Cont/ContEq1}). \texttt{kobsBoundsU\_mono\_C\(_0\)}: \texttt{KObsBoundsU} is monotone in
$C_0$. \texttt{masterCore\_exists\_C\(_0\)} is a copy of the master construction run with $C_0\ge1$, and
\texttt{thm\_master\_all\_cab} gives $M$ with $0<c_{ab}$. The existing library is unchanged.
\item \emph{Step 1} (\texttt{H02UV/Spin2/Cont/ContEq3}). \texttt{eq\_pair\_chain}, \texttt{eq\_abs\_sub\_le\_two}: the intermediate
point $Z=(Y_1,X_2)$, with one of the two orders good in each step, so a bound $B$ for good pairs gives $2B$.
\item \emph{Limits on $\{|x|\ge r\}$} (\texttt{Cont/AscoliOn}, \texttt{H02UV/Spin2/Cont/ContEq2}).
\begin{itemize}
\item \texttt{ClEquiOn}~$f\,r$ is \texttt{ClEqui} at one radius.
\item \texttt{ascoli\_asymp\_exists\_subseq\_on} handles one compact set; \texttt{\_annulus} handles $\{r_0<|x|\}$.
\item \texttt{thm\_limit\_of\_equiOn}: given \texttt{ClEquiOn} at $r$, a subsequence converges uniformly on
$\{\|x\|\ge r\}$. The limit is uniformly continuous there, nonnegative and bounded; it is invariant on $\{\|x\|>r\}$ and
obeys the short-distance law.
\item \texttt{thm\_limit\_of\_equiAnn}: given \texttt{ClEquiOn} for all $r>r_0$, the same holds on compacts of
$\{\|x\|>r_0\}$.
\end{itemize}
\end{itemize}
\begin{itemize}
\item \emph{Step 3, the chain part} (\texttt{H02UV/Spin2/Cont/ContEqC1--3}). \texttt{MasterFacts.eqc\_nbGB\_le}:
$\mathrm{GB}\le(N_h-j_Y+1)\,K\,R_l\,b\theta_{N_h}^2$ with $b\theta_{N_h}^2=\check g_{N_h}\vartheta_{N_h}$, for every envelope
$M_0\le R_l|\beta_{m_s}|$ (for the difference datum $R_l=\mathrm{cddRel}$). The inputs are $M_i\le K_eR_l\,b\theta_{m_s+i}$
(\texttt{cw4\_prod\_cmpD}, \texttt{cw2\_ratioD}, $|\beta_k|\le(1+3K_p)b\theta_k$ from $(H_\lambda)$),
$|Z_{n+1}|^{-1}\le e^{K_Z}\theta_n$, $\|\Gamma_k\|\le\ell^2|Z_{k+1}|^{-1}/16$, $s_k^2\le\frac98\ell^2|Z_{k+1}|^{-1}$,
$\|C_k(a,a)\|\le s_k^2/14$, and $\theta$ nondecreasing. The transients $k,w$ are bounded by $s_{m_s}^2M_0$ without division,
because $\Gamma_w\ge8c_N\cdot120B_k c\epsilon+c_B$. In addition, \texttt{eqc\_ctSig\_rel} gives $\sigma\le e^{c'}e^{c}R_l|\beta_{a,N_h}|$,
and \texttt{eqc\_varth\_ratio} gives $\vartheta_j\le K_g^{N_h-j}\vartheta_{N_h}$ with $K_g=\max(2,2K_1)$, from the $1/\check g$
recursion. \texttt{MasterFacts.eqc\_mid} (\texttt{H02UV/Spin2/Cont/ContEqC4}) is \texttt{fin\_T6\_mid} with conclusion
$|P(x)-P(y)|\le C_{\rm mid}\,\mathrm{cddRel}\,\vartheta_{N_h}^2$. The top case has no chain part left.
\item \emph{Degenerate cases and all pairs} (\texttt{H02UV/Spin2/Cont/ContEqC5--6}).
\begin{itemize}
\item The pair-dependent scalar hypotheses are collected in the named packages \texttt{EqcMidPkg} and \texttt{EqcTopPkg}.
\item \texttt{eqc\_mid\_pkg} and \texttt{eqc\_top\_pkg} turn them into bounds of the form $C\,\mathrm{cddRel}\,\vartheta_{N_h}^2$.
\item \texttt{eqc\_all\_pairs} assumes the packages for all good pairs in $\{J\le j(0,\cdot)\}$, with $m_s<J$ and $N\le J+C_\delta$.
It gives $|P(x)-P(y)|\le2\max(C_{\rm mid},C_{\rm top})\,\mathrm{cddRel}\,\vartheta_{N_h}^2$ for every $x$ whose
$L^{m_s-r}$-ball lies in that region and every $|y-x|_\infty\le L^{m_s-r}$.
\begin{itemize}
\item The boundary $j_Y=N_h$ versus $N_h+1$ is the case split.
\item The intermediate point lies in the region, so $m_s<j_X$ holds for it.
\item \texttt{eqc\_region} keeps the ball inside the region.
\item Small $m_s$ needs no separate bound: for fixed $\eta$, $m_s\to\infty$ as $N\to\infty$, so it is excluded by $N_0$. The
pointwise fallback is \texttt{eqc\_fallback}.
\end{itemize}
\end{itemize}
\end{itemize}

\emph{Part (i) is proved} (\texttt{H02UV/Spin2/Cont/ContEqS0--8}, \texttt{H02UV/Spin2/Cont/ContEqF1--5}).
\begin{itemize}
\item \emph{Master construction.} \texttt{eqs\_master}: a copy of the master construction with $C_{0,u}\ge1$ and
$C_{0,o}\ge1$ (\texttt{KObsBounds} is monotone in $C_0$), so $c_{ab}>0$ and $B>0$; $\Gamma_w>0$ and $c_N\ge1$ follow from
$c_{bn},c_N\ge1$. The short-distance chain is copied for a given $M$ (\texttt{eqs\_constsT}, \texttt{eqs\_allB},
\texttt{eqs\_shortX}, \texttt{eqs\_shortBound}, \texttt{eqs\_limitOn}, \texttt{eqs\_limitAnn}); \texttt{eqs\_bracketUp}:
$\check g_j\le4/\beta_R$ on the bracket.
\item \emph{Step 4 (scalar discharge).} \texttt{MasterFacts.eqs\_mid\_pkg}, \texttt{eqs\_top\_pkg} produce the packages
\texttt{EqcMidPkg}, \texttt{EqcTopPkg} for every admissible transient pair by explicit choices ($v_T=1+\Theta(C_\delta)$,
$m_0=m_0(C_\delta)$, $Q_0=(7/5)^n(\tau_w/(c\check g_{m_s})+\tau_k/(8B(c\check g_{m_s})^2))$, $mm$, $k_1$, $c_1$), from three
inputs: $\check g_{N_h},a_{m_s}\le g_M(C_\delta)$, $m^2L^{2N}\le\mu_M(C_\delta)$, and the waiting condition
$(3/4)^{J-m_s}K_W(C_\delta)\le\check g_{m_s}^2$. The waiting term uses $Z_{ab}/(c_{ab}r_{m_s})\le
(7/5)^{J-m_s}\Theta_b(m_s)K_2/\check g_{m_s}^2$ (\texttt{eqs\_zab\_le}, from $|\beta^{(y-x,y)}|\le3|\beta^{(0,y)}|$ and
$M_0\ge\check g|\beta^{(0,y)}|$) and $(8/15)^2\cdot\frac75\cdot\frac{81}{56}<\frac34$. \texttt{MasterFacts.eqs\_pairs}:
\texttt{eqc\_all\_pairs} with every scalar hypothesis discharged (shared ones from $g_M$, $\mu_M$, \texttt{M.jet},
\texttt{lastOut\_zeroT}, \texttt{topSeedD}).
\item \emph{Step 5 (modulus).} \texttt{MasterFacts.eqf\_equiOn}: with $C_\delta(r)=\lceil\log_L(1/r)\rceil+3$,
$g(r)=\min(g_A,g_U,g_M(C_\delta(r))/4)$ and $\mu(r)=\min(\mu_M(C_\delta(r)),g_M(C_\delta(r))/2,1)$, every admissible sequence
with $1/\beta_R\le g(r)$ and $m_N^2L^{2N}\le\mu(r)$ eventually satisfies \texttt{ClEquiOn}$(\mathrm{clF})\,r$. For given
$\varepsilon$ the gap $d=N_h-m_s$ and $r'$ are fixed (independent of $N$): $d$ makes the waiting condition hold and
$\check g_{m_s}\le\varepsilon'$ (\texttt{ct\_beta\_scale}: $|\beta_{m_s}-d\beta_L|\le T+a\log(2(d\beta_L+T+c'))$), and
$\delta=\frac12L^{-(d+r'+2)}$; then $|\mathrm{clF}(x)-\mathrm{clF}(y)|\le10K\,\mathrm{cddRel}\le
10KK_c(\check g_{m_s}+a_{m_s}+L^{-r'})$ (\texttt{eqf\_cddRel\_le}, \texttt{eqf\_scale}, the lattice geometry
\texttt{eqf\_supNorm\_latPt\_sub}, \texttt{eqc\_region}).
\item \emph{The theorem.} \texttt{Spin2.thm\_limit\_equi\_r}: for $L\ge L_0$ there are $M$, $g_0$, $N_0$, $g(r),\mu(r)>0$ such
that for $1/\beta_R\le\min(g_0,g(r))$ and every admissible sequence with $m_N^2L^{2N}\le\mu(r)$ eventually, a subsequence
of $\vartheta_\epsilon^{-2}P_\epsilon$ converges uniformly on $\{|x|\ge r\}$ (limit uniformly continuous, $\ge0$, bounded,
invariant, short-distance law). \texttt{thm\_limit\_equi\_rAnn}: the same on compacts of $\{|x|>r_0\}$ when the smallness
holds for all $r>r_0$. No hypotheses beyond the stated ranges.
\item \emph{Difference to the statement above.} The Lean $g(r)$, $\mu(r)$ are explicit but only their positivity is
exported; the $\Theta(r)^{-2}$, $\exp(-C\Theta(r)^2)$ forms are not checked in Lean. The modulus is not exported in the
$\omega_r(\eta)$ form; \texttt{ClEquiOn} is proved directly.
\end{itemize}

\begin{conjecture}[Short distances]\label{thm:short}
As $|x|\to0$,
\[
  G^{\rm ren}_{\beta_R}(x)=\Bigl(\frac{\beta(|x|)}{\beta_R}\Bigr)^{2/3}\bigl(1+O(\beta(|x|)^{-1})\bigr)
  =\beta_R^{-2/3}\Bigl(\frac{3}{2\pi}\log\frac1{|x|}\Bigr)^{2/3}\Bigl(1+O\Bigl(\frac{\log\log(1/|x|)}{\log(1/|x|)}\Bigr)\Bigr).
\]
Here $\beta(\ell)$ is the continuum running coupling with $\beta(1)=\beta_R$.
\end{conjecture}

\paragraph{Lean status (finite-$\epsilon$ form; \texttt{H02UV/Spin2/Cont/ContShort.lean}).}
\texttt{Spin2.thm\_short\_eps} (constants $M,g_0,K,K_c,N_0$ depending on $L$ only; $K$ the constant of
\texttt{thm\_multiscale\_final\_h0}): for $N\ge N_0$, $4m^2L^{2N}\le1/\beta_R\le g_0$, every $b$ in the
bracket, $x\ne0$ with $j_x\le N-2$, $\ell=|x|/L^N$, $\lambda=\log(1/\ell)$, $j=\lfloor N+\log_L\ell\rfloor$:
(A) $|\vartheta_\epsilon^{-2}P_\epsilon(x)-(\beta_j/\beta_R)^{2/3}|\le K(\beta_j/\beta_R)^{2/3}/\beta_j$
and, for $\lambda\ge0$, $|\beta_j-\beta_1\lambda|\le K_c(1+\beta_R+\log(1+\lambda))$;
(B) if $\lambda\ge0$ and $2K_c(1+\beta_R+\log(1+\lambda))\le\beta_1\lambda$, then
\[
\Bigl|\vartheta_\epsilon^{-2}P_\epsilon(x)-\Bigl(\frac{\beta_1\lambda}{\beta_R}\Bigr)^{2/3}\Bigr|\le
\Bigl(\frac{\beta_1\lambda}{\beta_R}\Bigr)^{2/3}\frac{K_c(1+\beta_R+\log(1+\lambda))+2K}{\beta_1\lambda},
\]
uniformly in $N\ge N_0$. \texttt{Spin2.thm\_short\_eps\_log}: for fixed $\beta_R$ there are
$\delta_0\in(0,1)$, $K'\ge0$ (depending on $\beta_R,L$) with
$|\vartheta_\epsilon^{-2}P_\epsilon(x)-(\beta_1\lambda/\beta_R)^{2/3}|\le K'(\beta_1\lambda/\beta_R)^{2/3}\log\lambda/\lambda$
for $\ell\le\delta_0$ (all such sample points, \texttt{CsPt}); and if
$\vartheta_\epsilon^{-2}P_\epsilon(x_k)\to G$ along samples with $\ell_k\to\ell\in(0,\delta_0)$, then $G$
satisfies the same bound at $\ell$ (every subsequential limit; existence of $G^{\rm ren}$ is not
claimed). The range is that of \texttt{thm\_multiscale\_final\_h0}: $j_x\le N-2$.

\paragraph{Lean status (all $x\ne0$; \texttt{H02UV/Spin2/Cont/ContShortX.lean}).} \texttt{Spin2.thm\_short\_epsX}: same constants
and prefix, every $b$ in the bracket and \emph{every} $x\ne0$ (from \texttt{thm\_multiscale\_final\_h0\_allB}), with
$\lambda=\log(1/\ell)>0$ automatically ($\ell=|x|/L^N<1$, \texttt{cx\_ell\_lt\_one}: $|x|\le\sqrt2|x|_\infty\le L^N/\sqrt2$) and
$k=\min(j,N-2)$: (A) $|\vartheta_\epsilon^{-2}P_\epsilon(x)-(\beta_k/\beta_R)^{2/3}|\le K(\beta_k/\beta_R)^{2/3}/\beta_k$ and
$|\beta_k-\beta_1\lambda|\le K_c(1+\beta_R+\log(1+\lambda))$ (for $j\ge N-1$: $\beta_1\lambda\le\beta_1\log L$,
$\beta_{N-2}\le5\beta_R$, \texttt{cx\_jphys\_top}); (B) the $\beta_1\lambda$ form as above; (C)
$\vartheta_\epsilon^{-2}P_\epsilon(x)\le(1+|K|)\bigl((\beta_1\lambda+K_c(1+\beta_R+\log(1+\lambda)))/\beta_R\bigr)^{2/3}$.
\texttt{Spin2.thm\_short\_eps\_logX}: \texttt{thm\_short\_eps\_log} for all sample points with $x\ne0$ (\texttt{CsPtX}, no
condition on $j_x$). Limit points: every limit $G$ from \texttt{thm\_limit\_of\_equi} (given asymptotic equicontinuity) satisfies this law for
$0<|x|<\delta_0$, with $\ell=|x|$ the Euclidean torus norm (\texttt{torusNorm}; the samples $\ell_N=|\lfloor L^Nx\rceil|/L^N\to|x|$).

\begin{conjecture}[Macroscopic form]\label{thm:macro}
For $|x|\ge\delta>0$,
$G^{\rm ren}_{\beta_R}(x)=1+\frac2{\beta_R}G_\T(x)+O(\beta_R^{-2})$ up to normalisation. Here
$G_\T$ is the Green function of $-\Delta$ on $\T$ without the zero mode. This is the
continuum version of $\theta^2+2\theta\lambda b_0W(x)$ in \texttt{cor:conn4pt}, and it matches
Expected result~\ref{thm:short} in the overlap region. Precisely: every subsequential limit $G$ of
Theorem~\ref{thm:limit-fixed} satisfies, for $|x|,|y|\ge\delta$,
\[
  \Bigl|G(x)-G(y)-\frac{2Z}{\beta_R}\bigl(G_\T(x)-G_\T(y)\bigr)\Bigr|\le\frac{C(\delta,L)}{\beta_R^2},
\]
where $Z=\lim Z_\epsilon$ along the subsequence and $Z_\epsilon=(\beta_{N-2}/\beta_R)^{2/3}\in[4^{-2/3},5^{2/3}]$.
At finite $\epsilon$ the statement is
$|\vartheta_\epsilon^{-2}(P_\epsilon(x)-P_\epsilon(y))-\frac{2Z_\epsilon}{\beta_R}(G_{\rm lat}(x)-G_{\rm lat}(y))|\le C(\delta)\beta_R^{-2}$.
\end{conjecture}

\begin{theorem}[Macroscopic form at order zero; Lean \texttt{Spin2.thm\_macro\_eps0}, \texttt{thm\_macro\_lim0}]\label{thm:macro0}
For $L\ge L_0$ there are $g_0>0$ and $N_0$, and for every $\delta>0$ a constant $C(\delta,L)$, such that the following holds.
\begin{itemize}
\item[(a)] For $N\ge N_0$, $4m^2L^{2N}\le1/\beta_R\le g_0$ and every $b$ in the bracket, the quotient
$\beta_{N-2}/\beta_R$ lies in $[\frac14,5]$, and every $x\ne0$ with $|x|/L^N\ge\delta$ satisfies
\[
  \bigl|\vartheta_\epsilon^{-2}P_\epsilon(x)-Z_\epsilon\bigr|\le\frac{C(\delta,L)}{\beta_R},\qquad Z_\epsilon=(\beta_{N-2}/\beta_R)^{2/3}.
\]
\item[(b)] Let $1/\beta_R\le g_0$, let the sequence be admissible, and let $\vartheta_\epsilon^{-2}P_\epsilon(\lfloor L^Nx\rceil)$
converge along a subsequence at $x$ and at $y$. If $\|x\|,\|y\|\ge\delta$, the limits satisfy
$|G(x)-G(y)|\le C(\delta,L)/\beta_R$ and $G(x)\ge4^{-2/3}-C(\delta,L)/\beta_R$. In particular this holds for every limit
of Theorem~\ref{thm:limit-fixed}, which \texttt{thm\_macro\_lim0} restates on the same master construction.
\end{itemize}
\end{theorem}

\begin{proof}
Use (A) of Theorem~\ref{thm:short} in the form \texttt{thm\_short\_epsX}. It gives
$|\vartheta_\epsilon^{-2}P_\epsilon(x)-X^{2/3}|\le KX^{2/3}/\beta_k$ with $X=\beta_k/\beta_R$ and $k=\min(j(|x|),N-2)$.
From $|x|/L^N\ge\delta$ we get $N-2-k\le|\log_L\delta|$. Summing the per-step band of Theorem~\ref{thm:beta} from $k$ to
$N-2$, with $\check g\le1$, gives $|\beta_k-\beta_{N-2}|\le(N-2-k)(\beta_L+C_{gu})+4c_1$, which has no $\beta_R$ in it.
On the bracket, $\check g_j\le4/\beta_R$ for $j\le N-2$ and $\check g_{N-2}\ge1/(5\beta_R)$. Hence
$\beta_k\ge\beta_R/4$, $Y=\beta_{N-2}/\beta_R\in[\frac14,5]$ and $|X-Y|\le\Delta(\delta)/\beta_R$. Then
$|X^{2/3}-Y^{2/3}|\le Y^{-1/3}|X-Y|\le2|X-Y|$ and $KX^{2/3}/\beta_k\le4|K|(1+X)/\beta_R$. This proves (a).
For (b), apply (a) at $\lfloor L^Nx\rceil$ and $\lfloor L^Ny\rceil$, which have physical distance $\ge\delta/2$ for large
$N$, and pass to the limit.
\end{proof}

\paragraph{Lean status.} \texttt{H02UV/Spin2/Macro/MacroA.lean}: \texttt{MasterFacts.mac\_beta\_diff} (the summed band),
\texttt{mac\_gap}, \texttt{mac\_eps} (part (a) for a given master construction). \texttt{H02UV/Spin2/Macro/MacroB.lean}:
\texttt{thm\_macro\_eps0} (part (a)) and \texttt{thm\_macro\_lim0} (the conclusion of \texttt{thm\_limit\_fixed\_betaR}
together with (b), with the same $M,g_0,N_0$; $C$ depends on $\delta$ and $L$ only). The first-order term
$\frac{2Z}{\beta_R}(G_\T(x)-G_\T(y))$ of Conjecture~\ref{thm:macro} is not proved; the missing estimates are listed in Section~\ref{sec:macro}.

\paragraph{Lean status (the continuum Green function $G_\T$; library \texttt{Cont}).}
In \texttt{Cont/TorusGreen.lean}, $\T=$ \texttt{UnitAddTorus (Fin 2)}. The heat kernel is
$p_t=\Theta_{4\pi t}$ with $\Theta_s(x)=\prod_i\texttt{cosKernel}(x_i,s)
=\prod_i\sum_{n}\cos(2\pi nx_i)e^{-\pi n^2s}$ (Mathlib). The Poisson form
$\Theta_s=s^{-1}\sum_{n\in\Z^2}e^{-\pi|x+n|^2/s}$ comes from \texttt{evenKernel\_functional\_equation}
(\texttt{b1\_cosK\_eq\_evenK}). Define
$G_\T(x)=\int_0^\infty(p_t(x)-1)\,dt$ (\texttt{torusGreen}, \texttt{torusGreen\_eq\_heat}).
The integral converges for $x\ne0$ (\texttt{torusGreen\_integrableOn}), and $G_\T$ is continuous on
$\T\setminus\{0\}$ (\texttt{torusGreen\_continuousOn}). It is invariant under $x\mapsto-x$, under
coordinate permutations and under reflections (\texttt{torusGreen\_neg}, \texttt{\_perm},
\texttt{\_swap}, \texttt{\_reflect}). It has the logarithmic singularity
$G_\T(x)+\frac1{2\pi}\log|x|_\T\to c$ as $x\to0$ (\texttt{torusGreen\_log\_singularity}, with
$|x|_\T$ the torus distance). The proof splits at $s=1$: for $s\ge1$, $|\Theta_s-1|\le Ce^{-\pi s}$;
for $s\le1$ and $|x_i|\le\frac34$, $0\le\Theta_s-s^{-1}e^{-\pi|x|^2/s}\le M$. It then uses
$\int_0^1e^{-\pi r^2/s}\frac{ds}s=-\log r^2+c_0+o(1)$ (\texttt{Cont/TorusLog.lean}).
On the lattice, the zero-mode-free Green function is $G_{\rm lat}=$ \texttt{fmat} of
$\lambda(k)^{-1}\mathbf 1_{k\ne0}$, with $(-\Delta)G_{\rm lat}=1-Q$, $G_{\rm lat}Q=0$ and
$\sum_xG_{\rm lat}(x)=0$ (\texttt{Cont/LatGreen.lean}).
The scaling convergence is proved (\texttt{Cont/LatScale.lean}). For $L\ge2$,
$G_{\Lambda_N}(\lfloor L^Nx\rceil)\to G_\T(x)$ uniformly on compact subsets of $\T\setminus\{0\}$
(\texttt{latGreen\_latPtN\_tendstoUniformlyOn}). Here $\lfloor L^Nx\rceil$ is the nearest lattice
point (\texttt{latPtN}). No rescaling is needed, because in $d=2$ the unit-spacing lattice Green
function is dimensionless. The underlying statement holds for every $n$-point torus: for
$\delta,\eta>0$ there is $n_0$ with $|G^{(n)}_{\rm lat}(\lfloor nx\rceil)-G_\T(x)|\le\eta$ for
$n\ge n_0$ and $|x|_\T\ge\delta$ (\texttt{latGreen\_latPt\_uniform}). Proof: with
$a(k)=n^2\lambda(k)/4\pi$, write $4\pi G_{\rm lat}=\int_0^\varepsilon(P_n(s)-1)\,ds+\sum_{k\ne0}
e^{-\varepsilon a(k)}\cos(k\cdot x)/a(k)$, where $P_n(s)=n^2p_{n^2s/4\pi}$ and $p_t=e^{t\Delta}(\cdot,0)$. Split $4\pi G_\T$ the same way, using the Fourier
series of $\Theta_s$. The long-time series converge uniformly in $x$ by Tannery's theorem, with
majorant $3e^{-4\varepsilon|k|^2/\pi}$ coming from $a(k)\ge4|k|^2/\pi$. For the short-time parts,
the off-diagonal heat-kernel bound is replaced by a second-moment bound in
\texttt{Cont/LatHeat.lean}. It uses $p_t\ge0$, $p_t(x)\le p_t(0)\le(C(\frac1n+\frac1{\sqrt t}))^2$
and $\sum_y\lambda(y)p_s(y)\le32\pi^2s/n^2$, which give
$\lambda(x)p_{2s}(x)\le C'(s/n^4+1/n^2)$. Hence
$|\int_0^\varepsilon(P_n-1)|=O(\varepsilon/\delta^2)$ uniformly in $n$.
\emph{Not yet formalised:} the weak identity $-\Delta G_\T=\delta-1$.

\begin{corollary}[Scaling window]\label{cor:window}
Let $\beta_\epsilon=A\log\frac1\epsilon$.
\begin{itemize}
\item[(a)] If $A>\frac3{2\pi}$, then $P_\epsilon(x)\to(1-\frac3{2\pi A})^{2/3}>0$ for fixed $x\ne0$.
  (Proved in Lean, \texttt{Spin2.thm\_cor\_window\_a}, \texttt{\_aX}.)
\item[(b)] If $A=\frac3{2\pi}$, Expected results~\ref{thm:limit}--\ref{thm:short} describe the scaling
  limit. (Remark; not a formal statement.)
\item[(c)] If $A<\frac3{2\pi}$, the flow leaves the perturbative regime, and this case is not covered.
  (Remark; not a formal statement.)
\end{itemize}
\end{corollary}

\paragraph{Lean (part (a)).} \texttt{Spin2.thm\_cor\_window\_a} (\texttt{H02UV/Spin2/Cont/ContWindow.lean},
pure limits in \texttt{H02UV/Spin2/Cont/ContWindowPure.lean}): there is $L_0$ such that for all $L\ge L_0$,
$A>\frac3{2\pi}$ and points $x_N\in\Lambda_N$ with $|x_N|/L^N\to\ell\in(0,\frac1{2L})$,
$P_{A\log L^N}[0\leftrightarrow x_N]\to(1-\frac3{2\pi A})^{2/3}$ ($\epsilon=L^{-N}$; $\frac3{2\pi}=\beta_1$,
the Lean step constant is $\beta_1\log L$). No other hypotheses. Proof: with $b_N=1/(A\log L^N)$ put
$\beta_R=1/b_N-(N-2)\beta_1\log L\approx(A-\frac3{2\pi})N\log L$ and $m^2=1/(4\beta_RL^{2N})$. By
Theorem~\ref{thm:beta} ($1/b_{\rm fin}(1/\beta')=\beta'+(N-2)\beta_1\log L+O(\log N)$ at $\beta'=2\beta_R,\beta_R/2$),
$b_N$ lies in the bracket $[b_{\rm fin}(1/(2\beta_R)),b_{\rm fin}(2/\beta_R)]$, where
Theorem~\ref{thm:multiscale} at $h=0$ holds for every $b$ (\texttt{fa\_main}). Along the $h=0$
trajectory from $b_N$ (in $D_\epsilon$), $|\beta_j-\beta_\epsilon+j\beta_1\log L|\le C_1+C_2\log\beta_\epsilon$
(\texttt{BetaAsymD}); with $j=j(\ell_N)=N+O(1)$ this gives $\beta_j/\beta_\epsilon\to1-\frac3{2\pi A}$ and
$K/\beta_j\to0$.

\emph{Every $\ell>0$} (\texttt{Spin2.thm\_cor\_window\_aX}, \texttt{H02UV/Spin2/Cont/ContWindowX.lean}): the same conclusion for
$|x_N|/L^N\to\ell>0$ with no upper bound on $\ell$ (realisable values: $\ell\le\sqrt2/2$). Here the bound of
\texttt{thm\_multiscale\_final\_h0\_allB} at $k=\min(j(\ell_N),N-2)$ is used; $|k-N|\le2+|\log_L\ell_N|$ still gives $k/N\to1$, so
$\beta_k/\beta_\epsilon\to1-\frac3{2\pi A}$ (at $k=N-2$: $\beta_{N-2}/\beta_\epsilon=\beta_R/\beta_\epsilon+O(1/N)\to1-\frac3{2\pi A}$)
(\texttt{a5w\_perNX}, \texttt{a5w\_limitX}).

\subsection{Heuristics}
\label{sec:heur}
After the substitution $\psi=\sqrt\beta\,\eta$, $\bar\psi=\sqrt\beta\,\xi$, the field has
dimension $[\psi]=0$ in $d=2$. So the mass $\psi\bar\psi$ is relevant, the gradient terms and the
quartic term $M=\psi\bar\psi(\nabla\psi)(\nabla\bar\psi)$ are marginal, and nothing else is
relevant or marginal (Section~\ref{sec:loc}).

The quartic coupling generates a field strength renormalisation $Z_j$ at first order. A bubble
renormalises the quartic coupling at second order. With $\beta_j=Z_j^2/b_j$, this gives
\begin{equation}\label{e:beta}
  \beta_{j+1}=\beta_j-\beta_1\log L+\frac{\beta_2}{\beta_j}\log L+O(\beta_j^{-2}),\qquad
  \beta_1=\frac{3}{2\pi}=\frac{2}{2\pi}+\frac{1}{2\pi}.
\end{equation}
Here $\frac2{2\pi}$ is the field strength contribution and $\frac1{2\pi}$ the bubble
contribution. Both are computed in Section~\ref{sec:pt-flow}.

The spin field has anomalous dimension $\gamma=\frac1{2\pi\beta}$ per unit of $\log$ length
(Section~\ref{sec:spin}). So the spin renormalisation is
$\vartheta_j^2=\exp(-\sum_k2\gamma_k\log L)=(\beta_j/\beta_\epsilon)^{2/3}$. The exponent is
$2\gamma\beta/\beta_1=\frac23$, which is the $O(n)$ value $\frac{n-1}{n-2}$ at $n=-1$.

The two-loop coefficient $\beta_2=\frac3{4\pi^2}$ is taken from the $O(n)$ model at $n=-1$ and
has \emph{not} been computed here. It only enters the $\log\log$ corrections, through
\[
  \beta_\epsilon=\beta_R+\beta_1\log\tfrac1\epsilon-\tfrac{\beta_2}{\beta_1}\log\tfrac{\beta_\epsilon}{\beta_R}+O(1).
\]

\subsection{Organisation and relation to \hr}
\label{sec:org}
The paper is organised like \hr{}, and identical parts are not repeated. The table lists the
correspondence.

\medskip
\begin{center}\small
\begin{tabular}{l|p{5.2cm}|p{5.2cm}}
here & \hr{} & change\\\hline
Section~\ref{sec:ward} & \texttt{sec:ward}, \texttt{sec:finvol} & $s_0=0$, tuning of $a_0$ only, finite $N$\\
Section~\ref{sec:frd} & \texttt{sec:frd} & dressed decomposition (\frd)\\
Section~\ref{sec:grassmann} & \texttt{sec:grassm-gauss-integr}, \texttt{sec:symmetries}, \texttt{sec:polymers}, \texttt{sec:norms}, \texttt{sec:Loc} & $d=2$ power counting, field scale $\ell_0$\\
Section~\ref{sec:pt} & --- & second-order map with finite-range remainder, coefficient formulas\\
Section~\ref{sec:rgmap} & \texttt{sec:step-def} & absorption step, second-order map\\
Section~\ref{sec:est} & \texttt{sec:step} & bound $O(\bar g^3)$\\
Section~\ref{sec:flow} & \texttt{sec:flow} & running marginal coupling, finite horizon\\
Section~\ref{sec:obs-flow} & \texttt{sec:obs-flow} & anomalous dimension of $\bar\psi\psi$\\
Section~\ref{sec:corr} & \texttt{sec:suscept}, \texttt{sec:obs} & multiscale factors\\
Section~\ref{sec:uv} & \texttt{prop:consistency} & comparison of $\T_\epsilon$ and $\T_{\epsilon/L}$
\end{tabular}
\end{center}
\medskip

\emph{Notation.} Here $Z_j$ is the field strength, as in \frd{}. The progressive integral called
$Z_j$ in \hr{} is denoted $\cZ_j$. The field strength increments are $\zeta_j$, with
$Z_{j+1}=Z_j(1+\zeta_j)$.

%======================================================================
\section{Fermionic formulation and Ward identities}
\label{sec:ward}

\subsection{Initial potential}
As in \hr{}, equation \texttt{e:quartic-rescale}, but with $s_0=0$, set
$\xi=\beta^{-1/2}\bar\psi$ and $\eta=\beta^{-1/2}\psi$. With a Gaussian mass $m^2>0$,
\begin{equation}\label{e:V0}
  \avg{F}_{\beta,h}=\frac{\E_C[e^{-V_0(\Lambda)}F]}{\E_C[e^{-V_0(\Lambda)}]},\qquad
  C=(-\Delta+m^2)^{-1},
\end{equation}
with
\begin{equation}\label{e:V0b}
  V_0=a_0\psi\bar\psi+b_0M,\qquad b_0=\frac1\beta,\qquad a_0=\frac{1+h}{\beta}-m^2 .
\end{equation}
There is no gradient coupling $s_0$ to tune: gradient terms are absorbed into the Gaussian at
every scale (Section~\ref{sec:rgmap}). The only tuned parameter is $a_0$.

\subsection{Finite volume and the Ward identity}
As in \texttt{sec:finvol}, $N$ is fixed and $m^2\downarrow0$. For $b_0\in(0,\bar b]$ and
$m^2\in(0,L^{-2(N-1)}]$, Section~\ref{sec:flow} gives a tuned value $a_0=a_0^c(b_0,m^2;N)$,
continuous in $m^2\ge0$, such that the flow stays perturbative. Define $(\beta,h)$ by
\eqref{e:V0b} with $a_0=a_0^c$. In \hr{} the argument of \texttt{cor:ward} gives $h(b_0,m^2)=m^2\beta(1+O(b_0))$:
the Ward identity $\avg{z_0}_{\beta,h}=h\sum_x\avg{\xi_0\eta_x}_{\beta,h}$ with $\chi\ge1/(2m^2)$. This needs
$m^2|\Lambda|\gg1$. \emph{Correction (fixed $N$, $m^2\downarrow0$): it fails.} Here $\chi\le|\Lambda|/b$, and in the window
$m^2\le c_0L^{-2N}$ the renormalisation group determines $\chi$ only up to $O(|\Lambda|)$ (through $|m^2Z_Nt_{\rm raw}-1|\le K_tm^2L^{2N}$).
The Ward identities hold for every $a_0$, the RG bounds are stable when the final condition $\mu_{N_h}=0$ is moved inside its
window, and $h$ is strictly monotone in that condition. So $h\downarrow0$ does not follow. The limit $h_0(b_0;N)$ is a finite-volume
shift of the critical point, set by the value of the final zero-mode coupling $|\Lambda|\tilde a_{N,N}$ (Log: ``$h\downarrow0$ by Ward identities'').
The difference from \hr{} is that $a_0^c$ depends on $N$: there is no infinite-volume limit, only a finite horizon.

\emph{Route (a1): the point $h=0$ in a window.} Instead of $h\downarrow0$, fix $h=0$ ($\mu_0=b-m^2$) and replace the final
condition $\mu_{N_h}=0$ by a window $|\mu_{N_h}|\le c_1\le\theta\check g_{N_h}$, $\theta=\frac34c_\mu$. With the sharp step-$(N-2)$
mass bound ($c_\nu$ without the $c_\mu$-term at order one, \texttt{glue\_nu\_sharp}) the $h=0$ point fits such a window (Log, size check).
In Lean: the shifted target $\mu_{N}=c'$ (\texttt{AdmW}); strictness, uniqueness and the $\check g$-cone for a common target; existence by
continuation in the target at fixed $b$ with the $\check g$-part of the signed cone, $|\delta\check g_N|\le p\,\delta\mu_N$
(\texttt{win\_interval}: initial data $[\mu_-,\mu_+]$ with targets $\mp c_1$, all in $D_\varepsilon$). Since $h=(\mu_0+m^2)/b-1$ is
affine, $b-m^2\in(\mu_-,\mu_+]$ once $h(\mu_+)\ge0$ and $h<0$ somewhere in the window. The Ward identity holds at either sign of $h$
given $\avg{Z_0}\ne0$ (\texttt{ward2\_of\_Zne}); $h(\mu_+)\ge0$ follows from $r\ge0$, $\chi>0$ at $\mu_+$ (named RG one-point input
\texttt{H0OnePtUp}: $r\approx M_N(1-q\check g_N/d)$), and $h<0$ somewhere follows from positivity ($h\ge0\Rightarrow0\le\chi\le|\Lambda|/b$)
once the dressed $\chi\approx|\Lambda|(1+k^0)/(Z_Nd)$ leaves $[0,|\Lambda|/b]$ in the window (named input \texttt{H0SusceptExit}).
Result \texttt{thm\_h0\_window} (one master construction). Along the $h=0$ family $b\mapsto\Re\check g_N$ is continuous, so
$\Re\check g_N=1/\beta_R$ is attained (\texttt{exists\_b\_h0}); uniqueness is open, since the $\check g$-cone is not forward invariant
without a common final condition.
Integration (\texttt{H02UV/Spin2/Master/Master}, \texttt{H02UV/Spin2/H0/H0WinConcl}, \texttt{MasterH0OnePt}, \texttt{MasterH0Shoot}): the master construction \texttt{MasterCore}
carries the windowed-flow hypotheses, the last two steps along every trajectory in $D_\varepsilon$ (\texttt{GlueD}) and the
trajectory clauses of \texttt{SpinConcl} in $D_\varepsilon$ (the existing proofs use only the domain, not $\mu_{N_h}=0$; they were
generalised in place). Hence \texttt{MagConcl}, \texttt{MultiConcl}, the zero-mode integration and the dressed susceptibility
$\chi=m^{-2}-m^{-4}X$ with the Ward identity hold along windowed trajectories at either sign of $h$
(\texttt{MasterCore.onePtWin}). The named inputs reduce exactly to the nonvanishing of the factors of $\avg{Z_0}$, $r\ge0$
and $X<m^2$ at $\mu_+$, and $X\notin[m^2-m^4|\Lambda|/b,m^2]$ somewhere in the window (\texttt{h0OnePtUp\_of\_X},
\texttt{h0SusceptExit\_of\_X}). Bracketing of $\Re\check g_N$ along the $h=0$ family by the admissible shooting values
(\texttt{MasterCore.h0Shoot}) gives existence of $b$ with $\Re\check g_N(b,b-m^2)=1/\beta_R$.
Windowed pair layer: the \texttt{MasterFacts} pair conclusions (\texttt{LoopConcl}, \texttt{RelConcl2G}, \texttt{WKConcl},
\texttt{GamBulkConcl}, \texttt{TopConcl}, the last-step facts, \texttt{PairInitConclW}, \texttt{LamConcl}, the $V$-step, the seed unit,
the joint sizes) and \texttt{thm\_pairAbove\_final} were moved to $D_\varepsilon$ forms (\texttt{H02UV/Spin2/Master/MasterAllFactsD}); hence (H$_\lambda$)
and \texttt{TwoPtSq} on the $h=0$ trajectory follow from the pair inputs. The chain factors $\prod_jc_j$, $\det_{N-1}$ of $\avg{Z_0}$
never vanish (Fourier symbol of $1+a_jA\Gamma_j$, \texttt{chainFactors\_ne}), so the window inputs reduce to $U=1+\tilde u\ne0$, $r\ge0$
and the position of $X$ at the window ends (\texttt{h0OnePtUp\_of\_U}, \texttt{h0SusceptExit\_of\_U}); since $U\propto d=\hat a_N+(1+k^0)m^2|\Lambda|$,
these need the quantitative last-two-step mass flow at $h=0$.
Mass flow (\texttt{H02UV/Step2/Flow/FlowWin2}, \texttt{H02UV/PtCorr2/Zero/WardNeg2}, \texttt{H02UV/Spin2/H0/H0MassFlow}, \texttt{MasterH0Ends}): the sharp
relation $|L^{2N}\nu_N/Z_N-L^4\mu_{N_h}|\le c_{\nu,\mathrm{last}}\check g_{N_h}$ through both last steps (\texttt{glue\_hist\_nu}),
$Z_N>0$ real, $|m^2Z_Nt_{\rm raw}-1|\le K_Tm^2|\Lambda|$ and the dressed susceptibility with the scale-dependent zero-mode
errors ($|e_j|\le Bm^2L^{2j}$, geometric sum, \texttt{suscept2\_dressed\_geo}) give, with $n=m^2|\Lambda|U=m^2|\Lambda|(1+k^0)+\hat a_Nw$,
$\chi=|\Lambda|p/(Z_Nn)+O(|\Lambda|\rho B(\rho+|\hat a_N|/|n|)/(L^2Z_N))$ and $\Re n=L^4\mu_{N_h}+O(\text{slack})$. If
$L^4c_1$ exceeds the explicit slack ($\approx(c_{\nu,\mathrm{last}}+2)\check g$ for $m^2|\Lambda|\le\check g$; $L^{-4}c_{\nu,\mathrm{last}}\approx1.09\ell^2<\frac34c_\mu\approx1.47\ell^2$),
then $U\ne0$ at both ends, $\Re\chi>0$ at $\mu_+$ and $\Re\chi<0$ at $\mu_-$ (\texttt{MasterCore.h0Mass}). Hence
\texttt{H0SusceptExit} is a theorem and \texttt{H0OnePtUp} follows from $r\ge0$ at $\mu_+$ alone (\texttt{MasterCore.h0Ends});
\texttt{thm\_multiscale\_h0E} is \texttt{thm\_multiscale\_h0} with these window inputs discharged (inputs: $r\ge0$ at $\mu_+$, the
slack inequality, the pair inputs). The slack inequality is discharged in \texttt{thm\_multiscale\_h0F}
(\texttt{H02UV/Spin2/H0/MasterH0Slack.lean}): at $\theta=\frac34c_\mu$, $c_1=\theta\check g^c$, the slack is monotone in $\check g$ and $m^2|\Lambda|$,
$\le(c_{\nu,\mathrm{last}}+1+X\check g)\check g$ for $m^2|\Lambda|\le\check g\le1$, the window condition gives $2p_W\theta<\frac14$, and
$5(c_{\nu,\mathrm{last}}+2)<3L^4c_\mu$ holds at the master $c_m=(48L^2+8)\ell^2/49+1$ for $L\ge49$, $c_\mu\varepsilon\le\frac1{16}$; the
$\varepsilon$-requirement $\varepsilon\le\frac45$, $\frac54X\varepsilon\le1$, $c_\mu\varepsilon\le\frac1{16}$ ($X=X(L,C_f,\ell,K_T,B,\rho)$) is part of the
master construction. New input: $m^2|\Lambda|\le\check g^c_{N_h}$ (compatible with the mass window, $\check g^c\asymp1/N$).
Input $r\ge0$ at $\mu_+$ (\texttt{H02UV/Spin2/H0/H0OnePt}, \texttt{H0OnePtM}, \texttt{H02UV/PtCorr2/Zero/OnePtWin}): reality needs no RG, since the
Ward identity $h'\chi=r$ with $h'$ real gives $h'\Re\chi=\Re r$, so $\Re\chi>0$, $\Re r\ge0$ force $h\ge0$, where positivity makes $r$ real
(\texttt{onept\_real\_of\_re}). The one-point chain gives exactly, at either sign of $h$, $r=\gamma_{a,N}+n_a\Omega$ (\texttt{onePt\_split\_win}), and
comparing the bulk normalisations of the observable and bulk chains $\mathrm{s}(\Omega)=U\,e^{u^{\rm obs}_N-u_N}$ with
$|u^{\rm obs}_N-u_N|\le L^4C_f\check g^2$ (the observable last step moves $\mathrm{Loc}\,S(\Lambda)$ into $u_N$). With
$|\gamma_{a,N}-\vartheta_{N_h}|\le(F'+c_L)\vartheta\check g$, $|\beta_{N-1}|/|Z_N|\le\frac{32}7e^{1/4}\vartheta\check g$ (sharp vertex bound),
$c2F_a\approx\beta$, $\varepsilon F_a=O(\beta/Z_N)$ (\texttt{LastOneRem}) and $\Re n\ge L^4\mu_{N_h}-S$: $|n_a\Omega|\le22\vartheta\check g/(L^4\mu_{N_h}-S)$.
At $\mu_+$, $L^4c_1-S\ge\frac54L^2\check g^c$ ($5(c_{\nu,\rm last}+2)+5L^2\le3L^4c_\mu$), so $\Re r\ge0$ (\texttt{MasterFacts.onePt\_re\_win},
\texttt{MasterFacts.h0EndsG}). \texttt{thm\_multiscale\_h0G} is \texttt{thm\_multiscale\_h0F} without the input $r\ge0$; in its place only
$O(\check g^c)$, $O(m^2|\Lambda|)$ smallness conditions and $N\ge3$ (the remainders of the two top steps come from
\texttt{TopConcl} and \texttt{gamRemNP\_le}, no \texttt{GamLastRem}).

\subsection{What the renormalisation group provides}
\label{sec:rgprovides}
In $d\ge3$, the finite-volume result \texttt{thm:psi4-full} has the following consequence.
Write $\kappa=b_0/(1+s_0)$, so that $\xi_0\eta_0=\kappa\,\bar\psi_0\psi_0$.

\begin{lemma}[Perfect square, $d\ge3$]\label{lem:square}
Let $h=0$ and drop the error terms $E_{00},E_{00xx}$ of \texttt{thm:psi4-full}. Then
\[
  \Pb[0\conn x]=\bigl(1-\kappa(\gamma-\lambda W_N(x))\bigr)^2 .
\]
\end{lemma}

\begin{proof}
Write $\tau=t_N|\Lambda_N|^{-1}/(1+\tilde u_{N,N})$. The Ward identity
$\kappa\avg{\bar\psi_0\psi_0}=1$ and \texttt{e:1pt-b-thm} give $\lambda\tau=\kappa^{-1}-\gamma$.
By \eqref{e:P} and \texttt{e:2pt-b-thm},
\[
  \Pb[0\conn x]=1-\kappa^2\bigl(\gamma^2-\lambda^2W^2+2\lambda\tau(\gamma-\lambda W)\bigr).
\]
Insert $\lambda\tau$. The result is $1-2\kappa(\gamma-\lambda W)+\kappa^2(\gamma-\lambda W)^2$.
\end{proof}

This is $(\theta+\kappa\lambda W)^2$ with $\theta=1-\kappa\gamma$, the main term of
\texttt{cor:conn4pt}. The combination $\gamma-\lambda W_N(x)=\sum_k(\lambda C_{k+1}(0,0)-\lambda^2C_{k+1}(0,x))$
receives no contribution from scales $L^k\gg|x|$ at leading order: the two observables have
coalesced there, and $\lambda=1+O(b_0)$.

In $d=2$ the $\bar\psi\psi$ vertex runs: $\hat\lambda_k=\lambda^{(2)}_k/\lambda^{(2)}_0$ grows
like $(\beta_\epsilon/\beta_k)^{1/3}$ (Section~\ref{sec:spin}). So the cancellation above the
coalescence scale is no longer automatic. It goes through the marginal coupling $r$ of the pair
observable (Section~\ref{sec:obs-flow}), which after coalescence multiplies a single composite
$\bar\psi\psi$. It is the renormalisation group form of $z_az_b\to z_a^2=1-2\xi_a\eta_a$ for
$a\to b$. The expected outcome is:

\begin{conjecture}[Renormalisation group input]\label{thm:rginput}
Let $j_x$ be the coalescence scale of $0$ and $x$, and let
\begin{equation}\label{e:PM}
  M_j=1-b_0\sum_{k<j}\hat\lambda_k\,\Gamma^{\rm int}_{k+1}(0,0).
\end{equation}
At $h=0$, $P(x)=M_{j_x}^2(1+O(\check g_{j_x}))$.
\end{conjecture}
\paragraph{Lean form.} \texttt{Spin2.RGInput}: along the trajectory $t$ from $(b,\mu_0)$ at mass $m^2$, for every
$x\ne0$ with $j_x=\texttt{Obs.coal}\,L\,0\,x\le N_h$,
$\bigl|P_{1/b}[0\conn x]-M_{j_x}^2\bigr|\le K_P|M_{j_x}|^2\check g_{j_x}$, with $M_j=\texttt{Mmag}$ as in \eqref{e:PM} for a vertex
sequence $\hat\lambda$ satisfying (H$_\lambda$) (\texttt{HLam}). \emph{Status: discharged on route (a1).} At $h=0$, on the
trajectory $(b,b-m^2)$ of Theorem~\ref{thm:multiscale} (fixed $m^2$ in its window), \texttt{TwoPtSq} with $r=0$, i.e.\ \texttt{RGInput}
(\texttt{rgInput\_of\_twoPtSq\_h0}), holds for $\hat\lambda=\mu_a/\mu_{a,0}$ with (H$_\lambda$) proved: \texttt{MasterFacts.twoPtSqD} with the pair
layer \texttt{fw\_pairAbove6\_MFD'}, used inside \texttt{Spin2.thm\_multiscale\_final\_h0}. It is not proved for general $h$ or along
$h\downarrow0$ (the route described next).

\emph{The field $h$ on the trajectory} (\texttt{H02UV/Spin2/Multiscale/RGInH.lean}, \texttt{H02UV/PtCorr2/Zero/TwoPtRep.lean}). The trajectory from
$(b,\mu_0^c(b))$ carries the model with $h=(\mu_0+m^2)/b-1$ (\texttt{hPar}), in general $h\ne0$; $\mu_0^c$ is fixed by the final
condition $\mu_{N_h}=0$. On the model side: $b\avg{\bar\psi_0\psi_x}=\Pb_{1/b,h}[0\conn x,\,0\not\leftrightarrow g]$
(\texttt{avg0\_two\_eq\_ghost}), which tends to $\Pb_{1/b}[0\conn x]$ as $h\to0$ (\texttt{tendsto\_ghostTwo}); with
$r=1-b\avg{\bar\psi_0\psi_0}=\Pb_{\beta,h}[0\leftrightarrow g]$, the Ward identity and positivity give $r=h\sum_x\Pb_{\beta,h}[0\conn x,\,0\not\leftrightarrow g]$, so
$h(1-r)\le r\le h|\Lambda|$ (\texttt{ward\_r\_bounds}). Thus $h\to0$ iff $r\to0$; the Ward identity alone does not determine $h$.
Exactly, $(\mathrm{rgProdO})_{ab}=\mu_a\mu_b\avg{\bar\psi_a\psi_b}\,(\mathrm{rgProdO})_\varnothing$ (\texttt{twopt\_rg\_repr}). The
two-point input at $h\ge0$ is \texttt{TwoPtSq}:
$|b\avg{\bar\psi_0\psi_x}-M_{j_x}^2|\le K_P|M_{j_x}|^2\check g_{j_x}+K_r|r|$; the shape of the $r$-term is that of Lemma~\ref{lem:square}
without its $h=0$ input, $P=M^2+2r(1-M)$ (\texttt{square\_r}). At $h=0$ it is \texttt{RGInput} (\texttt{rgInput\_of\_twoPtSq\_h0}).
\texttt{multiConclE\_of}, \texttt{multiConclH\_of} and \texttt{multiConclH0\_of} give, with the constants of Theorem~\ref{thm:beta}: on the trajectory,
$|\Pb_{1/b,h}[0\conn x,\,0\not\leftrightarrow g]-\vartheta_j^2|\le K\vartheta_j^2/\beta_j+K_rr$ (\texttt{MultiConclH}); and, at fixed $b$, for admissible
$(m_n^2,\mu_n)$ with $h_n\ge0$, $h_n\to0$, \texttt{TwoPtSq} and either (H$_\lambda$) or, for $\hat\lambda=\mu_a/\mu_{a,0}$, \texttt{Case2Rel} and
\texttt{Case2Loop}: for every $\delta>0$ some $n$ gives $|P_{1/b}[0\conn x]-\vartheta_j(m_n^2)^2|\le K\vartheta_j(m_n^2)^2/\beta_j(m_n^2)+\delta$
(\texttt{MultiConclH0}; for $\hat\lambda=\mu_a/\mu_{a,0}$, (H$_\lambda$) follows from \texttt{Case2Rel} and \texttt{Case2Loop} by \texttt{hLam\_case2}). Open: \texttt{TwoPtSq}, and the existence of such data with $h_n\to0$
(equivalently $b\avg{\bar\psi_0\psi_0}\to1$).

\emph{The $d=2$ square} (\texttt{H02UV/PtCorr2/Zero/TwoPtGhost.lean}, \texttt{H02UV/PtCorr2/Pair/TwoPtZZ.lean}, \texttt{H02UV/Spin2/Multiscale/TwoPtSq.lean}). With
$\tilde z_y=1-b\bar\psi_y\psi_y$: $\avg{z_az_x}_{\beta,h}=\Pb_{\beta,h}[a\leftrightarrow g,\,x\leftrightarrow g]-\Pb_{\beta,h}[a\conn x,\,a\not\leftrightarrow g]$
(\texttt{avg\_zz\_split}), so $b\avg{\bar\psi_0\psi_x}=-\avg{\tilde z_0\tilde z_x}+p$, $0\le p\le r$, and $\avg{\tilde z_x}=\avg{\tilde z_0}=r$
(\texttt{twopt\_eq\_neg\_zz}). $\tilde z_0\tilde z_x$ is the Case~(2) pair observable ($\gamma=1$, $\mu=-b$, no $ab$ couplings). Split the
observable RG product at the coalescence scale, $\Pi=\Pi^{<j_x}\Pi^{\ge j_x}$ (\texttt{rgProdO\_split}). Below $j_x$ the step constants
$e^{-u_{j+1}(\Lambda)}$ have no $\sigma_0\sigma_x$ part (Lemma~\ref{lem:coal}), so $\Pi^{<j_x}=c(1+\sigma_0A)(1+\sigma_xB)$ is rank one
(\texttt{prodLo\_rankOne}) with $A=\gamma_{0,j_x}=-u^{\rm obs}_{j_x,0}$, $B=\gamma_{x,j_x}$ (\texttt{prodLo\_sA\_sB}); $\Pi^{\ge j_x}=\omega_{\Gamma_{j_x}}(Z^{\rm obs}_{j_x})$
(\texttt{hist\_omegaO\_from}); write $E$ for its normalised $\sigma_0\sigma_x$ part. The one-point identity $r=A+A^{\ge}$ (the analogue of
$\kappa(\gamma+\lambda\tau)=1-r$) gives exactly (\texttt{ratio\_split}, \texttt{zz\_square\_split}, \texttt{twopt\_zz\_eq})
\[
  b\avg{\bar\psi_0\psi_x}=AB-r(A+B)-E+p .
\]
Hence \texttt{TwoPtSq} follows (\texttt{twoPtSq\_of\_zz\_flow}, with the constants of \texttt{FlowC2Concl}) from two named inputs:
\texttt{ZBelowU}, $|\gamma_{0,j_x}-M_{j_x}|,|\gamma_{x,j_x}-M_{j_x}|\le K_M|M_{j_x}|\check g_{j_x}$ (the Case~(2) $\gamma$ flow below coalescence;
the second-order terms must be the dressed $\Gamma^{\rm int}$ of \eqref{e:PM}), and \texttt{PairAboveW}, $|E|\le K_E|M_{j_x}|^2\check g_{j_x}+K_Rr$
(the pair part at coalescence: Pauli, the flow of the coalesced composite with the marginal $r$, the last step and the zero mode),
plus the bound $|M_j|\le M_b$. Not proved: \texttt{ZBelowU}, \texttt{PairAboveW}.
\emph{Correction and the $\gamma$ flow} (\texttt{H02UV/Obs2/Flow/GamChain2.lean}, \texttt{H02UV/Spin2/Multiscale/ZBelow*.lean}). At $j_x=0$ ($|x|_\infty<L/2$)
$\gamma_{0,0}=0\ne M_0=1$, so \texttt{ZBelowU} forces $K_M\ge1/b$ (\texttt{zBelowU\_coal\_zero}); the split is used only for
$j_x\ge1$ (\texttt{ZBelowU1}, \texttt{PairAboveW1}), and $j_x=0$ is the named input \texttt{TwoPtNear},
$|b\avg{\bar\psi_0\psi_x}-1|\le K_Nb+K_{RN}r$ (\texttt{twoPtSq\_of\_zz1\_flow}). Exactly (\texttt{gamA\_succ}, \texttt{uA\_succ}),
$\gamma_{0,k+1}=\gamma_{0,k}-\sum_B\epsilon(V_k(B)_0)-\beta_k\Gamma^{\rm int}_{k+1}(0,0)+\rho_k$ with $\beta_k=-\mu_{0,k}$, the step covariance being the
dressed $\Gamma^{\rm int}_{k+1}$ itself (\texttt{histCov\_gamInt}); hence $\gamma_{0,j}=M_j(\lambda^\gamma)$, $\lambda^\gamma_k=\beta_k/b-\rho_k/(b\Gamma^{\rm int}_{k+1}(0,0))$
(\texttt{uA\_eq\_mmag}). Comparing $M_j(\lambda^\gamma)$ and $M_j(\hat\lambda)$ with $\vartheta_j$ (\texttt{MagConcl}, \texttt{mmag\_near\_of}) needs no identification of the
pair vertex with $\hat\lambda$. \texttt{twoPtSq\_of\_gam\_flow}: with the master constants, \texttt{TwoPtSq} follows from (H$_\lambda$) for $\hat\lambda$ and for the
pair vertex (e.g.\ \texttt{Case2Rel}, \texttt{Case2Loop}, \texttt{pairHLam\_of}), \texttt{GamRemRel} ($|\rho_k|\le c_\gamma(\check g_k/b)^{1/3}(\check g_k+a_k)b|\Gamma^{\rm int}_{k+1}(0,0)|$),
\texttt{GamFormB} (the $\sigma_x$ mirror), \texttt{PairAboveW1}, \texttt{TwoPtNear} and the smallness $K(\hat\lambda)\check g_j\le\frac12$; $|M_k|\le M_b$ is derived from
\texttt{BetaAsym} (\texttt{varth\_le\_pure}). Not proved: \texttt{GamRemRel}, \texttt{GamFormB}, \texttt{PairAboveW1}, \texttt{TwoPtNear}.
\emph{The remainder $\rho_k$} (\texttt{H02UV/Obs2/Flow/GamRem.lean}, \texttt{H02UV/Spin2/Gam/GamRem.lean}). $\operatorname{loc}^a$ preserves $\epsilon$ and distributes it over the
blocks with total weight one, so the localisation cancels in $\sum_B$ and, exactly (\texttt{gamRem2\_split\_hist}),
$\rho_k=\sum_B\epsilon(e^{\mathcal L_{C_k}}Q_k(B)_0)+\sum_B\epsilon((\mathbb E\theta\mathcal W_k(B))_0)+\beta_kE_k$,
$E_k=\epsilon(F_{C_k}(\bar\psi_0\psi_0,V_{\kappa'_k}(\Lambda)))$ (bulk; the double contraction of $\bar\psi_0\psi_0$ with the quadratic part of $V_{\kappa'_k}$,
\texttt{eps\_FC\_pbp}). \texttt{gamRemRel\_of}: \texttt{GamRemRel} with $c_\gamma=(c_e+c_w)(1+3K_p)$ follows from (H$_\lambda$) for the pair vertex,
\texttt{GamBulkRel} ($|E_k|\le c_e(\check g_k+a_k)|\Gamma^{\rm int}_{k+1}(0,0)|$, bulk only) and \texttt{GamWKRel} (the $K$, $\mathcal W$ part
$\le c_w(\check g_k+a_k)|\beta_k||\Gamma^{\rm int}_{k+1}(0,0)|$, the $\epsilon$ analogue of \texttt{Case2Rel}); \texttt{twoPtSq\_of\_gam\_flow2} is
\texttt{twoPtSq\_of\_gam\_flow} with these two inputs in place of \texttt{GamRemRel}. Not proved: \texttt{GamBulkRel} (needs bounds on the couplings
$z',y',\nu'$ of $\kappa'_k$ and the gradient bounds of $C_k$), \texttt{GamWKRel}.
\texttt{GamBulkRel} is proved, at every point $a$ (\texttt{H02UV/Obs2/Flow/GamBulk.lean}, \texttt{H02UV/Spin2/Gam/GamBulk.lean}): by Wick and a summation by parts,
$E_k(a)=\nu'_k\sum_xC_k(a,x)C_k(x,a)+(y'_k+z'_k)\mathfrak l_{C_k}(a)$ (\texttt{eps\_FC\_pbp\_Vcoup}), $y'+z'=z_k+y_k-Z_k\zeta_k-g_k(c_{0,k}+\delta_k)$;
with the Gram bounds of $C_k$ at $\ell_k^2=\ell^2/(16|Z_{k+1}|)$ this gives $|E_k|\le K_b\check g_k/|Z_{k+1}|$ for $k\ge1$, and
$|Z_{k+1}\Gamma^{\rm int}_{k+1}(0,0)|\ge1$ from the value of $C'_{k+1}(0,0)$ if $\log L$ is large and $\epsilon$ small (\texttt{gamInt\_lower});
at $k=0$, $C_0=C_1$ has a nonnegative symbol, so $|C_0(x,y)|\le C_0(0,0)$ and $|E_0|\le K_0\check g_0|\Gamma^{\rm int}_1(0,0)|$.
\texttt{thm\_gam\_bulk}: \texttt{GamBulkRelA} along all admissible trajectories under the master constants. \texttt{twoPtSq\_of\_gam\_flow3}
(\texttt{H02UV/Spin2/Gam/GamBulk2.lean}): \texttt{TwoPtSq} from (H$_\lambda$) for $\hat\lambda$ and both pair vertices, \texttt{GamWKRel}, its mirror \texttt{GamWKRelS},
\texttt{PairAboveW1} and the smallness, with \texttt{GamBulkRel}, \texttt{GamFormB}, \texttt{TwoPtNear} discharged.
\texttt{Spin2.thm\_master\_all} (\texttt{MasterFacts.relG}, \texttt{MasterFacts.wk}; \texttt{H02UV/Spin2/Case2/Case2RelPairThm.lean}): under one set of master constants also
\texttt{RelConcl2G} and \texttt{WKConcl} ($|\mathrm{gamWK}_k|\le c_w(\check g_k+a_k)|\beta_k||\Gamma^{\rm int}_{k+1}(0,0)|$ for all
points, all $k<N_h$; from the block bounds of the relative recursion, \texttt{PtCorr2.obs\_hist\_rel2\_base},
$\sum_B\epsilon(e^{\mathcal L_C}Q(B)_a)=\sum_X(\epsilon(K(X)_a)+C(a,a)m_2(K(X)_a))$, $|\epsilon(e^{\mathcal L_C}\mathcal W_a)|\le Q_1\|\mathcal W_a\|$ and
$s_k^2\le\frac98\ell^2|\Gamma^{\rm int}_{k+1}(0,0)|$ from \texttt{gamInt\_lower}), giving \texttt{GamWKRel}, \texttt{GamWKRelS}, \texttt{GamWKTop}
(\texttt{wk\_inputs}). \texttt{Spin2.MasterFacts.twoPtSqD} (\texttt{H02UV/Spin2/H0/MasterH0Multi.lean}): \texttt{TwoPtSq} from (H$_\lambda$) for
$\hat\lambda$, \texttt{PairAboveW1} and the smallness only.
\emph{The pair part above coalescence} (\texttt{H02UV/PtCorr2/Pair/PairAbove.lean}, \texttt{H02UV/Spin2/Pair/PairAbove.lean}). With normalised components
$n_\bullet=s_\bullet/s_\varnothing$, $n_{ab}(XY)=n_{ab}X+n_{ab}Y+n_aXn_bY+n_bXn_aY$ (\texttt{nAB\_mul}). Split $\Pi^{\ge j_x}=X\Omega$ (\texttt{zHi\_eq}) into the constants
$X=\prod_{j_x\le i\le N_h}\mathrm{ch}_i\cdot\det e^{-u_N}$ and the zero-mode integral $\Omega=\omega_{cQ}(F_N)$, $c=t_{\rm raw}/|\Lambda|$. Exactly
(\texttt{zCst\_coup}, \texttt{pair\_E\_split}): $n_{ab}X=q_N+(\gamma_{a,N}-\gamma_{a,j_x})(\gamma_{b,N}-\gamma_{b,j_x})$ (the Pauli-adjusted accumulated pair
constant, $q=-\epsilon(u^{\rm obs}_{ab})$, $q_{j_x}=0$), $E=n_{ab}X+[n_{ab}\Omega+n_aXn_b\Omega+n_bXn_a\Omega]$ and $n_a\Omega=r-\gamma_{a,N}$ (one-point
closure). The zero mode acts as $s_\bullet(\Omega)=\epsilon(F_\bullet)-c\,c_2(F_\bullet)$ (\texttt{sX\_omegaCO\_cstC}); if
$c_2F_{ab}+n_aXc_2F_b+n_bXc_2F_a=\tilde\rho\,c_2F_a$, the bracket is $(\tilde x-\tilde\rho\,\epsilon F_a)/s_\varnothing(\Omega)+\tilde\rho\,n_a\Omega$
(\texttt{zm\_pauli\_elim}), so $c\approx(m^2|\Lambda|)^{-1}$ enters only through $n_a\Omega$. \texttt{pairAboveW1\_of\_parts(\_flow)}: \texttt{PairAboveW1}
with $K_E=K_1+K_2+K_3K_4$, $K_R=K_3M_b$ follows from \texttt{PairCstRel} ($|n_{ab}X|\le K_1|M_{j_x}|^2\check g_{j_x}$), \texttt{PairZeroRel}
($|\tilde\rho|\le K_3|M_{j_x}|\check g_{j_x}$, $|\tilde x-\tilde\rho\,\epsilon F_a|\le K_2|M_{j_x}|^2\check g_{j_x}|s_\varnothing(\Omega)|$) and \texttt{GamTopBd}
($|\gamma_{a,N}|\le K_4|M_{j_x}|$). Heuristically (merged picture, Pauli $(\bar\psi_a\psi_a)^2=0$) the increments of $n_{ab}X$ are $C_k(a,a)\tilde r_k$ with
$\tilde r=r+\mu_a(\gamma_b-\gamma_{b,j_x})+\mu_b(\gamma_a-\gamma_{a,j_x})$, $\tilde r_{j_x}=0$, sourced only by $C(a,a)-C(a,b)$, the $\delta$-terms of
$\operatorname{loc}^{ab}$ and $\chi$, and $\tilde\rho\approx\tilde r_N/\mu_{a,N}$. Not proved: \texttt{PairCstRel}, \texttt{PairZeroRel}, \texttt{GamTopBd}
(see below for the exact $ab$ one-step identity).
\emph{The pair chain} (\texttt{H02UV/Obs2/Pair/PairChain.lean}, \texttt{H02UV/Spin2/Pair/PairCst.lean}). Exactly, for any points and any step with $V_a=\alpha+\beta_a\bar\psi_a\psi_a$,
$V_b=\alpha'+\beta_b\bar\psi_b\psi_b$, no observable constants in $V$ and $C(a,a)=C(b,b)$ (\texttt{pair\_step\_id}): with the Pauli split
$\epsilon(e^{\mathcal L_C}F)=\epsilon F+C(a,a)m_2(F)+\mathfrak d_C(F)$, $\mathfrak d_C(F)=\epsilon(e^{\mathcal L_C}F)-\epsilon(e^{\mathcal L_{C(a,a)J}}F)$
(\texttt{eps\_expLC\_pauli}), the localisation cancellation for $\operatorname{loc}^{ab}$ (\texttt{eps\_total\_LS\_ab}) and
$\epsilon F_C(\bar\psi_a\psi_a,\bar\psi_b\psi_b)=-C(a,b)C(b,a)$ (\texttt{eps\_FC\_pbp\_pbp}),
\[
\Delta q+\Delta\gamma_aA_b+A_a\Delta\gamma_b+\Delta\gamma_a\Delta\gamma_b=-C(a,a)\tilde r+\beta_a\beta_b\bigl(C(a,a)^2-C(a,b)C(b,a)\bigr)-\mathfrak d+\mathrm{Rem},
\qquad \tilde r=r+\beta_aA_b+\beta_bA_a,
\]
$r=\sum_Bm_2(V(B)_{ab})$, $\beta_\bullet=\sum_Bm_2(V(B)_\bullet)=-\mu_\bullet$; Rem collects the $K$, $\mathcal W$ parts, $\epsilon F_C(V'_{ab},V'_\varnothing)$ and the
$\gamma$ remainders. Hence $n_{ab}X=\sum_{j_x\le k\le N}(-C_k(0,0)\tilde r_{j_x,k}+\mathrm{def}_k)$, the last step included (\texttt{nAB\_zCst\_sum}).
The $m_2$ flows (\texttt{betaA\_next}, \texttt{betaB\_next}, \texttt{rAB\_next}; $m_2$ of the $\delta$-terms of $\operatorname{loc}^{ab}$ vanishes,
\texttt{mass2\_Dt}) give exactly $\tilde r_{k+1}=(1-n_{a,k})\tilde r_k+\sigma_k$, $\beta_{a,k+1}=(1-n_{a,k})\beta_{a,k}-e_{a,k}$
(\texttt{rTil2\_succ}, \texttt{betaA\_hist\_succ}), $\tilde r_{j_x}=0$, with the source
$\sigma=\beta_a\beta_b(C(a,b)+C(b,a)-2C(a,a))+\beta_a\beta_bC(a,a)(n_a+n_b)+(n_a-n_b)\beta_bA_a-d_r+\mathfrak d_2-e_{ab}+f+\dots$
(\texttt{rSrc}): only defects, vertex corrections and remainders. A discrete Gronwall bound (\texttt{rel\_gronwall}, \texttt{rTil2\_rel\_bound})
gives $|\tilde r_k|\le(\sum s)\prod(1+\varepsilon)|\beta_{a,k}|$. \texttt{pairCstRel\_of\_parts}, \texttt{pairRTilRel\_of\_src}: \texttt{PairCstRel}
with $K_1=c_\sigma e^{c_\varepsilon}c_s+c_d$ follows from \texttt{PairSrcRel} ($|\sigma_k|\le s_k|\beta_{a,k+1}|$, $|e_{a,k}|\le\varepsilon_k|\beta_{a,k+1}|$,
$\sum s\le c_\sigma|M_{j_x}|\check g_{j_x}$, $\sum\varepsilon\le c_\varepsilon$), \texttt{PairVertSum} ($\sum_k|C_k(0,0)||\beta_{a,k}|\le c_s|M_{j_x}|$) and
\texttt{PairDefSum} ($\sum_k|\mathrm{def}_k|\le c_d|M_{j_x}|^2\check g_{j_x}$). Not proved: these three inputs, \texttt{PairZeroRel}.
\emph{The vertex term; the Pauli-reduced pair observable} (\texttt{H02UV/Obs2/Pair/PairRed.lean}, \texttt{H02UV/Spin2/Pair/PairRed.lean}). The source term
$\beta_a\beta_bC(a,a)(n_a+n_b)$, the terms $\beta_a'\rho_b+\beta_b'\rho_a$ and $e_{ab}$ are not defects. Each has relative size
$\vartheta_k\check g_k^2$ at every $k\ge j_x$, and $\sum_{k\ge j}\vartheta_k\check g_k^2\asymp\vartheta_N\check g_N\gg\vartheta_j\check g_j$
(ratio $(\beta_j/\beta_R)^{2/3}$). In the exact merged dynamics they cancel to all orders: the second-order map misses
$2\beta_a\beta_b(E-n(C+E))$ (\texttt{pauli\_closure\_pure}), which is carried by the third and higher cumulants
$F^{(p)}(V_a,V_b,V_\varnothing,\dots)$ in $K$. These are not Pauli-small, so \texttt{PairSrcRel} cannot be obtained from norm bounds with uniform
constants. The route is reorganised. $E=\omega_{\Gamma_j}(Z_j)_{ab}/\omega_{\Gamma_j}(Z_j)_\varnothing$ is the bulk expectation of the single linear insertion
$V_aV_b=\beta_a\beta_b\bar\psi_a\psi_a\bar\psi_b\psi_b$ (plus activities). It is run as a reduced history without $\sigma_a,\sigma_b$ parts. Then
$r'=(1-n_a)r-d_r+\mathfrak d_2-e_{ab}+f$ and $\Delta q=-C(a,a)r-\mathfrak d+X$, with no vertex and no $\gamma$-cross terms
(\texttt{rAB\_next\_red}, \texttt{pair\_step\_red}). The seed's defects are explicit:
$\epsilon(e^{\mathcal L_C}\bar\psi_a\psi_a\bar\psi_b\psi_b)=C(a,a)C(b,b)-C(a,b)C(b,a)$ and
$m_2(e^{\mathcal L_C}\bar\psi_a\psi_a\bar\psi_b\psi_b)=C(a,a)+C(b,b)-C(a,b)-C(b,a)$, both zero for $C$ constant on $\{a,b\}$
(\texttt{eps\_expLC\_pbp\_mul}, \texttt{mass2\_expLC\_pbp\_mul}, \texttt{pauli\_quartic\_const}). They are $\le2n\ell_k^4L^{-(k+1)}$ and
$\le2n\ell_k^2L^{-(k+1)}$ for $b$ at path distance $n$ (\texttt{norm\_detDef\_le}, \texttt{norm\_m2Def\_le}). Also $n_a=n_b$ along the history
(\texttt{nA2\_eq\_nB2\_hist}).
\emph{The reduced history} (\texttt{H02UV/Obs2/Map/SigmaScale.lean}, \texttt{H02UV/PtCorr2/Obs/ObsHistFrom.lean}, \texttt{PairReduce.lean},
\texttt{PairReduceChain.lean}, \texttt{H02UV/Spin2/Pair/PairRed2.lean}, \texttt{PairRed3.lean}). The elements without $\sigma_a,\sigma_b$ parts (the fixed
points of $\sigma_a\mapsto t\sigma_a$, $\sigma_b\mapsto t^{-1}\sigma_b$) form a subalgebra stable under every operation of the map, including
$\operatorname{loc}$ (\texttt{OAbS.step}). So the observable history started at scale $j$ from such a state $o'$ (\texttt{obsHistFrom}) keeps
none (\texttt{obsAbInv\_histFrom}), and $\omega_{\Gamma_k}(Z_k(o'))=\Pi'^{\ge k}$ (\texttt{hist\_omegaOF\_from}). If $Z_{j_x}(o')$ has the
$\varnothing$ and $\sigma_a\sigma_b$ components of $Z_{j_x}$, then $E=n_{ab}X'+n_{ab}\Omega'$ with the reduced constants $X'$ and zero mode
$\Omega'$, and $s_\varnothing X'=s_\varnothing X$ (\texttt{pair\_red\_ratio}); $n_{ab}X'=q'_N-q'_{j}=\sum_k(-C_k(0,0)r'_k+d'_k)+$ last step
(\texttt{nAB\_cstOF\_eq}, \texttt{qR\_eq\_sum}, \texttt{nAB\_zCstR\_sum}). The $r'$ flow has the same $n_{a,k}$ as the original $\beta_a$ flow
(\texttt{nA2\_fullV\_indep}), so the Gronwall bound is relative to $\beta_{a,k}$ (\texttt{pairRRelR\_of\_src}). \texttt{pairAboveW1\_of\_red\_parts}:
\texttt{PairAboveW1} from \texttt{PairSectRed} (the reduced state, PR1: regrouping the cross products of different polymer factors into
$K'_{ab}$ on connected polymers containing the coalescence region), \texttt{PairRedSeed} ($r'_j=0$), \texttt{PairSrcRelR} (sources
$-d_r+\mathfrak d_2-e_{ab}+f$, no vertex term), \texttt{PairVertSum}, \texttt{PairDefSumR}, \texttt{PairZeroRelR}
($c_2F'_{N,ab}=\tilde\rho\,c_2F_{N,a}$) and \texttt{GamTopBd}. The reduced state is constructed (\texttt{H02UV/Obs2/Pair/PairSect.lean}, \texttt{H02UV/PtCorr2/Pair/PairSectHist.lean}, \texttt{PairDGeom.lean}, \texttt{H02UV/Spin2/Pair/PairSect.lean}): for a connected $j$-polymer $D$ containing every block whose small-set neighbourhood contains $a$ or $b$ (the blocks meeting a ball of radius $8L^{j}+|a-b|_\infty$, \texttt{dxStd}), group the configurations $X$ of $Z_j=\sum_X I^{\Lambda\setminus X}K^X$ by the component $Y\supseteq D$ of $X\cup D$ (\texttt{sum\_dComp}); outside $Y$ all factors are bulk, so with $K'(Y)=\pi_{ab}\sum_{A:\,\mathrm{comp}_D(A)=Y}I^{Y\setminus A}K^A$, $V'=\mathcal W'=0$, the $\varnothing$ and $\sigma_a\sigma_b$ components of $Z_j$ are reproduced (\texttt{gZO\_sect}, \texttt{zinO\_sect}); the state is local and has no one-point parts (\texttt{pairSect\_hist}). So \texttt{PairSectRed} and \texttt{PairRedSeed} hold (\texttt{pairSectRed\_of}, \texttt{pairRedSeed\_of}), and \texttt{thm\_pairAbove\_red} gives \texttt{PairAboveW1} under the master constants of \texttt{thm\_case2\_pairL} from (H$_\lambda$) for $\hat\lambda$, \texttt{Case2TopStep} at both pair vertices, \texttt{GamLastRem(S)}, the smallness and the analytic inputs \texttt{PairSrcRelR}, \texttt{PairVertSum}, \texttt{PairDefSumR}, \texttt{PairZeroRelR}. \texttt{PairVertSum} is discharged (\texttt{H02UV/Spin2/Pair/PairPR5.lean}, \texttt{pairVertSum\_of}): $|C_k(0,0)|\le\frac{\ell^2}{16}|Z_{k+1}|^{-1}\le\frac{\ell^2e^{K_Z}}{14}\theta_k$, $|\beta_{a,k}|\le 2(1+3K_p)b\theta_k$ ((H$_\lambda$) for the pair vertex, \texttt{BetaLast}), $b\theta_k^2=\vartheta_k\check g_k$, and $\sum_{j\le k<N_h}b\theta_k^2\le\frac6\beta(1+K_S)\vartheta_j$ by \texttt{SpinSum} applied to the exact sequence $T_k=\frac\beta3b\theta_k^2$; with $\vartheta_j\le2|M_j|$ this gives $c_s=2\frac{\ell^2e^{K_Z}}{14}\,2(1+3K_p)(\frac6\beta(1+K_S)+2)$, and \texttt{thm\_pairAbove\_red'} drops the hypothesis. The $\varepsilon_k$ half of \texttt{PairSrcRelR} fails as stated: exactly $\mathfrak n^{(2)}_k-n_{a,k}=\sum_{i<k}\mathrm{tadInc}_2(i,k)$ (\texttt{nn2\_sub\_nA2h}), so $e_{a,k}=\mathrm{err}_k+\beta_{a,k}\sum_{i<k}\mathrm{tadInc}_2(i,k)$ is $O(\check g_k)$ relative and $\sum_k\varepsilon_k\asymp\sum_k\check g_k$ is not uniformly bounded. In the $\mathfrak n^{(2)}$ split, $r'_{k+1}=(1-\mathfrak n^{(2)}_k)r'_k+\tilde\sigma_k$ with $\tilde\sigma_k=\sigma'_k+(\mathfrak n^{(2)}_k-n_{a,k})r'_k$, and the $\varepsilon_k$ come from \texttt{Case2Rel}, \texttt{Case2TopStep} and $\sum_{k<N_h}\check g_k^2\le4/\beta$ (\texttt{GIncConcl}) (\texttt{H02UV/Spin2/Pair/PairPR5b.lean}, \texttt{pairRRelR\_of\_src2}). The $ab$ copy of the exact $\mu^{(2)}$ split along the reduced history is done (\texttt{H02UV/Obs2/Flow/MuChainAB.lean}, \texttt{H02UV/PtCorr2/Pair/PairChainAB.lean}): with $W'_{j,ab}=0$ the $ab$ tadpole functional unrolls (\texttt{wTadAB\_step}, \texttt{MuChainAB.wTadAB\_unroll}; the localised parts contribute only their Pauli tadpoles $m_2(e^{\mathcal L_\Gamma}\operatorname{loc}^{ab}S)-m_2(\operatorname{loc}^{ab}S)$), and splitting $V'(\Lambda)_{ab}=\epsilon+m_2(V'(\Lambda)_{ab})\bar\psi_a\psi_a+D$ ($D$ the gradient monomials of $\operatorname{loc}^{ab}$, $m_2(V'(\Lambda)_{ab})=r'+\mathfrak d_2$) gives exactly $\tilde\sigma_k=-\big(\sum_{i<k}(\hat r'_i-r'_k)\,\mathrm{tadInc}_2(i,k)+\mathrm{err}'_k\big)$, $\hat r'_i=r'_i$ for $i\ge j$ and $0$ below (\texttt{MuChainAB.srcAB\_eq}, \texttt{H02UV/Spin2/Pair/PairPR5d.lean}, \texttt{srcRT\_eq}); $\mathrm{err}'_k$ (\texttt{errAB}) collects the activity, $\mathcal W'_k(0)$, the Pauli and gradient defects, the far blocks and the memory terms (Pauli defect times tadpole, gradient tadpoles, tadpoles of $\operatorname{loc}^{ab}$, product increments). Since the drift contains $r'$, \texttt{PairSrcRelR2} is replaced by \texttt{PairErrRelR} ($\|\mathrm{err}'_k\|\le s_k\|\beta_{a,k+1}\|$, $\sum s_k\le c_e\|M_j\|\check g_j$), and the drift is handled by a relative Gronwall bound with memory (\texttt{H02UV/Obs2/Flow/MemGronwall.lean}, \texttt{rel\_gronwall\_geom}): for $\rho_k=r'_k/\beta_{a,k}$, $\beta_{a,k+1}(\rho_{k+1}-\rho_k)=\rho_k\mathrm{err}_k-D_k-\mathrm{err}'_k$, and $\hat r'_i-r'_k=(\rho_i-\rho_k)\beta_{a,i}+\rho_k(\beta_{a,i}-\beta_{a,k})$; with $|\mathrm{tadInc}_2(i,k)|\le c_t(2/L)^{k-i}\check g_i$, $c_t\check g\le1$ (\texttt{TadConcl}, \texttt{H02UV/Spin2/Case2/Case2RelPairLT.lean}), $|\beta_{a,i}|\le(64/49)^{k+1-i}|\beta_{a,k+1}|$ and $|\mathfrak n^{(2)}_k|\le B'\check g_k$, the increments $\rho_i-\rho_k$ are absorbed by the geometric kernel (column sums $\le\frac{64}{49}\cdot\frac1{16}$ for $L\ge49$) and the vertex increments give $O(\check g^2)$ Gronwall weights. This gives \texttt{PairRRelR} with $c_r=2c_e\exp(2((\frac87c_\delta G_2+\frac17)+\frac1{16}(\frac{64}{49}B'c_t+\frac87c_\delta)G_2+\frac{4}{49}))$ (\texttt{pairRRelR\_of\_err}). Open: \texttt{PairErrRelR}, \texttt{PairDefSumR}, \texttt{PairZeroRelR}. These need norm bounds for the $\sigma_a\sigma_b$ sector of the reduced history (a recursion from the nonzero initial activity $K'_{j}$; \texttt{RelRec} needs zero initial activity and $\beta_a\neq0$). All observable estimates of \texttt{Obs2} use single-source weights $w_{ab}=0$, so a weight $w_{ab}>0$ theory (with the Pauli-aware bound of $(1-\operatorname{loc})$ on the $ab$ sector) is missing. \emph{The $w_{ab}>0$ layer (started).} The barycentric $\operatorname{loc}^{ab}_{X,Y}$ vanishes for $X\cap\{a,b\}=\varnothing$, so $(1-\operatorname{loc}^{ab})$ contracts only on polymers meeting $\{a,b\}$; on the reduced history nothing else occurs: the $ab$ parts of $V'(B),\mathcal W'(B),K'(Y)$ vanish unless $B$ (resp.\ $Y$) meets $\{a,b\}$, since $\operatorname{Loc}_{X,Y}$ carries the factor $|Y\cap\{a,b\}|$, $S^\sharp(B)=\frac12F_C(V'_B,V'(\Lambda)^\varnothing)+\frac12F_C(\pi_{\rm obs}V'_B,V'(\Lambda))$ and the polymer formula only multiplies factors on subsets of $U$ (\texttt{H02UV/Obs2/Pair/AnchAB.lean}, \texttt{OAnchS.step}; \texttt{H02UV/PtCorr2/Pair/ObsHistAnch.lean}, \texttt{obsAnchInv\_histFrom}, \texttt{obsAnchInv\_redSt}). On $\mathrm{abSub}$ the observable norm is $w_{ab}\|F_{ab}\|$ for any weights and the product rule is linear, $\|\pi_{\rm obs}(FG)\|\le\|F_\varnothing\|\|\pi_{\rm obs}G\|+\|\pi_{\rm obs}F\|\|G_\varnothing\|$; the two-scale bound, the bound of $\operatorname{Loc}$ and, on anchored elements, $\|\pi_{\rm obs}(F-\operatorname{Loc}_{X,X}F)\|_{T_{j+1}(\ell')}\le CL^{-1}\|\pi_{\rm obs}F\|_{T_j(\ell)}$ hold, and $\|\bar\psi_a\psi_a\bar\psi_b\psi_b\|_{T_h(\ell)}\le\ell^4(|a-b|_1/h)^2$ (\texttt{H02UV/Obs2/Est/EstAB.lean}). Missing: the $ab$ copies of the block and step bounds (\texttt{EstBlock}, \texttt{BlockO1--6}, \texttt{LinBoundO}, \texttt{StepKO}, \texttt{WObs1--3}) and the recursion from $K'_{j}$. \emph{The $w_{ab}>0$ layer (cont.).} The linear part contracts on the $ab$ sector of anchored activities: $A^{(|\mathcal B_{j+1}(U)|-4)_+}\|\pi_{\rm obs}\mathcal L(K)(U)\|_{T_{j+1}(\ell')}\le C_0(L^{-1}+A^{-\eta})k$ for any weights, if $K(X)\in\mathrm{anc}(X)$ (\texttt{H02UV/Obs2/Est/LinBoundAB.lean}, \texttt{Llin2O\_ab\_bound(L)}; $(1-\operatorname{loc}^{ab})$ is applied only on small sets meeting $\{a,b\}$, counted at $a$ and at $b$). At coalescence, for weights $\ell_{ab}\le\ell_a\ell_b$, $\|\pi_{ab}\,\mathrm{zLoc}\,Y\|\le\sum_{A\in\mathcal S_Y}\prod_{B\subseteq Y\setminus A}\|I(B)\|\prod_{X\in\mathrm{Comp}A}\|K(X)\|$, $Y\setminus A\subseteq D$, $|\mathcal S_Y|\le2^{|\mathcal B(D)|}$, so $\|I(B)\|\le\iota$ and $A^{|\mathcal B(X)|}\|K(X)\|\le1$ give $A^{(|\mathcal B(Y)|-4)_+}\|K'(Y)_{ab}\|\le M_aM_b(2\iota A)^{|\mathcal B(D)|}$ in the weights $(M_a^{-1},M_b^{-1},(M_aM_b)^{-1})$ (\texttt{H02UV/PtCorr2/Pair/PairInitAB.lean}, \texttt{kRed\_ab\_poly\_le}); the factor bounds need the $ab$ parts of the unreduced state at $j_x$ relative to $M_aM_b$, which no estimate provides (all unreduced estimates use $w_{ab}=0$): open. In $\mathrm{err}'$: $\mathrm{farABm}=0$ exactly on anchored inputs, only the blocks of $a$ and $b$ carry $ab$ parts, $|\mathrm{qAB}|\le2Q_1\ell^{-2}\sup_B\|Q'(B)_{ab}\|$, $|\mathcal W'(0)\text{-tadpole}|\le2\ell^{-2}\sup_B\|\mathcal W'(B)_{ab}\|$, $|\mathfrak d_2|$ by \texttt{pauli\_lip\_loc} (\texttt{H02UV/Obs2/Est/ErrAB.lean}); hence \texttt{PairErrRelR} follows from the named input \texttt{PairErrSizesR} (the $ab$ sizes of $Q'_k,\mathcal W'_k$, $\mathfrak d_2$, the gradient tadpole and the memory terms relative to $\|\beta_{a,k+1}\|$, summed $\le c_e\|M_j\|\check g_j$; \texttt{H02UV/Spin2/Pair/PairErrAB.lean}, \texttt{pairErrRelR\_of\_sizes}, anchoring of \texttt{oRed} free). The seed defects are summable above coalescence: $\sum_{k\ge j}|\det\text{-}\mathrm{def}_k|\le2c_0\bar\ell^4/(L-1)$ for $|a-b|_1\le c_0L^{j}$ (\texttt{H02UV/Spin2/Pair/PairDefAB.lean}); \texttt{PairDefSumR} still needs the $ab$ sizes of the reduced history. \emph{The $w_{ab}>0$ layer: initial bound, block count, defects.} The hypotheses of \texttt{kRed\_ab\_poly\_le} split into one-point norms (weights $(M_a^{-1},M_b^{-1},0)$, those of the existing estimates) and the $ab$ parts relative to $M_aM_b$ (\texttt{H02UV/PtCorr2/Pair/PairInitAB2.lean}, \texttt{oTj\_wM\_eq}, \texttt{kRed\_ab\_poly\_le\_split}, \texttt{redHist\_K\_ab\_le}); by Lemma \texttt{lem:coal} the block potentials have no $ab$ part at $j\le j_{ab}$ (\texttt{inV\_ab\_eq\_zero\_of\_coal}). The block count of the covering polymer is $|\mathcal B_{j_x}(D_x)|\le(L+18)^2$ (\texttt{H02UV/PtCorr2/Pair/PairDGeom2.lean}, \texttt{nb\_ballPoly\_coal\_le}; \texttt{H02UV/Spin2/Pair/PairDefAB2.lean}, \texttt{nb\_dxStd\_le}), so $A^{(|\mathcal B(Y)|-4)_+}\|K'_{j_x}(Y)_{ab}\|\le M_aM_b(2\iota A)^{(L+18)^2}$. Open: $\|\mathcal W_{j_x}(B)_{ab}\|$, $A^{|\mathcal B(X)|}\|K_{j_x}(X)_{ab}\|\lesssim M_aM_b$ for the unreduced history; below coalescence $V_{ab}=\mathcal W_{ab}=0$ and $K_{ab}=0$ on small sets, so $\operatorname{Loc}$ never sees an $ab$ part and the step chain in the full norm with $w_{ab}=w_aw_b$ (a third copy of \texttt{stepKO\_ts}, for inputs with one-point parts) should close; not done. On the reduced history $\rho_a=0$, $\mathrm{farAB}=0$, so $\mathrm{hDefR}_k=-\mathfrak d_k+\mathrm{wkAB}_k+\epsilon F_{C_k}(V'_{ab},V'_\varnothing)$ (\texttt{hDefR\_eq}), $|\mathrm{wkAB}|\le2Q_1(\bar q+\bar w)$ (\texttt{norm\_wkAB2\_le}), and with the seed split $\mathfrak d_k=M_{4,k}\det\text{-}\mathrm{def}_k+O(s^p_k)$, a lattice path $0\to x$ of $\le2|x|_\infty<L^{j_x+1}$ steps (\texttt{exists\_path\_sum}) and \texttt{seed\_detDef\_sum\_le}, \texttt{PairDefSumR} follows from the named sizes \texttt{PairDefSizesR} (\texttt{pairDefSumR\_of\_sizes}; the seed contributes $\bar M_4\cdot2L\bar\ell^4/(L-1)$, the last step is a size). \emph{The $w_{ab}>0$ layer: block and step bounds.} The one-point block and step chain is copied to inputs in $\mathrm{abSub}$ (\texttt{H02UV/Obs2/Est/EstNormAB.lean}, \texttt{EstBlockAB}, \texttt{BlockOAB1--6}, \texttt{Step2OAB}, \texttt{Step2OLAB}, \texttt{Step3OAB}, \texttt{Step3OLAB}, \texttt{StepKOWAB}, \texttt{KObsAB}, \texttt{WObsAB}): with all four components local the two-scale bounds hold for any weights; on $\mathrm{abSub}$, $F_C(\mathbb E\theta\pi_{\rm obs}A,\mathbb E\theta\pi_{\rm obs}B)=0$ and the product rule is linear, and the Cauchy estimates use the four pairings $F_\varnothing,F_a,F_b,F_{ab}$. This gives \texttt{stepKO\_tsWAB}, \texttt{kobsK\_stepWAB} (rescaling by $(\bar g/m)^{1/2}$, since $\|\pi_{\rm obs}F\|_{w_\lambda}=\lambda^2\|\pi_{\rm obs}F\|_w$ on $\mathrm{abSub}$) and, at the raw level for anchored states, \texttt{H02UV/PtCorr2/Pair/KObsStepAB.lean}, \texttt{kobs\_raw\_stepAB}, with the constants of \texttt{kobs\_raw\_step}. For $\mathcal W'$ the localisation acts after the expectation on the $(j+1)$-block, and $\operatorname{loc}^{ab}_X$ depends on $X$; the needed contraction $\|\mathbb E\theta G-\operatorname{loc}^{ab}_{X,X}\mathbb E\theta G\|_{T_{j+1}(\ell')}\le KL^{-1}\|G\|_{T_j(\ell)}$ ($G\in\mathcal N(B^\square)$ anchored at $B\subseteq X$, $a,b\in B^\square$) is \texttt{EthSubLocAB}, and \texttt{H02UV/PtCorr2/Pair/WObsStepAB.lean}, \texttt{wobs\_raw\_stepAB} is conditional on it. It holds (\texttt{H02UV/Obs2/Est/WObsAB2.lean}, \texttt{ethSubLocAB}): the range of $\operatorname{loc}^{ab,t}_X$ is spanned by $1$, $\bar\psi_u\psi_v$ ($u,v\in\{a,b\}$), $\bar\psi_a\psi_a\bar\psi_b\psi_b$, which is the whole $U(1)$-invariant part of $\mathcal N(\{a,b\})$ (\texttt{H02UV/Obs2/Loc/LocABRepro.lean}, \texttt{eq\_proj0\_of\_ab}), so for every $X$ meeting $\{a,b\}$, $\operatorname{loc}^{ab}_X$ is the identity on it (\texttt{locABtX\_repro}, also for $a=b$); $\mathbb E\theta$ preserves $\mathcal N(\{a,b\})$ and $U(1)$ invariance, so with $G=(G-\operatorname{loc}^{ab}_BG)+\operatorname{loc}^{ab}_BG$ the second part gives $(1-\operatorname{loc}^{ab}_X)\mathbb E\theta\operatorname{loc}^{ab}_BG=0$ exactly, and the first part is bounded by $(1+C_b)C_cL^{-1}\|G\|_{T_j(\ell)}$ (contraction of $1-\operatorname{loc}^{ab}_B$ into $T_{j+1}(2\ell)$, two-scale bound of $\mathbb E\theta$, boundedness of $\operatorname{loc}^{ab}_X$ at scale $j+1$). Hence \texttt{wobs\_raw\_stepAB2} (\texttt{H02UV/PtCorr2/Pair/KObsHistAB.lean}) is unconditional. \emph{The $w_{ab}>0$ layer: the recursion from the seed.} Along $\mathrm{obsHistFrom}\,j_x(\mathrm{redSt})$ the reduced state is local, $U(1)$-invariant and anchored (\texttt{H02UV/PtCorr2/Obs/ObsSymFrom.lean}, \texttt{obsSymInv\_histFrom}, \texttt{obsSymInv\_redSt}), so \texttt{kobs\_raw\_stepAB}, \texttt{wobs\_raw\_stepAB2} apply at every step $k\ge j_x$ of a regular history (\texttt{H02UV/PtCorr2/Pair/KObsHistAB.lean}, \texttt{kobs\_hist\_stepAB}, \texttt{wobs\_hist\_stepAB}). With the $V'$-sizes $\|V'_k(B)\|_w\le m_k$ as input ($m_k\le2m_{k+1}$), the seed $A^{(|\mathcal B(X)|-4)_+}\|K'_{j_x}(X)\|_w\le k_0$ with $8k_0\le64B\,r_{j_x}^2c_Wm_{j_x}$ and $\mathcal W'_{j_x}=0$, a joint induction gives $\|\mathcal W'_k(B)\|_w\le c_Wr_km_k$ ($c_W=4C_1+1$) and $A^{(|\mathcal B_k(X)|-4)_+}\|K'_k(X)\|_w\le(8\bar\theta)^{k-j_x}k_0+\sum_{j_x\le i<k}(8\bar\theta)^{k-1-i}B\,r_i^2c_Wm_i$, $8\bar\theta\le\frac29$ (\texttt{H02UV/PtCorr2/Pair/PairRecAB.lean}, \texttt{obs\_histFromAB}, \texttt{kseqF\_eq\_sum}); $\|Q'_k(B)\|_w\le C\cdot(\text{that})$ and, in the weights $(M_a^{-1},M_b^{-1},(M_aM_b)^{-1})$, $\|F_{ab}\|\le M_aM_b\|F\|_w$ (\texttt{H02UV/PtCorr2/Pair/PairSizesAB.lean}). \emph{Finding.} The seed $k_0\asymp(2\iota A)^{|\mathcal B(D_x)|}$ is $O(1)$, not $O(\check g)$, relative to $M_aM_b$ (it contains $V_aV_b\propto\beta_a\beta_b\bar\psi_a\psi_a\bar\psi_b\psi_b$); after one step its localisation enters $\mathrm{qAB}$ through the tadpole $m_2(e^{\mathcal L_C}\bar\psi_a\psi_a\bar\psi_b\psi_b)$, of size $|\beta_a\beta_b|\cdot O(L^{-1})$ summed, so the norm route cannot give $\sum_k\ell_k^{-2}\hat s^q_k\le c_e\|M_j\|\check g_j$ in \texttt{PairErrSizesR}: the quartic seed must be separated exactly (as in \texttt{PairDefSumR}) before the crude bounds apply. \emph{PR6 (the zero mode of the reduced history).} At the last scale $\mathcal P_N(\Lambda)=\{\varnothing,\Lambda\}$ and $K_N(\varnothing)=1$, so $F_N=\prod_B\tilde I_N(B)+K_N(\Lambda)$ with $\tilde I_N(B)=e^{-V_N(B)}(1+\mathcal W_N(B))$ (\texttt{H02UV/PtCorr2/Pair/PairZero.lean}, \texttt{finO\_last}). Restricted to constant fields $(\epsilon,c_2)$ is multiplicative in the dual-number sense, and the $\sigma_a$ component (and, without one-point parts, the $\sigma_a\sigma_b$ component) of a product is a derivation (\texttt{deriv\_prod\_of}); hence exactly (\texttt{finO\_last\_a}, \texttt{finO\_last\_ab}) $c_2(F_N)_a=\beta_NP_0+R_a$, $c_2(F'_N)_{ab}=r'_NP'_0+R'_{ab}$, $\epsilon(F_N)_a=\sum_B\epsilon\mathcal W_N(B)_a\prod_{B'\ne B}\alpha_{B'}+\epsilon K_N(\Lambda)_a$ (same for $ab$), $\alpha_B=1+\epsilon\mathcal W_N(B)_\varnothing$, $P_0=\prod_B\alpha_B$, $\beta_N,r'_N$ the couplings produced by the last step, and $R_a$, $R'_{ab}$ contain only the observable parts of $\mathcal W_N$, $K_N$ (so without $\sigma_a\sigma_b$ activity $c_2F'_{ab}=r'_NP_0$, $\epsilon F'_{ab}=0$: \texttt{finO\_last\_ab\_noAct}). With the named last-step remainders \texttt{PairZeroRemR} $e_1,e_2,e_3$ ($P_0\ne0$, $|c_2(F_N)_a-\beta_{a,N-1}P_0|\le\frac12|\beta_{a,N-1}||P_0|$, $|c_2(F'_N)_{ab}-r'_{N-1}P_0|\le e_1|M|\check g|\beta_{a,N-1}||P_0|$, $|\epsilon(F'_N)_{ab}|\le e_2|M|^2\check g|s_\varnothing(\Omega)|$, $|\epsilon(F_N)_a|\le e_3|M||s_\varnothing(\Omega)|$), \texttt{PairRRelR} $c_r$ and $\beta_{a,N-1}\ne0$ (\texttt{hBeta\_rel\_of}) give \texttt{PairZeroRelR} with $\tilde\rho=c_2F'_{ab}/c_2F_a\approx r'_{N-1}/\beta_{a,N-1}$, $K_3=2(c_r+e_1)$, $K_2=e_2+K_3e_3$ (\texttt{H02UV/Spin2/Pair/PairPR6.lean}, \texttt{zero\_rel\_pure}, \texttt{pairZeroRelR\_of}), so \texttt{PairZeroRelR} can be replaced by \texttt{PairZeroRemR}. Open: \texttt{PairZeroRemR} ($e_1,e_2$ need the $w_{ab}>0$ bounds of $\mathcal W'_N,K'_N$ and the last-step source of $r'$; the one-point items need last-step bounds of $\mathcal W_{N,a},K_{N,a}$; the $\epsilon$ items are relative to $|s_\varnothing(\Omega)|=|\epsilon F_\varnothing-c\,c_2F_\varnothing|$, for which only $s_\varnothing(\Omega)\ne0$ is known). \emph{PR6 without $s$.} The bulk normaliser $s=s_\varnothing(\Omega)$ is $1+\tilde u_{N,N}$ for the bulk final object, and with the Ward coefficient $X=\tilde a\pi_N/(1+\tilde u)+S$ exactly $(1+\tilde u)(1-t_{\rm raw}(X-S)/\pi_N)=P_0+\epsilon K_N(\Lambda)$; a lower bound $|s|\gtrsim|P_0|$ is equivalent to a bound on $|t_{\rm raw}\tilde a/(1+\tilde u)|$, which is not available. It is not needed: the $\Omega$-part of the pair constant is $(x-\tilde\rho x_a)/s+\tilde\rho\,n_a\Omega$ with $n_a\Omega=D/s$, $D=\epsilon F_a-c\,c_2F_a$, so $|x-\tilde\rho x_a|\le K_2|M|\check g|D|$ gives $|\cdot|\le(K_2+K_3)|M|\check g|n_a\Omega|$, and $n_a\Omega=r-\gamma_{a,N}$ is fixed by the one-point closure; $s$ cancels (\texttt{H02UV/Spin2/Pair/ZeroNorm.lean}, \texttt{PairZeroRelW}, \texttt{pairAboveW1\_of\_redW}, $K_E=K_1+(K_2+K_3)K_4$, $K_R=(K_2+K_3)M_b$). The remainders \texttt{PairZeroRemW} $e_1,e_2$ are those of \texttt{PairZeroRemR} with the $\epsilon$ items relative to the zero-mode vertex: $|\epsilon(F'_N)_{ab}|\le e_2|M|\check g\,|c||\beta_{a,N-1}||P_0|$, $|\epsilon(F_N)_a|\le\frac14|c||\beta_{a,N-1}||P_0|$ (then $|D|\ge\frac14|c\beta_{a,N-1}P_0|$; \texttt{zero\_rel\_pureW}, \texttt{pairZeroRelW\_of}: $K_2=4e_2+2(c_r+e_1)$, $K_3=2(c_r+e_1)$). Open: \texttt{PairZeroRemW}. \emph{PR6, items (a), (e).} The $\varnothing$-part of $\mathcal W_N(B)$ of the observable last step is $S(B)-\operatorname{Loc}_{\Lambda,B}S(\Lambda)$, while the bulk final object keeps $S(B)$; since $\epsilon(\operatorname{Loc}_{X,Y}F)=\epsilon(F)|Y|/|X|$, the bounds of \texttt{GlueOut} give $|P_0-1|\le L^2we^{L^2w}$, $w=(1+L^2)C_f\check g_{N_h}^2$, hence $P_0\neq0$, $|P_0|\ge\frac12$ for $L^2w\le\frac14$ (\texttt{H02UV/PtCorr2/Last/LastObsA.lean}, \texttt{eps\_Wp\_last\_le}; \texttt{H02UV/Spin2/Pair/ZeroNormA.lean}, \texttt{p0\_last\_le}, \texttt{zP0\_ne\_of}). For (e) the master construction is changed once: $\rho$ is in addition below the $\bar\zeta$ of \texttt{prop\_last\_d2} (\texttt{H02UV/Step2/Flow/FlowMasterHypLZ.lean}, \texttt{H02UV/Spin2/Master/Master.lean}, exporting \texttt{ZeroTOk}), so $|m^2Z_Nt_{\rm raw}-1|\le K_tm^2L^{2N}$ along every history (\texttt{lastOut\_zeroT}) and $|c|\ge1/(2m^2|Z_N|L^{2N})$ if $K_tm^2L^{2N}\le\frac12$ (\texttt{zCc\_ge}). The one-point activity of the last step has the natural size $|\epsilon F_a|\lesssim\ell^2|\beta_{a,N-1}||P_0|/|Z_N|$ (the vertex against $C_N(a,a)$), while the zero-mode vertex is $|c\beta_{a,N-1}P_0|\approx|\beta_{a,N-1}P_0|/(m^2|Z_N|L^{2N})$: at the top of \texttt{mRange} ($m^2L^{2N}=L^2$) the zero mode does not dominate and (e) is not available from norm bounds. \emph{Deviation: mass window.} The pair layer assumes $K_tm^2L^{2N}\le\frac12$ and $8K_Am^2L^{2N}\le1$, i.e.\ $m^2\le c_0L^{-2N}$ instead of $m^2\le L^{-2(N-1)}$; the $h\downarrow0$ route (small masses $m_n^2\to0$) is unaffected. Remaining named inputs: \texttt{LastOneRem} $K_A$ (item (b) and $|\epsilon F_a|\le K_A|\beta_{a,N-1}||P_0|/|Z_N|$; needs a bound of $K^{\rm obs}_{N_h+1}$ and the last observable step at $C_N$) and \texttt{PairZeroRemAB} (items (c), (d)). \emph{PR6, item (b).} At the last scale the bulk inputs are translation covariant and $C_N$ is translation invariant, and every block is a translate of the block $0$; so the bulk block scalars do not depend on $B$, and $\epsilon\mathcal W_N(B)_\varnothing=\epsilon S(B)-\epsilon S(\Lambda)|B|/|\Lambda|=0$: $P_0=1$ exactly (\texttt{H02UV/PtCorr2/Last/LastOneTI.lean}, \texttt{eps\_Wp\_TI}). With $\sum_B\epsilon\mathcal W_N(B)_a=\sum_Bm_2\mathcal W_N(B)_a=0$ the block remainders cancel: $c_2(F_N)_a=\beta_N+c_2K_N(\Lambda)_a$, $\epsilon(F_N)_a=\epsilon K_N(\Lambda)_a$ (\texttt{remCa\_of\_TI}, \texttt{finO\_last\_TI}), and $\beta_N=\beta-\beta\bar\nu-\sum_Bm_2(\mathbb E\theta\mathcal W(B))_a$ (\texttt{betaN\_last}), with $|\bar\nu|=|m_2F_{C_N}(\bar\psi_a\psi_a,V'(\Lambda)_\varnothing)|\le C_f(\ell^2/8+\ell^4/128)\check g_{N_h}$ and the $\mathcal W$ tadpole $\le c_BQ_1w\,s^{-2}$ (\texttt{H02UV/Spin2/Gam/LastOneRem.lean}, \texttt{mass2\_bulk\_last\_le}, \texttt{mass2\_EthO\_W\_le}). The block bound of $\mathcal W^{\rm obs}_{N_h+1}$ keeps its factor $\check g_{N_h}$, $w=c_W\check g_{N_h}s_{N_h}^2|\beta_{N_h}|$ (\texttt{H02UV/Spin2/Gam/LastOneM.lean}, \texttt{LastObsConclW}, \texttt{MasterFacts.lastW}). Hence \texttt{LastOneRem} $K$ follows from the named \texttt{LastKRem} $K$ ($|c_2K_N(\Lambda)_a|\le K\check g_{N_h}|\beta_{a,N-1}|$, $|\epsilon K_N(\Lambda)_a|\le K|\beta_{a,N-1}|/|Z_N|$) and $(C_f(\ell^2/8+\ell^4/128)+\frac94c_BQ_1c_W+K)\check g_{N_h}\le\frac12$ (\texttt{lastOneRem\_of}); these replace \texttt{LastOneRem} and the smallness of $P_0-1$ in the pair layer. Open: \texttt{LastKRem} (no estimate of the observable activity at $j+1=N$; an observable copy of the crude bound of $K_N(\Lambda)$ is needed) and \texttt{PairZeroRemAB}. \emph{PR6, item (b): the activity of the last step.} The map with sources at $j+1=N$ contains $\operatorname{loc}_{\Lambda,B}S(\Lambda)$, and $\Lambda$ is not small at any scale $<N$ (the lifted coordinates wrap around the torus), so $K_N(\Lambda)$ itself has no norm bound. But only its $\sigma_a$ jets enter: $\epsilon$ and $c_2$ of the $\varnothing$ and $\sigma_a$ components form a ring homomorphism $JT$ into $(\mathbb C[\delta]/\delta^2)[\sigma]/\sigma^2$ (\texttt{H02UV/PtCorr2/Last/LastObsK2.lean}), and $F_N=e^{u(\Lambda)}\mathbb E\theta Z_{N-1}=e^{u(\Lambda)}e^{-u'(\Lambda)}(\prod_BG'(B)+K'(\Lambda))$ (\texttt{finG\_eq}) with the reference map without localisation, $u'(B)=\epsilon(V'_B)$, $G'(B)=e^{-(V'_B-u'(B))}(1+S(B))$, $K'(\Lambda)$ the generic polymer map with $J=0$ (\texttt{H02UV/Obs2/Est/LastKO1.lean}, \texttt{stepN\_ZVK\_last}); the localisation enters only through $\sum_B\epsilon(\operatorname{loc}_{\Lambda,B}S(\Lambda))$ and the exact totals of its $\sigma_a$ part. In the weighted norm $|z_{00}|+\alpha|z_{01}|+\beta|z_{10}|+\alpha\beta|z_{11}|$ (submultiplicative, \texttt{LastObsK3.lean}, \texttt{jet\_main}) this gives $|c_2(F_N)_a+\sum_Bc_2V_N(B)_a|\le c_J(c_\rho+4/r_0^2)\tau M/\ell'^2$ and $|\epsilon(F_N)_a|\le c_J(c_\rho+4/r_0^2)\tau M$ for bulk input sizes $\le\tau$, observable block sizes $\le M$, $\|\mathcal W^{\rm obs}\|\le\tau M$, $\|K^{\rm obs}\|\le\tau M/r_0$ (\texttt{LastObsK5.lean}, \texttt{finG\_jets\_le}). The input $K'(\Lambda)$ is bounded by a crude bound of the reference map with sources at the last scale (no range of $C$, no $\operatorname{loc}$; \texttt{H02UV/Obs2/Est/LastKO2.lean}, \texttt{KNO\_crude\_ts}) and a one-parameter Cauchy estimate in the bulk scale (the $\sigma_a$ part vanishes at zero bulk inputs, \texttt{KNO\_ea\_zero}): $\|K'(\Lambda)_a\|\le\tau M/r_0^2$, $\|K'(\Lambda)_\varnothing\|\le\tau/r_0$ (\texttt{H02UV/PtCorr2/Last/LastObsK1.lean}, \texttt{KNO\_a\_le}). Along the trajectory the bulk inputs are $O(\check g_{N_h})$ by \texttt{GlueOut} (\texttt{LastObsK6.lean}, \texttt{bulk\_last\_in}), $\mathcal W^{\rm obs}_{N_h+1}$ by \texttt{LastObsConclW}, and $K^{\rm obs}_{N_h+1}\le c_K\check g_{N_h}s_{N_h}^2|\beta_{N_h}|$ by one $K^{\rm obs}$ step at $N_h$ (\texttt{LastObsK7.lean}, \texttt{kobs\_hist\_stepL}; exported by the master construction, \texttt{H02UV/Spin2/Gam/LastL1M.lean}, \texttt{LastObsConclK}, \texttt{MasterFacts.lastK}, whose requirement on $A$ also contains the regulator requirement $A/2\ge A_0$ of the jet bound). So \texttt{LastKRem} $K_A$ holds with $K_A=c_J(c_\rho+4/r_0^2)c_\tau\frac98c_A(1+\ell^2)$ (\texttt{H02UV/Spin2/Gam/LastKRem1.lean}, \texttt{lastKRem\_of}) from the named input \texttt{LastL1} $c_A$, $\sum_B|m_2V_{N_h+1}(B)_a|\le c_A|\beta_{a,N_h+1}|$ (the $\ell^1$ size of the shares one step past the horizon of the relative recursion; open) and smallness conditions on $\check g_{N_h}$; these replace \texttt{LastKRem} in the pair layer. \texttt{LastL1} is discharged with $c_A=3$ (\texttt{H02UV/PtCorr2/Last/LastL1.lean}, \texttt{H02UV/Spin2/Gam/LastL1M.lean}): the relative recursion at horizon $N_h$ gives $D_{N_h}\le\frac34|\beta_{N_h}|$ and the $\ell^1$ error $E_{N_h}\le c_\delta r_{N_h}^2|1-\mathfrak n_{N_h}||\beta_{N_h}|$ (\texttt{RelRec.E\_top}); with $D_{N_h+1}\le|1-\mathfrak n_{N_h}|D_{N_h}+2E_{N_h}$, $|\beta_{N_h+1}|\ge|1-\mathfrak n_{N_h}||\beta_{N_h}|-E_{N_h}$ and $c_\delta r_{N_h}^2\le\frac18$ this gives $\sum_B|m_2V_{N_h+1}(B)_a|\le3|\beta_{a,N_h+1}|$ whenever $|\mathfrak n^{(2)}_{N_h}|\le\frac17$ (\texttt{obs\_hist\_rel2\_topA}, exported as \texttt{LastL1Concl} by the master construction, \texttt{MasterFacts.lastL1}); so the pair layer does not need \texttt{LastL1}, using the $\mathfrak n$ half of \texttt{Case2TopStep}. \emph{The $w_{ab}>0$ layer: the seed and the unreduced chain.} Along the trajectory $M_{4,k}\asymp\beta_a\beta_b\asymp(b\theta_k)^2$ and $\ell_k^4\asymp\theta_k^2$, so the seed defects are $|M_{4,k}|\,|C_{aa}C_{bb}-C_{ab}C_{ba}|_k\asymp(\vartheta_k\check g_k)^2L^{-(k-j_x)}$, summed $\lesssim|M_{j_x}|^2\check g_{j_x}^2$: the seed is of the target size, but only with per-scale weights; uniform bounds $\bar M_4\bar\ell^4\asymp(\vartheta_{N_h}\check g_{N_h})^2$ are top-dominated. The named sizes of \texttt{PairDefSumR} therefore use per-scale seed weights (\texttt{PairDefSizesR3}, \texttt{pairDefSumR\_of\_sizes3}, \texttt{seed\_geom2\_le}). In the norm chain the quartic seed is of the size of $V$ in the $ab$ sector, so it has to be split off exactly before the crude bounds apply. For the unreduced history (one-point parts present) the block and step chain holds in the full norm with $w_{ab}\le w_aw_b$ once the $V$ sizes are halved, $\|V^\varnothing\|\le r/2$, $\|V^{\rm obs}\|\le m_V/2$, $m_V\le r$; this absorbs the products $F_C(V^{\rm obs},V^{\rm obs})$ with the constants of the one-point chain (\texttt{H02UV/Obs2/Est/EstBlockU.lean}, \texttt{BlockOU2--6}, \texttt{Step2OLU}, \texttt{Step3OLU}, \texttt{StepKOWU}, \texttt{KObsU}). The localisation bound holds for all four components (\texttt{EstLocU}), and the linear part needs only $K(X)_{ab}=0$ on small $X$ (\texttt{LinBoundU}). Along the history below coalescence ($\texttt{ObsCoalInv}$, \texttt{obsHist\_coalInv\_valid}) this gives $A^{(|\mathcal B_j(X)|-4)_+}w_{ab}\|K_j(X)_{ab}\|\le c_2r_jm_j^2$ for $j\le j_{ab}$ from the one-point sizes (\texttt{H02UV/PtCorr2/Obs/ObsHistU.lean}, \texttt{kab\_hist}). The initial reduced activity then follows without the $ab$ activity hypothesis (\texttt{H02UV/PtCorr2/Pair/PairInitU.lean}, \texttt{redHist\_K\_ab\_le\_U}). Open (wiring): the one-point sizes in the weights $(M_a^{-1},M_b^{-1},0)$ from \texttt{obs\_hist\_rel2\_base} and its $a\leftrightarrow b$ mirror, and $\|\mathcal W_{j_{ab}}(B)_{ab}\|\lesssim M_aM_b$ at the last step, which follows from the block bound \texttt{BHOU.Wp\_o}. \emph{The $w_{ab}>0$ layer: the linear regime.} On $\mathrm{abSub}$ the quadratic term of the $K$ step contributes $C_1(\bar g/m)k^2\le C_1c_k\bar gk$ to the $ab$ activity once $k\le c_km$, so the a priori condition $k\le c_k\bar g^2m$ can be dropped with unchanged conclusion (\texttt{H02UV/Obs2/Pair/SeedAB.lean}, \texttt{kobsK\_stepWABL}; \texttt{H02UV/PtCorr2/Pair/SeedSplit.lean}, \texttt{kobs\_raw\_stepABL}, \texttt{kobs\_hist\_stepABL}, \texttt{obs\_histFromABL}: seed $8k_0\le c_kc_Wm_{j_0}$), and $\sum_kk_k\le2(k_0+B\sum_kr_k^2m_k)$ (\texttt{kseqF\_sum\_le}). Since $M_a=2s^2|\beta_a|\asymp\ell^2|M_j|\check g_j$, an $O(1)$ seed in the weights $(M_a^{-1},M_b^{-1},(M_aM_b)^{-1})$ is of the target size of \texttt{PairErrSizesR} per scale and summable along the geometric decay. Open: per-scale weights for the $r'$ part of $V'$, the $r'$-proportional part of the $\mathcal W'$ tadpole (of relative size $\check g_k$, not summable), the $V'$ sizes. \emph{The $w_{ab}>0$ layer: the initial bound under master constants.} With $\gamma_a=\gamma_b=0$ (vertex \texttt{zV0}; the constants decouple from scale $1$ on, \texttt{obsHist\_zV0}) the $\sigma_b$ components at points $(0,x)$ are the $\sigma_a$ components at $(x,0)$ (\texttt{mirror\_b\_eq}), so \texttt{obs\_hist\_rel2\_baseX} at both points gives the one-point sizes in the weights $(M_a^{-1},M_b^{-1},0)$, $M_a=2s_{j_x}^2|\beta_{a,j_x}|$, $M_b$ its mirror: with $x_j=2s_j^2|\beta_{a,j}|/M_a+2s_j^2|\beta^{\rm mir}_{a,j}|/M_b$ and $m_j=\max(c_Wx_j,1)$, $\|V_j\|\le m_j$, $\|\mathcal W_j\|\le r_jm_j$, $K_j\le8Br_j^2m_j$ below coalescence; $m_j\le2m_{j+1}$ since $s_j^2\le\frac65s_{j+1}^2$ and $|\beta_j|\le\frac{64}{49}|\beta_{j+1}|$ (\texttt{bt2\_ratio}, from \texttt{Case2Rel} and \texttt{Case2Loop}), and $m_{j_x}=2c_W$. With $\|K_{j_x}\|\le c_x\check g^3$ and the constants of \texttt{kab\_hist} (requirements \texttt{kabAReq}, \texttt{kabCond}) this gives $A^{(|\mathcal B(Y)|-4)_+}\|K'_{j_x}(Y)_{ab}\|\le M_aM_b(2\iota A)^{(L+18)^2}$ along every admissible trajectory, $1\le j_x\le N_h$ (\texttt{H02UV/Spin2/Pair/PairInitSizes.lean}, \texttt{pairInit\_traj}; \texttt{PairInitConclU}, \texttt{MasterFacts.pairInitU}), conditional only on the block bound $\|I_{j_x}(B)\|_{w_M}\le\iota$ of the unreduced input. That block bound holds with $\iota_W=e^{7/2}(1+e^{L^2C_W(1+16c_W^2)})$: $\|\mathcal W_{j+1}(B)\|_w\le C_W(\bar g+m^2)$ for unreduced inputs ($\bar g\le m$), by \texttt{BHOU.Wp\_o} in the rescaled weights $w_\lambda$, $\lambda=\bar g/m$, and $\lambda^2\|F\|_w\le\|F\|_{w_\lambda}$ (\texttt{H02UV/Obs2/Est/WObsU.lean}, \texttt{wpO\_fullWU}, \texttt{Wnext2O\_full\_le}; \texttt{H02UV/PtCorr2/Obs/WObsStepU.lean}, \texttt{wobs\_hist\_fullWU}), $V_{j_x,ab}=0$, $\|V_{j_x}\|_{w_M}\le\frac32+2$, and $\|e^{-V}(1+\mathcal W)\|\le e^{\|V\|}(1+\|\mathcal W\|)$ (\texttt{H02UV/Spin2/Pair/PairInitW.lean}, \texttt{pairInit\_trajW}; \texttt{PairInitConclW}, \texttt{MasterFacts.pairInitW}): the initial $ab$ activity bound is unconditional. \emph{The $w_{ab}>0$ layer: the $\mathcal W'$ tadpole, per-scale weights.} The term $\mathrm{wTadAB}_k(\mathcal W'_k,0)=\sum_Bm_2(\mathcal W'_k(B)_{ab})$ of $\mathrm{err}'_k$ has no linear part: \texttt{wTadAB\_step} at $\Gamma=0$ gives exactly $\mathrm{wTadAB}_{k}(\mathcal W'_{k},0)=\mathrm{farABm}_{k-1}+\sum_{B_+}m_2\big((\prod_{B\subseteq B_+}(1+\mathcal W'_{k-1}(B))-1-\sum_B\mathcal W'_{k-1}(B))_{ab}\big)$ (\texttt{H02UV/Obs2/Flow/MuChainAB2.lean}, \texttt{wTadAB\_zero\_succ}, \texttt{MuChainAB.errAB\_eq2}), since $\tau^{ab}(0)$ and the Pauli tadpoles of $\operatorname{loc}^{ab}$ vanish at $\Gamma=0$ and $\sum_Bm_2\operatorname{Loc}S(B)=\sum_Bm_2S(B)-\mathrm{farABm}$. On anchored inputs only the blocks of $a,b$ contribute and the product term is linear in the $ab$ size $o$ times the bulk size $e$ of $\mathcal W'_{k-1}$: $\le2L^2o((1+e)^{L^2}-1)\ell_k^{-2}$ (\texttt{norm\_wProdAB0\_le}), i.e.\ $O(\check g^3)|r'_{k-1}|$ relative, summable; no new drift term is needed. The relative Gronwall bound with memory allows errors $\|E_k\|\le s_k\|\beta_{k+1}\|+\sum_{j\le i\le k}\eta_{ik}\|r_i\|$ at the price of $\sum_k\sum_i(\frac{64}{49})^{k+1-i}\eta_{ik}$ in the exponent (\texttt{H02UV/Obs2/Flow/MemGronwall2.lean}, \texttt{rel\_gronwall\_memR}, \texttt{rel\_gronwall\_geomR}); hence \texttt{PairRRelR} from \texttt{PairErrRelR2} (\texttt{H02UV/Spin2/Pair/PairErrAB3.lean}, \texttt{pairRRelR\_of\_err2}, constant with $+c_a$ in the exponent) and \texttt{PairErrRelR2} from the named sizes \texttt{PairErrSizesR3} (as \texttt{PairErrSizesR}, with the product form of the $\mathcal W'$ mass and $r'$-proportional parts allowed; \texttt{pairErrRelR2\_of\_sizes3}, \texttt{pairRRelR\_of\_sizes3}). On $\mathrm{abSub}$ a change of the weights $(M_a^{-1},M_b^{-1},(M_aM_b)^{-1})$ is a rescaling, so \texttt{obs\_histFromABL} runs with per-scale weights $M_{\bullet,k}$ (\texttt{H02UV/PtCorr2/Pair/SeedSplit2.lean}, \texttt{oTo\_wM\_rescale}, \texttt{obs\_histFromABL\_wM}: condition $\lambda_jm_j\le2\lambda_{j+1}m_{j+1}$, $\lambda_j=M_{a,j}M_{b,j}/M_{a,j_0}M_{b,j_0}$); \texttt{kseqF} is linear in the $V$ sizes (\texttt{kseqF\_add\_m}), so the $r'$ part of the activity enters $\mathrm{err}'$ with $\eta_{ik}\propto\theta^{k-1-i}Br_i^2$ and weight $\le2(\frac{64}{49})^2B\sum r_i^2$ (\texttt{geom\_kernel\_weight\_le}). Open: the $V'$ sizes ($r'$, $M_4$, gradient part) along the reduced history (an $ab$ copy of the coupling step), \texttt{memAB}, \texttt{pdefM2}, the gradient terms, i.e.\ \texttt{PairErrSizesR3} and \texttt{PairDefSizesR3}. \emph{The $w_{ab}>0$ layer: the $ab$ coupling.} Along the reduced history $V'_k(B)_{ab}$ is a $U(1)$-invariant element of $\mathcal N(\{a,b\})$ (\texttt{H02UV/PtCorr2/Pair/VObsStepAB1.lean}, \texttt{obsVabInv\_histFrom}), so $V'_k(B)_{ab}=\sum_{u,v\in\{a,b\}}c_{uv}\bar\psi_u\psi_v+M_4\bar\psi_a\psi_a\bar\psi_b\psi_b=m_2\bar\psi_a\psi_a+(c_{ab}+c_{bb})\bar\psi_a\delta+(c_{ba}+c_{bb})\bar\delta\psi_a+c_{bb}\bar\delta\delta+M_4\bar\psi_a\psi_a\bar\delta\delta$, $\delta=\psi_b-\psi_a$, with the norm bounded by the coefficients times powers of $|e|_1\ell/h$ (\texttt{H02UV/Obs2/Flow/VObsAB1.lean}, \texttt{tNorm\_of\_ab\_le}). The Pauli defects are exact: $\mathfrak d_2=M_4(C_{aa}+C_{bb}-C_{ab}-C_{ba})$, $\mathfrak d-M_4(C_{aa}C_{bb}-C_{ab}C_{ba})=\sum c_{uv}(C_{uv}-C_{aa})$ (\texttt{pdefM2\_hVR\_eq}, \texttt{pdefAB2\_hVR\_sub}), and with the Gram representation both defects are of second order, $\le(n\ell/L^{k+1})^2$ resp.\ $2\ell^2(n\ell/L^{k+1})^2$ (\texttt{H02UV/Obs2/Flow/VObsAB2.lean}); this is needed because the barycentric $\operatorname{loc}^{ab}$ produces $M_4\asymp(L^k/|x|)^2$ relative to the norm. Hence \texttt{PairDefSizesR3} uses the second-order seed weight (\texttt{H02UV/Spin2/Pair/PairDefAB4.lean}, \texttt{pairDefSumR\_of\_sizes3}, \texttt{seed\_geom2\_le}). The total quartic coefficient obeys $M_4'=M_4-\sum_BM_4(Q(B)_{ab})-\sum_BM_4(\operatorname{Loc}S(B)_{ab})$ with $M_4(\operatorname{loc}^{ab}_XF)$ the barycentric moment (\texttt{H02UV/Obs2/Flow/VObsAB3.lean}, \texttt{m4Tot\_succ}, \texttt{cfAB4\_locABtX}). If $a,b$ stay in different blocks (at all scales for some $x$), the per-block masses are not controlled by the total $r'$; but the exact $r'$ flow holds verbatim for the mass $r'_k(X)=\sum_{B\in\mathcal B_k(X)}m_2(V'_k(B)_{ab})$ of any region $X$ (a union of blocks of the top scale), with the same factor $1-\mathfrak n^{(2)}_k$ and the same tadpole increments, the block sums of the error restricted to $X$ (\texttt{H02UV/Obs2/Pair/ShareAB1.lean}, \texttt{SFO\_ab\_red}, \texttt{rABX\_next\_red}, \texttt{MuChainAB.rABX\_flow}; \texttt{H02UV/PtCorr2/Pair/ShareABHist1.lean}, \texttt{rRX\_flow}). The relative Gronwall bound therefore holds uniformly in $X$, and with $X=B$ a $k$-block, $|m_2(V'_k(B)_{ab})|\le c_r|M_j|\check g_j|\beta_{a,k}|$ for all $j_x\le k\le N_h+1$ (\texttt{H02UV/Spin2/Pair/PairShareAB1.lean}, \texttt{pairRRelRX\_of\_err2}, \texttt{pairBlockMass\_of\_err2}; named input \texttt{PairErrRelRX}, the region error). The $w_M$ block sizes of $V'$ follow from the mass, gradient and $M_4$ coefficients (\texttt{H02UV/PtCorr2/Pair/ShareABHist2.lean}, \texttt{oTo\_wM\_le\_coeffs}, \texttt{hVo\_of\_coeffs}). The $M_4$ and gradient coefficients of $V'_{ab}$ (\texttt{H02UV/Obs2/Flow/VObsAB4.lean}, \texttt{H02UV/Obs2/Flow/VObsAB4X.lean}, \texttt{H02UV/PtCorr2/Pair/VObsStepAB2.lean}): $|a-b|_\infty^2\le|e|^2$ (\texttt{supNorm\_sq\_le\_eN2}), hence $|M_4(\operatorname{loc}^{ab}_XF)|\le K_{M4}\ell^{-4}(L^j/|a-b|_\infty)^2\|F\|_{T_j(\ell)}$ (\texttt{norm\_cfAB4\_locABtXY\_le}), which bounds the two sources of $M_4'=M_4-\sum_BM_4(Q(B))-\sum_BM_4(\operatorname{Loc}S(B))$ by $\sum_{X\supseteq B}\|K(X)_{ab}\|$ at scale $j$ and $\|S(B_+)_{ab}\|$ at scale $j+1$ (\texttt{norm\_cfAB4\_Q2O\_le}, \texttt{norm\_cfAB4\_LSO\_le}). With $M_4=0$ at $j_x$ and $L^{j_x}\le2|a-b|_\infty$, the second-order seed weights of \texttt{PairDefSizesR3} are summable, $\sum_k|M_{4,k}|\,2\ell_k^2(L^{j_x+1}\ell_k/L^{k+1})^2\le16\sum_k\sigma_k$ (\texttt{m4\_weighted\_sum\_le}, \texttt{m4\_hist\_weighted\_sizes}), and per region/block $|M_4(V'_k(B))|\le K_{M4}\sum_{j_x\le i<k}(\ell_i^{-4}(L^i/|a-b|)^2\kappa^Q_i(B)+\ell_{i+1}^{-4}(L^{i+1}/|a-b|)^2\kappa^S_i(B))$ (\texttt{m4X\_succ}, \texttt{cfAB4\_hVR\_block\_sizes}). Under $e^{\mathcal L_C}$, $\mathcal L_C(\bar\psi_a\psi_a\bar\psi_b\psi_b)=C_{bb}\bar\psi_a\psi_a+C_{aa}\bar\psi_b\psi_b-C_{ab}\bar\psi_b\psi_a-C_{ba}\bar\psi_a\psi_b$ (\texttt{LopC\_m4}), so $c_{uv}(e^{\mathcal L_C}F)=c_{uv}(F)+M_4(F)\kappa_{uv}$ (\texttt{cfAB2\_expLC\_of\_ab}, \texttt{gradCf\_expLC\_of\_ab}) and the gradient part flows by \texttt{gradAB\_VpO}, \texttt{cf2X\_succ}; the gradient part of $\operatorname{loc}^{ab,t}_X$ is $D_t$ (\texttt{gradAB\_locABt}) and contracts, $\|\cdot\|_{T_{j+1}}\le K_D L^{-1}\|F\|_{T_j}$ (\texttt{tNorm\_gradAB\_locABt\_le}). The seed of \texttt{obs\_histFromABL\_wM} is $k_0=(2\iota_WA)^{(L+18)^2}$ from \texttt{PairInitConclW} (\texttt{H02UV/Spin2/Pair/PairSeedW.lean}, \texttt{pairSeed\_hK0}). \emph{The joint induction} (\texttt{H02UV/Spin2/Pair/PairJoint2--3.lean}, \texttt{H02UV/Obs2/Est/ErrABX.lean}): the region error $\mathrm{err}'_k(X)$ is bounded by the same per-block sizes as the total one (\texttt{norm\_qABX\_le}, \texttt{norm\_wProdABX0\_le}); the region Gronwall bound holds for one region from the errors below its top scale (\texttt{regionMass\_of\_err}), so one induction over the scale $n$ with the invariant ``$V'$ block sizes on $[j_x,n)$ and region masses $\le c_r\|M_j\|\check g_j|\beta_{a,k}|$ up to $n$'' breaks the circularity (\texttt{pairJoint\_x}): from a $V$-step (masses up to $n$ $\Rightarrow$ $V'$ sizes at $n$) and an $E$-step (the size items at the scales $<M\le n$ from the $V'$ sizes on $[j_x,n)$, sequences fixed in advance) one gets \texttt{PairErrRelRX}, all region masses, and at $X=\Lambda$ \texttt{PairRRelR} (\texttt{pairRelRX\_of\_joint}, \texttt{pairRRelR\_of\_relRX}, \texttt{pairRRelR\_of\_jointSizes}); the two steps form the named input \texttt{PairJointSizes}. Ingredients for the steps: the gradient part seen from $a$ of a $b$-based localisation is $D_t+m_2(F)(\bar\psi_b\psi_b-\bar\psi_a\psi_a)$ (\texttt{gradAB\_swap}, \texttt{gradAB\_locABt\_swap}, \texttt{H02UV/Obs2/Flow/GradAB1--3.lean}), a scale-weighted gradient size contracting by $L^{-1}$ per step (\texttt{gSzE\_scale\_le}) with sources bounded by $K$, $S$ norms (\texttt{gSzE\_locABtXY\_le}, \texttt{gSzE\_VpO\_le}) and its per-block flow along the history (\texttt{H02UV/PtCorr2/Pair/GradABHist2.lean}, \texttt{gSzE\_hVR\_block\_le}, \texttt{gSrcX\_le}), \texttt{hVo\_of\_coeffsP} (per-block base point); the gradient tadpole $|m_2F_C(D,V'_\varnothing)|\le2\ell^{-2}\|D\|\|V'_\varnothing\|$ (\texttt{norm\_gradTad\_le}, \texttt{norm\_gradTadX\_le}); the memory terms by \texttt{pauli\_lip\_loc} with $I=(i,k)$: $|\mathrm{mem}(i,k)|\le|\mathfrak d_{2,i}||\mathrm{tadInc}_2(i,k)|+1280K^2(8/7L)^{k-i}s_{i+1}^{-2}(f_i+2\lambda_i+2w_i)$ (\texttt{H02UV/PtCorr2/Pair/MemAB1.lean}, \texttt{memAB\_hist\_le}, \texttt{memABX\_hist\_le}). Open: discharging \texttt{PairJointSizes} (locality of the memory arguments, covariance differences in the $M_4$ gradient source, instantiation of \texttt{obs\_histFromABL\_wM} along the induction) and \texttt{PairDefSizesR3}. Since then (\texttt{H02UV/Obs2/Flow/GradAB4.lean}, \texttt{H02UV/PtCorr2/Pair/GradABHist3.lean}, \texttt{H02UV/Spin2/Pair/PairJoint4.lean}): no first-order covariance differences are needed, $\mathrm{gSz}(\mathcal L_Cm_4)\le c(4\ell g+g^2)$ from $|C|\le c$ (\texttt{gSzE\_LopM4\_le}); $|e_{\{a\}}|_1\le d|a-b|_\infty$ (\texttt{eN1\_single\_le}); the field scale $s_j$ is not monotone, so the gradient flow is redone with $s_{i+1}\le qs_i$, contraction $q^4/L$ (\texttt{gSzE\_scale\_le\_q}, \texttt{gSzE\_hVR\_block\_leq}); the $\gamma_k$ input of \texttt{hVo\_of\_coeffsP} for a $k$-block is bounded by the per-scale sizes $\mu_i$ ($M_4$), $\kappa^Q_i$, $\kappa^S_i$ (\texttt{gradM4Sz\_hVR\_block\_sizes}, \texttt{m4X\_hVR\_mu\_le}, \texttt{m4SrcX\_le\_sizes}). From \texttt{MasterFacts}: \texttt{pairRRelR\_of\_jointSizes\_MF} ($\texttt{PairJointSizes}\Rightarrow\texttt{PairRRelR}$ for \texttt{oRed}\,\dots\,\texttt{dxStd} along every admissible trajectory; $V'=\mathcal W'=0$ for \texttt{oRed}, \texttt{oRed\_V\_ab\_zero}) and \texttt{pairDefSumR\_of\_sizes3\_MF}. On the joint route, under \texttt{MasterFacts}, the inputs \texttt{PairSrcRelR2}, \texttt{PairDefSumR} of the pair layer are replaced by \texttt{PairJointSizes}~$c_e,c_a,c_\sigma$ (with $\mathrm{crErr}_2(c_e,c_\delta,4/\beta,c_t,B'_+,c_a)\le c_\sigma$) and \texttt{PairDefSizesR3}. The requirements of the $ab$ chain with per-scale weights (\texttt{obs\_histFromABL\_wM}: \texttt{KObsBoundsABL}, \texttt{WObsBoundsAB}, $\theta_0\le1/72$, $K_w/L\le1/12$, and the $\varepsilon$-conditions with $c_k=1$) are part of the master construction (\texttt{H02UV/Spin2/Master/AbwReq.lean}, \texttt{abwAReq}, \texttt{abwEReq} in \texttt{allAReq}/\texttt{allEReq}), and \texttt{MasterCore.abwHist} runs \texttt{obs\_histFromABL\_wM} along every admissible trajectory; the second-order part has $\|S_j(B_+)_{ab}\|_{T_{j+1}(s_{j+1})}\le M_aM_bL^2s_O(c\check g_j,m_V,m_W)$ and vanishes unless $B_+$ meets $\{a,b\}$ (\texttt{H02UV/PtCorr2/Pair/SXOStepAB.lean}, \texttt{sxo\_hist\_ab\_le}, \texttt{MasterCore.sxoHist}): the $\kappa^S_j$ of the $\gamma$ and $M_4$ bounds. Still open: the $V$-step and $E$-step themselves (see \texttt{Log.md}). \emph{The $E$-step} (\texttt{H02UV/PtCorr2/Pair/MemAB2.lean}, \texttt{H02UV/PtCorr2/Pair/PairEStep1--2.lean}, \texttt{H02UV/Spin2/Pair/PairEStep1--4.lean}, \texttt{H02UV/Spin2/Pair/PairDefSeed.lean}): the memory terms only use the covariances up to $N_h$ (\texttt{memABX\_hist\_leL}; $\texttt{ValidAt}$ below $N_h$ and $\texttt{ValidAtL}$ at $N_h$, \texttt{MasterCore.pairTraj}); the three memory arguments are localised near $a$ ($\mathrm{Loc}\,S_i(B)_{ab}\in\mathcal N(\{a,b\})$, the product parts in $(B_+)^\square$, and, by the finite range of $C_i$, $F_C(D,V'_\varnothing)=F_C(D,V'^{\rm near}_\varnothing)$ with only the blocks within $L^{i+1}/2$ of $\{a,b\}$: \texttt{FC\_pair\_eq\_near}, \texttt{memSz\_of\_sizes}); $|\beta_{a,k}|\le\frac{64}{49}|\beta_{a,k+1}|$ (\texttt{hBeta\_ratio}) converts the memory bounds to $s^m_k|\beta_{a,k+1}|$ with $\eta^m=0$ (\texttt{mem\_conv\_le}); the gradient tadpole uses the near part of $V'_\varnothing$; $|\mathfrak d_2(X)|\le|M_4(X)|(n\ell_k/L^{k+1})^2$. The locality item $Q'_k(B)_{ab}\in\mathcal N(B^\square)$ of \texttt{PairErrItemsX} can fail at $k=j_x$ and is not needed, since $Q'_k(B)_{ab}\in\mathcal N(\{a,b\})$ (\texttt{norm\_qABX\_leP}): \texttt{PairErrItemsXP}, \texttt{PairJointSizesP}, \texttt{pairRRelR\_of\_jointSizesP\_MF}. The $E$-step (\texttt{pairEStep\_of\_sizes}, \texttt{pairJointSizesP\_MF}): $\texttt{VB}\,n$ and the masses below $n$ give the items up to $n$ from the per-scale sizes \texttt{PairESizes} (the $ab$ sizes of $K_k$ on small sets, of $\mathcal W'_{k-1}$, $|M_4|$, the gradient part, the near $V'_\varnothing$, $\mathfrak d_2$, $\mathrm{tadInc}_2$, $\mathrm{Loc}\,S$ and the product parts); all trajectory facts are discharged for \texttt{oRed}. The seed term of \texttt{PairDefSizesR3}: with $s_k\le(9/8)^{k-i}s_i$ instead of monotonicity it is $\le16\sum_kK_{M4}((9/8)^4\kappa^Q_k+\kappa^S_k)$ (\texttt{m4\_hist\_weighted\_sizes\_q}, \texttt{pairDefSeed\_MF}). \emph{The $V$-step} (\texttt{H02UV/PtCorr2/Pair/PairVStep1--4.lean}, \texttt{H02UV/Spin2/Pair/PairVReq.lean}, \texttt{H02UV/Spin2/Master/MasterAllVStep.lean}, \texttt{H02UV/Spin2/Pair/PairVStep1--2.lean}): with the weights $M_{a,j}=2s_j^2|\beta_{a,j}|$, $M_b=\|M_{j_x}\|\check g_{j_x}$ ($M_{a,j}\le2M_{a,j+1}$) the mass part of $\|V'_n(B)\|_{wM}$ is exactly $c_r/2$ by \texttt{PairMassUpTo}\,$n$. The first-order coefficients of $\mathcal L_Cm_4$ are covariance \emph{differences}, $|C_i(a,a)-C_i(b,a)|\le|w|\ell_C^2/L^{i+1}$ along a lattice path (\texttt{covDiff\_le\_path}); with the uniform bound $|C_i|\le c_i$ alone the $M_4$ part of the gradient would grow like $L^{i-j_x}$ (\texttt{gradM4Sz\_hVR\_block\_sizesD}). Only blocks near $a$, $b$ contribute to $\kappa^Q_i\le2c_Nc_S\cdot8\,M_{a,i}M_b\,k_i$ and $\kappa^S_i\le2L^2M_{a,i}M_bL^2s_O$ (\texttt{kappaQ\_le\_count}, \texttt{kappaS\_le\_count}, anchoring), and the gradient sum is $\le4M_{a,n}M_b((2K_G+3q^2+2K_{M4})Q+(2K_G+3+2K_{M4})S)$ (\texttt{gamma\_budget}, \texttt{gradV\_hist\_le}). With $K'$ sizes $\le\mathrm{vck}\,c_W\bar m/8$ (\texttt{kseqF\_leL}) and the new $\varepsilon$-requirement $16\,\mathrm{vA2}(L,C_{1w})c\varepsilon\le1$, $32B(c\varepsilon)^2\le\mathrm{vck}$ (\texttt{abwEReqV}, appended to \texttt{allEReq}), the gradient part is $\le\bar m/2$: \texttt{MasterCore.vstep}, \texttt{pairVStep\_MF} (the $V$-step of the joint induction with $\texttt{VB}=\texttt{PairVB}\,\bar m$ for $c_r\le\bar m$, $8c_u(2\iota_WA)^{(L+18)^2}\le\mathrm{vck}\,c_W\bar m$). The only unit relation is at the seed and is a named input: \texttt{PairUnitSeed}\,$c_u$, $2s_{j_x}^2|\beta^{\rm mir}_{a,j_x}|\le c_u\|M_{j_x}\|\check g_{j_x}$. It is discharged (\texttt{H02UV/Spin2/Pair/PairUnit.lean}, \texttt{pairUnitSeed\_MF}): with $j=j_x=n+1$, $u=(\check g_n/b)^{1/3}\le v=(\check g_j/b)^{1/3}$, $s_j^2=\ell^2|Z_j|^{-1}\le\ell^2e^{K_Z}u$ (the $\log|Z|$ bound of \texttt{SpinConcl}), $|\beta^{\rm mir}_{a,j}|\le\frac{81}{64}|\beta^{\rm mir}_{a,n}|\le\frac{81}{64}(1+3K_p)bu$ ((H$_\lambda$) for $\mu_a/\mu_{a,0}$ at the mirror points, vertex $-b$, via \texttt{hLam\_case2}) and $\|M_j\|\check g_j\ge\frac12(b/\check g_j)^{1/3}\check g_j=\frac12bv^2$ (\texttt{MagConcl} and the smallness of $K(K_\lambda)\check g_j$), so $c_u=\frac{81}{16}\ell^2e^{K_Z}(1+3K_p)$, independent of $\lambda$, $K_\lambda$, $x$. \texttt{thm\_pairAbove\_final} (\texttt{H02UV/Spin2/Pair/PairFinal.lean}): the joint route (input \texttt{PairJointSizesP}) with the $V$-step and the seed unit discharged; the pair inputs are \texttt{PairESizes}\,$c_\sigma$ (with $\texttt{VB}=\texttt{PairVB}\,\bar m$) and the two sums, \texttt{PairDefSizesR3P}, \texttt{PairZeroRemAB}; constants $0<\bar m$, $c_\sigma\le\bar m$, the seed condition. \emph{Constant loop} (\texttt{H02UV/Spin2/Pair/PairLoop.lean}): $\mathrm{crErr2}(c_e)=c_eX$, $X=2\exp(2(\dots))$ depending on $L$, $c_a$; if the $E$-step gives $c_e=\Theta_0k_0+\Theta_1\bar m$ (seed part plus $\bar m$-proportional part) with $X\Theta_1\le\frac12$, every $\bar m\ge\max(2X\Theta_0k_0,\ \text{seed bound})$ closes the loop with $c_\sigma=\bar m$ (\texttt{pairLoop\_close}, \texttt{pairLoop\_exists}). Second-order per-scale sizes are fine: $\sum_{j\le k<n}\check g_k^2\le\frac4\beta(\check g_n-\check g_j)\le\frac{4\varepsilon}{\beta}$ (\texttt{sum\_gsq\_Ico\_le}), giving an $\varepsilon$-requirement $X\,C_E\frac{4\varepsilon}{\beta_L}\le\frac12$; first-order sizes $C(L)\check g_k\bar m$ without decay in $k-j_x$ are not summable uniformly in $b$ ($\sum_{j_x\le k<N_h}\check g_k\asymp\beta^{-1}\log(\check g_{N_h}/\check g_{j_x})$). With the present bound $\kappa^S\lesssim M_aM_bL^2s_O(c\check g,\bar m,c_Wc\check g\bar m)$ the $S$-driven items ($s^p$ through $\mu$, $s^m$ through $\mathrm{Loc}\,S$) are of the latter type; closing the loop needs their first-order part to be of gradient type (factor $L^{-(k-j_x)}$) or part of the exact one-step factor (as $\mathfrak n$ for $\beta_a$). \emph{Resolved} (\texttt{H02UV/Obs2/Pair/SFirst1.lean}, \texttt{H02UV/PtCorr2/Pair/SFirst1.lean}, \texttt{H02UV/Spin2/Pair/PairLoop2--5.lean}): $M_4(\mathrm{Loc}\,S_i)$ is genuinely first order, but its mass effect cancels exactly. $\mathrm{Loc}\,S_i(B)_{ab}\in\mathcal N(\{a,b\})$ is symmetric, so the tadpole of the localised parts removed from $\mathcal W'$ is linear in the covariance, $\mathrm{pdefLS}_i(\Gamma)=M_4(\mathrm{Loc}\,S_i)\,(\Gamma_{aa}+\Gamma_{bb}-\Gamma_{ab}-\Gamma_{ba})$. On the $V'$ side $M_{4,k}=-\sum_{i<k}(M_4(Q_i)+M_4(\mathrm{Loc}\,S_i))$. Hence $-\mathfrak d_{2,k}(1-n_k)-\sum_iM_4(\mathrm{Loc}\,S_i)\,\mathrm{m2D}(C_k)=(\sum_iM_4(Q_i))\,\mathrm{m2D}(C_k)(1-n_k)-n_k(\sum_iM_4(\mathrm{Loc}\,S_i))\,\mathrm{m2D}(C_k)$ (\texttt{errABX2\_eq3}, \texttt{errABX2\_eq3\_hR}): what remains of $S$ carries the vertex factor $n_k=O(\check g_k)$. The items are reorganised accordingly (\texttt{PairErrItemsXP3}, \texttt{PairESizes3}; no $\mathrm{Loc}\,S$ size in the memory). All remaining sizes are seed-decaying or of convolution type $\sum_{i<k}q^{k-i}\check g_i\check g_k$, summable by \texttt{sfirst\_conv\_le}. The loop closes (\texttt{pairLoop2\_close}) under the $\varepsilon$-requirement $X\,C_E(24c^2/\beta_L+2c)\varepsilon\le\frac12$. \emph{$E$-step sizes} (\texttt{H02UV/Spin2/Pair/PairJoint6*.lean}, \texttt{H02UV/PtCorr2/Pair/PairESz1.lean}, \texttt{H02UV/Spin2/Pair/PairESz2--4.lean}): the joint route takes \texttt{PairJointSizesP} and \texttt{PairDefSizesR3P} (the false locality item $Q'_k(B)_{ab}\in\mathcal N(B^\square)$ dropped; the $\varepsilon$-part of $\mathrm{wkAB2}$ is bounded for $F\in\mathcal N(\{a,b\})$, \texttt{norm\_eps\_expLC\_le\_pair}). Of \texttt{PairESizes}: $|1-n_{a,k}|\le3$ ($n_a=\mathfrak n_k-\sum_{i<k}\mathrm{tadInc}_2$, both bulk), the tadpole kernel $c_t(2/L)^{k-i}\check g_i$ (\texttt{TadConcl}) and the Pauli-defect sizes $\mu_i(|w|s_i/L^{i+1})^2$ are discharged (\texttt{pairESizes\_of\_R\_MF}); the per-block $\mathcal W'_{i+1}(B)$ (bulk part $\le c_W(L)c\check g_i$, $ab$ part $\le M_{a,i}M_b(s_O+C_{lo}L^2s_O)$), $\mathrm{Loc}\,S_i(B)_{ab}$ and the product parts of $\mathcal W_{i+1}(B_+)$ are exposed from the block bounds (\texttt{wpo\_raw\_stepAB}, \texttt{wpo\_hist\_ab\_le}) and inserted along the $\mathcal W$ chain (\texttt{pairWSizes\_MF}, \texttt{pairESizesR\_of\_K\_MF}); \texttt{pairJointSizesP\_of\_VRK} gives \texttt{PairJointSizesP} from \texttt{PairESizesK} (the $K$ sizes $\kappa^Q$, $|M_4|$, the $T_{k+1}(4s_k)$ sizes of $\mathrm{dABX}$ and of $V'^{\rm near}_\varnothing$, and $F_C(\mathrm{dABX},V'^{\rm near}_\varnothing)$) and the first summability condition. \emph{$\texttt{PairESizes3}$ discharged} (\texttt{H02UV/PtCorr2/Pair/PairESz2.lean}, \texttt{H02UV/Spin2/Master/MasterVHist.lean}, \texttt{H02UV/Spin2/Pair/PairESz5--8.lean}): the history $K$, $S$, bulk-$V$ and per-scale gradient sizes are exported as \texttt{MasterCore} lemmas (\texttt{vhistK} up to $k=N_h$ via \texttt{ValidAtL}, \texttt{vhistS}, \texttt{vhistE}, \texttt{vhistG}); with the scale change $T_k(s_k)\to T_{k+1}(\ell')$, $\ell'\le qs_k$, of the gradient form (factor $q^4/L$, \texttt{gradM4Sz\_scale\_le}), the bound of $\mathrm{dABX}$ from per-block sizes (\texttt{Tj\_dABX\_le}), the bulk $V'$ near the pair ($\le 2(L+39)^2c^4\,c\check g_k$, \texttt{Tj\_VprO\_e\_le}, \texttt{card\_near\_le}) and the bilinear bound of $F_C$, every item of \texttt{PairESizes3} holds with explicit per-scale sequences (\texttt{pairESizes3\_MF}), and \texttt{pairJointSizesP3\_of\_VS} gives \texttt{PairJointSizesP3} from the first summability condition alone (the per-scale bookkeeping $s_k\le\mathrm{seed}_k+C_E\bar m[\dots]$ for \texttt{pairLoop2\_sum\_le} is not yet in Lean; now in Lean: \texttt{H02UV/Spin2/Pair/PairLoop6a--g.lean}, \texttt{pl6\_item\_le} proves $s_k\le V C_E\,\mathrm{NF}(k)+[k=N_h]VC_N\bar m c\varepsilon$ with explicit $C_E=C_1+\dots+C_5$, \texttt{pl6\_sum\_le}/\texttt{pairJointSizesP3\_closed} discharge the first summability condition with $c_e=2C_Ek_0+\Theta\bar m$, $\Theta=C_E(6c^2G_2+2c\varepsilon)+C_Nc\varepsilon$, $G_2=4\varepsilon/\beta_L+\varepsilon^2$, under $C_{1k}(3+16c_x)c\varepsilon\le\frac1{72}$ and $(1+c_W(L)c\varepsilon)^{L^2}\le2$; \texttt{pairAbove\_final4\_MF} (\texttt{H02UV/Spin2/Pair/PairFinal4.lean}) closes the loop at an explicit $\bar m$ under $X\Theta\le\frac12$; the three $\varepsilon$-requirements are now part of the master construction: \texttt{pairEReq} (\texttt{H02UV/Spin2/Pair/PairEReq.lean}) states them at one set of $\varepsilon$-independent witnesses (those of \texttt{c2RelReqU} and \texttt{abwEReq}, canonical $C_{lo},C_{lb}$), \texttt{masterCore\_exists} takes these witnesses as the exported constants, and \texttt{pairAbove\_final4\_MF'} holds without them). The history items of \texttt{PairDefSizesR3P} are discharged (\texttt{H02UV/Spin2/Pair/PairDef5a--e}, \texttt{PairDef5D}, \texttt{PairDef5eD}, \texttt{H02UV/PtCorr2/Pair/PairDef5S}): the absolute weight bound $2s_k^2|\beta_{a,k}|\le c_u\|M_{j_x}\|\check g_k$ (\texttt{vwA\_le\_mag\_MF}, the scale-$k$ form of \texttt{pairUnitSeed\_MF}) and $|\beta_{a,k+1}|\le\frac{21}{16}|\beta_{a,k}|$ turn the relative $ab$ sizes into absolute ones with an extra factor $\check g_k$; the $K$, $S$ (now also at $N_h$, \texttt{sxo\_hist\_ab\_leL}) and $\mathcal W$ sizes of the reduced history give the $Q'$, $\mathcal W$ items and the seed term (\texttt{pairDefSizesR3P\_of\_NS\_MFC}), and the gradient Pauli defect $\sum_B\sum_{u,v}c_{uv}(C_{uv}-C_{aa})$ is $\le\frac2L\mathrm{gradM4Sz}$ per block (\texttt{pd5\_gradDef\_block\_le}, \texttt{pairDefSizesNS\_of\_NF\_MFC}); the $V$-bound input $\texttt{PairVB}(N_h+1)$ comes from the region masses (\texttt{pairVB\_top\_MF}). \texttt{pairAbove\_final5\_MF'} (\texttt{H02UV/Spin2/Pair/PairFinal5.lean}) is \texttt{pairAbove\_final4\_MF'} with \texttt{PairDefSizesR3P}~$c_d$ replaced by \texttt{PairDefSizesNF}~$c_d$ (only the $\varepsilon$-tadpole $\varepsilon F_{C_k}(V'_{k,ab},V'_{k,\varnothing})$ and the last-step defect remain; constant $c_d+\texttt{pd5Cp}+\texttt{pd5CS}$). The last pair inputs \texttt{PairDefSizesNF}, \texttt{PairZeroRemAB} are discharged (\texttt{H02UV/Spin2/Pair/PairDef6a}, \texttt{PairDef6b(D)}, \texttt{H02UV/PtCorr2/Last/LastAB1--4}, \texttt{H02UV/Spin2/Pair/PairZeroAB1--4}, \texttt{H02UV/Spin2/Pair/PairFinal6.lean}): the $\varepsilon$-tadpole is $\le2(64s_k^2|r'_k|+2\cdot8^4L^{-1}\gamma_k)e_k$ with only $V'^{\rm near}_\varnothing$ entering (finite range, $e_k=2(L+39)^28^4c\check g_k$), hence second order by the weight bound (\texttt{pd6\_epsTad\_le}, \texttt{pairDefSizesNF\_of\_last\_MFC}); at the last scale the map $\sigma_a\sigma_b\mapsto\sigma_a$ is multiplicative on elements without one-point parts and commutes with the reference map without localisation, so $(F'_N)_{ab}$ is the $\sigma_a$-part of $F_N$ for the swapped inputs (\texttt{sw\_finG}) and the jet bound \texttt{LastJetBd} of the master construction applies (\texttt{finG\_jets\_ab}; no new requirement on $A$); the last-step defect is bounded crudely by the input sizes (\texttt{lastDef\_ab\_le}); the $ab$ sizes at $N_h+1$ come from one more step of the $ab$ chain (\texttt{MasterCore.vstepTopD}, \texttt{pairTopSizes\_MF}), and \texttt{pairZeroAB\_of} gives \texttt{PairZeroRemAB}~$e_1\,1$ and the last-step defect. \texttt{pairAbove\_final6\_MF'} and its windowed form \texttt{fw\_pairAbove6\_MFD'} have no named pair inputs; new hypothesis: the mass window $2\cdot\frac{81}{16}c_J(c_\rho+4/r_0^2)c_S\ell^2\,mL^{2N}\le1$ of item~(d) (the analogue of $8K_AmL^{2N}\le1$), folded into $g_0$ in Theorem~\ref{thm:multiscale}.
\emph{One master construction.} All master variants are consolidated (\texttt{H02UV/Spin2/Master/MasterAll*.lean}): \texttt{masterCore\_exists} runs the flow construction once, with the conjunction of all requirements on $A$ and $\varepsilon$ (\texttt{allAReq}, \texttt{allEReq}, including the step-$N_h$ requirement \texttt{topReq}); \texttt{thm\_master\_all} gives one set of constants with all exports of the earlier variants (a structure \texttt{MasterFacts}), together with the per-scale relative bounds up to $N_h$ (\texttt{RelBaseConcl}) and the $\ell^1$ error $E_{N_h}$ of the step $N_h$ (\texttt{TopRelConcl}). The hypothesis \texttt{Case2TopStep} is discharged by \texttt{TopConcl} (\texttt{MasterFacts.top}).
\texttt{GamTopBd} is reduced (\texttt{H02UV/Obs2/Flow/GamTop.lean}, \texttt{H02UV/Spin2/Gam/GamTop.lean}): the $\gamma$ step holds for any covariance
(\texttt{eps\_uXO\_univ\_a}), so at the last step $\gamma_{a,N}=\gamma_{a,N-1}-\beta_{N-1}C_N(0,0)+\rho_N$ ($C_N=$ \texttt{stCovN}, \texttt{gamNa\_eq});
$\gamma_{a,N_h}=M_{N_h}(\lambda^\gamma)$ with $\lambda^\gamma$ satisfying (H$_\lambda$) once $\rho_k$ is relatively small for all $k<N_h$; then
$|\gamma_{a,N}|\le\vartheta_{N_h}(1+K+c_L)\le2(1+K+c_L)|M_{j_x}|$, using that $\check g$ is nondecreasing along admissible trajectories
(\texttt{GMonoConcl}, from (F1) of the master construction). \texttt{gamTopBd\_flow}: \texttt{GamTopBd} with $K_4=2(1+K(K_p+c_\gamma)+c_L)$ from
(H$_\lambda$) for $\hat\lambda$ and the pair vertex, \texttt{GamBulkRel}, \texttt{GamWKTop} (\texttt{GamWKRel} for all $k<N_h$) and
\texttt{GamLastRel}: $|\beta_{N_h}\Gamma_{N_h}-\rho_{N_h}|+|\beta_{N_h+1}C_N(0,0)-\rho_N|\le c_L\vartheta_{N_h}\check g_{N_h}$ (the natural size
$b\theta_k/|Z_{k+1}|=\vartheta_k\check g_k$).
\texttt{GamLastRel} is reduced (\texttt{H02UV/Spin2/Gam/GamLast.lean}): the step covariance at $N_h$ (\texttt{ValidAtL}) and the last covariance
\texttt{stCovN} (\texttt{ValidLast}, now exported by the master construction: \texttt{H02UV/Step2/Flow/FlowMasterHypL.lean}, \texttt{H02UV/Spin2/Master/MasterL.lean},
\texttt{LastValidConcl} in \texttt{H02UV/Spin2/Master/Master.lean}) are \texttt{GoodCov2}, so $|\Gamma_{N_h}|,|C_N(0,0)|\le\ell^2/(16|Z|)$, and $1/|Z_{N_h+1}|\le e^{K_Z}\theta_{N_h}$,
$|Z_N|\ge\frac78|Z_{N_h+1}|$; with $b\theta_k^2=\vartheta_k\check g_k$ this gives \texttt{GamLastRel} with $c_L=2(1+c_r)K_\beta\ell^2e^{K_Z}/14$
(\texttt{gamLastRel\_flow}) from the named inputs \texttt{BetaLast} ($|\beta_{N_h}|,|\beta_{N_h+1}|\le K_\beta b\theta_{N_h}$) and
\texttt{GamLastRem} ($|\rho_{N_h}|\le c_r|\beta_{N_h}||\Gamma_{N_h}|$, $|\rho_N|\le c_r|\beta_{N_h+1}||C_N(0,0)|$).
The point $x$ (\texttt{H02UV/Spin2/Gam/GamTopB.lean}): \texttt{prjO\,false} of the last increment is \texttt{prjO\,true} of the mirrored one, so
$\gamma_{b,N}$ is $\gamma_{a,N}$ of the mirrored flow (\texttt{gamNb\_eq}); \texttt{gamTopBdB\_flow} gives $|\gamma_{b,N}|\le K_4'|M_{j_x}|$ from
the mirrored inputs (\texttt{GamRemTopS}, \texttt{BetaLastS}, \texttt{GamLastRemS}); \texttt{gamTopBd\_both\_flow}:
both bounds with the master constants.
Inputs discharged (\texttt{H02UV/Spin2/Gam/GamLastBd.lean}, \texttt{H02UV/Spin2/Case2/Case2RelPairL.lean}):
$\beta_{N_h}=\beta_{N_h-1}(1-\mathfrak n_{N_h-1})-\mathrm{err}_{N_h-1}$ is covered by \texttt{Case2Rel}/\texttt{Case2Loop} at $k=N_h-1$, and
$\beta_{N_h+1}=\beta_{N_h}(1-\mathfrak n_{N_h})-\mathrm{err}_{N_h}$ (the chain identity with horizon $N_h+1<N$) by the named input
\texttt{Case2TopStep} ($|\mathfrak n_{N_h}|\le\frac18$, $|\mathrm{err}_{N_h}|\le\frac18|\beta_{N_h}(1-\mathfrak n_{N_h})|$; \texttt{Case2Rel} extended by one step,
\texttt{case2TopStep\_of\_rel}); with (H$_\lambda$) at $N_h-1$ and $\theta_{N_h-1}\le\theta_{N_h}$ this gives \texttt{BetaLast}, \texttt{BetaLastS} with
$K_\beta=2(1+3K_p)$ (\texttt{bt2\_last\_le}). \texttt{GamRemTopS} follows from \texttt{GamBulkRelA} at $a=x$ and the mirrored $K/\mathcal W$ bound of
\texttt{wk\_inputs} (\texttt{gamRemTopS\_of}). \texttt{thm\_case2\_pairL} gives \texttt{RelConcl2G}, \texttt{LoopConcl}, \texttt{WKConcl} and \texttt{GamBulkConcl} through \texttt{thm\_flow\_masterL}
($\epsilon$-requirement \texttt{c2RelReq}$\wedge$\texttt{gbEReq}), together with \texttt{GMonoConcl} and \texttt{LastValidConcl};
with these constants, \texttt{GamTopBd}$\wedge$\texttt{GamTopBdB} hold along admissible trajectories given only
(H$_\lambda$) for $\hat\lambda$, \texttt{Case2TopStep} at both pair vertices, \texttt{GamLastRem(S)} and the smallness of $K(K_\lambda)\check g_j$.
The step $N_h$ (\texttt{H02UV/Obs2/Flow/RelRecTop.lean}, \texttt{H02UV/PtCorr2/Obs/ObsHistTop.lean}, \texttt{H02UV/PtCorr2/Obs/ProdBd2Top.lean}, \texttt{H02UV/Spin2/Case2/Case2Top.lean},
\texttt{H02UV/Spin2/Case2/Case2TopMaster.lean}, \texttt{H02UV/Spin2/Gam/GamTopFull.lean}): at $j=N_h$ only \texttt{ValidAtL} holds; the lemmas of the relative recursion use
\texttt{ValidAt} at $j$ only through \texttt{ValidCore}, $\rho\in[0,\frac18]$, \texttt{RawInv}, the a priori bounds and $|m^2|L^{2j}\le1$, so their copies
at $j=N_h$ hold (\texttt{kobs\_EK2L}, \texttt{gamWK2\_leL}, \texttt{prodBd2\_histL}, \texttt{C2Ctx.tadInc2\_leU}, \texttt{C2Ctx.gamBulk\_le\_posU}).
The relative recursion closed at horizon $N_h$ plus one error bound (\texttt{RelRec.E\_top}; no summability up to $N_h+1$ is needed) gives
$|\mathrm{err}_{N_h}|\le c_\delta r_{N_h}^2|\beta_{N_h}(1-\mathfrak n_{N_h})|$ (\texttt{obs\_hist\_rel2\_top}). The bound
$|\mathfrak n_{N_h}|\le(4L^2+3)c\,\check g_{N_h}$ follows from $\mathfrak n_j=\nu'_j\,\mathrm{rowS}_{C_j}-g'_j\mathfrak l_{C_j}+\sum_{k<j}\mathrm{tadInc}_{kj}$
(\texttt{C2Ctx.frakN2\_top\_le}). It is crude, since the couplings at $N_h+1$ are not controlled. It is made $\le\frac18$ by the new requirement
$(4L^2+3)c\,\epsilon\le\frac18$ (\texttt{topReq}). Hence \texttt{Case2TopStep} at every point configuration and $|\rho_{N_h}|\le c_\rho|\beta_{N_h}|/|Z_{N_h+1}|$
(\texttt{TopConcl}, \texttt{thm\_case2\_top}: a copy of \texttt{thm\_case2\_pairL} with $\epsilon$-requirement
\texttt{c2RelReq}$\wedge$\texttt{gbEReq}$\wedge$\texttt{topReq}).
Deviation: \texttt{GamLastRem(S)} is replaced by the form relative to $1/|Z_{N_h+1}|\le e^{K_Z}\theta_{N_h}$, because no lower bound on $C_N(0,0)$ is available.
\texttt{gamLastRel\_ofZ} then gives \texttt{GamLastRel} with $c_L=2K_\beta(\ell^2e^{K_Z}/14+c_re^{K_Z})$.
\texttt{thm\_gamTopBd\_case2'}: \texttt{GamTopBd}$\wedge$\texttt{GamTopBdB} hold given only (H$_\lambda$) for $\hat\lambda$, the last-step remainders
\texttt{GamRemNZ}, \texttt{GamRemNZS} ($|\rho_N|\le c_r|\beta_{N_h+1}|/|Z_{N_h+1}|$) and the smallness.
\emph{The last step} (\texttt{H02UV/PtCorr2/Obs/ObsHistLast.lean}, \texttt{H02UV/Spin2/Gam/GamRemLast.lean}, \texttt{H02UV/Spin2/Gam/GamTopFullL.lean}). At $j=N_h+1=N-1$ the map has no $Q$ (\texttt{Q2O} vanishes for $j+1=N$), so $\rho_N$ has no activity term and
no $\Delta$-bound of the last step is needed: $\rho_N=\sum_B\varepsilon(e^{\mathcal L_{C_N}}\mathcal W^{\rm obs}_{N_h+1}(B)_a)+\beta_{N_h+1}\,\varepsilon F_{C_N}(\bar\psi_a\psi_a,V'_N(\Lambda)_\emptyset)$
(\texttt{gamRemG\_last\_split}). The bulk part is $V_{\kappa_N}(\Lambda)$ plus a constant (\texttt{VprLO\_e\_last}, \texttt{rawStepN\_VN}), so
$\varepsilon F=\nu_N\sum_xC_N(a,x)C_N(x,a)+(y_N+z_N)\mathfrak l_{C_N}(a)$; on the torus at the last scale the sums have $L^{2N}$ terms
($|\sum CC|\le L^{2N}\ell_C^4$, $|\mathfrak l_C|\le2\ell_C^4$, $\ell_C^2=\ell^2/(16|Z_N|)$), and the final couplings of \texttt{GlueOut}
($|L^{2N}\nu_N/Z_N|,|z_N/Z_N|,|y_N/Z_N|\le C_f\check g_{N_h}$) with $|Z_N|\ge\frac78|Z_{N_h+1}|$ give $|\varepsilon F|\le C_f\ell^4\check g_{N_h}/|Z_{N_h+1}|$
(\texttt{eps\_bulk\_last\_le}). The $\mathcal W$ term is $\le c_BQ_1w$ (\texttt{eps\_EthO\_W\_le}, \texttt{GoodCov2} of $C_N$ at $s_{N_h+1}$), where
$w=(3C_1+1)s_{N_h}^2|\beta_{N_h}|$ comes from one $\mathcal W^{\rm obs}$ step at $N_h$ (\texttt{wobs\_hist\_stepL}, \texttt{ValidAtL}) fed by the
relative recursion at horizon $N_h$ (\texttt{obs\_hist\_rel2\_baseX}, which also gives $\sum_B|m_2(V_{N_h}(B)_a)|\le2|\beta_{N_h}|$). With
$|\beta_{N_h}|\le2|\beta_{N_h+1}|$ (\texttt{Case2TopStep}) and $s_{N_h}^2\le\frac98\ell^2/|Z_{N_h+1}|$: $|\rho_N|\le c_r|\beta_{N_h+1}|/|Z_{N_h+1}|$ with
$c_r=2c_BQ_1(3C_1+1)\frac98\ell^2+C_f\ell^4$ (\texttt{gamRemNP\_le}), for both points (\texttt{gamRemN}, \texttt{gamRemNS}).
The master construction adds the block bound \texttt{LastObsConcl} (\texttt{MasterFacts.lastO}), and
\texttt{MasterFacts.fw\_gamTopBdD} (\texttt{H02UV/Spin2/Final/FinalWire.lean}): \texttt{GamTopBd} holds along every trajectory in $D_\varepsilon$ given only (H$_\lambda$) for $\hat\lambda$ and the smallness.
\emph{The $\sigma_x$ mirror and $j_x=0$} (\texttt{H02UV/Obs2/Map/SwapO.lean}, \texttt{H02UV/Obs2/Map/SwapHist.lean}, \texttt{H02UV/PtCorr2/Zero/ForestNear.lean},
\texttt{H02UV/Spin2/Multiscale/TwoPtNear.lean}, \texttt{H02UV/Spin2/Gam/GamFormB.lean}). In Case~(2) ($\tau=\mathrm{id}$) the projection $F\mapsto F_\emptyset+\sigma_aF_i$ ($i\in\{a,b\}$)
is an algebra endomorphism commuting with every operation of the map, and $\operatorname{loc}^a,\operatorname{loc}^b$ are the same operator; so it
commutes with the whole observable step and history with the points changed from $(a,b)$ to $(p_i,\cdot)$ (\texttt{prj\_obsStep},
\texttt{prj\_obsHist}). For the symmetric couplings of $\tilde z_0\tilde z_x$ the $\sigma_x$ part of the pair flow is the $\sigma_a$ part of the flow
with points $(x,0)$ (\texttt{uB\_eq\_uMir}); with $C_k(x,x)=\Gamma^{\rm int}_{k+1}(0,0)$ (\texttt{histCov\_gamInt\_at}) the chain of the $0$-part applies
verbatim (\texttt{gamFormB\_of}: \texttt{GamFormB} from (H$_\lambda$) for the mirrored pair vertex and \texttt{GamRemRelS}, the mirror of \texttt{GamRemRel}).
\texttt{TwoPtNear} is proved at model level (\texttt{twoPtNear\_of}, $K_N=L$, $K_{RN}=1$): $1-b\avg{\bar\psi_0\psi_x}=P_{\beta,h}[\neg(0\leftrightarrow x,0\not\leftrightarrow g)]
\le P[0\leftrightarrow g]+P[0\not\leftrightarrow x]$, $P[0\leftrightarrow g]=r$, and $P_w[u\not\leftrightarrow v]\le 1/w_{uv}$ (adding the edge $uv$ to a forest
with $u\not\leftrightarrow v$ is injective), so along a nearest-neighbour path $P[0\not\leftrightarrow x]\le2|x|_\infty/\beta<Lb$, uniformly in $N$.
\texttt{twoPtSq\_of\_gam\_flowB}: \texttt{twoPtSq\_of\_gam\_flow} without \texttt{GamFormB} and \texttt{TwoPtNear}, with (H$_\lambda$) for the mirrored pair
vertex and \texttt{GamRemRelS} added. Not proved: \texttt{GamRemRelS} (same analysis as \texttt{GamRemRel} at $a=x$).

Section~\ref{sec:spin} shows that $M_j=\vartheta_j(1+O(\check g_j))$ with
$\vartheta_j=(\beta_j/\beta_\epsilon)^{1/3}$. Together with Expected result~\ref{thm:rginput},
this gives Theorem~\ref{thm:multiscale} (on route (a1) Expected result~\ref{thm:rginput} is proved at $h=0$; see there).

%======================================================================
\section{The dressed finite range decomposition}
\label{sec:frd}

This section collects everything about covariances. References of the form ``F-x.y'' are to
\frd{}.

\subsection{The standard decomposition in $d=2$}
We use the decomposition of \texttt{sec:frd} (Lean: \texttt{FRD}, with
\texttt{stdPt := PtF.prePt.toPtFamily}):
\[
  (-\Delta+m^2)^{-1}=\sum_{k=1}^{N-1}C_k+C_N+t_NQ_N,\qquad\rng(C_k)<\tfrac12L^k\ (k<N).
\]
\begin{itemize}
\item It is a standard decomposition (Definition~F-2.1), in every $d\ge2$
  (Proposition~F-2.2; Lean: \texttt{DressedFRD.frd\_stdFull}). For $L\ge4$ it is also
  non-degenerate (Definition~F-2.5, Proposition~F-2.6).
\item In $d=2$ it is scale invariant: $\hat e_j(p)=E(L^j|p|)+O(|p|^2+m^2)$ with an explicit
  profile $E$ (proof of Proposition~\ref{prop:d2}).
\item Consequently $C_{j+1}(0,0)$ and $C_{j+1}(0,e_1)$ equal $\frac1{2\pi}\log L+O(L^{-j}+m^2L^{2j})$,
  and $C_{j+1}(x,y)=c_*(L^{-j}(x-y))+O(L^{-j})$. There is no factor $L^{-(d-2)j}$.
\end{itemize}

\subsection{The obstruction}
The gradient coefficient of the effective potential is not small: $Z_j\to0$ (Section~\ref{sec:spin}).
So the gradient term must be moved into the Gaussian with the exact identity
\begin{equation}\label{e:absorb}
  \E_\Gamma\bigl[e^{-z(\psi,A\bar\psi)}F\bigr]=\det(1+zA\Gamma)\,\E_{\Gamma(1+zA\Gamma)^{-1}}[F],
  \qquad A=-\Delta+m^2 .
\end{equation}
For $\Gamma=Z^{-1}\sum_{k>j}C_k$, the right-hand covariance is a Neumann series in
$E_j=A\sum_{k\le j}C_k$, which does not have finite range. Rescaling the field alone multiplies
covariances by scalars and renormalises nothing; see \frd{}.

\subsection{The dressed decomposition}
Theorem~F-3.1 constructs, for increments $|\zeta_i|\le\bar\zeta(L)$ and
$Z_{j+1}=Z_j(1+\zeta_j)$, covariances $\Gamma_j$ (remaining) and $\Gamma^{\rm int}_{j+1}$
(integrated at step $j$) with the following properties.
\begin{enumerate}[label=(D\arabic*)]
\item\label{D1} \emph{Exact splitting:}
  $\Gamma_j(1+z_jA\Gamma_j)^{-1}=\Gamma^{\rm int}_{j+1}+\Gamma_{j+1}$, with $z_j=\zeta_jZ_j$ and
  $\Gamma_0=A^{-1}$.
\item\label{D2} \emph{Finite range:} $\rng(\Gamma^{\rm int}_{j+1})<\frac12L^{j+1}$.
\item\label{D3} \emph{Closeness:} $\|Z_{j+1}\Gamma^{\rm int}_{j+1}-C_{j+1}\|_{(j)}\le K|\zeta|_{\le j}$,
  hence Gram-admissibility. The first-order term is $\zeta_jC_{\le j}(1-E_j)$ (Theorem~F-3.1(v)).
\item\label{D4} \emph{Positivity:} for real $\zeta$, $Z_{j+1}\Gamma^{\rm int}_{j+1}$ is
  positive semidefinite up to a remainder of relative size $\bar\zeta^{L/256}$
  (Theorem~F-3.1(vii)).
\item\label{D5} \emph{Adaptedness:} $\Gamma^{\rm int}_{j+1}$ and $\Gamma_{j+1}$ depend
  holomorphically on $\zeta_0,\dots,\zeta_j$ only, and are real for real $\zeta$.
\end{enumerate}
The hypotheses are $|\zeta_j|\le\bar\zeta(L)$ and $m^2\le L^{-2(N-1)}$. The first holds because
$\zeta_j=O(g\log L)$, provided $\beta_R\ge C(L)$. The second is the regime $m^2\downarrow0$.

\subsection{The one-step identity}
\label{sec:onestep-id}
\begin{lemma}[One step]\label{lem:onestep}
For every element $F$ of the Grassmann algebra and $j<N$,
\[
  \E_{\Gamma_j}\bigl[e^{-z_j(\psi,A\bar\psi)}F\bigr]=\det(1+z_jA\Gamma_j)\;
  \E_{\Gamma_{j+1}}\bigl[\E_{\Gamma^{\rm int}_{j+1}}\theta F\bigr].
\]
\end{lemma}

\begin{proof}
The Grassmann identity \eqref{e:absorb} follows by completing the square, or from
$\E_\Gamma\theta=e^{\cL_\Gamma}$. Combine it with \ref{D1} and the convolution property
$\E_{C_1+C_2}=\E_{C_2}\circ\E_{C_1}\theta$. All three are algebraic and hold for arbitrary
complex covariances. In the Lean library \texttt{Grassmann}, $\cL_C$ and $\E_C$ are defined for
real matrices only; the identity needs the complex versions.
\end{proof}

\subsection{Normalised covariances}
In fields normalised at scale $j+1$, $\psi=Z_{j+1}^{-1/2}\psi^{(j+1)}$, the integrated and the
accumulated covariances are
\begin{equation}\label{e:Cw}
  C'_{j+1}=Z_{j+1}\Gamma^{\rm int}_{j+1},\qquad
  w'_j=\sum_{k\le j}\frac{Z_{j+1}}{Z_k}C'_k,\qquad w'_{j+1}=w'_j+C'_{j+1}.
\end{equation}
For $|\zeta|\le\bar\zeta$, the factors $Z_{j+1}/Z_k$ are at most $(1+\bar\zeta)^{j+1-k}$, which
is compensated by the decay of $\hat C_k$ at momenta $|p|\gtrsim L^{-k}$. The field scale is the constant $\ell_0$. The Gram bound
$\|\E_{C'}\theta F\|_{T_{j+1}(\ell_{j+1})}\le\|F\|_{T_j(\ell_j)}$ follows from
\ref{D3}, Lemma~F-1.1(ii) and the Gram--Hadamard inequality. It replaces \texttt{e:E-contract}
of \hr{}, which uses positive semidefiniteness.

\subsection{Universal coefficients}
Proposition~\ref{prop:d2}: uniformly in $N$, $m^2\le L^{-2(N-1)}$ and $|\zeta|\le\bar\zeta$,
\begin{equation}\label{e:univ}
  C'_{j+1}(0,0),\ C'_{j+1}(0,e_1),\ \mathfrak l(w'_{j+1})-\mathfrak l(w'_j),\
  \mathfrak B(w'_{j+1})-\mathfrak B(w'_j)\ =\ \frac{\log L}{2\pi}+O(L^{-j}+m^2L^{2j}+|\zeta|_{\le j}).
\end{equation}
Here $\mathfrak l(w)=(w,-\Delta w)$, and $\mathfrak B$ is the bubble functional of
Section~\ref{sec:pt-coef}. Only the telescoped differences are universal.

\subsection{Estimates for the continuum limit}
\begin{itemize}
\item \emph{Physical units.} For fixed $k$, $C_{N-k}(x/\epsilon,y/\epsilon)\to c_*(L^k(x-y))$
  as $\epsilon\to0$, with derivatives. The same holds for $C'_{N-k}$, given convergence of the
  increments $\zeta_{N-k}$ (Section~\ref{sec:uv}).
\item \emph{Green function.} $\sum_{k<N}C_k(0,x)=\frac1{2\pi}\log\frac{L^N}{|x|}+c+O(|x|^{-1})$.
  This is the $d=2$ analogue of \texttt{FRD.GreenAsymp}.
\end{itemize}

%======================================================================
\section{Grassmann integration, symmetries, polymers, norms}
\label{sec:grassmann}

\subsection{Identical parts}
The following are as in \hr{} and are not repeated:
\begin{itemize}
\item Grassmann Gaussian integration: $\E_\Gamma\theta=e^{\cL_\Gamma}$, the semigroup property,
  and factorisation for separated supports. Finite range is \ref{D2}.
\item Symmetries: $U(1)$ and lattice automorphisms. The $\Gamma^{\rm int}$ are in $\cT$, hence
  automorphism invariant.
\item Blocks, polymers, small sets, $X^\square$, $\bar X$, and the factorisation property of
  $K$ (\texttt{sec:polymers}).
\item The $\Phi_j(\ell)$ and $T_j(\ell)$ norms, the product property, monotonicity, and the
  doubling bound. However, \texttt{e:normcontract} fails in $d=2$; it is replaced in
  Section~\ref{sec:normcontract}.
\end{itemize}

\subsection{Field scale}
Normalised fields are used throughout, so $\ell_j=\ell_0$ for all $j$, instead of
$\ell_0L^{-\frac12(d-2)j}$. The constant $\ell_0$ is chosen so that
$\|C'_{j+1}\|_{\Phi_{j+1}(\ell_0)}\le1$ in the Gram sense, which is possible by \ref{D3}. In bare units this is
$\ell_j=\ell_0|Z_j|^{-1/2}$.

Precisely, by Corollary~F-3.2(a), $\ell_0^2\ge(c_1+K\bar\zeta)L^{d-2+q}$ gives
$\|\Gamma^{\rm int}_{j+1}\|_{\Phi_{j+1}(\ell_{j+1})}\le1$ with
$\ell_j=\ell_0|Z_j|^{-1/2}L^{-(d-2)j/2}$. This is the \texttt{norm} field of
\texttt{GoodCov}. The \texttt{range} and \texttt{inv} fields are \ref{D2} and
$\Gamma^{\rm int}\in\cT$. The \texttt{psd} field is replaced by Gram-admissibility
(Definition~F-1.2). For a non-degenerate decomposition, Theorem~F-3.1(vii) says more: the
Gram representation is the usual psd one for $\Gamma^+_{j+1}$, plus the representation of
Lemma~F-1.1(ii) for a remainder of relative size $\bar\zeta^{L/256}$.

\subsection{Localisation in $d=2$}
\label{sec:loc}
$\Loc_X$ and $\Loc_{X,Y}$ are defined as in \hr{} (\cite{MR3332939}) with field dimension
$[\psi]=0$ and maximal dimension $d_+=2$. Proposition \texttt{prop:Loc} holds with $d=2$; its
proof in \cite{MR3332939} does not depend on $d$.

The space $\cV$ of formal local polynomials is the same as in Definition \texttt{def:cV}:
\[
  V=y(\nabla\psi)(\nabla\bar\psi)+\tfrac z2\bigl((-\Delta\psi)\bar\psi+\psi(-\Delta\bar\psi)\bigr)+\nu\,\psi\bar\psi+g\,M .
\]
Now $M$ is marginal (dimension $2=d$), not irrelevant.
\begin{itemize}
\item By nilpotency and $U(1)$ invariance, every other local monomial of dimension $\le2$ is
  zero, or is odd under lattice reflections.
\item In particular $\psi\bar\psi\partial_i\psi\partial_j\bar\psi$ has only the $\delta_{ij}$
  part.
\item On the torus, $y$ and $z$ multiply the same quadratic form. So
  $\zeta=y+z$ is the coefficient of $(\psi,-\Delta\bar\psi)$.
\end{itemize}

\paragraph{The quartic part of $\Loc$}
Since $M$ is marginal, $\Loc_X$ must reproduce the quartic test polynomials of dimension $2$.
Concretely (Lean: \texttt{Loc/Loc4Def.lean}), with lifted coordinates $\ell(x)=r(x-b_X)$,
\[
  \mu^{(4)}_{ik}(F)=\sum_{x_0,\dots,x_3}F_{(x_0,+)(x_1,-)(x_2,+)(x_3,-)}\,\ell(x_0)_i\,\ell(x_1)_k ,
  \qquad
  M^{ik}_x=(\nabla_i\psi_x)(\nabla_k\bar\psi_x)\,\psi_x\bar\psi_x ,
\]
and $\Loc_{X,Y}=\Loc^{(2)}_{X,Y}+\Loc^{(4)}_{X,Y}$, where $\Loc^{(2)}$ is the localisation of
\hr{} (constant and quadratic part) and
$\Loc^{(4)}_{X,Y}F=\sum_{i,k}|X|^{-1}\mu^{(4)}_{ik}(F)\sum_{x\in Y}M^{ik}_x$. The pairing
matrix is diagonal: $\mu^{(4)}_{i'k'}(M^{ik}_x)=\delta_{ii'}\delta_{kk'}$ for $x\in X$ (patch
condition), so $\mu^{(4)}(F-\Loc_XF)=0$. $\Loc^{(4)}$ has coefficients only on sequences of
length $4$ with two fields of each spin, maps into $\cN(Y^\square)$, and satisfies
\texttt{prop:Loc}(i) (Lean: \texttt{exists\_Loc4XY\_bound}, \texttt{exists\_Loc2XY\_bound}).
For lattice-symmetric $F$ only the trace $\sum_i\mu^{(4)}_{ii}$ survives, which is the
coefficient of $M$. The Lean localisation \texttt{Loc.LocX} is $\Loc^{(2)}$ only; the
$d=2$ localisation is \texttt{Loc.Loc2X}.

\subsection{Replacement of \texttt{e:normcontract}}
\label{sec:normcontract}

\paragraph{What fails}
In \hr{}, \texttt{e:T-monotone} and \texttt{e:E-contract} combine into \texttt{e:normcontract}:
\[
  \|F\|_{T_{j+1}(\ell_{j+1})}\le\|F\|_{T_{j+1}(2\ell_{j+1})}\le\|F\|_{T_j(\ell_j)},\qquad
  \|\E_{C_{j+1}}\theta F\|_{T_{j+1}(\ell_{j+1})}\le\|F\|_{T_j(\ell_j)} .
\]
The middle inequality uses $2\ell_{j+1}\le\ell_j$, i.e.\ $L^{(d-2)/2}\ge2$. The doubling map
$\theta$ costs a factor $2$ in the field scale (\texttt{e:theta-bd}), and in $d\ge3$ the
decrease of $\ell_j$ pays for it. In $d=2$, $\ell_j=\ell_0$, and the inequality is false: for
$F=\psi_x\bar\psi_x$,
\[
  \|F\|_{T_{j+1}(2\ell_0)}=4\ell_0^2>\ell_0^2=\|F\|_{T_j(\ell_0)} .
\]
For the same reason, the contraction factor of \texttt{prop:Loc}(ii),
$\bar CL^{-d}L^{-d_+}$ with $d_+=\frac{d-2}2\wedge1$ in \hr{}, gives no gain in $d=2$:
$d_+=0$. In the Lean development, \texttt{Tj\_succ\_le} and \texttt{Tj\_EC\_theta\_le} (file
\texttt{Grassmann/Torus.lean}) are \texttt{e:normcontract} and assume $d\ge3$.

\paragraph{The mechanism in $d=2$}
The loss is a factor $2$ per field. It is compensated by the Pauli principle: an
antisymmetric kernel with many fields of one type in one block, tested against a function
that is smooth at the larger scale $L^{j+1}$, is small. So the loss is at most a constant
\emph{per block}, not per field. Per block it is absorbed by the large-set regulator. On small
sets, it is absorbed by the derivative gain of $1-\Loc$.

Let $p=p_\Phi$ be the number of derivatives in $\Phi_j(\ell)$ (\texttt{def:Phi}); in the Lean
development $p=2d$. Both lemmas below only use Taylor expansions of order $\le2$ and the
antisymmetry of the kernels: a test function that does not depend on $x_i$ (a \emph{frozen}
variable) has vanishing pairing with $F$ as soon as two frozen variables of equal spin lie in
the same block.

\paragraph{Partition of unity}
Let $R=L^j$ and let $b_R$ be the discrete B-spline of order $p$ on the torus:
$b^{(0)}=\mathbf 1_{[0,R)^d}$, $b^{(m+1)}=R^{-d}\,\mathbf 1_{[0,R)^d}\ast b^{(m)}$, $b_R=b^{(p)}$.
For grid points $Rq$, $q\in\{0,\dots,L^{N-j}-1\}^d$, put $\chi_q(x)=b_R(x-Rq)$. Then
$\sum_q\chi_q=1$, $\chi_q$ is supported in $Rq+[0,(p+1)R)^d$, and
$\sum_q|\nabla^\beta\chi_q|\le(2/R)^{|\beta|}$ for $|\beta|\le p$ (the identity
$\nabla(\mathbf 1\ast f)=(\nabla\mathbf 1)\ast f$ and $\sum|\nabla_e\mathbf 1_{[0,R)^d}|\le2R^{d-1}$).
For $Y\subset\Lambda$ let $\cB_j(Y)$ be the set of $q$ with $\chi_q\not\equiv0$ on $Y$. Then
$|\cB_j(Y)|\le(p+1)^d\,|\{j\text{-blocks meeting }Y\}|$.

\begin{lemma}[Doubling with Pauli gain]\label{lem:pauli}
Let $C_\chi=2^pd(3p+1)$, $C_*=(4\cdot2^pC_\chi)^2$ and $L\ge L_0=4\cdot2^p(p+1)^{2d}C_\chi$.
For $j<N$, $Y\subset\Lambda$, $F\in\cN(Y)$ and $\ell>0$,
\[
  \|F\|_{T_{j+1}(2\ell)}\le C_*^{\,|\cB_j(Y)|}\,\|F\|_{T_j(\ell)} .
\]
Consequently, if $C'$ is Gram-admissible with $\|C'\|_{\Phi_{j+1}(\ell)}\le1$ in the sense of
the Gram bound of Section~\ref{sec:frd}, then
$\|\E_{C'}\theta F\|_{T_{j+1}(\ell)}\le C_*^{|\cB_j(Y)|}\|F\|_{T_j(\ell)}$.
\end{lemma}

\begin{proof}
Let $g$ be in the unit ball of $\Phi_{j+1}(2\ell)$ and fix the field degree $n$.
\begin{itemize}
\item \emph{One-variable operators.} For a block $q$ with centre $c_q=Rq$ and a variable $i$,
  let $P^{(i)}_qH=\chi_q(x_i)\,H|_{x_i=c_q}$ (frozen) and
  $R^{(i)}_qH=\chi_q(x_i)\,(H-H|_{x_i=c_q})$ (remainder); the spin of $x_i$ is kept. For an
  assignment $\alpha=(q_i,o_i)_{i\le n}$, $o_i\in\{P,R\}$, let $G_\alpha$ be the composition of
  these operators over $i=1,\dots,n$, applied to $g$. Since $\sum_q\chi_q=1$ and $P+R=\chi$,
  $g=\sum_\alpha G_\alpha$ on $Y^n$.
\item \emph{Pauli cancellation.} Call $\alpha$ inadmissible for the spins $\sigma$ if two
  frozen variables $i\ne k$ have $q_i=q_k$ and $\sigma_i=\sigma_k$. Then $G_\alpha$ factorises
  as $\prod_{\text{frozen}}\chi_{q_i}(x_i)\cdot K(\text{frozen points moved to the centres})$,
  and the involution exchanging the canonical such pair $(i,k)$ shows
  $\sum_zF_zG_\alpha(z)=-\sum_zF_zG_\alpha(z)$. So $\ip Fg=\ip Fh$ with $h$ the sum of the
  admissible $G_\alpha$. An admissible $\alpha$ has at most $2|\cB_j(Y)|$ frozen variables.
\item \emph{Product bounds.} By the discrete Leibniz rule, $|\nabla^\beta G_\alpha|$ is bounded
  by $(2\ell)^n\prod_iw_{o_i}(x_i,\beta_i)$ with weights summing over $q$ to
  $\le2^pC_\chi R^{-|\beta_i|}$ (frozen) and $\le2^p(p+1)^{2d}C_\chi L^{-1}R^{-|\beta_i|}$
  (remainder; the path from $c_q$ to $x_i$ has length $\le d((p+1)R+2p)$ on
  $\operatorname{supp}\chi_q$, and $g$ varies at scale $LR$).
\item \emph{Counting.} With $A_P=2^pC_\chi$ and $A_R=2^p(p+1)^{2d}C_\chi/L\le\frac14$,
  $2^n\sum_{\alpha\text{ adm}}\prod_iw\le(4A_P)^{2|\cB|}(4A_R)^{n-2|\cB|}\le C_*^{|\cB|}$.
\end{itemize}
The consequence follows from \texttt{e:theta-bd} and the Gram bound.
\end{proof}

\begin{lemma}[Localisation contraction in $d=2$]\label{lem:loc2}
There are $C_P$ and $L_0$ such that for $L\ge L_0$, $j<N$, $\ell'\le2\ell$, $X\in\cS_j$ and
$U(1)$-invariant $F\in\cN(X^\square)$,
\[
  \|(1-\Loc_X)F\|_{T_{j+1}(\ell')}\le C_PL^{-3}\|F\|_{T_j(\ell)} ,
\]
where $\Loc_X$ is the localisation of Section~\ref{sec:loc}, \emph{including the quartic
marginal part} $\Loc^{(4)}_X$.
\end{lemma}

\begin{proof}
Let $G=(1-\Loc_X)F$; by \texttt{prop:Loc}(i), $\|G\|_{T_j(\ell)}\le(1+\bar C)\|F\|_{T_j(\ell)}$.
Split the pairing with a test function $g$ (antisymmetrised, in the unit ball of
$\Phi_{j+1}(\ell')$) by the field degree $n$.
\begin{itemize}
\item $n\in\{0,1,3\}$: the coefficients of $G$ vanish ($\Loc$ reproduces the constant; odd
  degrees vanish by $U(1)$).
\item $n=2$: as in \hr{}: $G$ annihilates the Taylor polynomial of degree $2$ of $g$ at the base
  point of $X$, and the remainder of order $3$ gains $L^{-3}$, times $(\ell'/\ell)^2\le4$.
\item $n=4$: the coefficients $G_{(x_0,+)(x_1,-)(x_2,+)(x_3,-)}$ are antisymmetric under
  $x_0\leftrightarrow x_2$ and $x_1\leftrightarrow x_3$. Hence $G$ annihilates the constant,
  all linear terms, and the quadratic terms in two points of equal spin. The mixed terms
  $\ell(x_a)_i\ell(x_b)_k$ are $\pm$ the quartic pairings $\mu^{(4)}_{ik}(G)$, which vanish by
  the definition of $\Loc^{(4)}_X$. (Every $U(1)$-invariant spin pattern is a permutation of
  $(+,-,+,-)$.) So again only the Taylor remainder of order $3$ contributes: $L^{-3}$, times
  $(\ell'/\ell)^4\le16$.
\item $n\ge5$: $G_n=F_n$. Use the decomposition of Lemma~\ref{lem:pauli} with a \emph{single}
  block: $\chi\equiv1$, centre the base point $b$ of $X$, distance $\|r(x-b)\|_1$. At most one
  variable of each spin is frozen, so at least $n-2\ge3$ variables carry a remainder, each
  $O(L^{-1})$ on $X^\square$. After the localisation of the test function at scale $L^j$
  (constant $C_l$ per variable), the bound is
  $(4C)^2(4CA/L)^{n-2}\le(4C)^2(4CA)^3L^{-3}$ for $L\ge4CA$, uniformly in $n$ (here
  $C=C_l+1$, $A=d(2^{d+2}+4d)+2$, and $(\ell'/2\ell)^n\le1$ absorbs the field scale).
\end{itemize}
\end{proof}

\paragraph{Remark}
With the localisation of \hr{} (constant and quadratic part only) the bound fails at $n=4$:
the antisymmetry gives only $L^{-2}$ (\texttt{exists\_long\_bound}), and the test functions
$(x_0-x_2)_i(x_1-x_3)_k$ are marginal in $d=2$. With $O(L^2)$ small sets per block, $L^{-2}$
gives no contraction. This is why $\Loc^{(4)}_X$ is part of $\Loc_X$ in $d=2$.

\paragraph{How each use of \texttt{e:normcontract} in \hr{} is replaced}
\begin{enumerate}[label=(U\arabic*)]
\item \emph{Local polynomials:} $\E\theta V_j(B)$ and $\E\theta Q(B)$, in \texttt{E:vj}.
  These have degree $\le4$, and $W_j(B)$ has degree $\le8$. So directly,
  $\|\E_{C'}\theta P(B)\|_{T_{j+1}(\ell_0)}\le2^{\deg P}\|P(B)\|_{T_j(\ell_0)}$. The constant
  $2^{\deg P}$ is harmless. But $\|V_{j+1}\|\le\|V_j\|+O(\|K_j\|)$ (\texttt{E:vj}) is no longer
  used: $V_{j+1}$ is controlled coupling by coupling through the explicit flow
  (Section~\ref{sec:pt-flow}).
\item \emph{Small sets} (\texttt{lem:smallset-K}). Use Lemma~\ref{lem:loc2} in place of
  \texttt{e:Loc-contract}. With the $O(L^2)$ small sets $X$ with $\bar X=U$, this gives
  $O(L^{-1})\|K_j\|_j$. So the contraction exponent is $\kappa_0=1$, replacing $d_+$.
\item \emph{The block factors $\delta I$} (\texttt{lem:smallset-V}, \texttt{lem:smallset-V1}).
  There $2\ell_{j+1}\le\ell_j$ was used for $\theta(e^{-V_j(B)}-1)$. Now
  $\|\theta(I^\flat_j(B)-1)\|_{T_{j+1}(\ell_0)}\le O(\|V^\flat_j(B)\|_{T_{j+1}(2\ell_0)}+\|W^\flat_j(B)\|_{T_{j+1}(2\ell_0)})$,
  and the bounded degree gives $O(2^8)(\|V_j\|_j+\|W_j\|_j)$. This is absorbed in the
  smallness $\bar g\le\epsilon_0(L)$.
\item \emph{Large sets} (\texttt{e:largeset2}--\texttt{e:largeset3}). By
  Lemma~\ref{lem:pauli},
  $\|\E_{C'}\theta K_j(X)\|_{T_{j+1}(\ell_0)}\le C_*^{|\cB_j(X)|}\|K_j(X)\|_{T_j(\ell_0)}$.
  The extra $C_*^{|\cB_j(X)|}$ is absorbed by choosing $A\ge A_0(L)$ with $A_0\ge 4C_*$ in
  \texttt{e:Achoice}. Then $A^{-\eta}$ becomes $(C_*/A)^{\eta'}$.
\end{enumerate}
The Gaussian contraction \texttt{e:E-contract2} for Gram-admissible, possibly non-psd, $C'$
is the Gram--Hadamard bound of Section~\ref{sec:frd}; it does not depend on $d$. The ordering
of the steps is unchanged: first $\theta$ and $\E$ at scale $j+1$, then the change of scale.
All constants $C_*,C_P$ depend only on $p_\Phi$. They are uniform in $j$, $N$, $L\ge L_0$,
and in the dressing.

%======================================================================
\section{Second-order perturbation theory with dressed covariances}
\label{sec:pt}

This section collects the algebra and the coefficient formulas that the second-order
renormalisation group of Section~\ref{sec:rgmap} needs when the covariances are the normalised
dressed covariances \eqref{e:Cw}. Everything up to Section~\ref{sec:pt-coef} holds in any
dimension. Section~\ref{sec:pt-d2} evaluates the coefficients in $d=2$. By Theorem~F-3.1,
$\rng(w'_j)<\frac12L^j$ and $\|C'_{j+1}-C_{j+1}\|_{(j)}\le K|\zeta|_{\le j}$.

\subsection{Truncated expectations}
For $\Gamma\in\cT_\C$ put $\cL_\Gamma=\sum_{x,y}\Gamma(x,y)\partial_{\psi_y}\partial_{\bar\psi_x}$,
so that $\E_\Gamma\theta=e^{\cL_\Gamma}$ on the Grassmann algebra. This is the definition
used in \hr{} (equation \texttt{e:fWick}), and it requires no positivity. Define
\[
  F_\Gamma(A,B)=e^{\cL_\Gamma}\bigl((e^{-\cL_\Gamma}A)(e^{-\cL_\Gamma}B)\bigr)-AB ,
\]
so that $F_\Gamma(\E_\Gamma\theta A,\E_\Gamma\theta B)=\E_\Gamma\theta(AB)-(\E_\Gamma\theta A)(\E_\Gamma\theta B)$
is the truncated expectation.

\begin{lemma}\label{lem:F}
For $\Gamma,\Gamma_1,\Gamma_2\in\cT_\C$ and even $A,B$:
\begin{enumerate}[label=(\roman*)]
\item $F_\Gamma$ is bilinear and symmetric, and $F_0=0$.
\item $F_{\Gamma_1+\Gamma_2}(e^{\cL_{\Gamma_2}}A,e^{\cL_{\Gamma_2}}B)
  =e^{\cL_{\Gamma_2}}F_{\Gamma_1}(A,B)+F_{\Gamma_2}(e^{\cL_{\Gamma_2}}A,e^{\cL_{\Gamma_2}}B)$.
\item If $A$ depends only on fields in $X$, $B$ only on fields in $Y$, and
  $\operatorname{dist}_\infty(X,Y)>\rng(\Gamma)$, then $F_\Gamma(A,B)=0$.
\item $e^{\cL_\Gamma}$ does not enlarge supports: if $A$ depends only on fields in $X$,
  then so does $e^{\cL_\Gamma}A$.
\end{enumerate}
\end{lemma}

\begin{proof}
(ii): Since $\cL_{\Gamma_1}$ and $\cL_{\Gamma_2}$ commute,
$e^{\cL_{\Gamma_1+\Gamma_2}}=e^{\cL_{\Gamma_2}}e^{\cL_{\Gamma_1}}$. Both
sides of (ii) then equal
$e^{\cL_{\Gamma_1+\Gamma_2}}\bigl((e^{-\cL_{\Gamma_1}}A)(e^{-\cL_{\Gamma_1}}B)\bigr)-(e^{\cL_{\Gamma_2}}A)(e^{\cL_{\Gamma_2}}B)$.
(iii): $\cL_\Gamma(A'B')=(\cL_\Gamma A')B'+A'(\cL_\Gamma B')$ plus terms with
a factor $\Gamma(x,y)$, $x\in X$, $y\in Y$, which vanish. (iv) is clear.
\end{proof}

\subsection{The second-order map}
\label{sec:pt-map}
Let $\Loc_x$ be the localisation of \hr{} (Proposition \texttt{prop:Loc}) with
image $\C\oplus\mathcal V$, $\mathcal V\ni(z,y,a,b)$. A \emph{second-order remainder} is a
family $W=(W(x))_{x\in\Lambda}$ with $W(x)$ even. We write
$W(X)=\sum_{x\in X}W(x)$, $V(\Lambda)=\sum_xV_x$, and
$F(V,V)(x)=F(V_x,V(\Lambda))$.

\begin{definition}\label{def:Vpt}
Given a local $V$, a remainder $W$ and a covariance $C'$, let $V'=\E_{C'}\theta V$ and
\begin{align*}
  u_{\rm pt}+V_{\rm pt}(x)&=V'_x-\Loc_x\Bigl[\tfrac12F_{C'}(V'_x,V'(\Lambda))+\E_{C'}\theta\,W(x)\Bigr],\\
  W_{\rm pt}(x)&=(1-\Loc_x)\Bigl[\tfrac12F_{C'}(V'_x,V'(\Lambda))+\E_{C'}\theta\,W(x)\Bigr].
\end{align*}
\end{definition}

\begin{proposition}\label{prop:Vpt}
Replace $V$ by $tV$ and $W$ by $t^2W$. Then, as formal power series in $t$,
\[
  \E_{C'}\theta\bigl[e^{-V(\Lambda)}(1+W(\Lambda))\bigr]
  =e^{-u_{\rm pt}|\Lambda|-V_{\rm pt}(\Lambda)}\bigl(1+W_{\rm pt}(\Lambda)\bigr)+O(t^3).
\]
Moreover, suppose $W(x)$ depends only on fields within distance $R$ of $x$, and
$\rng(C')<R'$. Then $W_{\rm pt}(x)$ depends only on fields within distance
$\max(R,R'+1)$ of $x$.
\end{proposition}

\begin{proof}
By the cumulant expansion,
$\E\theta e^{-tV}=\exp\bigl(-tV'+\tfrac{t^2}2F_{C'}(V',V')\bigr)+O(t^3)$ and
$\E\theta(e^{-tV}t^2W)=t^2\E\theta W+O(t^3)$. Split $\frac12F+\E\theta W$
into $\Loc$ and $1-\Loc$. Locality follows from Lemma~\ref{lem:F}(iii),(iv). The range of
$V_x$ is $1$.
\end{proof}

With $W_0=0$, the remainder $W_j$ after $j$ steps is defined recursively by
Definition~\ref{def:Vpt}. By Proposition~\ref{prop:Vpt}, $W_j(x)$ depends only on fields
within distance $\frac12L^j+1$ of $x$. This is the finite-range analogue of the coordinate
$W_j$ of Bauerschmidt--Brydges--Slade, Chapter~5. There $W$ has the closed form
$\frac12(1-\Loc)F_{w_j}(V,V)$. In our setting, $V$ and the normalisation change at each step,
so the recursive form is more convenient. The two agree to second order:

\begin{lemma}\label{lem:frozen}
If $W=\frac12(1-\Loc)F_w(V,V)$, then
\[
  V_{\rm pt}=V'-\tfrac12\Loc F_{w+C'}(V',V')+\tfrac12\E_{C'}\theta\,\Loc F_w(V,V).
\]
\end{lemma}

\begin{proof}
By Lemma~\ref{lem:F}(ii),
$\frac12F_{C'}(V',V')=\frac12F_{w+C'}(V',V')-\frac12\E_{C'}\theta F_w(V,V)$. Insert this in
Definition~\ref{def:Vpt}, and use
$\Loc\E\theta F-\Loc\E\theta(1-\Loc)F=\Loc\E\theta\Loc F$.
\end{proof}

\paragraph{Gradient terms and absorption}
In the renormalisation group with dressed covariances, the global gradient term of $V_j$ is
absorbed into the Gaussian before $V_{\rm pt}$ is formed (Section~\ref{sec:onestep}, step (A)). So
Definition~\ref{def:Vpt} is applied to $V^\flat$ without gradient terms. The remainder
$W_j$ is carried along unchanged by the absorption, since absorption changes only the
Gaussian. The field rescaling $\psi'\mapsto(1+\zeta_j)^{-1/2}\psi'$ acts on $W_j$ as on
any polynomial.

\subsection{Coefficient formulas}
\label{sec:pt-coef}
Write
\[
  M_x=\psi_x\bar\psi_x(\nabla\psi)_x(\nabla\bar\psi)_x,\qquad
  (\nabla\psi)_x(\nabla\bar\psi)_x=\tfrac12\sum_{e\in\mathcal E_d}(\nabla^e\psi)_x(\nabla^e\bar\psi)_x .
\]
For real $w\in\cT$ of finite range define the functionals
\[
  w^{(1)}=\hat w(0)=\sum_xw(0,x),\qquad
  \mathfrak l(w)=\sum_xw(0,x)\,((-\Delta)w)(x,0)=\avint\lambda\,\hat w^2\,dp,
\]
and $\mathfrak B(w)$ by
\[
  \tfrac12\Loc_xF_w(M_x,M(\Lambda))=-\mathfrak B(w)\,M_x+(\text{monomials of degree}\le2).
\]
$\mathfrak B(w)$ is well defined, since $M$ is the only quartic monomial in $\mathcal V$.

\begin{lemma}[First order]\label{lem:first}
For $\Gamma\in\cT_\C$, up to a constant,
\[
  \E_\Gamma\theta M_x=M_x-\Gamma(0,0)(\nabla\psi\nabla\bar\psi)_x+\Delta\Gamma(0)\,\psi_x\bar\psi_x
  -\frac{\Delta\Gamma(0)}{2d}\cdot\frac12\bigl((-\Delta\psi)_x\bar\psi_x+\psi_x(-\Delta\bar\psi)_x\bigr),
\]
where $\Delta\Gamma(0)=\sum_e(\Gamma(0,e)-\Gamma(0,0))$. On the torus the total gradient
coefficient ($y+z$) generated by $gM$ is therefore $-g\,\Gamma(0,e_1)$, and the mass shift is
$g\,\Delta\Gamma(0)$.
\end{lemma}

\begin{proof}
This is the Wick computation in the proof of \texttt{thm:flow} of \hr{}. For the
total gradient coefficient, use $-\Gamma(0,0)+(\Gamma(0,0)-\Gamma(0,e_1))=-\Gamma(0,e_1)$.
\end{proof}

\begin{lemma}[Second order]\label{lem:second}
Let $w\in\cT$ have finite range, and let $\Gamma\in\cT$.
\begin{enumerate}[label=(\roman*)]
\item $\Loc_xF_w\bigl((\psi\bar\psi)(\Lambda),M_x\bigr)=4w^{(1)}M_x+(\text{degree}\le2)$.
\item For an observable $\sigma_a\bar\psi_a\psi_a$ with commuting $\sigma_a$,
  $\Loc_aF_w(\sigma_a\bar\psi_a\psi_a,M(\Lambda))=-\mathfrak l(w)\,\sigma_a\bar\psi_a\psi_a+(\text{constant})$.
\item For an observable $\sigma_a\bar\psi_a$ with anticommuting $\sigma_a$,
  $F_w(\sigma_a\bar\psi_a,M(\Lambda))$ has no component in $\C\sigma_a\bar\psi_a$, and
  $\Loc_aF_w(\sigma_a\bar\psi_a,(\psi\bar\psi)(\Lambda))=w^{(1)}\sigma_a\bar\psi_a$.
\end{enumerate}
Consequently, let $V=\nu\psi\bar\psi+gM-\mu^{(1)}_a\sigma_a\bar\psi_a-\mu^{(2)}_a\sigma'_a\bar\psi_a\psi_a$,
with $W$ as in Lemma~\ref{lem:frozen} for the accumulated covariance $w$. Write
$w_+=w+C'$, $\nu'=\nu+g\Delta C'(0)$, and $\mathfrak n=\nu'w_+^{(1)}-\nu w^{(1)}$. Then
\begin{align*}
  g_{\rm pt}&=g+g^2\bigl(\mathfrak B(w_+)-\mathfrak B(w)\bigr)-4g\,\mathfrak n,
  &(y+z)_{\rm pt}&=-g\,C'(0,e_1)+O_2,\\
  \mu^{(2)}_{\rm pt}&=\mu^{(2)}\bigl(1+g(\mathfrak l(w_+)-\mathfrak l(w))-2\mathfrak n\bigr),
  &\mu^{(1)}_{\rm pt}&=\mu^{(1)}(1-\mathfrak n)+O_2,
\end{align*}
up to terms of third order in $(g,\nu w^{(1)},\mu)$. Here $O_2$ denotes explicit second-order
terms $-\frac12\Loc[F_{w_+}(V',V')-\E\theta F_w(V,V)]$ in the respective monomial. They enter
only the two-loop coefficients.
\end{lemma}

\begin{proof}
Apply Lemma~\ref{lem:frozen} and read off the coefficients. Each statement is a single Wick
contraction (for (i) and (iii)) or a double contraction (for (ii) and $\mathfrak B$) between
two local monomials. (i) and the second part of (iii) express the substitution
$\psi\mapsto(1-\nu w*)\psi$ on each of the four legs of $M$, respectively on the single leg
$\bar\psi_a$; at zero momentum $w*$ acts as $w^{(1)}$. In (iii), $\sigma_a\bar\psi_a$ has a
single field, so a nonzero $F_w$ requires exactly one line. The remaining cubic polynomial has
no $\sigma_a\bar\psi_a$ component. The signs were checked by computer algebra
(\texttt{scripts/pt-checks/masscross.py}, \texttt{obs.py}).
\end{proof}

\subsection{Values in two dimensions}
\label{sec:pt-d2}

\begin{proposition}\label{prop:d2}
Let $d=2$, and let $(C_k)$ be the decomposition of \texttt{sec:frd}. Uniformly in $N$, in
$0<m^2\le L^{-2(N-1)}$, and in $|\zeta|_{\le j}\le\bar\zeta$, for $j+1<N$:
\begin{align*}
  C'_{j+1}(0,0),\ C'_{j+1}(0,e_1)&=\tfrac1{2\pi}\log L+O(L^{-j}+m^2L^{2j}+|\zeta|_{\le j}),\\
  \mathfrak l(w'_{j+1})-\mathfrak l(w'_j)&=\tfrac1{2\pi}\log L+O(L^{-j}+m^2L^{2j}+|\zeta|_{\le j}),\\
  \mathfrak B(w'_{j+1})-\mathfrak B(w'_j)&=\tfrac1{2\pi}\log L+O(L^{-j}+m^2L^{2j}+|\zeta|_{\le j}).
\end{align*}
\end{proposition}

\begin{proof}
First let $\zeta=0$, so that $C'_k=C_k$ and $w'_j=C_{\le j}$. By \frd{}, Appendix~A,
$\hat e_j(p)=1-G(\frac12L^j,\hat u(p))$, and $G(T,u)=\sum_i\varpi_iG^0(T/2^i,u)$ with
F-(A.4). For small $u$, $u/(c_\psi x^2)=c_0^{-1}(1+O(u))$, and $x=x(\hat u(p))$ is
$|p|/(2\sqrt2\,\pi)\,(1+O(|p|^2+m^2))$. So $\hat e_j(p)=E(L^j|p|)+\epsilon_j(p)$, with the
scaling profile
\[
  E(y)=1-\sum_i\varpi_i\,c_0^{-1}\Psi_1\bigl(y/(2^{i+1}2\sqrt2\,\pi)\bigr),
\]
which satisfies $E(y)=O(y^2)$ at $0$, $E(\infty)=1$, and has rapidly decaying $yE'(y)$. The
error $\epsilon_j$ is $O(|p|^2+m^2)$ times a function of $L^j|p|$ that decays rapidly. The
terms with $T/2^i\le2$, which carry the weights $\varpi_i\le\varpi_{J}$ with $2^J\asymp L^j$, are
smaller than any power of $L^{-j}$. In the continuum approximation $\avint dp\to\int d^2p/(2\pi)^2$
and $\lambda(p)\to|p|^2$, the three quantities are
\[
  \frac1{4\pi}\int_0^\infty\frac{ds}{s}\bigl(E_{j+1}-E_j\bigr),\qquad
  \frac1{4\pi}\int_0^\infty\frac{ds}{s}\bigl(E_{j+1}^2-E_j^2\bigr),\qquad
  \frac1{4\pi}\int_0^\infty\frac{ds}{s}\bigl(\Phi_{j+1}^2-\Phi_j^2\bigr),
\]
with $s=|p|^2$, $E_j=E(L^j\sqrt s)$ and $\Phi_j=\Phi(L^j\sqrt s)$, where $\Phi(y)=E(y)-\frac12yE'(y)$.
For $\mathfrak B$ this uses the continuum bubble formula: for a radial
$\hat w(p)=f(|p|^2)$, the $M$-coefficient of $-\frac12F_w(M,M)$ has density
$-f(s)\,(s^3f'(s))'$ in $\int d^2p/(2\pi)^2$. This follows from the double Wick contraction
(computer-algebra check: \texttt{scripts/pt-checks/general.py}). After an integration by
parts, the boundary terms at $s=\infty$ cancel in the difference, and
$s^3f'^2=s^{-1}\Phi^2$. By Frullani's formula,
$\int_0^\infty\frac{ds}{s}(h(L^{j+1}\sqrt s)-h(L^j\sqrt s))=2\log L\,(h(\infty)-h(0))$. Here
$h(\infty)=1$ and $h(0)=0$ for $h=E,E^2,\Phi^2$, since $E(y)\sim cy^2$ gives $\Phi(0)=0$. Each
quantity is therefore $\frac{2\log L}{4\pi}=\frac{\log L}{2\pi}$.
\begin{itemize}
\item \emph{Lattice and mass corrections.} These are localised at $|p|\asymp L^{-j}$ and give
  $O(L^{-j}+m^2L^{2j})$.
\item \emph{The dressing.} $C'_{j+1}-C_{j+1}=O(|\zeta|_{\le j})$ in $\|\cdot\|_{(j)}$ by
  Theorem~F-3.1(iv). The factors $Z_{j+1}/Z_k=1+O((j+1-k)\bar\zeta)$ in $w'_j$
  multiply $\hat C_k$, which at $|p|\asymp L^{-j}$ is $O(L^{-2(j-k)})$ relative to $\hat C_j$.
  So $w'_j$ differs from $C_{\le j}$ by $O(|\zeta|_{\le j})$ in the region that contributes.
\end{itemize}
\end{proof}

\paragraph{Formalisation (continuum part; Lean: \texttt{H02UV/Step2/D2/D2Frullani.lean}, \texttt{H02UV/Step2/D2/D2Continuum.lean}).}
A \emph{profile} is a function $h$ that is continuous and bounded on $(0,\infty)$, with
$|h-h_0|\le Ky^\alpha$ for $y\le1$ and $|h-h_\infty|\le Ky^{-\alpha}$ for $y\ge1$.
\begin{itemize}
\item \emph{Frullani} (\texttt{Profile.frullani}, \texttt{frullani\_trunc}). For profiles,
  $\int_0^\infty(h(bt)-h(at))\,dt/t=(h_\infty-h_0)\log(b/a)$. The truncation error is at most
  $(\sup_{[a\varepsilon,b\varepsilon]}|h-h_0|+\sup_{[aR,bR]}|h-h_\infty|)\log(b/a)$. The same holds in
  the variable $s=t^2$.
\item \emph{The three continuum values} $\frac{\log L}{2\pi}$:
  \begin{itemize}
  \item \texttt{continuum\_C}: $\int\frac{d^2p}{(2\pi)^2}(E_{j+1}-E_j)/|p|^2$;
  \item \texttt{continuum\_l}: $\int\frac{d^2p}{(2\pi)^2}|p|^2(\hat w_{j+1}^2-\hat w_j^2)$;
  \item \texttt{continuum\_B}: the bubble density $-f(s^3f')'$. It goes through the integration by
    parts \texttt{bubble\_ibp},
    $-\int_\varepsilon^Rf_c(s^3f_c')'=[E\Phi]_{c\sqrt\varepsilon}^{c\sqrt R}+\int_\varepsilon^R\Phi^2ds/s$,
    and the value is a limit of truncations, since the single-scale terms diverge.
  \end{itemize}
\end{itemize}
\paragraph{Formalisation (lattice and torus; Lean: \texttt{H02UV/Step2/D2/D2Lattice.lean}, \texttt{H02UV/Step2/D2/D2Torus.lean}).}
Let $\lambda(p)=\sum_i(2-2\cos p_i)$, $Q=[-\pi,\pi]^2$, $I_\alpha=\int_{\R^2}(1+|p|)^{-\alpha}d^2p$.
\begin{itemize}
\item \emph{Lattice$\to$continuum} (\texttt{lattice\_radial}). If $|h(y)|+|yh'(y)|\le K(1+y)^{-\alpha}$ with
  $\alpha>2$ and $h(y)/y\in L^1(0,\infty)$, then for $c>0$,
  $\bigl|\int_Q\frac{h(c\sqrt\lambda)}{\lambda}-\int_{\R^2}\frac{h(c|p|)}{|p|^2}\bigr|\le
  (\frac{\pi^6}{256}+\frac1{\pi^2})\frac{I_\alpha K}{c^2}$. The proof uses $0\le|p|^2-\lambda\le|p|^4/4$ and
  $\lambda\ge4|p|^2/\pi^2$ on $Q$.
\item \emph{Gradient direction} (\texttt{lattice\_cos}, \texttt{continuum\_cos}): the integrals
  weighted by $1-\cos p_1$ are $O(I_\alpha Kc^{-2})$.
\item \emph{Mass} (\texttt{lattice\_mass}): if $|g'(t)|\le Mc^4(1+c\sqrt t)^{-\beta}$ with $\beta>2$, then
  $|\int_Q(g(\lambda+m^2)-g(\lambda))|\le\frac{\pi^2}{4}I_\beta Mm^2c^2$.
\item \emph{Values} (\texttt{lattice\_C}, \texttt{lattice\_l}, \texttt{lattice\_Ce1}): for a profile $E$
  with $|1-E|+|yE'|\le K(1+y)^{-\alpha}$, $\alpha>2$, and $\hat e_k=E(L^k\sqrt\lambda)$, the lattice
  values $\avint_Q\hat C_{j+1}$, $\avint_Q\cos p_1\,\hat C_{j+1}$ and
  $\mathfrak l(w_{j+1})-\mathfrak l(w_j)$ are $\frac{\log L}{2\pi}+O_{K,\alpha}(L^{-2j})$.
\item \emph{Torus $=$ Brillouin zone} (\texttt{Cj\_eq\_integral\_BZ}): if $r\in\Z^2$ lifts $x-y$
  and $L^j/2+|r|_\infty<n$, then $C_j(x,y)=\int_Q\hat c_j(p)\cos(p\cdot r)\frac{d^2p}{(2\pi)^2}$ exactly.
  The symbol is $\hat c_{(a,b]}=(\hat e_b-\hat e_a)/(\lambda+m^2)$ (\texttt{chat\_eq\_ehat}).
\item \emph{Combined} (\texttt{Cj\_succ\_diag\_sub\_log}): $|C_{j+1}(0,0)-\frac{\log L}{2\pi}|\le
  CL^{-2j}+\int_Q|\hat c_{j+1}-(E_{j+1}-E_j)/\lambda|\frac{d^2p}{(2\pi)^2}$.
\end{itemize}
\paragraph{Formalisation (profile, values, dressing; Lean: \texttt{H02UV/Step2/D2/D2Profile.lean},
\texttt{D2Ehat}, \texttt{D2Cov}, \texttt{D2Dress}, \texttt{D2Lfun}, \texttt{D2Lval}).}
This is for the concrete family $P_c=$ \texttt{FRD.PtF.prePt.toPtFamily}.
\begin{itemize}
\item \emph{Exact form of $\hat e_T$} (\texttt{ehat\_repr}). For $u=\hat u(p)\in(0,2]$,
  $x=\frac1{2\pi}\arccos(1-u)$ and $T>0$:
  $\bigl|\hat e_T(p)-1+\frac{\sin^2\pi x}{(\pi x)^2}(1-H(Tx))\bigr|\le Cu\,T^{-2}$, where
  $H(z)=1-\sum_i\varpi_ic_0^{-1}\Psi_1(z/2^i)$ and $\Psi_1(z)=\int_z^\infty r\psi(r)\,dr$.
\item \emph{The profile} (\texttt{Ep\_profile}, \texttt{Ep\_decay}, \texttt{Phi\_Ep\_le}). Let
  $E(y)=H(y/(4\sqrt2\pi))$. Then $E$ is a profile from $0$ to $1$,
  $|1-E|+|yE'|\le K(1+y)^{-3}$, and $|E-\frac12yE'|\le My^4(1+y)^{-4}$.
\item \emph{(R1)} (\texttt{chat\_profile\_error}). For $0\le m^2\le1$:
  $\int_Q|\hat c_{j+1}-(E_{j+1}-E_j)/\lambda|\le AL^{-2j}+Bm^2(L^{2j}+L^{2(j+1)})$.
\item \emph{Values} (\texttt{Cj\_succ\_diag\_value}, \texttt{Cj\_succ\_e1\_value}). For $j\ge1$:
  $C_{j+1}(0,0),\,C_{j+1}(0,e_1)=\frac{\log L}{2\pi}+O(L^{-2j}+m^2L^{2j}(1+L^2))$.
\item \emph{Dressed} (\texttt{Cprime\_values}). For $L\ge7$ and $|\zeta|\le\bar\zeta$:
  $C'_{j+1}(0,0),\,C'_{j+1}(0,e_1)=\frac{\log L}{2\pi}+O(L^{-2j}+m^2L^{2j}(1+L^2)+|\zeta|_{\le j})$.
  This uses \texttt{Concl.close} and, in $d=2$, $|w(x,0)|\le\|w\|_{(j)}$ (\texttt{norm\_kern\_le\_snorm}).
\item \emph{$\mathfrak l$} (\texttt{lfun\_cov\_eq\_integral\_BZ}, \texttt{lfun\_value\_sum}). If $2T+1<n$, then
  $\mathfrak l(C_{(0,T]})=\int_Q\lambda\hat c_{(0,T]}^2\frac{d^2p}{(2\pi)^2}$ exactly, by Parseval. Hence
  $\mathfrak l(w_{j+1})-\mathfrak l(w_j)=\frac{\log L}{2\pi}+O(L^{-2j}+m^2L^{2j}(1+L^2))$ for $j\ge1$.
\end{itemize}
All constants are independent of $L$, $j$, $m$ and $n$. The lattice error is $O(L^{-2j})$, better
than the $O(L^{-j})$ of the statement. Not yet formalised:
\begin{itemize}
\item the momentum form of $\mathfrak B$ from \texttt{FlowSeq.bubble\_explicit}, and its value;
\item the dressing of $\mathfrak l$ and $\mathfrak B$, which needs the localised argument at
  $|p|\asymp L^{-j}$; the bound from \texttt{Concl.close} alone is not uniform in $j$;
\item the assembled statement.
\end{itemize}
\emph{Remark on the proof.} For $j+1<N$ the torus average $\avint$ can be replaced
\emph{exactly} by $\int_{[-\pi,\pi]^2}$. The reason is finite range: $C_{j+1}$ and $w_j$ are
polynomials in $\Delta$ of range $<L^N/2$, so their torus kernels equal those on $\Z^2$. Without
this, the Riemann-sum error of the torus is not uniform in $N$ for $j$ close to $N$.

\begin{remark}
The three values in Proposition~\ref{prop:d2} are universal: they do not depend on the profile
$E$. Individual single-scale sums, such as $\mathfrak l(C'_{j+1})$, are not universal; only
the telescoped differences are. This is why the second-order map must carry $W$,
rather than use $\frac12F_{C'_{j+1}}$ alone.
\end{remark}

\begin{corollary}[One-loop flow in $d=2$]\label{cor:oneloop}
In the setting of Lemma~\ref{lem:second}, with $C'=C'_{j+1}$ and $w=w'_j$, absorb the new
gradient coefficient $\zeta_{j+1}=(y+z)_{\rm pt}$ and renormalise, so that
$g_{j+1}=g_{\rm pt}/(1+\zeta_{j+1})^2$. Then
\[
  g_{j+1}=g_j+\frac{3\log L}{2\pi}\,g_j^2-4g_j\mathfrak n_j+O\bigl(g_j^3+g_j^2(L^{-j}+|\zeta|_{\le j})\bigr),
  \qquad \zeta_{j+1}=-\frac{\log L}{2\pi}g_j+O(g_j^2+g_jL^{-j}).
\]
For $\beta_j=1/g_j$ this gives $\beta_{j+1}=\beta_j-\beta_1\log L+O(\beta_j^{-1})$ with
$\beta_1=\frac3{2\pi}$, up to the telescoping term $4\mathfrak n_j/g_j$. That term is removed by
the change of variables $\check g_j=g_j(1+4\nu_jw'^{(1)}_j)$, as in
Bauerschmidt--Brydges--Slade. The field-strength part contributes $\frac2{2\pi}$ and the bubble
$\frac1{2\pi}$.
\end{corollary}

\begin{proof}
$g_{\rm pt}=g+g^2\frac{\log L}{2\pi}-4g\mathfrak n$ by Lemma~\ref{lem:second} and
Proposition~\ref{prop:d2}. Also $(1+\zeta_{j+1})^{-2}=1+2g\frac{\log L}{2\pi}+O(g^2)$.
\end{proof}

%======================================================================
\section{The renormalisation group map}
\label{sec:rgmap}

\subsection{Coordinates}
At scale $j$, in the fields $\psi^{(j)}$, the progressive integral is written as
\begin{equation}\label{e:coords}
  \cZ_j=e^{-u_j|\Lambda|}\sum_{X\in\cP_j}I_j(\Lambda\setminus X)\,K_j(X),\qquad
  I_j(Y)=\prod_{B\in\cB_j(Y)}e^{-V_j(B)}\bigl(1+W_j(B)\bigr).
\end{equation}
The remaining Gaussian has covariance $Z_j\Gamma_j$.
\begin{itemize}
\item $V_j=(z_j,y_j,\nu_j,g_j)\in\cV$.
\item $W_j(B)=\sum_{x\in B}W_j(x)$ is the second-order remainder of Section~\ref{sec:pt-map}. Each
  $W_j(x)$ depends only on fields within distance $\frac12L^j+1$ of $x$, so
  $W_j(B)\in\cN(B^\square)$.
\item $K_j\in\mathcal K_j$ is as in \texttt{def:cKbulk}.
\end{itemize}
The initial data are $V_0$ from \eqref{e:V0b}, $W_0=0$, $K_0=1_{X=\varnothing}$, $Z_0=1$ and
$\Gamma_0=A^{-1}$.

\subsection{One step}
\label{sec:onestep}
The map $(V_j,W_j,K_j)\mapsto(u_{j+1},V_{j+1},W_{j+1},K_{j+1})$ and $Z_j\mapsto Z_{j+1}$
consists of four parts.

\paragraph{(A) Absorption}
Set $\zeta_j=y_j+z_j$ and $Z_{j+1}=Z_j(1+\zeta_j)$. Let $V_j^\flat=(0,0,\nu_j-\zeta_jm^2,g_j)$,
so that $(V_j-V_j^\flat)(\Lambda)=\zeta_j(\psi^{(j)},A\bar\psi^{(j)})$. Since $V_j-V_j^\flat$ is
local,
\[
  I_j(\Lambda\setminus X)=e^{-\zeta_j(\psi^{(j)},A\bar\psi^{(j)})}\;e^{(V_j-V^\flat_j)(X)}\;I^\flat_j(\Lambda\setminus X),
\]
where $I^\flat_j$ is $I_j$ with $V^\flat_j$. Put $K^\flat_j(X)=e^{(V_j-V_j^\flat)(X)}K_j(X)$,
which again has the factorisation property. Apply Lemma~\ref{lem:onestep} to the remaining Gaussian. By
\ref{D1}, it becomes the Gaussian with covariance $\Gamma^{\rm int}_{j+1}+\Gamma_{j+1}$, in bare
fields. The determinant goes into $u_{j+1}$, and $W_j$ is unchanged.

\paragraph{(B) Normalisation}
Rewrite everything in $\psi^{(j+1)}=(1+\zeta_j)^{1/2}\psi^{(j)}$. A monomial of degree $2n$ is
multiplied by $(1+\zeta_j)^{-n}$:
\[
  \nu^\flat=\frac{\nu_j-\zeta_jm^2}{1+\zeta_j},\qquad g^\flat=\frac{g_j}{(1+\zeta_j)^2},
\]
and similarly for $W_j$, $K^\flat_j$ and the observables (Section~\ref{sec:obs-flow}). In
$\psi^{(j+1)}$ the integrated covariance is $C'_{j+1}$.

\paragraph{(C) Second-order potential}
Define $Q(B)$ and $J(B,X)$ from $K^\flat_j$ by \texttt{e:Q-def}--\texttt{e:J-def} of \hr{}. Then,
with Definition~\ref{def:Vpt} applied to $(V^\flat,W^\flat_j,C'_{j+1})$,
\begin{align}
  u_{j+1}|B|+V_{j+1}(B)&=V_{\rm pt}(B)-\E_{C'_{j+1}}\theta Q(B),\label{e:V+}\\
  W_{j+1}(x)&=W_{\rm pt}(x).\label{e:W+}
\end{align}
Explicitly, with $V'=\E_{C'}\theta V^\flat$,
\begin{align*}
  V_{\rm pt}(x)&=V'_x-\Loc_x\bigl[\tfrac12F_{C'}(V'_x,V'(\Lambda))+\E_{C'}\theta\,W^\flat_j(x)\bigr],\\
  W_{\rm pt}(x)&=(1-\Loc_x)\bigl[\tfrac12F_{C'}(V'_x,V'(\Lambda))+\E_{C'}\theta\,W^\flat_j(x)\bigr],
\end{align*}
where $F_C(A,B)=e^{\cL_C}\bigl((e^{-\cL_C}A)(e^{-\cL_C}B)\bigr)-AB$. If $j+1=N$, set $Q=J=0$.

\paragraph{(D) Polymer activity}
$K_{j+1}$ is defined by \texttt{def:K+} with the following replacements:
\begin{itemize}
\item $e^{-V_j(B)}$ by $I^\flat_j(B)=e^{-V^\flat(B)}(1+W^\flat_j(B))$;
\item $e^{-(u+V)_{j+1}(B)}$ by $e^{-u_{j+1}|B|}e^{-V_{j+1}(B)}(1+W_{j+1}(B))$;
\item $K_j$ by $K^\flat_j$, and $C_{j+1}$ by $C'_{j+1}$.
\end{itemize}

\begin{proposition}\label{prop:ZVK}
With these definitions,
$\cZ_{j+1}=\E_{\Gamma^{\rm int}_{j+1}}\theta\,\cZ_j^{(A)}$, where $\cZ_j^{(A)}$ is $\cZ_j$ after
the absorption (A). The coordinates satisfy \eqref{e:coords} at scale $j+1$, and $K_{j+1}$ has
the factorisation property and is automorphism invariant.
\end{proposition}

\begin{proof}
For (A) and (B), use \eqref{e:absorb} and the change of variables. For (C) and (D), the proof
of \texttt{prop:ZVK} applies verbatim: it only uses the block-product structure of $I$, the
identity $\sum_XJ(B,X)=0$, and $\E\theta$ as a homomorphism. The finite range of
$\Gamma^{\rm int}_{j+1}$ (\ref{D2}) is used as before.
\end{proof}

\subsection{Consistency}
The flows on tori $\Lambda_{N_1}\subset\Lambda_{N_2}$ with the same lattice spacing coincide
up to scale $N_1-1$ for polymers that do not wrap. This holds as in \texttt{prop:consistency},
since $C'_{j+1}$ depends only on $(\zeta_i)_{i\le j}$ and not on $N$ (for $j+1<N$). The
continuum limit uses instead the comparison of different lattice spacings
(Section~\ref{sec:uv}).

%======================================================================
\section{Estimates for the renormalisation group map}
\label{sec:est}

The norms $\|V\|_j$ and $\|K\|_j$ are \texttt{e:V-norm}--\texttt{e:K-norm} of \hr{} with
$\ell_j=\ell_0$. In $d=2$, $\|V(B)\|_{T_j(\ell_0)}\asymp|z|+|y|+L^{2j}|\nu|+|g|$. Let
$\bar g_j=\max(|g_j|,L^{2j}|\nu_j|,|\zeta_j|/\log L)$.

\paragraph{The activity $K_{j+1}$}
Let $K^{\rm pf}_{j+1}$ be the output of the polymer formula (D) (Brydges--Slade's $K^{(3)}$).
With $V'_B=\E_{C'}\theta V(B)$ and $F=F_{C'}$, its second-order part is
\begin{equation}\label{e:K2}
  K^{(2)}(U)=\sum_{B:\,\bar B=U}\Bigl(-\tfrac12F\bigl(V'_B,V'(\Lambda\setminus B)\bigr)\Bigr)
  +\sum_{\{B,B'\}:\,\overline{B\cup B'}=U}F(V'_B,V'_{B'}) ,
\end{equation}
i.e.\ $K^{\rm pf}_{j+1}=K^{(2)}+O(t^3)$ along $(tV,t^2W,0)$. It is nonzero
(Brydges--Slade's $h_{\rm lead}$, \emph{A single RG step} (5.18)), but $\sum_UK^{(2)}(U)=0$. Let
$\tilde I_{j+1}(B)=e^{-V_{j+1}(B)}(1+W_{j+1}(B))$, so that
$I_{j+1}(B)=e^{-u_{j+1}(B)}\tilde I_{j+1}(B)$. For a $(j+1)$-block $B_+$ let
$\tilde I_{j+1}(B_+)=\prod_{B\subset B_+}\tilde I_{j+1}(B)$. Split $K^{(2)}$ over the $(j+1)$-blocks:
\[
  J(B_+,U)=\sum_{X\in\cP_j:\,\bar X=U}\kappa(X,B_+)\quad\text{if $U$ is small,}\qquad J(B_+,U)=0
  \text{ otherwise},
\]
where $\kappa(X,B_+)$ is the part of the summand of \eqref{e:K2} for $X$ that comes from the
$j$-blocks of $X$ in $B_+$. Then $\sum_{B_+\subset U}J(B_+,U)=K^{(2)}(U)$ for connected $U$, and
$\sum_{U\supset B_+}J(B_+,U)=0$. The second identity uses that $F(V'_B,V'_{B'})=0$ when $B$ and
$B'$ lie in non-touching $(j+1)$-blocks. Put $\bar J(B_+,U)=\tilde I^U_{j+1}J(B_+,U)$ and
$M(U)=K^{\rm pf}_{j+1}(U)-\tilde I^U_{j+1}K^{(2)}(U)$.

The activity of the step is the output of the change of variables of Brydges--Slade
(Prop.~4.1, (D.13); their Map~4), at scale $j+1$:
\[
  K_{j+1}(W)=\sum_{(\mathcal X,U_M)}\tilde I_{j+1}^{W\setminus(U_M\cup X_{\mathcal X})}
  \prod_{Y\in{\rm Comp}(U_M)}M(Y)\prod_{(B_+,U)\in\mathcal X}\bar J(B_+,U).
\]
The sum runs over sets $\mathcal X$ of pairs ($U$ small, $U\supset B_+$, distinct and pairwise
non-touching $U$'s) and $(j+1)$-polymers $U_M$ not touching $X_{\mathcal X}=\bigcup U$, such that
$U_M\cup\bigcup_{(B_+,U)\in\mathcal X}B_+^\square=W$. For connected $W$,
$K_{j+1}(W)=M(W)+E(W)$, where $E(W)$ collects the terms with two or more small factors; the terms
with one pair and $U_M=\emptyset$ cancel. The intermediate activities $K^{\rm pf}_{j+1}$ and $M$
appear only in the proof.

\begin{theorem}[Step estimates]\label{thm:step}
Let $d=2$, $j+1<N$, and fix a field scale $\ell>0$ (the same at scales $j$ and $j+1$). Let
$V=(V(B))_B$ and $W=(W(B))_B$ be families of even, $U(1)$-invariant elements indexed by the
$j$-blocks, with $V(B),W(B)\in\cN(B^\square)$ and $\epsilon(V(B))=0$, and let
$K_j\in\mathcal K_j$ (locality and $U(1)$ invariance on connected polymers). Let $C'=C'_{j+1}$
be a complex covariance with a Gram representation satisfying the Gram bound at scale $j+1$ and
field scale $\ell$, and with range $<\frac12L^{j+1}$. Write $\|V\|_j=\sup_B\|V(B)\|_{T_j(\ell)}$,
and let $(u_{j+1},V_{j+1},W_{j+1},K_{j+1})$ be the map of the Lean formalisation below, with
$K_{j+1}$ as above. For a $(j+1)$-block $B_+$ put
$V_{j+1}(B_+)=\sum_{B\subset B_+}V_{j+1}(B)$ and
$W_{j+1}(B_+)=\prod_{B\subset B_+}(1+W_{j+1}(B))-1$.
\begin{enumerate}
\item[(i)] (Exactness.) $\E_{C'}\theta Z_j=Z_{j+1}$, where
  $Z_j=\sum_XI_j^{\Lambda\setminus X}\prod_{Y\in{\rm Comp}(X)}K_j(Y)$,
  $I_j(B)=e^{-V(B)}(1+W(B))$, and
  \[
    Z_{j+1}=\sum_{U\in\cP_{j+1}}I_{j+1}^{\Lambda\setminus U}e^{-u_{j+1}(U)}
    \prod_{Y\in{\rm Comp}(U)}K_{j+1}(Y).
  \]
  So the next step only uses $K_{j+1}$ on connected polymers.
\item[(ii)] (Coupling constants.) If $L\ge L_0$, then for every $j$-block $B$
  \[
    \bigl\|(u+V)_{j+1}(B)-V'_B+\Loc_{B_+,B}S(B_+)\bigr\|_{T_{j+1}(\ell)}
    =\|\E_{C'}\theta Q(B)\|_{T_{j+1}(\ell)}\le C\|K_j\|_j .
  \]
\item[(iii)] (Polymer activity.) There are $C_0,\eta>0$ and $L_0$ such that for $L\ge L_0$
  there is $A_0(L)$, and for $A\ge A_0$ there are $\epsilon_0(L,A)>0$ and $C(L,A)$, such that
  the following holds. If $\|V\|_j\le\bar g$, $\|W\|_j\le\bar g^2$, $\bar g\le\epsilon_0$ and
  $\|K_j\|_j\le\epsilon_0$, then for every connected $(j+1)$-polymer $U$
  \[
    A^{(|\cB_{j+1}(U)|-4)_+}\|K_{j+1}(U)\|_{T_{j+1}(\ell)}\le
    C_0(L^{-1}+A^{-\eta})\|K_j\|_j+C(L,A)\bigl(\bar g^3+\bar g\|K_j\|_j+\|K_j\|_j^2\bigr).
  \]
  In particular, if $\|K_j\|_j\le\bar g^3$ then
  $\|K_{j+1}\|_{j+1}\le O(L^{-1}+A^{-\eta})\|K_j\|_j+C(L,A)\,\bar g^3$.
\item[(iv)] (Closure.) If $L\ge L_0$, then $K_{j+1}\in\mathcal K_{j+1}$: for every connected
  $(j+1)$-polymer $U$, $K_{j+1}(U)\in\cN(U^\square)$ (at scale $j+1$) is $U(1)$-invariant.
  Moreover $V_{j+1}$ and $W_{j+1}$ are families as in the hypotheses, at scale $j+1$:
  $V_{j+1}(B_+),W_{j+1}(B_+)\in\cN((B_+)^\square)$ are even and $U(1)$-invariant, and
  $\epsilon(V_{j+1}(B_+))=0$. If $\|V\|_j\le v$, $\|W\|_j\le w$ and $\|K_j\|_j\le k$ with
  $v,w,k\le1$, then
  \[
    \|V_{j+1}\|_{j+1}\le C(L)(v+w+k),\qquad\|W_{j+1}\|_{j+1}\le C(L)(v^2+w).
  \]
\end{enumerate}
(Lean, package \texttt{Step2}: (i) \texttt{stepK\_ZVK}, (ii) \texttt{step2\_V}, (iii)
\texttt{stepK\_d2}, and \texttt{stepK} for every $d\ge2$ with $L^{d-3}$ in place of $L^{-1}$, (iv)
\texttt{InK2\_Kplus}, \texttt{Vnext2\_inF}, \texttt{Wnext2\_inF}, \texttt{$\varepsilon$\_Vnext2},
\texttt{coords\_bound}. Here $K_{j+1}$ is \texttt{Step2.Kplus}, $K^{\rm pf}_{j+1}$ is
\texttt{Step2.Knext2}, and $M$ is \texttt{Step2.Knew}.)
\end{theorem}


\emph{Proof of (iii).} The proof has four steps.
\begin{itemize}
\item By Cauchy estimates (Lean formalisation below),
  $A^{(|\cB_{j+1}(U)|-4)_+}\|K^{\rm pf}_{j+1}(U)-K^{(2)}(U)\|\le
  C_0(L^{-1}+A^{-\eta})\|K_j\|_j+C(\bar g^3+\bar g\|K_j\|_j+\|K_j\|_j^2)$
  (\texttt{Step2.step2}).
\item With $\|\tilde I^U_{j+1}-1\|=O(\bar g+\|K_j\|_j)$ and $\|K^{(2)}(U)\|=O(\bar g^2)$, the same
  bound holds for $M$ (\texttt{Step2.step3}). It even holds with the stronger large-set regulator
  $\lambda A$ in place of $A$, for a fixed $\lambda=\lambda(d)$ (\texttt{Step2.step3L}); the slack
  in Lemma \texttt{lem:comb1} absorbs $\lambda^{|\cB_{j+1}(U)|}$.
\item Each term of $E(W)$ contains two small factors: two $\bar J=O(\bar g^2)$, or one $\bar J$
  and one $M$. The number of terms grows exponentially in $|\cB_{j+1}(W)|$; the stronger
  regulator on $M$ absorbs it. Hence
  $A^{(|\cB(W)|-4)_+}\|E(W)\|\le C(L,A)(\bar g^4+\bar g^2\|M\|)$ (\texttt{Step2.Kout\_sub\_le},
  as in Brydges--Slade (D.17)--(D.20), where $a_{\rm out}<a_{\rm in}$ plays the role of
  $\lambda$).
\item Combining the three steps gives (iii).
\end{itemize}

\emph{Proof of (i) and (iv).}
\begin{itemize}
\item The identity is (D.4) of Brydges--Slade, \texttt{Step2.reblock\_ZVK}. It applies because
  $K^{\rm pf}_{j+1}$ factorises over non-touching polymers.
\item The block factor $I_{j+1}(B)$ is not local at scale $j$, because $W_{j+1}(B)$ reaches the
  near blocks. It enters the expectation only through the inclusion, so it contains no
  fluctuation field $\zeta$. The expectation factorises over separated $\zeta$-supports
  (\texttt{Step2.ECC\_mul\_of\_sepZ}): if $F_i$ depends on all $\psi$ but only on $\zeta_x$,
  $x\in X_i$, and $C'$ vanishes between the disjoint sets $X_1,X_2$, then
  $\E_{C'}(F_1F_2)=\E_{C'}F_1\,\E_{C'}F_2$.
\item The factorisation of the output $K_{j+1}$ itself is \texttt{Step2.Kout\_union}.
\item Locality and symmetry follow from a membership principle for the polymer formula and for
  (D.13). It uses that $\E_{C'}$ maps $\cN(X\sqcup X)$ to $\cN(X)$ and commutes with the $U(1)$
  action. The range of $C'$ enters through $F(V'_B,V'(\Lambda))$: only blocks near $B$
  contribute, and they lie in $(B_+)^\square$.
\end{itemize}

\subsection{Contraction of the remainder $W$}
\label{sec:W-contract}
Theorem~\ref{thm:step}(iv) gives $\|W_{j+1}\|_{j+1}\le C(L)(v^2+w)$, whose constant grows
under iteration. The part of $W_{j+1}$ linear in $W$ is
$(1-\Loc_{B_+})\sum_{B\subset B_+}\E\theta\,W(B)$. It has the same form as the linear part of
$K_{j+1}$ on a single block, and for translation-covariant data it contracts by
Lemma~\ref{lem:loc2} (Proposition~\ref{prop:W-contract}). We keep the recursive $W$ of Definition~\ref{def:Vpt} rather than the
closed form $\frac12(1-\Loc)F_{w_j}(V,V)$ of Brydges--Slade. The closed form appears not to be
consistent to second order with the absorption (A). The absorbed gradient term
$G=\zeta_{j+1}(\psi,A\bar\psi)$ is first order in $\bar g$ (by \eqref{e:zflow}). So the closed form
at scale $j+1$, built from $V^\flat_{j+1}$, misses the cross terms
$\frac12(1-\Loc)[2F_{w'_{j+1}}(V'^\flat,G)+F_{w'_{j+1}}(G,G)]$, which are $O(\bar g^2\log L)$.
(This is a heuristic computation, not checked in detail.) The recursive form is exactly what
each step produces, whatever is absorbed. In $d=2$ the closed form would also not give a
uniform bound for free, since the Gram constant of $w'_j$ grows linearly in $j$.

\paragraph{Translation covariance}
For $a\in\Z^d$ let $\tau_a$ be the automorphism of $\cN$ induced by $x\mapsto x+a$. A block
family $F=(F(B))_B$ at scale $j$ is \emph{translation covariant} if
$F(B+a)=\tau_aF(B)$ for all $a\in L^j\Z^d$ (Lean \texttt{FamCovG}). The proposition below uses only
the weaker form $F(B)=\tau_{a_B}F(B_0)$, where $a_B$ is the corner of $B$ and $B_0$ the block at the
origin (\texttt{FamCov}; \texttt{FamCovG.toFamCov}). A polymer activity $K$ is translation
covariant if $K(X+a)=\tau_aK(X)$ for $a\in L^j\Z^d$ (\texttt{KCov}). A covariance is translation
invariant if $C'(x+a,y+a)=C'(x,y)$ for all $a\in\Z^d$ (\texttt{TInv}).

A polynomial on $X$ is \emph{uniform} if it has the form
$\kappa+\sum_Mc_MM(X)+\sum_{i,k}e_{ik}M^{(4)}_{ik}(X)$ with constants $\kappa,c_M,e_{ik}$, where
$M(X)=\sum_{x\in X}M_x$ runs over the monomials of $\Loc^{(2)}$ and $M^{(4)}_{ik}(X)$ over those of
$\Loc^{(4)}$ (\texttt{unif}).

\begin{lemma}[Localisation and translations; proved in Lean]\label{lem:loc-tr}
Let $L\ge L_0$, $j+1<N$, and let $X$ be a small $j$-polymer, $F\in\cN(X^\square)$.
\begin{enumerate}
\item[(i)] $\tau_a\Loc_{X,Y}F=\Loc_{X+a,Y+a}\tau_aF$ for all $a\in\Z^d$ and $Y$.
  The quartic pairings $\mu^{(4)}$ do not depend on the base point.
\item[(ii)] $\Loc_XF$ is uniform on $X$, with coefficients $\kappa=F_\varnothing$,
  $c=c(F)$ and $e_{ik}=\mu^{(4)}_{ik}(F)/|X|$.
\item[(iii)] For every patch $X'$, $\Loc_{X'}$ is the identity on uniform polynomials on $X'$.
\end{enumerate}
(Lean: \texttt{Step2.Loc2XY\_transl}, \texttt{Step2.mu4\_base\_indep}, \texttt{Step2.mu4\_transl},
\texttt{Step2.Loc2X\_eq\_unif}, \texttt{Step2.Loc2X\_unif}.)
\end{lemma}

\begin{proof}
(i) For $\Loc^{(2)}$ this is \texttt{prop:Loc}(iv). For $\Loc^{(4)}$: a shift of the base point
changes $\ell(x_0)_i\ell(x_1)_k$ by terms that do not depend on $x_0$ (resp.\ $x_1$). They
antisymmetrise to $0$ under $x_1\leftrightarrow x_3$ (resp.\ $x_0\leftrightarrow x_2$), the
positions of equal spin, so $\mu^{(4)}$ does not change. (ii) This holds by definition.
(iii) This is the moment identity: $A_{X'}$ is invertible, and
$\mu^{(4)}(M^{(4)}_{ik}(X'))=|X'|\delta$.
\end{proof}

\begin{proposition}[Contraction of $W$; proved in Lean]\label{prop:W-contract}
Let $d\ge2$. In the setting of Theorem~\ref{thm:step}, assume in addition that $W$ is
translation covariant at scale $j$ and that $C'$ is translation invariant. There are $C_0$
(independent of $L$) and $L_0$ such that, for $L\ge L_0$, there is $C(L)$ with the following
property. If $\|V\|_j\le v$ and $\|W\|_j\le w$ with $v,w\le1$, then
\[
  \|W_{j+1}\|_{j+1}\le C_0L^{d-3}w+C(L)\,(v^2+w^2),
\]
that is, $C_0L^{-1}w+C(L)(v^2+w^2)$ in $d=2$. Moreover, if $j+2<N$ and $V$, $W$, $K_j$ are
translation covariant at scale $j$, then $V_{j+1}$, $W_{j+1}$ and $K_{j+1}$ are translation
covariant at scale $j+1$. Here $C_0=(1+\bar C)(1+16\bar C)C_P$, with $\bar C$ the constant of
\texttt{prop:Loc}(i) and $C_P$ that of Lemma~\ref{lem:loc2}.

(Lean, package \texttt{Step2}: \texttt{W\_contract}, \texttt{W\_contract\_d2},
\texttt{Wnext2\_cov}, \texttt{Vnext2\_cov}, \texttt{Kplus\_cov}.)
\end{proposition}

\begin{proof}
For each $j$-block $B$ write $W(B)=A_B+P_B$ with $P_B=\Loc_BW(B)$, and put
$R_B=\E\theta P_B$. Expanding $\prod(1+W_{j+1}(B))-1$, we get
$W_{j+1}(B_+)=\sum_{B\subset B_+}W_{j+1}(B)+R$, where $\|R\|\le n^2\delta^2(1+\delta)^n$,
$n\le L^d$ and $\delta=O_L(v^2+w)$ (\texttt{tn\_prod\_sub\_one\_sub}). Since
$\sum_{B\subset B_+}\Loc_{B_+,B}=\Loc_{B_+}$,
\[
  \sum_{B\subset B_+}W_{j+1}(B)=(1-\Loc_{B_+})(S_V+X_1+U),
\]
where:
\begin{itemize}
\item $S_V=\sum_B\frac12F(V'_B,V'(\Lambda))$;
\item $X_1=\sum_B\bigl(\E\theta A_B+(1-\Loc_B)R_B\bigr)$;
\item $U=\sum_B\Loc_BR_B$.
\end{itemize}

\emph{The term $U$.} By covariance, $R_B$ and hence $\Loc_BR_B$ is the translate of the
corresponding object at $B_0$ (Lemma~\ref{lem:loc-tr}(i), and $\E\theta$ commutes with
$\tau_a$). By Lemma~\ref{lem:loc-tr}(ii), all $\Loc_BR_B$ are uniform with the same
coefficients. Hence $U$ is uniform on $B_+$, and $\Loc_{B_+}U=U$ by
Lemma~\ref{lem:loc-tr}(iii) (\texttt{loc\_U}).

\emph{The terms of $X_1$.} Both gain $L^{-3}$ by Lemma~\ref{lem:loc2}:
\begin{itemize}
\item $\|\E\theta A_B\|_{T_{j+1}(\ell)}\le\|A_B\|_{T_{j+1}(2\ell)}\le C_PL^{-3}w$, exactly as for
  the linear part of $K_{j+1}$;
\item $\|(1-\Loc_B)R_B\|_{T_{j+1}(\ell)}\le C_PL^{-3}\|R_B\|_{T_j(\ell)}$, with
  $\|R_B\|_{T_j(\ell)}\le\|P_B\|_{T_j(2\ell)}\le16\|P_B\|_{T_j(\ell)}\le16\bar Cw$. The first step
  is the Gram bound at scale $j$, which follows from the one at scale $j+1$. The second holds
  because $P_B$ has degree $\le4$ (\texttt{tNorm\_two\_le\_of\_deg}).
\end{itemize}
There are at most $L^d$ blocks, so $\|X_1\|\le L^dC_PL^{-3}(1+16\bar C)w$. Together with
$\|S_V\|\le C(L)v^2$ and $\|\Loc_{B_+}\|\le\bar C$ this gives the bound.

\emph{Covariance.} Covariance of the outputs follows from that of the inputs, since the
following commute with $\tau_a$ for $a\in L^{j+1}\Z^d$: $\Loc$ (Lemma~\ref{lem:loc-tr}(i)),
$F_{C'}$, $\E_{C'}\theta$, the polymer formula (\texttt{AE\_gKnext}), the splitting of $K^{(2)}$,
and the change of variables (\texttt{AE\_Kout}). The condition $j+2<N$ makes $(j+1)$-blocks
regions for $\Loc$ at scale $j+1$. It can be replaced by $j+1<N$ and $L>2^{d+4}+4$: the
neighbourhood of a small $k$-set has diameter $\le2^{d+2}L^k<L^N/4$ for every $k<N$
(\texttt{region\_small\_L}; the covariance statements for $j+1<N$ are in \texttt{H02UV/Step2/Last/Last*.lean}).
\end{proof}

Without translation covariance, the uniform term $U$ would be replaced by a sum of uniform
polynomials with block-dependent coefficients, whose size is of order $\|W\|_j$. The
contraction would then (heuristically) only hold over two steps.

\begin{corollary}[Uniform bound on $W$; proved in Lean, \texttt{Flow2.w\_uniform}]\label{cor:W-uniform}
Let $\theta\le\frac14$, $C>0$, $M\in\mathbb N$, and let $v_k,w_k\ge0$ satisfy, for $k<M$, $w_0=0$,
$w_{k+1}\le\theta w_k+C(v_k^2+w_k^2)$, $v_k^2\le2v_{k+1}^2$, and $v_k\le\epsilon$ for $k\le M$, with
$\epsilon\le1/(8C)$. Then $w_k\le8C\,v_k^2$ for all $k\le M$. Along the flow, with $v_k=\bar g_k$
and $\theta=C_0L^{-1}$, this gives $\|W_k\|_k\le8C(L)\bar g_k^2$, uniformly in $k$.
\end{corollary}

\begin{proof}
By induction: if $w_k\le8Cv_k^2$, then $w_k^2\le64C^2\epsilon^2v_k^2\le v_k^2$, and so
$w_{k+1}\le2Cv_k^2+Cv_k^2+Cv_k^2=4Cv_k^2\le8Cv_{k+1}^2$.
\end{proof}

\begin{conjecture}[Remaining step estimates]\label{thm:step-rest}
In the setting of Theorem~\ref{thm:step}:
\begin{itemize}
\item the analogue of (iii)--(iv) for the last step $j+1=N$;
\item analyticity of the map in $(\zeta_0,\dots,\zeta_j)$, and continuity in $m^2\ge0$.
\end{itemize}
\end{conjecture}

\begin{proposition}[Holomorphy and Lipschitz bounds of the step; proved in Lean]\label{prop:step-lip}
Let $d\ge2$, $L\ge L_0$, $A\ge A_0(L)$, $j+1<N$, and let $\kappa_0=C_0(L^{d-3}+A^{-\eta})$ and $C_1$ be the
constants of Theorem~\ref{thm:step}(iii).
\begin{enumerate}
\item[(i)] (Holomorphy.) If the entries of $C(e)$ are holomorphic in $e$ and $V(e),W(e),K(e)$ are
  coefficient-wise holomorphic with $\epsilon(V(e)(B))=0$, then $K_{j+1}$, $W_{j+1}$ and $V_{j+1}$ are
  coefficient-wise holomorphic in $e$.
\item[(ii)] (Lipschitz in $(V,W,K)$.) Let $C$ satisfy \texttt{GoodCov2}$_j(\ell)$, and let
  $0<e\le\epsilon_1$. Suppose $\|V\|_j,\|V'\|_j\le e$, $\|W\|_j,\|W'\|_j\le e^2$ and
  $\|K\|_j,\|K'\|_j\le e^2$. Then, for connected $U$,
  \[
    \|K_{j+1}(x')(U)-K_{j+1}(x)(U)\|_{j+1,U}\le4(\kappa_0+8C_1e)\bigl(e\|\delta V\|_j+\|\delta W\|_j
      +\|\delta K\|_j\bigr).
  \]
\item[(iii)] (The $W$-map.) $W_{j+1}=\mathrm{Src}_j(V,V)+T_jW+R_j(V,W)$, where $R_j$ is the
  product remainder. For $C$ translation invariant, $W,W'$ translation covariant and
  $0<e\le\frac12$ (same domain as (ii)),
  \[
    \|W_{j+1}(x')-W_{j+1}(x)\|_{j+1}\le C_0\tfrac{L^d}{L^3}\|\delta W\|_j+8C_2e\|\delta V\|_j
      +72C_3e^2\bigl(e\|\delta V\|_j+\|\delta W\|_j\bigr).
  \]
\item[(iv)] (Lipschitz in the covariance.) Let $\zeta\mapsto C(\zeta)$, $\zeta\in\C^m$, be holomorphic
  on $\{|\zeta_i|<\rho_1\}$, and let \texttt{GoodCov2}$_j(\ell)$ hold on $\{|\zeta_i|\le\rho\}$ with
  $\rho<\rho_1$. Let $|\zeta_i|,|\zeta'_i|\le\rho/2$ and $|\zeta'_i-\zeta_i|\le\delta$. If $\Phi$ is any of
  $K_{j+1}(U)$, $R_j(B_+)$ or $T_jW(B_+)$, at fixed $(V,W,K)$, then
  $\|\Phi(C(\zeta'))-\Phi(C(\zeta))\|\le8B\delta/\rho$. Here $B$ is the corresponding bound:
  \begin{itemize}
  \item $\kappa_0k+C_1(g^3+gk+k^2)$ for $K_{j+1}$;
  \item $C_3(v^2+w)^2$ for $R_j$;
  \item $C_0L^{d-3}w$ for $T_j$, where translation invariance is also required.
  \end{itemize}
\end{enumerate}
\end{proposition}
\noindent\emph{Lean.} The files are \texttt{H02UV/Step2/Map/HolC.lean}, \texttt{H02UV/Step2/Lip/Lip.lean} and \texttt{H02UV/Step2/Lip/LipW.lean}:
\begin{itemize}
\item (i): \texttt{Kplus\_holC}, \texttt{Wnext2\_holC}, \texttt{Vnext2\_holC}, based on
  \texttt{Hol.ECC\_var} and \texttt{HolE.FCE\_var}. The point is that $\E_C$ and $e^{\cL_C}$ are finite sums of
  powers of $\cL_C$, which is linear in $C$.
\item (ii): \texttt{Kplus\_lip}, via the abstract \texttt{lip\_of\_bound}.
\item (iii): \texttt{Wnext2\_decomp}, \texttt{RemW\_lip}, \texttt{Src\_lip}, \texttt{Wnext2\_lip}.
\item (iv): \texttt{Kplus\_lip\_cov}, \texttt{RemW\_lip\_cov}, \texttt{TrL\_lip\_cov}, and the generic
  \texttt{Tj\_lip\_cov}.
\end{itemize}
\emph{Proof.} Every estimate is a one-variable Cauchy estimate along a complex line
$z\mapsto x+z(x'-x)$, or $z\mapsto C(\zeta+z(\zeta'-\zeta))$, applied to the pairings with test
functions:
\begin{itemize}
\item the general estimate is $|\varphi(1)-\varphi(0)|\le M/(R-1)$ (\texttt{norm\_taylor\_le}, and
  \texttt{norm\_taylor\_le\_on} for a closed disc);
\item the radius is chosen so that the line stays in the domain of the corresponding bound. For (ii)
  it is $R=e^2/(e\|\delta V\|+\|\delta W\|+\|\delta K\|)$; points far apart are handled by the bound
  itself.
\end{itemize}
In (iii) the source and the remainder are treated separately. A single line would lose the
smallness of the $W$-coefficient, because $v^2$ dominates the bound. The dressed covariances
satisfy the hypotheses of (iv) with $m=j+1$, $\rho_1=\bar\zeta$ and any $\rho<\bar\zeta$, and with the field scale
$\ell=(c_S+K\bar\zeta)^{1/2}L^{-j(d-2)/2}L^{p_\Phi}$ (\texttt{dressed\_family}, \texttt{H02UV/Step2/Lip/CdrHol.lean}).
The inputs are \texttt{DressedFRD.thm\_main\_hol}, Proposition~\ref{prop:C5}, and translation invariance of
\texttt{kmat}.
\emph{Continuity in $m^2$} (\texttt{H02UV/Step2/FlowCont/MsqCont.lean}):
\begin{itemize}
\item the symbols $\hat C_k$ are continuous in $m^2>0$ (\texttt{frdC\_continuousAt});
\item at fixed $\zeta$, $\Gamma^{\rm int}_{j+1}(p)$ is continuous in $m^2>0$
  (\texttt{MsqCont.Gint\_mem}, the rational structure of \texttt{DressedFRD.Data.Hol} with $m^2$ in
  place of $\zeta$);
\item $C'_{j+1}$ is continuous on $m^2\in(0,L^{-2(N-1)}]$ (\texttt{Cdr\_continuousOn});
\item $K_{j+1}$ is a continuous function of the covariance (\texttt{Kplus\_continuous\_cov}).
\end{itemize}
At $m^2=0$ the zero mode of $C_N$, which enters $\Gamma_j$, is singular. Not done: the last step
$j+1=N$.

\paragraph{Lean formalisation}
The formalisation (\texttt{H02UV/Step2/**/*.lean}) uses the following concrete choices.
\begin{itemize}
\item \emph{Coordinates.} $V$ and $W$ are families of even elements indexed by the $j$-blocks,
  $V(B),W(B)\in\cN(B^\square)$, $U(1)$-invariant, $\epsilon(V(B))=0$; there are no coupling
  constants. This covers $V=V^\flat$ in coupling form, and steps (A), (B) are applied before.
  The block factor is $I_j(B)=e^{-V(B)}(1+W(B))$.
\item \emph{Covariance.} $C'$ is a complex matrix with a Gram representation satisfying the
  Gram bound at scale $j+1$ with field scale $\ell$, and range $<\frac12L^{j+1}$. The field
  scale $\ell$ is the same at scales $j$ and $j+1$.
\item \emph{Localisation.} $\Loc=\Loc^{(2)}+\Loc^{(4)}$ (Section~\ref{sec:loc}) in $Q$, $J$ and
  $V_{\rm pt}$. The pointwise $\Loc_x$ of Definition~\ref{def:Vpt} is replaced by the
  localisation on the $(j+1)$-block $B_+\supset B$: with
  $S(B)=\frac12F(V'_B,V'(\Lambda))+\E\theta\,W(B)$ and $S(B_+)=\sum_{B\subset B_+}S(B)$,
  \[
    (u+V)_{j+1}(B)=V'_B-\E\theta Q(B)-\Loc_{B_+,B}S(B_+),\qquad
    W_{j+1}(B)=S(B)-\Loc_{B_+,B}S(B_+).
  \]
  So $-V^{(2)}_{j+1}(B)+W_{j+1}(B)=S(B)$ exactly, which is what the second-order computation
  uses. At scale $j+1$, $I_{j+1}(B_+)=\prod_{B\subset B_+}I_{j+1}(B)$, i.e.
  $W_{j+1}(B_+)=\prod_{B\subset B_+}(1+W_{j+1}(B))-1=(1-\Loc_{B_+})S(B_+)+O(W^2)$.
\item \emph{Proof of the third-order bound for $K^{\rm pf}_{j+1}-K^{(2)}$.} $K^{\rm pf}_{j+1}$ is entire in $(V,W,K)$, and bounded on a
  ball by a crude version of the estimates of \hr{} (with Lemma~\ref{lem:pauli} in place of
  \texttt{e:normcontract}). Cauchy estimates along $t\mapsto(tV,t^2W,0)$, $s\mapsto(0,0,sK)$
  and in $(s,t)$ jointly then reduce the bound to two exact statements:
  \begin{itemize}
  \item the first and second $t$-derivatives at $0$ are $0$ and $2K^{(2)}$;
  \item the linear part in $K$ at $V=W=0$ contracts by $O(L^{-1}+A^{-\eta})$, by
    Lemma~\ref{lem:loc2} on small sets and Lemma~\ref{lem:pauli} on large sets.
  \end{itemize}
\item \emph{Status (done).} All of Theorem~\ref{thm:step}, in the package \texttt{Step2}:
  \begin{itemize}
  \item intermediate results for $K^{\rm pf}_{j+1}$ and $M$: \texttt{H02UV/Step2/Map/Main.lean},
    \texttt{H02UV/Step2/Map/Reapp.lean}, \texttt{H02UV/Step2/Map/Lambda.lean};
  \item locality, symmetry and coordinates: \texttt{H02UV/Step2/Map/Mem.lean}, \texttt{H02UV/Step2/Map/Coord.lean};
  \item factorisation: \texttt{H02UV/Step2/Map/ZetaFact.lean}, \texttt{H02UV/Step2/Map/Fact.lean};
  \item the change of variables: \texttt{H02UV/Step2/Map/Reblock.lean}, \texttt{H02UV/Step2/Map/ReblockBound.lean},
    \texttt{H02UV/Step2/Map/Split.lean};
  \item the theorem: \texttt{H02UV/Step2/Map/StepFinal.lean}.
  \end{itemize}
  The ingredients of the Cauchy argument are \texttt{Knext2\_expansion}, \texttt{Knext2\_lin},
  \texttt{Llin2\_bound}, \texttt{Knext2\_crude}, \texttt{Knext2\_hol} and the lemmas of
  \texttt{H02UV/Step2/Map/Cauchy.lean}. The constant $C(L,A)$ is not tracked.
\item \emph{Not formalised yet.} Expected result~\ref{thm:step-rest}.
\end{itemize}

\emph{What is new compared with \texttt{thm:step}.}
\begin{itemize}
\item The third-order bound. As in Bauerschmidt--Brydges--Slade, Chapter~10, the second-order
  terms of $K_{j+1}$ cancel because $V_{\rm pt}$ and $W_{\rm pt}$ are exact to second order
  (Proposition~\ref{prop:Vpt}).
\item The local factor $e^{(V_j-V_j^\flat)(X)}$ has norm $e^{O(\bar g_j\log L)|X|}$. It is
  absorbed by the large-set regulator $A^{(|\cB(X)|-4)_+}$ if $\bar g_j\log L\ll\log A$.
\item The Gram bound of Section~\ref{sec:frd} replaces \texttt{e:E-contract}.
\item \texttt{e:normcontract} fails in $d=2$. It is replaced by Lemmas~\ref{lem:pauli}
  and~\ref{lem:loc2} (Section~\ref{sec:normcontract}). The contraction factor is
  $L^{-\kappa_0}$ with $\kappa_0=1$, and the large-set regulator must satisfy $A\ge4C_*$.
\end{itemize}

%======================================================================
\section{The coupling-constant layer: statements}
\label{sec:coupling}

This section states the results that connect the block-family map of Theorem~\ref{thm:step} to the
coupling constants $(\zeta,\nu,g)$, the absorption steps (A)--(B), and the covariances of
Section~\ref{sec:frd}. Together with Proposition~\ref{prop:fh-flow} they feed into Theorem~\ref{thm:flow}. All statements in this section are \emph{targets}. Each is stated in the
form in which it will be formalised.

Three issues came up while preparing these statements. They are addressed below.
\begin{itemize}
\item \emph{Exact lattice symmetry} (Section~\ref{sec:coupling-sym}). The quartic monomials
  $M^{ik}_x=(\nabla_i\psi)(\nabla_k\bar\psi)\psi\bar\psi$ of $\Loc^{(4)}$ use forward
  differences, whereas the quadratic basis of $\Loc^{(2)}$ uses central differences. For
  $\E\theta$ with a lattice-symmetric $C'$, contracting $\nabla_i\psi_x$ with $\bar\psi_x$ in $M^{ii}_x$
  produces $(C'(e_i,0)-C'(0,0))\,\psi_x(\nabla_i\bar\psi)_x$. Its odd part $\psi\partial_i\bar\psi$
  is a relevant monomial outside the symmetric coupling space. Such terms would be generated
  at every step. The quartic basis must therefore be replaced by the symmetric one, matching
  $M_x=\psi_x\bar\psi_x(\nabla\psi\nabla\bar\psi)_x$ of Section~\ref{sec:pt-coef}.
\item \emph{History dependence} (Section~\ref{sec:coupling-flow}). The covariance
  $C'_{j+1}$ depends on $\zeta_0,\dots,\zeta_j$ (\ref{D5}). So the step map is not a function of the
  current state alone, as Proposition~\ref{prop:fh-flow} assumes.
\item \emph{Reality} (Section~\ref{sec:coupling-flow}). Proposition~\ref{prop:fh-flow} is about
  real coordinates. The coupling constants are real if all data are real.
\end{itemize}

\subsection{Symmetric localisation and lattice symmetries}
\label{sec:coupling-sym}
Let $\mathfrak G_j$ be the group of lattice automorphisms generated by the translations by
$L^j\Z^d$, the coordinate reflections and the coordinate permutations about the centre of the
$j$-block at the origin. These preserve $j$-blocks (Lean: \texttt{AffAut}, \texttt{IsGen},
\texttt{PresBlocks}, \texttt{reflCenter}, \texttt{permCorner}). For $i,k\in\{1,\dots,d\}$ let
\[
  \hat M^{ik}_x=\psi_x\bar\psi_x\,m^{ik}_x,\qquad
  m^{ik}_x=\begin{cases}(\partial_i\psi)_x(\partial_k\bar\psi)_x,&i\ne k,\\
  \frac12\sum_{s=\pm}(\nabla_{se_i}\psi)_x(\nabla_{se_i}\bar\psi)_x,&i=k,\end{cases}
\]
which is the stencil \texttt{m2} times $\psi_x\bar\psi_x$. Let $\widehat\Loc^{(4)}$ and
$\widehat\Loc=\Loc^{(2)}+\widehat\Loc^{(4)}$ be defined as $\Loc^{(4)}$ and $\Loc$, with
$\hat M^{ik}$ in place of $M^{ik}$.

\begin{lemma}[C0, symmetric localisation; proved in Lean]\label{lem:C0}
\begin{enumerate}
\item[(i)] $\mu^{(4)}_{i'k'}(\hat M^{ik}_x)=\delta_{ii'}\delta_{kk'}$ for $x$ in a patch $X$. Hence
  $\mu^{(4)}(F-\widehat\Loc_XF)=0$.
\item[(ii)] \texttt{prop:Loc}(i), Lemma~\ref{lem:loc2} and Lemma~\ref{lem:loc-tr} hold for
  $\widehat\Loc$, with the same proofs.
\item[(iii)] $\tau\widehat\Loc_{X,Y}F=\widehat\Loc_{\tau X,\tau Y}\tau F$ for every
  $\tau\in\mathfrak G_j$, small $X$ and $F\in\cN(X^\square)$.
\item[(iv)] Theorem~\ref{thm:step} and Proposition~\ref{prop:W-contract} hold verbatim for the
  map with $\widehat\Loc$ in place of $\Loc$.
\end{enumerate}
\end{lemma}

\emph{Lean.} The symmetric localisation is \texttt{Loc.Loc2sXY} ($\hat M^{ik}$ is
\texttt{Loc.mono4s}).
\begin{itemize}
\item (i) is \texttt{Loc.mu4\_mono4s\_patch}.
\item (ii) consists of \texttt{Loc.exists\_Loc2sXY\_bound}, \texttt{Loc.loc2s\_contract} and
  \texttt{Step2.Loc2XY\_transl}.
\item (iii) is \texttt{Step2.Loc2sXY\_aut}, for every \texttt{IsGen} lattice map and $F$
  supported in a region of $X$.
\end{itemize}
The \texttt{Step2} map now uses $\widehat\Loc$ throughout. All of its theorems (Theorem~\ref{thm:step},
Proposition~\ref{prop:W-contract}) were re-checked for it, which gives (iv).

All quartic monomials are combinations of
$g^{ee'}_x=(\nabla_e\psi)_x(\nabla_{e'}\bar\psi)_x\psi_x\bar\psi_x$ with $e,e'\in\mathcal E_d$
(\texttt{gmono4}), so the coefficient, locality and norm properties of
\texttt{Loc4Def}/\texttt{Loc4Est} carry over (\texttt{Loc/Loc4Sym.lean},
\texttt{Loc/Loc2Sym.lean}). For (iii), every \texttt{IsGen} map acts on the unit directions as a
signed permutation $(\pi,s)$ (\texttt{isGen\_sperm}). Then
$\tau\hat M^{ik}_x=s_is_k\hat M^{\pi i,\pi k}_{\tau x}$ and
$\mu^{(4)}_{\pi\kappa,\pi\kappa'}(\tau F)=s_\kappa s_{\kappa'}\mu^{(4)}_{\kappa\kappa'}(F)$; the
shift of the base point drops out by antisymmetry.

\begin{proposition}[C1, lattice-symmetry covariance; proved in Lean]\label{prop:C1}
In the setting of Theorem~\ref{thm:step} with $\widehat\Loc$ and $j+2<N$, let
$\tau\in\mathfrak G_{j+1}$. Assume that $C'$ is $\tau$-invariant, and that $V$, $W$ and $K_j$ are
covariant under $\mathfrak G_j$ ($F(\tau B)=\tau F(B)$, $K(\tau X)=\tau K(X)$). Then
$V_{j+1}$, $W_{j+1}$ and $K_{j+1}$ are covariant under $\mathfrak G_{j+1}$.
\end{proposition}
\emph{Lean} (\texttt{H02UV/Step2/Map/SymCov.lean}): \texttt{Wnext2\_latsym}, \texttt{Vnext2\_latsym},
\texttt{Kplus\_latsym}. The formal version is slightly stronger:
\begin{itemize}
\item the hypotheses are covariance under all \texttt{IsGen} maps preserving $j$-blocks
  (\texttt{FamSym}, \texttt{KSym}) and invariance of $C'$ under all \texttt{IsGen} maps
  (\texttt{CSym});
\item the conclusion is covariance under all \texttt{IsGen} maps preserving $(j+1)$-blocks.
\end{itemize}
These maps include $\mathfrak G_{j+1}$ (\texttt{presBlocks\_reflCenter},
\texttt{presBlocks\_permCorner}). The step also needs $\epsilon(V(B))=0$.

\subsection{Coupling constants}
\label{sec:coupling-coords}
For $\kappa=(z,y,\nu,g)\in\C^4$ let
\[
  V_{\kappa,x}=z\,(\nabla\psi\nabla\bar\psi)_x
  +\tfrac y2\bigl((-\Delta\psi)_x\bar\psi_x+\psi_x(-\Delta\bar\psi)_x\bigr)+\nu\,\psi_x\bar\psi_x+g\,\hat M_x,
  \qquad V_\kappa(X)=\sum_{x\in X}V_{\kappa,x},
\]
with $(\nabla\psi\nabla\bar\psi)_x=\sum_im^{ii}_x$ and $\hat M_x=\sum_i\hat M^{ii}_x$. These are
uniform polynomials in the sense of Lemma~\ref{lem:loc-tr} (\texttt{unif}).

\paragraph{The symplectic symmetry}
Under lattice symmetries alone the coefficients of $(\Delta_i\psi)\bar\psi$ and
$\psi(\Delta_i\bar\psi)$ in an invariant uniform polynomial can differ. The coupling form
$V_\kappa$ has them equal. They are swapped by the symplectic involution
$J\psi=\bar\psi$, $J\bar\psi=-\psi$ (old \texttt{Step}: \texttt{Jsym}, \texttt{memHyp\_jinv}),
under which $V_0$, the localisation and $\E_{C'}\theta$ for symmetric $C'$ are covariant. So
Proposition~\ref{prop:C2}(ii) also assumes $J$-invariance of $W$ and $K_j$, and includes
$J$-invariance of the outputs (through the membership principle \texttt{GMem} of
\texttt{H02UV/Step2/Map/Mem.lean}, with the $J$-equaliser in place of $\cN(X)$).

\begin{proposition}[C2, closure of the coupling space; proved in Lean]\label{prop:C2}
Let $C'$ be symmetric ($C'(x,y)=C'(y,x)$) and invariant under all lattice symmetries. Then
$C'(x,x)=c_0$ and $C'(x+e,x)=C'(x,x+e)=c_0+\delta$ for all $x$ and unit directions $e$, with
$c_0=C'(0,0)$ and $\delta=C'(e_1,0)-C'(0,0)$ (\texttt{covProps}).
\begin{enumerate}
\item[(i)] (Lemma~\ref{lem:first}, exact form.) For every $X$, $\E_{C'}\theta\,V_\kappa(X)=V_{\kappa'}(X)+c$
  for a constant $c\in\C$, where
  \[
    \kappa'=(z-gc_0,\ y-g\delta,\ \nu+2d\,g\delta,\ g).
  \]
  In particular $z'+y'=z+y-g\,C'(e_1,0)$ and $\nu'=\nu+g\,\Delta C'(0)$
  (\texttt{expLC\_Vcoup}; Wick formulas \texttt{LopC\_two}, \texttt{LopC\_four},
  \texttt{expLC\_four}).
\item[(ii)] Let $V(B)=V_\kappa(B)$ on every $j$-block, and let $V$, $W$, $K_j$ satisfy the
  hypotheses of Theorem~\ref{thm:step}. Assume further that they are covariant under the lattice
  symmetries (\texttt{FamSym}, \texttt{KSym}) and $J$-invariant (\texttt{FamJ}, \texttt{KJ}), and
  that $j+2<N$. Then there is one $\kappa_{j+1}\in\C^4$ such that
  \[
    V_{j+1}(B)=V_{\kappa_{j+1}}(B)\ \text{for every $j$-block $B$, and}\quad
    V_{j+1}(B_+)=V_{\kappa_{j+1}}(B_+)\ \text{for every $(j+1)$-block $B_+$.}
  \]
  The constants go into $u_{j+1}$ (\texttt{CoupHyp.Vp\_coup}, \texttt{CoupHyp.Vnext2\_coup}).
  Explicitly, $\kappa_{j+1}=\kappa'-\kappa_Q'-\kappa_S$. Here $\kappa_Q$ and $\kappa_S$ are the
  coupling vectors of $Q(B)$ and $\Loc_{B_+}S(B_+)$, and $'$ is the map of (i).
\item[(iii)] (Coordinates; proved in Lean.) $\zeta_{j+1}=z_{j+1}+y_{j+1}$ is the coefficient of
  $(\psi,-\Delta\bar\psi)$ in $V_{j+1}(\Lambda)$ (\texttt{Vcoup\_grad}). If $\|K_j\|_{j,A}\le r$,
  then $\kappa_Q=(z_Q,y_Q,\nu_Q,g_Q)$ satisfies, with $h=L^j$ and $|B|=L^{dj}$,
  \[
    |z_Q|,|y_Q|\le C\,c_SC_l\,\ell^{-2}h^2\frac r{|B|},\qquad
    |\nu_Q|\le C\,c_SC_l\,\ell^{-2}\frac r{|B|},\qquad
    |g_Q|\le C\,c_SC_l\,\ell^{-4}h^2\frac r{|B|}
  \]
  (\texttt{exists\_coup\_bound}, \texttt{BH.coupQ\_le}). For $d=2$ this is
  $|\zeta_Q|+L^{2j}|\nu_Q|+|g_Q|\le C(L)\|K_j\|$ in the dimensionless couplings. So
  $|\kappa_{j+1}-(\kappa'-\kappa_S)|\le C(L)\|K_j\|_j$.
\end{enumerate}
\end{proposition}
\noindent\emph{The initial potential.} A block family $V(B)=V_\kappa(B)$, in particular $V_0$,
satisfies the hypotheses \texttt{FamSym} and \texttt{FamJ}. The reason is that
$\tau V_\kappa(Y)=V_\kappa(\tau Y)$ for every lattice automorphism $\tau$ generated by
translations, reflections and permutations, and $JV_\kappa(Y)=V_\kappa(Y)$
(\texttt{Step2.autAlg\_Vcoup}, \texttt{Jsym\_Vcoup}, \texttt{famSym\_Vcoup}, \texttt{famJ\_Vcoup}).
The outputs keep these properties by C1 and Proposition~\ref{prop:C2}.

\emph{Proof of (ii).}
\begin{itemize}
\item $V'_B=\E\theta V_\kappa(B)$ is $V_{\kappa'}(B)+c$ by (i).
\item $Q(B)=\sum_{X\supset B}\Loc_{X,B}K(X)$ is uniform on $B$. Its coefficients are invariant
  under the reflections and permutations fixing $B$ (\texttt{AE\_Q2\_j}) and under $J$. So they
  are coupling coefficients (\texttt{unif\_sym\_coeff}, \texttt{coup\_of\_sym}). By translation
  covariance they are the same for all $B$ (\texttt{Q2\_coup}).
\item $\Loc_{B_+,B}S(B_+)$ is uniform on $B$ with the coefficients of $\Loc_{B_+}S(B_+)$, which
  are symmetric by Proposition~\ref{prop:C1} for the stabiliser of $B_+$ (\texttt{LS\_coup}).
\item Coefficients of uniform polynomials are unique (\texttt{unif\_coeff}). The reflections act
  by $(-1)^{\#\text{derivatives in that direction}}$ and the permutations by relabelling
  (\texttt{reflCenter\_unif}, \texttt{permCorner\_unif}), and $J$ exchanges
  $(\Delta\psi)\bar\psi\leftrightarrow\psi\Delta\bar\psi$ (\texttt{Jsym\_unif}).
\end{itemize}
$J$-invariance of the step is \texttt{Vnext2\_J}, \texttt{Wnext2\_J}, \texttt{Kplus\_J}. It
needs $\E_{C'}\theta J=J\E_{C'}\theta$ for symmetric $C'$ (\texttt{ECC\_Jsym2}) and
$\widehat\Loc J=J\widehat\Loc$ (\texttt{Loc2sXY\_Jsym}, using
$\mu^{(4)}_{\kappa\kappa'}(JF)=\mu^{(4)}_{\kappa'\kappa}(F)$).

\subsection{Absorption and normalisation}
\label{sec:coupling-abs}
For a complex covariance $\Gamma$ write $\E_\Gamma[F]=\epsilon(\E_\Gamma\theta F)$ for the
Gaussian expectation, that is, convolution followed by $\psi=0$. Let
$Q=(\psi,A\bar\psi)=\sum_x\psi_x(A\bar\psi)_x$ with $A=-\Delta+m^2$, and let $\sigma_\lambda$ be the
field rescaling $\psi\mapsto\lambda\psi$, $\bar\psi\mapsto\lambda\bar\psi$.

\begin{lemma}[C3, Gaussian identities; proved in Lean]\label{lem:C3}
For complex matrices $\Gamma,\Gamma_1,\Gamma_2$ and $z\in\C$ with $\det(1+zA\Gamma)\ne0$:
\begin{enumerate}
\item[(i)] $\E_{\Gamma_1+\Gamma_2}\theta=\E_{\Gamma_2}\theta\circ\E_{\Gamma_1}\theta$;
\item[(ii)] $\E_\Gamma[e^{-zQ}F]=\det(1+zA\Gamma)\,\E_{\Gamma(1+zA\Gamma)^{-1}}[F]$ for all
  $F\in\cN$ (this is \eqref{e:absorb}; Lemma~\ref{lem:onestep});
\item[(iii)] $\E_\Gamma[\sigma_\lambda F]=\E_{\lambda^2\Gamma}[F]$, and
  $\sigma_\lambda\E_\Gamma\theta=\E_{\lambda^2\Gamma}\theta\,\sigma_\lambda$.
\end{enumerate}
\end{lemma}
(Lean: (i) \texttt{Grassmann.expLC\_add}; (ii) \texttt{Grassmann.omegaC\_exp\_mul\_det},
\texttt{Grassmann.omegaC\_onestep}; (iii) \texttt{Step2.omegaC\_sig}, \texttt{Step2.expLC\_sig};
here $\E_C\theta$ acts as $e^{\mathcal L_C}$, so (iii) reads
$e^{\mathcal L_\Gamma}\sigma_\lambda=\sigma_\lambda e^{\mathcal L_{\lambda^2\Gamma}}$.)

\begin{proposition}[C3', the full step; algebraic part proved in Lean]\label{prop:C3}
Let $(\Gamma_j,\Gamma^{\rm int}_{j+1},Z_j)$ be as in Section~\ref{sec:frd}, with $z_j=\zeta_jZ_j$
and $|\zeta_i|\le\bar\zeta$. Suppose that for some $j+1<N$,
\[
  \E_{\Gamma_0}\bigl[e^{-V_0(\Lambda)}\bigr]=e^{-U_j}\,
  \E_{Z_j\Gamma_j}\Bigl[\sum_XI_j^{\Lambda\setminus X}\prod_{Y\in{\rm Comp}(X)}K_j(Y)\Bigr]
\]
in fields normalised at scale $j$, with $I_j(B)=e^{-V_{\kappa_j}(B)}(1+W_j(B))$. Define
\begin{itemize}
\item $V^\flat=V_{(0,0,\nu_j-\zeta_jm^2,g_j)}$ and $K^\flat(X)=e^{(V_{\kappa_j}-V^\flat)(X)}K_j(X)$;
\item the rescaled data $\sigma_{(1+\zeta_j)^{-1/2}}$ applied to $V^\flat$, $W_j$ and $K^\flat$;
\item the \texttt{Step2} map with $C'=C'_{j+1}=Z_{j+1}\Gamma^{\rm int}_{j+1}$.
\end{itemize}
Then the same representation holds at scale $j+1$, with
$U_{j+1}=U_j+u_{j+1}|\Lambda|-\log\det(1+z_jA\Gamma_j)$ and the outputs
$(\kappa_{j+1},W_{j+1},K_{j+1})$ of Proposition~\ref{prop:C2}. Moreover $K^\flat$ satisfies
the hypotheses of Theorem~\ref{thm:step} with
$\|K^\flat\|_{j,A'}\le\|K_j\|_{j,A}$ for $A'=Ae^{-C|\kappa_j|}$.
\end{proposition}
\noindent\emph{Lean.} The algebraic part is \texttt{Step2.omegaC\_full\_step}
(\texttt{H02UV/Step2/Map/FullStep.lean}), in the following form. Let $V(B)=V^\flat(B)+\Delta(B)$ with
$\Delta(B)=V_{(z,y,\zeta m^2,0)}(B)$, $\zeta=z+y$, so that $\kappa^\flat=\kappa-(z,y,\zeta m^2,0)
=(0,0,\nu-\zeta m^2,g)$. Let $\Gamma(1+\zeta A\Gamma)^{-1}=\Gamma_1+\Gamma_2$ with
$\det(1+\zeta A\Gamma)\ne0$, and let $\lambda s=1$. Then
\[
  \E_\Gamma\bigl[Z_j(V,W,K)\bigr]=\det(1+\zeta A\Gamma)\,
  \E_{\lambda^2\Gamma_2}\Bigl[\sum_{U\in\cP_{j+1}}I_{j+1}^{\Lambda\setminus U}\,
  e^{-u_{j+1}(U)}K_{j+1}(U)\Bigr],
\]
where $Z_j(V,W,K)=\sum_XI^{\Lambda\setminus X}K(X)$, and $(I_{j+1},u_{j+1},K_{j+1})$ is the
output of the \texttt{Step2} map (\texttt{stepK\_ZVK}) with $C'=\lambda^2\Gamma_1$ applied to
$(\sigma_sV^\flat,\sigma_sW,\sigma_sK^\flat)$, $K^\flat(X)=e^{\Delta(X)}K(X)$,
$\Delta(X)=\sum_{B\in\cB_j(X)}\Delta(B)$. The steps are: summation by parts
$V_{(z,y,\zeta m^2,0)}(\Lambda)=\zeta Q$ (\texttt{Vcoup\_grad}); the factorisation
$Z_j(V^\flat+\Delta,W,K)=e^{-\Delta(\Lambda)}Z_j(V^\flat,W,K^\flat)$ (\texttt{gZ\_factor},
using that $\Delta(X)$ is additive over the components of $X$, \texttt{Kext\_bpow\_mul});
Lemma~\ref{lem:C3}(ii), and then (iii) with $\lambda=(1+\zeta_j)^{1/2}$
(\texttt{omegaC\_gZ\_step}). The choice $\lambda^2\Gamma_1=C'_{j+1}$, $\lambda^2\Gamma_2=Z_{j+1}\Gamma_{j+1}$
comes from the DressedFRD split. The bound on $K^\flat$ is \texttt{Step2.Knorm\_Kflat\_le}:
$\|K^\flat\|_{j,Ae^{-\delta}}\le e^{2^d\delta}\|K\|_{j,A}$ for $\sup_B\|\Delta(B)\|_j\le\delta$. Here
$\delta$ is bounded by the couplings via \texttt{Step2.Tj\_Vcoup\_le}. The identification with the
DressedFRD matrices is C5.

\subsection{Second-order coefficients}
\label{sec:coupling-coef}
\begin{proposition}[C4, the second-order part of $W$; exact for $W^{\rm lin}$, see the correction]\label{prop:C4}
Define $W^{(2)}_j$ recursively by $W^{(2)}_0=0$ and
\[
  W^{(2)}_{j+1}(B)=(1-\widehat\Loc_{B_+,B})\Bigl[\tfrac12F_{C'_{j+1}}\bigl(V'_B,V'(\Lambda)\bigr)
  +\E_{C'_{j+1}}\theta\,\sigma_{(1+\zeta_j)^{-1/2}}W^{(2)}_j(B)\Bigr],\qquad V'=\E\theta V^\flat_j,
\]
which is an explicit quadratic function of $\kappa_0,\dots,\kappa_j$. Along the flow of
Proposition~\ref{prop:C3}, with $\|K_k\|_k\le\bar g_k^3$ and $\bar g_k$ slowly varying as in
Corollary~\ref{cor:W-uniform},
\[
  \|W_j-W^{(2)}_j\|_j\le C(L)\,\bar g_j^3 .
\]
\end{proposition}
\noindent\emph{Correction.} $W_j$ is \emph{not} exactly quadratic in the couplings. The
\texttt{Step2} block recursion
\[
  W'_{j}(B)=S(B)-\Loc_{B_+,B}S(B_+),\qquad S(B)=\tfrac12F_{C'}(V'_B,V'(\Lambda))+\E_{C'}\theta W(B),
\]
does not involve $K_j$. It is linear in $W$, with a source that is quadratic in $V$
(\texttt{Step2.Wp\_split}):
\[
  W'_j(B)=W'_j|_{W=0}(B)+\E_{C'}\theta W(B)-\Loc_{B_+,B}\E_{C'}\theta W(B_+).
\]
But the remainder at the next scale is $W_{j+1}(B_+)=\prod_{B\subset B_+}(1+W'_j(B))-1$
(\texttt{Wnext2}). So
\[
  W_{j+1}(B_+)=\sum_{B\subset B_+}W'_j(B)
    +\sum_{\substack{A\subset\cB(B_+)\\|A|\ge2}}\prod_{B\in A}W'_j(B),
\]
and the products are of order $\ge4$ in the couplings. Define the \emph{linearised} remainder
$W^{\rm lin}$ by dropping the products, so that $W^{\rm lin}_{j+1}(B_+)=\sum_{B\subset B_+}W'_j(B)$
computed from $\sigma_{s_j}W^{\rm lin}_j$. Then $W^{(2)}_j=W^{\rm lin}_j$ exactly, when
$W^{(2)}$ is formed with the actual couplings of the flow. The difference
$W_j-W^{\rm lin}_j$ consists of the transported product terms, which are of order $\ge4$.
They enter $\Loc S_j$ linearly (\texttt{FlowSeq.LS\_sub\_lin}), and bounding them is part of
C4$'$(d). A further difference appears when the past couplings $\kappa_k$, $k<j$, are replaced
by the current ones (the frozen form of Lemma~\ref{lem:frozen}). Proposition~\ref{prop:C4b}
treats this; that difference turns out to be of second order in the $g\nu$ term.

\paragraph{Setting for C4$'$.}
Run the flow of Proposition~\ref{prop:C3}. At step $k$ the \texttt{Step2} map receives:
\begin{itemize}
\item the potential $V^{\rm in}_k=\sigma_{s_k}V^\flat_k=V_{\kappa^{\rm in}_k}$, with
  $s_k=(1+\zeta_k)^{-1/2}$ and
  $\kappa^{\rm in}_k=(0,0,\nu^{\rm in}_k,g^{\rm in}_k)=(0,0,s_k^2(\nu_k-\zeta_km^2),s_k^4g_k)$;
\item the remainder $\sigma_{s_k}W_k$;
\item the activity $\sigma_{s_k}K^\flat_k$;
\item the covariance $C'_{k+1}$.
\end{itemize}
Let $\kappa'_k=(z'_k,y'_k,\nu'_k,g'_k)$ be the image of $\kappa^{\rm in}_k$ under the first-order
map of Proposition~\ref{prop:C2}(i) with $C'_{k+1}$. So $g'_k=g^{\rm in}_k$,
$z'_k+y'_k=-g^{\rm in}_kC'_{k+1}(0,e_1)$ and $\nu'_k=\nu^{\rm in}_k+g^{\rm in}_k\Delta C'_{k+1}(0)$.
By Proposition~\ref{prop:C2}(ii),
\begin{equation}\label{e:kappa-exact}
  \kappa_{k+1}=\kappa'_k-(\kappa_{Q,k})'-\kappa_{S,k},
\end{equation}
where $\kappa_{S,k}$ is the coupling vector of $\Loc_{B_+,B}S_k(B_+)$. By C2(iii),
$|(\kappa_{Q,k})'|\le C(L)\|K_k\|_k$ in the scaled units of C2(iii).

\paragraph{Line expansion.}
For even $A,B$ write $F_C(A,B)=\sum_{n\ge1}F^{(n)}_C(A,B)$, where $F^{(n)}_C$ collects the terms of
$e^{\cL_C}(AB)-e^{\cL_C}A\,e^{\cL_C}B$ with exactly $n$ contractions between $A$ and $B$ and no
contraction inside $A$ or inside $B$. For quartic $A,B$ the term $F^{(n)}_C(A,B)$ has degree
$8-2n$. With $M_x=\psi_x\bar\psi_x(\nabla\psi)_x(\nabla\bar\psi)_x$ (the monomial of $g$) define, for
a $j$-block $B$ and $C_1,C_2\in\cT_\C$,
\begin{align*}
  -\mathfrak B_B(C)&:=\text{the $g$-coefficient of }
    \Loc_{B_+,B}\tfrac12\textstyle\sum_{B'\subset B_+}F^{(2)}_C\bigl(M(B'),M(\Lambda)\bigr),\\
  -\mathfrak B_B(C_1;C_2)&:=\text{the $g$-coefficient of }
    \Loc_{B_+,B}\tfrac12\textstyle\sum_{B'\subset B_+}\cL_{C_2}F^{(1)}_{C_1}\bigl(M(B'),M(\Lambda)\bigr),\\
  \mathfrak c_B(C)&:=\text{the $g$-coefficient of }
    \Loc_{B_+,B}\tfrac12\textstyle\sum_{B'\subset B_+}\bigl[F_C\bigl((\psi\bar\psi)(B'),M(\Lambda)\bigr)
    +F_C\bigl(M(B'),(\psi\bar\psi)(\Lambda)\bigr)\bigr].
\end{align*}
Here ``the $g$-coefficient'' is the last component of the coupling vector of a polynomial
$u|B|+V_\kappa(B)$. Such a polynomial is obtained when $C,C_1,C_2$ are lattice-symmetric and
symmetric, by C1 and C2. By translation invariance these numbers do not depend on $B$, and we drop
the index.

\begin{proposition}[C4$'$, the exact second-order coefficients; proved in Lean]\label{prop:C4b}
In this setting, for $j+2<N$:
\begin{enumerate}
\item[(a)] (Quadratic structure; proved in Lean for $W^{\rm lin}$.) Let $\kappa^{\rm lin}_{S,j}$
  be the coupling vector of $\Loc_{B_+,B}S_j(B_+)$, with $S_j$ formed from $\sigma_{s_j}W^{\rm lin}_j$.
  There are bilinear maps $\beta_{j,k}:\C^4\times\C^4\to\C^4$, $k\le j$, depending only on
  $C'_{k+1},\dots,C'_{j+1}$ and $s_{k+1},\dots,s_j$ (and on the block $B$), such that
  \[
    \kappa^{\rm lin}_{S,j}=\sum_{k\le j}\beta_{j,k}(\kappa^{\rm in}_k,\kappa^{\rm in}_k).
  \]
  Explicitly:
  \begin{itemize}
  \item $\beta_{j,j}(\kappa,\lambda)$ is the coupling vector of
    $\Loc_{B_+,B}\sum_{B'\subset B_+}\frac12F_{C'_{j+1}}(V'_\kappa(B'),V'_\lambda(\Lambda))$, where
    $V'=\E_{C'_{j+1}}\theta V$;
  \item for $k<j$, $\beta_{j,k}(\kappa,\lambda)$ is the coupling vector of
    $\Loc_{B_+,B}\sum_{B'\subset B_+}\E_{C'_{j+1}}\theta\,\sigma_{s_j}T_{j,k}\,{\rm Src}_k(\kappa,\lambda)(B')$.
    Here ${\rm Src}_k(\kappa,\lambda)(B_+)=\sum_{B\subset B_+}(1-\Loc_{B_+,B})\frac12F_{C'_{k+1}}(V'_\kappa(B),
    V'_\lambda(\Lambda))$, and $T_{j,k}$ is the composition of the transports
    $W\mapsto\sum_{B\subset B_+}\bigl(\E\theta\sigma W(B)-\Loc_{B_+,B}\E\theta\sigma W(B_+)\bigr)$
    at the scales $k+1,\dots,j-1$.
  \end{itemize}
  The unrolled recursion is $W^{\rm lin}_j=\sum_{k<j}T_{j,k}{\rm Src}_k$.

  \emph{Lean} (\texttt{H02UV/Step2/Coup/CoupQuad.lean}):
  \begin{itemize}
  \item \texttt{FlowSeq.Wlin\_eq} (the unrolled recursion);
  \item \texttt{FlowSeq.kappaS\_quad} and \texttt{kappaS\_quad'} (the decomposition, with the
    coefficient vector \texttt{cvec4});
  \item \texttt{FlowSeq.beta\_congr} (the dependence only on the stated data).
  \end{itemize}
  Since constants drop out of $F$ and $V'_\kappa=V_{\kappa'}+{\rm const}$ (Proposition~\ref{prop:C2}(i)),
  $\beta_{j,k}(\kappa,\lambda)=\beta^q_{j,k}(\kappa',\lambda')$, where $\beta^q_{j,k}$ is obtained by
  replacing $V'_\kappa$ with $V_{\kappa'}$, and $\kappa\mapsto\kappa'$ is the linear first-order shift
  of $C'_{k+1}$. So $q_{j,k}(\kappa')=\beta^q_{j,k}(\kappa',\kappa')$, and
  $\kappa^{\rm lin}_{S,j}=\sum_{k\le j}q_{j,k}(\kappa'_k)$ (\texttt{H02UV/Step2/Coup/CoupQuad2.lean}:
  \texttt{FC\_const\_left}, \texttt{SrcF\_shift}, \texttt{FlowSeq.beta\_eq\_betaQ},
  \texttt{FlowSeq.kappaS\_quadQ}).
\item[(b)] (The $g$-component; proved in Lean.) Write $C'=C'_{j+1}$,
  $C'^{(1)}=\sum_xC'(x,0)$, $\nu'_j=\nu^{\rm in}_j+2d\,g^{\rm in}_j\delta_j$ (the first-order image),
  $-\mathfrak B(C'):=q_{j,j}(e_g)_g$ and $b_{j,k}:=\beta_{j,k}(e_g,e_g)_g$ for $k<j$. Then, exactly,
  \[
    (\kappa^{\rm lin}_{S,j})_g=\sum_{k<j}(g^{\rm in}_k)^2\,b_{j,k}-(g^{\rm in}_j)^2\,\mathfrak B(C')
      +4\,C'^{(1)}\,g^{\rm in}_j\,\nu'_j ,
  \]
  and therefore
  \[
    g_{j+1}=g^{\rm in}_j+(g^{\rm in}_j)^2\mathfrak B(C')-4C'^{(1)}g^{\rm in}_j\nu'_j
      -\sum_{k<j}(g^{\rm in}_k)^2b_{j,k}-(\kappa_{Q,j})'_g-(\kappa_{S,j}-\kappa^{\rm lin}_{S,j})_g .
  \]
  In particular the quadratic couplings $(z,y,\nu)$ enter the $g$-flow only through the single
  term $4C'^{(1)}g\nu'$ of the current step; the gradient couplings $z',y'$ do not enter at all, and
  the past steps contribute only through $g_k^2$.

  \emph{Lean:} \texttt{FlowSeq.kappaS\_g} and \texttt{FlowSeq.kappaS\_g'} (\texttt{H02UV/Step2/Coup/CoupG.lean})
  for lattice-symmetric symmetric covariances with \texttt{GoodCov2}, $d\ge1$, $2^{d+3}+8d<L$ and
  $j+1<N$. The ingredients are:
  \begin{itemize}
  \item \texttt{betaQ\_self\_g}: no terms quadratic in the quadratic couplings at the current step;
  \item \texttt{beta\_past\_g}: past steps contribute only through $g_k^2$ (degree counting, the
    vanishing quartic pairing of $(1-\Loc)$, \texttt{mu4\_wpsum}, and base-point independence);
  \item \texttt{qjj\_g}, from (c).
  \end{itemize}
\item[(c)] (Wick evaluation; proved in Lean.)
  \begin{itemize}
  \item \texttt{FC\_sub\_Xc\_degLe}: for $A$ of degree $\le2$ and $B$ of degree $\le4$, $F_C(A,B)$ and
    the single cross contraction
    $X_C(A,B)=\sum C(x,y)\bigl[(\partial_{y+}A)(\partial_{x-}B)-(\partial_{x-}A)(\partial_{y+}B)\bigr]$
    differ by a polynomial of degree $\le2$.
  \item \texttt{Xc\_pm\_$\iota$Multi}, \texttt{mu4\_$\iota$Multi}, \texttt{mu4\_Xc\_gmono4}: $X_C(\iota p\,\iota q,\cdot)$
    acts on a product of four vectors as the derivation of $u\mapsto\langle p,u\rangle_Cq+\langle u,q\rangle_Cp$,
    and the quartic pairing of a product of four vectors is the determinant of their columns.
    This gives a closed formula for each pair $\psi_s\bar\psi_t$ (\texttt{Phi}).
  \item The lattice sums (\texttt{H02UV/Step2/Coup/WickSum.lean}, \texttt{CovH}): vanishing first moments by
    reflection symmetry, row sums, and telescoping. Results: the mass cross term gives
    $4C^{(1)}|X|$ in both summation orders (\texttt{mu4\_mass\_univ\_X}, \texttt{mu4\_mass\_X\_univ}),
    and all gradient cross terms vanish (\texttt{mu4\_m2\_*}, \texttt{mu4\_a2\_*},
    \texttt{mu4\_b2\_*}).
  \item Hence $\beta^q_{j,j}(e_g,e_\nu)_g+\beta^q_{j,j}(e_\nu,e_g)_g=4C^{(1)}$
    (\texttt{betaQ\_g\_nu}) and the $gz$, $gy$ cross coefficients vanish (\texttt{betaQ\_g\_z},
    \texttt{betaQ\_g\_y}). This is Lemma~\ref{lem:second}(i) at the current step, with $w=C'$, and
    the claim that the cross terms with the gradient couplings have no zero-momentum quartic part.
  \end{itemize}
  \begin{itemize}
  \item \emph{The past coefficients} (\texttt{FlowSeq.beta\_past\_eq}, \texttt{H02UV/Step2/Coup/CoupPast6.lean}):
    for $k<j$,
    \[
      b_{j,k}=\Bigl(\prod_{i=k+1}^{j}s_i^6\Bigr)\,\frac{1}{|B_+|}\,
      \mu_4\Bigl(\cL_{C'_{j+1}}\tfrac12F_{C'_{k+1}}\bigl(V_g(B_+),V_g(\Lambda)\bigr)\Bigr),
    \]
    with $V_g=\sum_iM^{ii}$ the quartic coupling monomial and $\mu_4$ the quartic pairing
    (based at $B_+$, direction $i_0$). Only the degree-$6$ part of the step-$k$ source reaches $g$
    at step $j$. The reasons:
    \begin{itemize}
    \item $\sigma_s$ multiplies the degree-$6$ part by $s^6$ (\texttt{degLe\_sig\_sub}, which uses
      that even elements have no odd coefficients);
    \item the intermediate $e^{\cL}$ and $\Loc$ only change parts of degree $\le4$;
    \item $\mu_4\circ\cL$ kills degree $\le4$.
    \end{itemize}
  \item \emph{The telescoping identity} (\texttt{mu4\_FC\_add\_cov}): for $A,B$ of degree $\le4$ and
    $A'=e^{\cL_{C_2}}A$, $B'=e^{\cL_{C_2}}B$,
    \[
      \mu_4F_{C_1+C_2}(A',B')=\mu_4F_{C_1}(A,B)+\mu_4\cL_{C_2}F_{C_1}(A,B)+\mu_4F_{C_2}(A',B').
    \]
    This is the exact form of $\mathfrak B(w+C')=\mathfrak B(w)+2\mathfrak B(w;C')+\mathfrak B(C')$.
  \end{itemize}
  \emph{Consequence for the frozen comparison:}
  \begin{itemize}
  \item With frozen couplings and $s_i=1$, the current coefficient $-\mathfrak B(C'_{j+1})$ and the
    past coefficients $b_{j,k}$ combine, by the telescoping identity, into
    $\mathfrak B(w'_{j+1})-\mathfrak B(w'_j)$. Here $\mathfrak B$ is evaluated on the primed
    arguments $V'_g=e^{\cL}V_g$.
  \item Replacing $V'_g$ by $V_g$ changes $\mu_4F$ by the cross terms of $V_g$ with the quadratic
    (tadpole) part $V'_g-V_g$. By (c) these equal $4C^{(1)}$ times the mass-type tadpole
    coefficient $g\Delta C'(0)$, which is the term $-4g(\nu'-\nu)w'^{(1)}_j$ of \eqref{e:gflow}.
  \end{itemize}
  So \eqref{e:gflow} is the frozen form of (b). The remaining differences are of the size of the
  coupling variation $|g_k-g_j|$, of $1-s_i$, and of third order.
\item[(d)] (Estimates; proved in Lean, \texttt{H02UV/Step2/Coup/Est4.lean}.) Let $d\ge2$. There are $C_0,K\ge0$
  and $L_0$ such that for $L\ge L_0$ there are $C_2,C_3\ge0$ with the following properties. Assume
  \texttt{GoodCov2}$_i(\ell)$ and translation invariance of $C'_{i+1}$, and $|s_i|\le1$, on the
  relevant ranges. Write $\theta=C_0L^{d-3}$ (precisely $C_0L^d/L^3$).
  \begin{itemize}
  \item (Transport, \texttt{TrL\_contract'}.) For $W$ in \texttt{InF}$_i$ and
    translation covariant (\texttt{FamCovG}$_i$),
    $\|T_iW\|_{i+1}\le\theta\|W\|_i$. Hence $\|T_{k+1+n,k}W\|_{k+1+n}\le\theta^n\|W\|_{k+1}$,
    and locality and covariance are preserved (\texttt{TrIt\_bound}).
  \item (Source, \texttt{Src\_bound}.) If $\|V\|_k\le v\le1$ and $\epsilon(V)=0$, then
    $\|\mathrm{Src}_k(V,V)\|_{k+1}\le C_2v^2$.
  \item (Coefficients, \texttt{LpartL\_coef\_bound}.) If $\|F\|_j\le w$, then the coupling vector
    $u=(z,y,\nu,g)$ of $\Loc_{B_+,B}\sum_{B'}\E\theta\sigma_sF(B')$ satisfies
    \[
      |u_a|\le \frac{K}{\ell^2\,|B_+|}\,\ell^{-2[a=g]}\,(L^{j+1})^{2[a\ne\nu]}\,L^dQ_1\,w .
    \]
  \item (Past coefficients, \texttt{FlowSeq.beta\_bound}.) For $k<j$ and $\|V_k(u)\|_k\le v\le1$, every
    component of $q_{j,k}(u)$ obeys the bound of the previous item with
    $w=\theta^{\,j-1-k}C_2v^2$. In particular $b_{j,k}$ decays like $\theta^{\,j-k-1}$. The base is
    $C_0L^{d-3}$: the earlier form $C(L)L^{-(j-k)}$ is not uniform in $j-k$.
  \item (Product remainder.) Let $W_{k+1}=\mathrm{Wnext2}_k(V_k,\sigma_{s_k}W_k)$, $W_0=0$
    (\texttt{FlowSeq.Wact}). Then the following hold.
    \begin{itemize}
    \item (\texttt{Wact\_sub\_Wlin\_succ}.)
      $W_{k+1}-W^{\rm lin}_{k+1}=R_k+T_k(W_k-W^{\rm lin}_k)$, where
      $R_k=\prod_B(1+W'(B))-1-\sum_BW'(B)$.
    \item (\texttt{prodRem\_bound}.) $\|R_k\|_{k+1}\le C_3(v_k^2+w_k)^2$, where $v_k$ and $w_k\le1$
      bound $\|V_k\|_k$ and $\|W_k\|_k$.
    \item (\texttt{Wact\_sub\_Wlin\_bound}.) Suppose $\theta\rho\le\frac12$ and
      $(v_k^2+w_k)^2\le\rho\,(v_{k+1}^2+w_{k+1})^2$ for $k<M$. Then for all $k\le M$,
      \[\|W_k-W^{\rm lin}_k\|_k\le2\rho\,C_3(v_k^2+w_k)^2.\]
    \item (\texttt{kappaS\_sub\_lin\_bound}.) Under the same hypotheses, with $M+1<N$,
      $\kappa_{S,M}-\kappa^{\rm lin}_{S,M}$ is the coupling vector of
      $\Loc_{B_+,B}\sum_{B'}\E\theta\sigma(W_M-W^{\rm lin}_M)(B')$ (\texttt{LS\_sub\_lin}).
      It obeys the coefficient bound with $w=2\rho C_3(v_M^2+w_M)^2$.
    \end{itemize}
  \end{itemize}
  With $v_k\asymp\bar g_k$ and $w_k\lesssim\bar g_k^2$ (from C6), this gives
  $|\kappa_{S,j}-\kappa^{\rm lin}_{S,j}|\le C(L)\bar g_j^4$ in the scaled units of C2(iii).
  The inputs $v_k,w_k$ and the slow variation of $\bar g_k$ are hypotheses here; they are supplied
  by the flow.
\end{enumerate}
\end{proposition}

\noindent\emph{Comparison with \eqref{e:gflow} (correction of an earlier remark).} The exact
current-step $g\nu$ term is $-4C'^{(1)}_{j+1}g\nu'$, involving only the current covariance. In
\eqref{e:gflow} the term $-4g\mathfrak n_j=-4g(\nu'w'^{(1)}_{j+1}-\nu w'^{(1)}_j)$ contains in
addition $-4g(\nu'-\nu)w'^{(1)}_j=-4g^2\Delta C'_{j+1}(0)\,w'^{(1)}_j$. An earlier version of this
remark called this a second-order discrepancy. That was premature: in the exact recursion the
corresponding term sits in the past coefficients $b_{j,k}$. There, $\cL_{C'_{j+1}}$ contracts two
legs of the same $M$ in the degree-$6$ part of the step-$k$ source (a tadpole $\propto\Delta C'_{j+1}(0)$),
and the line $C'_{k+1}$ then acts on the remaining leg as in the mass computation, giving
$\propto\Delta C'_{j+1}(0)\,C'^{(1)}_{k+1}$. Summed over $k$ with frozen couplings, this has the form
of the extra term. So \eqref{e:gflow} is consistent with (b) in the frozen approximation; see
the consequence paragraph in (c).

\paragraph{The bubble coefficient (proved in Lean).} \texttt{H02UV/Step2/Coup/Bubble.lean} states:
\begin{itemize}
\item for homogeneous quartic $A,B$, only the terms with exactly two lines reach the quartic
  pairing (\texttt{mu4\_FC\_quartic}):
  \[
    \mu_4F_C(A,B)=\mu_4c_2(A,B),\qquad
    c_2=\tfrac12\cL^2(AB)-\cL(\cL A\cdot B)-\cL(A\cdot\cL B)+\tfrac12\cL^2A\cdot B+\cL A\cdot\cL B
      +\tfrac12A\cdot\cL^2B ;
  \]
\item by the Leibniz rules for even and odd factors (\texttt{LopC\_mul\_even}, \texttt{LopC\_mul\_odd}),
  $c_2=\frac12X_2(A,B)$ (\texttt{c2\_eq\_X2}, \texttt{LopC\_Xc}), where
  \[
    X_2(A,B)=\sum_{x,y}C(x,y)\bigl[X_C(\partial_{\bar\psi_x}A,\partial_{\psi_y}B)
      -X_C(\partial_{\psi_y}A,\partial_{\bar\psi_x}B)\bigr]
  \]
  is the double cross contraction and $X_C$ the single one;
\item hence (\texttt{FlowSeq.bubble\_eq})
  \[
    -\mathfrak B(C')=q_{j,j}(e_g)_g=\frac{\mu_4X_2\bigl(V_g(B_+),V_g(\Lambda)\bigr)}{4|B_+|}.
  \]
\end{itemize}
So $\mathfrak B$ is the quadratic form of the double contraction, with no tadpole terms (those
are in $b_{j,k}$; see (c)). The proof splits $e^{-\cL}A\,e^{-\cL}B$ into homogeneous parts of
degrees $8,6,4,2,0$ (\texttt{IsHom}). The pairing $\mu_4$ only sees the part of degree $4$, and
$\cL$ lowers degrees by $2$. 
\emph{Explicit form} (\texttt{H02UV/Step2/Coup/BubbleX.lean}). For products of vectors
$\iota^p=\iota p_0\cdots\iota p_n$:
\begin{itemize}
\item $\partial_a\iota^p=\sum_r(-1)^rp_r(a)\,\iota^{p\setminus r}$ (\texttt{D\_$\iota$Multi});
\item with the antisymmetric kernel $K_C(u,w)=\langle u,w\rangle_C-\langle w,u\rangle_C$, where
  $\langle u,w\rangle_C=\sum_{x,y}C(x,y)u(y,+)w(x,-)$ (\texttt{Xc\_$\iota$Multi}, \texttt{X2\_$\iota$Multi}),
  \[
    X_C(\iota^p,\iota^q)=\sum_{r,s}(-1)^{r+s}K_C(p_r,q_s)\,\iota^{p\setminus r}\iota^{q\setminus s},\qquad
    X_2(\iota^p,\iota^q)=-\sum_{r,s}(-1)^{r+s}K_C(p_r,q_s)\,X_C(\iota^{p\setminus r},\iota^{q\setminus s}).
  \]
\end{itemize}
Hence, for four vectors on each side (\texttt{mu4\_X2\_$\iota$Multi}),
\[
  \mu_4X_2(\iota^p,\iota^q)=-\sum_{r,s,r',s'}(-1)^{r+s+r'+s'}K_C(p_r,q_s)\,K_C(p'_{r'},q'_{s'})\,
  \det\bigl[\mathrm{col}(p''),\mathrm{col}(q'')\bigr].
\]
Here $p'=p\setminus r$ and $p''=p'\setminus r'$, and similarly for $q$. The columns are those of
\texttt{mu4\_$\iota$Multi}, i.e.\ first moments of lifts. With
$\hat M^{ii}_x=\frac12\sum_{e=\pm e_i}\iota^{v(e,x)}$, the legs $v(e,x)$ being those of $\hat M$
(\texttt{v4g}), this gives (\texttt{FlowSeq.bubble\_explicit})
\[
  -\mathfrak B(C')=\frac1{16|B_+|}\sum_{x\in B_+}\sum_{y\in\Lambda}\sum_{i,i'}\sum_{e=\pm e_i}\sum_{f=\pm e_{i'}}
    \mu_4X_2\bigl(\iota^{v(e,x)},\iota^{v(f,y)}\bigr),
\]
an explicit finite sum of products of two kernels $K_{C'}$ on the legs, times $4\times4$
determinants of lifts. The three line configurations correspond to the terms with $K_C\ne0$:
a leg of type $\psi$ is contracted with a leg of type $\bar\psi$. Not formalised: the universal
values of Proposition~\ref{prop:d2}.

\paragraph{Status of the proof.} (a)--(d) and the reduction of $\mathfrak B$ to $X_2$ are in Lean.
The universal values of $\mathfrak B(w+C')-\mathfrak B(w)$ and $C'(0,e_1)$ (Proposition~\ref{prop:d2})
are a separate, analytic task. They are not part of (c).

\subsection{Admissibility of the covariances, and the flow}
\label{sec:coupling-flow}
\begin{proposition}[C5, \texttt{GoodCov2} along the flow; proved in Lean except reality]\label{prop:C5}
Let $d\ge2$, $L\ge7$, and let $|\zeta_i|\le\bar\zeta(L)$ for all $i$, with $j+1<N$. Then
$C'_{j+1}=Z_{j+1}\Gamma^{\rm int}_{j+1}$ satisfies \texttt{GoodCov2} at scale $j$ with field scale
\[
  \ell=(c_S+K|\zeta|_{\le j})^{1/2}\,L^{-j(d-2)/2}\,L^{p_\Phi}
\]
(for $d=2$, $\ell=(c_S+K|\zeta|_{\le j})^{1/2}L^{p_\Phi}$). It is invariant under all lattice
automorphisms generated by translations, reflections and coordinate permutations, and it is
symmetric, $C'(x,y)=C'(y,x)$. For real $\zeta$ it is real.
\end{proposition}
\noindent\emph{Lean.} \texttt{Step2.prop\_C5} (\texttt{H02UV/Step2/Map/DCov.lean}), for the decomposition
\texttt{DressedFRD.frdC} with $q=2p_\Phi$. Correction to the earlier statement: the factor
$L^{p_\Phi}$ is needed. \texttt{cor\_gram} bounds $\Gamma^{\rm int}_{j+1}$ with derivative scale
$L^j$, while \texttt{GoodCov2} at scale $j$ uses $L^{j+1}$ (\texttt{GramBound.rescale}). This is the
same loss as in the old \texttt{goodCov\_Cfrd} ($\ell_0^2\ge KL^{d-2+2p_\Phi}$). The proof:
\begin{itemize}
\item the Gram part: \texttt{kmat\_isGram}, \texttt{kmat\_gramBound}, \texttt{cor\_gram};
\item the range: \texttt{Concl.range};
\item the symmetry: the symbol of $\Gamma^{\rm int}_{j+1}$ is invariant under signed permutations of
  momenta (\texttt{symSub}, \texttt{DData.Gint\_symMem}, the analogue of \texttt{Gint\_mem}), and
  this gives \texttt{CSym} (\texttt{kmat\_CSym}) and the symmetry (\texttt{kmat\_symm}, the
  reflection $x\mapsto-x$).
\end{itemize}
Reality of the matrix is not formalised (the symbol is real and even by
\texttt{DressedFRD.frd\_Gint\_real}); \texttt{Step2} does not use it.
The Gram part is \texttt{DressedFRD.cor\_gram} with \texttt{kmat\_gramBound} and
\texttt{gram\_rep}. The range is \ref{D2}. The invariance follows because $\Gamma^{\rm int}$ is
a Fourier multiplier with a lattice-symmetric symbol. Reality is \texttt{frd\_Gint\_real}.

\begin{proposition}[C6, history-dependent finite-horizon flow with growing $g$; proved in Lean]\label{prop:C6}
Let $\Lambda>1$, $\beta>0$, $0\le\delta\le\frac12$, $0\le\kappa<1$, $c_g,c_m,c_T>0$, $A,a\ge0$, and let
$G_j,m_j,T_j$ depend continuously and causally on the history $s_{\le j}=(s_0,\dots,s_j)$. With
$D_\epsilon=\{0<g\le\epsilon,\ |\mu|\le c_\mu g,\ \|x\|\le c_xg^3\}$, $c_\mu=\frac{2c_m}{\Lambda-1}$,
$c_x=\frac{2c_T}{1-\kappa}$ (\texttt{InD3}), call a history $s_{\le j}$ \emph{regular} if
$s_i\in D_\epsilon$ for $i\le j$ and its steps satisfy the band (F1),
$|g_{i+1}-g_i-\beta g_i^2|\le\delta\beta g_i^2+c_gg_i^3$ for $i<j$ (\texttt{InDh}, \texttt{RegH}). Assume,
for regular histories:
\begin{enumerate}
\item[(F1)] $|G_j-g-\beta g^2|\le\delta\beta g^2+c_gg^3$ (the coupling \emph{grows}; $\delta$ is a
  relative band around the one-loop coefficient);
\item[(F2)] $|m_j|\le c_mg$;
\item[(F3)] $\|T_j\|\le\kappa\|x\|+c_Tg^3$;
\item[(L)] for two regular histories, with $M\ge0$ bounding $|\delta g_i|+|\delta\mu_i|+\|\delta x_i\|$
  for $i<j$,
  \begin{align*}
    |G_j(s_{\le j})-G_j(s'_{\le j})|&\le(1+a)|\delta g_j|+a|\delta\mu_j|+A\|\delta x_j\|+aM,\\
    |m_j(s_{\le j})-m_j(s'_{\le j})|&\le A|\delta g_j|+a|\delta\mu_j|+A\|\delta x_j\|+aM,\\
    \|T_j(s_{\le j})-T_j(s'_{\le j})\|&\le\kappa\|\delta x_j\|+a(|\delta g_j|+|\delta\mu_j|)+aM.
  \end{align*}
\end{enumerate}
Let $\beta_+=(1+\delta)\beta+c_g$ and $b_{\max}=(1/\epsilon+N\beta_+)^{-1}$. If $0<\epsilon\le\epsilon_0^u$,
where $\epsilon_0^u=\min\{\epsilon_0(\beta,\kappa,c_g),\ c_m/(c_\mu\beta_+)\}$, and $a\le a_0^{\rm h}$ with
\[
  a_0^{\rm h}=\min\Bigl\{1,\frac1{K+Q},\frac{\Lambda-1}{2(AK+AQ+2+K+Q)}\Bigr\},\qquad
  Q=\frac4{1-\kappa},\quad K=\frac{2(4+AQ)}{\Lambda-1},
\]
then:
\begin{enumerate}
\item[(i)] for every $b\in(0,b_{\max}]$ there is a unique $\mu_0$ with $s_j\in D_\epsilon$ for all $j\le N$
  and $\mu_N=0$;
\item[(ii)] along this flow $|g_{j+1}^{-1}-g_j^{-1}+\beta|\le2(1+\beta)\delta\beta+C g_j$ with
  $C=2(\beta^2+(1+\beta)c_g)$;
\item[(iii)] $b\mapsto\mu_0(b)$ and $b\mapsto g_N(b)$ are continuous on $(0,b_{\max}]$.
\end{enumerate}
\end{proposition}
\noindent\emph{Lean.} \texttt{Flow2.Hist.fh\_flow\_hist} (\texttt{H02UV/Flow2/History.lean}). The maps act
on sequences $t=(t_i)_{i\in\mathbb N}$ and are \emph{causal}: $G_j(t)$ depends only on $t_0,\dots,t_j$.
The trajectory satisfies $s_{j+1}=(G_j(s),\Lambda\mu_j+m_j(s),T_j(s))$ (\texttt{traj\_succ}).
\begin{itemize}
\item The hypotheses are only needed along regular histories, since trajectories are regular by
  (F1) (\texttt{admissible\_InDh}, \texttt{trajc\_InDh}). This is what makes the RG bounds
  available: along a regular history $g_i\le(6/5)^{j-i}g_j$, so quantities built from the past
  (the gradient couplings of the previous step, $W^{\rm lin}_j$) are bounded in terms of the current
  $g_j$.
\item Existence uses the flow with $\mu$ clamped to $|\mu|\le c_\mu g$. It stays in $D_\epsilon$ by the
  invariant $1/g_i\ge1/b-i\beta_+$ (\texttt{trajc\_InDh}, from $G\le g+\beta_+g^2$ and
  \texttt{inv\_step}), and $1/b\ge1/\epsilon+N\beta_+$.
\item The clamp saturation propagates because $c_\mu G\le c_\mu g+c_mg$ (\texttt{cmu\_G\_le}, using
  $c_\mu\beta_+\epsilon\le c_m$).
\item In the uniqueness argument the $\mu$-differences increase along the trajectory. The induction
  carries $|\delta g_i|+|\delta\mu_i|+\|\delta x_i\|\le(1+p+q)|\delta\mu_k|$ for $i<k$, with $p=Ka$ and
  $q=Qa$.
\end{itemize}
For the RG, the history enters through $\zeta_0,\dots,\zeta_{j-1}$ and $W^{\rm lin}_j$, which are
functions of the previous coordinates. The band $\delta$ absorbs the non-universal parts of the
second-order coefficient: the lattice and mass corrections $O(L^{-j}+m^2L^{2j})$ and the
telescoping term $4\mathfrak n_j/g$, which are $O(1)$ against $\beta=\frac{3\log L}{2\pi}$, and the
dressing correction, which is $O(\bar\zeta^{L/256}\epsilon)$ (\ref{D3} with
\texttt{DressedFRD.Concl.first}). Proposition~\ref{prop:fh-flow} is the earlier version with
decreasing $g$; the $\bbH^{0|2}$ coupling grows towards the infrared
($\beta_{j+1}=\beta_j-\beta_1\log L$), so Proposition~\ref{prop:C6} is the one used.

\paragraph{Order of formalisation}
The planned order is:
\begin{enumerate}
\item C0, then C1;
\item C2 and C3 (algebraic);
\item C5 (it only uses \texttt{DressedFRD});
\item C6 (abstract);
\item C4 and C4'.
\end{enumerate}
The connection to Theorem~\ref{thm:flow} then verifies the hypotheses of
Proposition~\ref{prop:C6} for $s_j=(g_j,L^{2j}\nu_j,(W_j,K_j))$. The Lipschitz bounds (L) need
the analyticity part of Expected result~\ref{thm:step-rest}.

\subsection{The reduction of Theorem~\ref{thm:flow} to Proposition~\ref{prop:C6}}
\label{sec:reduction}
This subsection describes how the step estimates feed into Proposition~\ref{prop:C6}. Fix $d=2$.

\paragraph{Raw fields and causal field scales.} The covariances are not rescaled. The step at
scale $j$ uses the raw covariance $\Gamma^{\rm int}_{j+1}=Z_{j+1}^{-1}C'_{j+1}$, which is exact by
the splitting \ref{D1} (\texttt{omegaC\_full\_step\_raw}, \texttt{split\_kmat}). By
Proposition~\ref{prop:C5} and \texttt{GoodCov2.smul}, $\Gamma^{\rm int}_{j+1}$ satisfies
\texttt{GoodCov2} at the Gram scale $s_{j+1}\ell_C$, where $s_j=|Z_j|^{-1/2}$ (\texttt{raw\_family}).
Field norms at scale $j$ are taken at the \emph{causal field scale} $\ell_j=s_j\ell$ with $\ell=2\ell_C$.
\begin{itemize}
\item \emph{The two-scale step.} The step reads its input at $\ell_j$ and writes its output at
  $\ell_{j+1}=\ell_j|1+\zeta_j|^{-1/2}$. The covariance has Gram scale $\ell_{j+1}/2$, and
  $\ell_{j+1}+\ell_{j+1}/2\le2\ell_j$ for $|\zeta_j|\le\frac18$ (\texttt{scale\_ok}). The bounds of
  \texttt{Step2} hold in this setting, because $\E_C\theta$ only needs $\ell'+\ell_C\le2\ell$:
  \begin{itemize}
  \item \emph{Gaussian rescaling} (\texttt{H02UV/Step2/Map/ThetaScale.lean}): $\E_Cx=\E_{\lambda^{-2}C}\sigma_\lambda x$,
    where $\sigma_\lambda$ scales the fluctuation field (\texttt{ECC\_sigR}). Also
    $\|\theta_\lambda F\|_{T(\ell)}\le\|F\|_{T((1+\lambda)\ell)}$ for $\theta_\lambda F=F(\psi+\lambda\zeta)$,
    $\lambda\ge0$ (\texttt{tNorm\_thetaL\_le}).
  \item \emph{The two-scale expectation bound} (\texttt{H02UV/Step2/Map/TwoScale.lean}): if $C$ has a Gram bound
    at $\ell_C$ and $\ell'+\ell_C\le2\ell$, then
    $\|\E_C\theta F\|_{T_{j+1}(\ell')}\le q^{|\mathcal B_j(Y)|}\|F\|_{T_j(\ell)}$ (\texttt{EthE\_succ\_le2}).
  \item \emph{The doubled norm} (\texttt{H02UV/Step2/Map/TNorm2.lean}):
    $\|x\|_{T(\ell'\sqcup\ell_C)}=\|\sigma_{\ell_C/\ell'}x\|_{T(\ell')}$ (\texttt{T2}). It is
    submultiplicative and satisfies:
    \begin{itemize}
    \item $\|\E_Cx\|_{T(\ell')}\le\|x\|_{T(\ell'\sqcup\ell_C)}$ (\texttt{tNorm\_ECC\_le\_T2});
    \item $\|\theta F\|_{T(\ell'\sqcup\ell_C)}\le\|F\|_{T(\ell'+\ell_C)}$ (\texttt{T2\_theta\_le});
    \item $\|F\|_{T(\ell'\sqcup\ell_C)}\le\|F\|_{T(\ell')}$ on the first copy (\texttt{T2\_incl\_le}).
    \end{itemize}
  \item \emph{The bound chain}, generalised to input $\ell$, output $\ell'$ and Gram scale $\ell_C$
    with $\ell'+\ell_C\le2\ell$:
    \begin{itemize}
    \item the theorems are \texttt{Llin2\_bound\_ts}, \texttt{Knext2\_crude\_ts}, \texttt{step2\_ts},
      \texttt{step2\_V\_ts}, \texttt{step3\_ts}, \texttt{stepK\_ts}, \texttt{coords\_bound\_ts},
      \texttt{W\_contract\_ts} (this one also with $\ell_C\le\ell$), \texttt{Kplus\_lip\_ts},
      \texttt{Kplus\_lip\_cov\_ts}, and the $\lambda A$ versions in \texttt{H02UV/Step2/Map/Lambda.lean};
    \item \texttt{stepK\_tsW} also has the output weight $\mu A$, $\mu\ge1$, which absorbs the
      factor $e^{2^d\delta}$ of $K^\flat$;
    \item the single-scale statements are the specialisations $\ell'=\ell_C=\ell$.
    \end{itemize}
  \end{itemize}
\item \emph{Field rescaling} (\texttt{H02UV/Step2/Map/FieldScale.lean}). For $s>0$,
  $\|\sigma_sF\|_{T(\ell)}=\|F\|_{T(s\ell)}$ (\texttt{tNorm\_sig\_eq}). The rescaling $\sigma_s$ preserves the
  admissible block families and activities (\texttt{sigF\_mem}, \texttt{sigK\_mem}), and
  $\|\sigma_sF\|_{j}=\|F\|_{j}$ at $s\ell$ (\texttt{Fnorm\_sigF}, \texttt{Knorm\_sigK}).
\end{itemize}

\paragraph{Coordinates.} The state at scale $j$ consists of the raw couplings
$(z_j,y_j,\nu_j,g_j^{\rm raw})$, $W_j$, $K_j$, the field strength $Z_j=\prod_{i<j}(1+\zeta_i)$ and
the history $\zeta_{<j}$. The coordinates of Proposition~\ref{prop:C6} are $s_j=(g_j,\mu_j,x_j)$:
\begin{itemize}
\item $g_j=\mathrm{Re}\,(g^{\rm raw}_j/Z_j^2)$ and $\mu_j=\mathrm{Re}\,(L^{2j}\nu_j/Z_j)$;
\item $x_j=\sigma_{s_j}(D_j,K_j)$ with $D_j=W_j-W^{\rm lin}_j$, in the max norm of
  $\|\cdot\|_j\times\|\cdot\|_{j,A}$ at the fixed field scale $\ell$. So $\|x_j\|$ is the norm of
  $(D_j,K_j)$ at $\ell_j$.
\end{itemize}
Everything else ($z_j,y_j$, $W^{\rm lin}_j$, $Z_j$, $\zeta_{<j}$) is a causal function of
$s_0,\dots,s_{j-1}$, recomputed from the previous steps.

\paragraph{One step (Lean: \texttt{H02UV/Step2/Map/RGMap.lean}).}
\begin{itemize}
\item \texttt{rawHist} reconstructs the raw state from $t_0,\dots,t_{j-1}$.
\item \texttt{rawStep} absorbs $\zeta_j=\mathrm{clamp}_\rho((z_j+y_j)/Z_j)$, where \texttt{clampDisk} is
  the continuous radial clamp and $\rho\le\bar\zeta\le\frac18$. With $a=Z_j\zeta_j$ and
  $y_c=Z_j\,\mathrm{clamp}_\rho(y_j/Z_j)$ it absorbs \emph{both} gradient couplings through
  $\Delta=V_{(a-y_c,y_c,am^2,0)}$. This is allowed because the full step only uses
  $\Delta(\Lambda)=a(\psi,A\bar\psi)$ (summation by parts, \texttt{omegaC\_full\_step\_raw}). It leaves
  $V^\flat=V_{(z_j+y_c-a,\,y_j-y_c,\,\nu-am^2,\,g)}$, whose gradient part vanishes when the clamps are
  inactive, and $K^\flat=e^{\Delta}K$ (\texttt{Kflat}). Absorbing only $z_j+y_j$ in the $z$-slot would
  leave $(-y_j,y_j)$ in $V^\flat$, and $y_j$ would accumulate over all scales.
\item It then applies the \texttt{Step2} map with $\Gamma^{\rm int}_{j+1}(\zeta)$ and reads off
  $g_{j+1}$, $\mu_{j+1}$, $D_{j+1}$, $K_{j+1}$, with $x_{j+1}=\sigma_{s_{j+1}}(D_{j+1},K_{j+1})$.
\item Along every history $|\zeta_i|\le\rho$, $Z_j=\prod_{i<j}(1+\zeta_i)\ne0$ and
  $W^{\rm lin}_j\in$ \texttt{BSub} (\texttt{RawInv}, \texttt{rawInv\_hist}). The step hypotheses hold
  whenever the parameters are valid (\texttt{ValidAt}: geometry, $\rho\le\frac18$, and the
  covariance properties at Gram scale $s_{j+1}\ell/2$ for admissible dressed histories, which
  \texttt{raw\_family} supplies).
\item The flow data \texttt{rgData} ($G_j=g_{j+1}$, $m_j=\mu_{j+1}-L^2\mu_j$, $T_j=x_{j+1}$) are
  causal (\texttt{causal\_G}, \texttt{causal\_m}, \texttt{causal\_T}).
\end{itemize}

\paragraph{The hypotheses of Proposition~\ref{prop:C6}.}
\begin{enumerate}
\item[(F3), $K$.] \texttt{stepK\_tsW} at input $\ell_j$, output $\ell_{j+1}$, Gram scale $\ell_{j+1}/2$,
  with inputs $(V^\flat,W,K^\flat)$:
  \begin{itemize}
  \item $\|V^\flat\|_{\ell_j}\le c(|g|+|\mu|+|z_j/Z_j|+|y_j/Z_j|)$ (\texttt{Tj\_Vcoup\_le}; the
    normalisation by $Z_j$ is exactly the field scale $\ell_j$);
  \item $\|W\|\le\|W^{\rm lin}\|+\|D\|=O(\bar g^2)$;
  \item $\|K^\flat\|_{j,Ae^{-\delta}}\le e^{2^d\delta}\|K\|_{j,A}$ (\texttt{Knorm\_Kflat\_le}).
  \end{itemize}
\item[(F3), $D$.] $D_{j+1}=R_j+T_jD_j$ (\texttt{Wact\_sub\_Wlin\_succ}), with the contraction of $T_j$
  and $R_j=O(g^4)$.
\item[(LT).] From the Lipschitz bounds of Proposition~\ref{prop:step-lip}, two-scale.
\item[(F2), (Lm).] $\mu_{j+1}=L^2\mu_j+m_j$, $m_j=O(g)$ from the first-order coefficient and the
  bounds C2(iii), C4$'$(d).
\item[(F1), (LG).] $g_{j+1}=g_j+\beta_jg_j^2+O(g_j^3)$, where $\beta_j$ is the sum of the field-strength
  part $2C'_{j+1}(0,e_1)$, the bubble difference and the telescoping term. By
  Proposition~\ref{prop:d2}, $|\beta_j-\beta|\le\delta\beta$ with $\beta=\frac{3\log L}{2\pi}$ and
  $\delta=O(1/\log L)+O(\bar\zeta^{L/256}\epsilon)$, which is the band of Proposition~\ref{prop:C6}.
\end{enumerate}
The abstract part of this reduction is \texttt{H02UV/Flow2/Reduction.lean}: (F3) and (LT) follow from
bounds of the form
\begin{itemize}
\item $\|K_{j+1}\|\le\kappa_1\|K\|+C(g^3+g\|K\|+\|K\|^2)$ and $\|D_{j+1}\|\le\kappa_1\|D\|+Cg^4$
  (\texttt{RGBound});
\item Lipschitz bounds $\kappa_1\|\delta x\|+C\epsilon(\|\delta x\|+|\delta g|+|\delta\mu|+M)$ (\texttt{RGLip}).
\end{itemize}
Then \texttt{F3\_of\_rg}, \texttt{LT\_of\_rg}, \texttt{hhyp\_of\_rg} and \texttt{flow\_of\_rg} give the
conclusions of Proposition~\ref{prop:C6}, with (F1), (F2), (LG), (Lm), causality and continuity as
hypotheses.

\paragraph{Lean status of the estimates (\texttt{H02UV/Step2/Flow/FlowBounds.lean}, \texttt{CoupFlow},
\texttt{FlowStep}, \texttt{FlowRG}, \texttt{FlowHist}).}
\begin{itemize}
\item \emph{Field rescaling} (\texttt{H02UV/Step2/Map/FieldScale.lean}): besides \texttt{tNorm\_sig\_eq}, the
  Lipschitz bound in the scale, $\|\sigma_aF-\sigma_bF\|_{T(\ell)}\le\frac{|a-b|}\tau\|F\|_{T(\ell)}$ for
  $0\le a,b\le1-\tau$ (\texttt{tNorm\_sig\_sub\_le}), which is used for the dependence on the history
  through $s_j$.
\item \emph{The couplings of $V_{j+1}$} (\texttt{coupP\_bound}, \texttt{rawStep\_coup}):
  $\kappa_{j+1}=\kappa^{\flat\prime}-\kappa_Q'-\kappa_L$, where $'$ is the first-order shift with
  $c_0=C(x,x)$, $\delta=C(x+e,x)-C(x,x)$. Here $|\kappa_Q|\lesssim\|K^\flat\|$ (\texttt{coupQ\_le}) and
  $|\kappa_L|\lesssim v^2+w$ (\texttt{coupL\_le}). The Gram bound gives $|c_0|\le\ell_C^2$ and
  $|\delta|\le\ell_C^2/(2L^{2(j+1)})$ (\texttt{cov\_c0\_delta\_le}, from the first-difference profile,
  \texttt{norm\_cov\_diff\_le}).
\item \emph{The marginal coordinate and the normalisation.} The raw state carries the accumulated zero
  mode $w^{(1)}_j=\sum_{k<j}\sum_xC_k(x,0)$, and the marginal coordinate is $\check g=g(1+n)$ with
  $n=\mathrm{clamp}_{1/8}(4\nu w^{(1)})$. The couplings of step $j$ are normalised by
  $Z_{j+1}=Z_j(1+\zeta_j)$, the outputs by $Z_{j+2}$; so the field strength generated in a step enters
  $\check g_{j+1}$ without lag (with a lag the $g^2$-coefficient at $j=0$ would be $\beta/3$).
\item \emph{Normalised one-step bounds} (\texttt{step\_coup\_norm}, \texttt{rg\_step\_bounds} $=$
  \texttt{StepBounds}). With inactive clamps ($|z/Z|,|y/Z|\le\eta$), $g_r=\check g/(1+n)$ and
  $v=c_M\ell^2(2|\mu|+2\eta)+8\ell^4\check g$:
  \begin{itemize}
  \item $|z'/Z_{j+1}|,|y'/Z_{j+1}|\le E_z=\frac{\ell^2}{16}|g_r|+O(\|K\|+v^2+w)\le\frac{8\ell^2}{49}\check g+\dots$;
  \item $|\nu'-L^2\mu|\le\frac{16}7L^2E_z|\mu|+3L^2\eta+\frac{\ell^2}7|g_r|+O(L^2\|K\|+v^2+w)$;
  \item $\|W^{\rm lin}_{j+1}\|\le C_2v^2+\frac{C_0}L\|W^{\rm lin}_j\|$;
  \item $\|D_{j+1}\|\le\frac{C_0}L\|D_j\|+C_D(v^2+w)^2$;
  \item $\|K_{j+1}\|\le8C_0(L^{-1}+(A/2)^{-\eta})\|K_j\|+C_K(r^3+r\|K\|+\|K\|^2)$.
  \end{itemize}
  All norms are at the causal field scales. \texttt{stepK\_tsW} (output weight $2\cdot\frac A2$) absorbs
  $e^{4\delta}$ from $K^\flat$.
\item \emph{A priori bounds along regular histories} (\texttt{apriori}): $|z_j/Z_j|,|y_j/Z_j|\le C_zg_j$ and
  $\|W^{\rm lin}_j\|\le C_wg_j^2$, with $C_z=\frac{16\ell^2}{49}$. Hence the clamps are inactive on regular
  histories, and $V^\flat$ has no gradient part.
\item \emph{(Kb), (Yb)}: \texttt{RGBound} holds for \texttt{rgData} (\texttt{rgBound\_rg}), with
  $\kappa_1=8C_0(L^{-1}+(A/2)^{-\eta})$ and a constant independent of $c_T$.
\item \emph{(F2)} holds for \texttt{rgData} (\texttt{F2\_rg}), with
  $c_m\ge3L^2C_z+\frac{8\ell^2}{49}+O(\epsilon)$.
\item \emph{(F1)} (\texttt{H02UV/Step2/Flow/FlowF1*.lean}). Along a history, $\mathcal W^{\rm lin}$ is the $W^{\rm lin}$
  of a \texttt{FlowSeq} with $s\equiv1$ (\texttt{histFS\_Wlin}), so \texttt{FlowSeq.kappaS\_g} gives the
  exact $g$-component of $\Loc S$ (\texttt{hist\_LS\_g}); the $D$-part is $O(\|D\|)$ (\texttt{LS\_D\_bound}).
  In $P=\kappa'_g(1+n')/Z_{j+1}^2$ the term $4C^{(1)}g(\nu+4g\delta)$ cancels exactly against
  $(1+n')/(1+n)$ (\texttt{P\_identity}); the remaining $g^2$-terms are $4g_r\,a\,m^2w^{(1)}$ and $16g_r^2Z^2\delta w^{(1)}$
  (\texttt{rP\_bound}). They must be bounded \emph{sharply}: the Gram constant of \texttt{GoodCov2} scales
  like $L^{p_\Phi}=L^4$ (derivatives at the next scale), so entry bounds derived from it cannot enter the
  $g^2$-coefficient. With the sharp zero modes $|Z_{j+1}w^{(1)}_j|\le M_0L^{2j}$ and first differences
  $|Z_{j+1}\delta_j|\le M_2L^{-2(j+1)}$ they are $O(M_0M_2L^{-2})+O(\log L\cdot m^2L^{2j})$. With $(1+\zeta_{j+1})^{-2}$ (\texttt{F1\_alg}):
  \[
    \bigl|\check g_{j+1}-\check g_j-\mathrm{Re}\bigl(2Z_{j+1}C_j(e_1,0)-Z_{j+1}^2(q_{jj}+\textstyle\sum_{k<j}b_{jk})\bigr)\check g_j^2\bigr|
    \le e\,\check g_j^2+c_3\check g_j^3
  \]
  (\texttt{step\_F1}, \texttt{F1\_core}), and along regular histories (\texttt{F1\_rg}) the band
  holds with $\delta\beta\ge e+E_0$, given (HV) $|\mathrm{Re}(\dots)-\beta|\le E_0+E_1\check g_j$ with
  $E_0$ independent of $L$, and (HD) decay of the normalised past coefficients. By
  $\sum_{k<j}b_{jk}+q_{jj}=-(\mathfrak B(w+C_j)-\mathfrak B(w))+16\delta_j\hat w(0)$
  (\texttt{FlowSeq.past\_sum\_eq}, \texttt{tadT\_eq}), the bubble value (\texttt{bub\_line3}), the bilinear
  decay $|\mathfrak B(C'_k,C'_j)|\le KL^{-2(j-k)}$ (\texttt{bubX\_scale}) and line~1 of
  Proposition~\ref{prop:d2}, (HV) holds with $\beta=\frac{3\log L}{2\pi}$.
\item \emph{Continuity} of \texttt{rgData} (\texttt{H02UV/Step2/FlowCont/FlowCont.lean}, \texttt{rgData\_cont}), given
  continuity of $\Gamma^{\rm int}$ in $\zeta$ (\texttt{covCont\_of\_small}).
\item \emph{Lipschitz dependence of the dressing on the past}
  (\texttt{DressedFRD/LipDecay.lean}, \texttt{H02UV/Step2/Lip/CdrLip.lean}): for every $\rho\in(0,\frac12]$,
  $\|Z_{j+1}\Gamma^{\rm int}_{j+1}(\zeta)-Z'_{j+1}\Gamma^{\rm int}_{j+1}(\zeta')\|_{(j)}\le K\sum_{k\le j}\rho^{j-k}|\zeta_k-\zeta'_k|$
  (\texttt{lip\_decay}, \texttt{prop\_C5\_lip}). The decay comes from the contraction of the weighted tails, not
  from Cauchy estimates on single pieces.
\item \emph{Separate constants.} Proposition~\ref{prop:C6} is stated with separate constants $c_g$ (F1),
  $c_m$ (F2 and $c_\mu$) and $c_T$ (F3 and $c_x$). The remainder constant of (F3) depends on the
  $\mu$-domain through $|V^\flat|$, so with a single $c$ the conditions would be circular.
\item \emph{Assembly} (\texttt{H02UV/Step2/Flow/FlowMain.lean}, \texttt{flow\_rg}): the conclusions of
  Proposition~\ref{prop:C6} (\texttt{flow\_of\_rg}) for \texttt{rgData}, with causality, continuity,
  (Kb)/(Yb), (F2) and (F1) discharged, and with the one-loop inputs (HV), (HD), (HS) and the Lipschitz
  bounds as hypotheses. \texttt{validAt\_concrete} (\texttt{H02UV/Step2/Flow/FlowValid.lean}) gives the step
  hypotheses \texttt{ValidAt} for the dressed decomposition of FRD when $\ell\ge\ell_0(L)$ and
  $\rho\le\bar\zeta$.
\item \emph{Lipschitz bounds} (\texttt{H02UV/Step2/Lip/FlowLip*.lean}, \texttt{lip\_rg}). The coordinates are
  normalised by a \emph{complex} root $u_j=\prod_{i<j}(1+\zeta_i)^{1/2}$ ($x_j=\sigma_{u_j^{-1}}(D_j,K_j)$);
  with the real scale $|Z_j|^{-1/2}$ the phase $u_j/|u_j|$ accumulates over the history and the
  sensitivity of $T_j$ to the past grows like $j\check g_j^3$, which is not uniform in $N$. With the
  complex root, the normalised step is phase-free (equivariance \texttt{Kplus\_sig} etc.), Lipschitz in
  all inputs (\texttt{Kplus\_lip\_joint}, \texttt{WlinNext\_lip\_joint}, \ldots), and depends on the past
  only through the dressed covariances, with geometric decay (\texttt{CvH\_sub\_good}). This gives
  (LG), (Lm), (LT) with $a=C(\epsilon+\vartheta)$, under $m^2L^{2(j+2)}\le\vartheta$.
\item \emph{The flow theorem} (\texttt{H02UV/Step2/Flow/FlowThm.lean}, \texttt{thm\_flow}). For $L\ge L_0$ there are
  $\ell,A,\rho,\vartheta,\kappa,c_g,c_m,c_T,\epsilon$ such that for all $N$, $0<m^2\le L^{-2(N-1)}$ and
  every horizon $N_h$ with $N_h+2\le N$ and $m^2L^{2(N_h+1)}\le\vartheta$, the conclusions of
  Proposition~\ref{prop:C6} hold for \texttt{rgData} with $\beta=\frac{3\log L}{2\pi}$, $\delta=\frac12$:
  existence and uniqueness of the admissible $\mu_0$ for $b\in(0,b_{\max}]$, the band
  $|1/\check g_{j+1}-1/\check g_j+\beta|\le2(1+\beta)\delta\beta+C\check g_j$, and continuity in $b$.
  \texttt{flow\_of\_rg2} (\texttt{H02UV/Flow2/Reduction.lean}) separates the constant of (Kb)/(Yb)
  ($\le c_T$) from the Lipschitz constant (which depends on $c_T$).
\item \emph{The sharp marginal flow} (\texttt{H02UV/Step2/Flow/FlowF1S1}--\texttt{S3}, \texttt{H02UV/Step2/Flow/FlowSharp1}--\texttt{3}).
In the exact expansion of $\check g'$ (\texttt{P\_identity}) the term $g_r\cdot16g_rZ^2\delta w^{(1)}$ equals
$16\check g^2Z^2\delta w^{(1)}+16(g_r^2-\check g^2)Z^2\delta w^{(1)}$, and the second piece is $O(\check g^3)$ since
$|Z^2\delta w^{(1)}|\le\ell^4/32$. Moving the first into the $g^2$ coefficient
(\texttt{P\_identity2}, \texttt{F1\_core2}, \texttt{step\_F1s}) makes the coefficient
$\Re(2c-(Z^2q+S_B)+T_j)$ with the dressed tadpole $T_j=16(C'_{j+1}(e_1)-C'_{j+1}(0))\hat w'_j(0)$
(\texttt{htad\_eq}; $T_0=0$). With \texttt{HV\_sharp} (\texttt{HVs\_hist}), for $j\ge1$ (\texttt{F1\_rg\_sharp}):
\[
  |\check g_{j+1}-\check g_j-\beta\check g_j^2|\le\bigl(C_0(L^{-2j}+m^2L^{2(j+1)})+C_mm^2L^{2j}\bigr)\check g_j^2
    +c_g\check g_j^3,\qquad\beta=\tfrac{3\log L}{2\pi},
\]
with $C_0$ absolute. Along admissible trajectories (\texttt{thm\_flow\_sharp}, via \texttt{ginv\_sharp}),
$|1/\check g_{j+1}-1/\check g_j+\beta|\le C(L^{-2j}+m^2L^{2(j+1)})+C\check g_j$ for $1\le j<N_h$. Summing
(\texttt{thm\_flow\_sharp\_sum}), for $1\le n\le N_h$:
\[
  \Bigl|\frac1{\check g_n}-\frac1b+\beta n\Bigr|\le C(1+2\vartheta)+(1+\beta)\beta+C'\sum_{j<n}\check g_j .
\]
The step $j=0$ only contributes the band constant $(1+\beta)\beta$; a sharp constant needs
$\lim_T(\mathfrak B(w_T)-\frac{\log T}{2\pi})$. The sum $\sum_{j<n}\check g_j$ is bounded by a bootstrap
Let $B=1/b$, $K_0=C(1+2\vartheta)+(1+\beta)\beta$ and $K_1=C'$.
If $B\ge6(K_0+4K_1/\beta)$, strong induction gives $\check g_j\le6/B$, so $\sum_{j<n}\check g_j\le4/\beta$.
Otherwise $N_h<4(K_0+4K_1/\beta)/\beta$, and $\check g_j\le\epsilon$ gives the bound. Both cases use
$b\le b_{\max}$, i.e.\ $1/b\ge\frac32\beta N_h$. Hence,
for $0<b\le b_{\max}$ and $1\le n\le N_h$,
\[
  \Bigl|\frac1{\check g_n}-\frac1b+\beta n\Bigr|\le C_{\rm fl}
    =K_0+K_1\bigl(4+6\epsilon(K_0+4K_1/\beta)\bigr)/\beta ,
\]
with $C_{\rm fl}$ depending on $L$ only. This is (i) on the horizon allowed by $b_{\max}$. There
$\beta_\epsilon/(\beta_\epsilon-\beta n)\le3$, so the $O(\log(\beta_\epsilon\check g_n))$ is $O(1)$. The genuinely
logarithmic regime lies beyond $b_{\max}$; it is covered by the sharp range below.

\emph{The sharp range} (\texttt{H02UV/Flow2/RangeKeep}, \texttt{H02UV/Flow2/RangeSharp}, \texttt{H02UV/Flow2/RangeFlow},
\texttt{H02UV/Step2/Flow/FlowRangeHyp}, \texttt{H02UV/Step2/Flow/FlowRange}). The range $b\le b_{\max}$,
$1/b_{\max}=1/\epsilon+N_h\beta_+$, uses only the band: it forces $1/\check g_{N_h}\ge1/\epsilon+N_h(2\delta\beta+c_g)$, so
$\check g_{N_h}\to0$ as $N_h\to\infty$. In the proofs of Proposition~\ref{prop:C6} the range enters only through the a priori
bound that every clamped trajectory from $b$ stays a regular history in $D_\epsilon$ (\texttt{KeepD}); existence,
uniqueness, continuity and shooting hold on any set of such $b$ (\texttt{fh\_flow\_histK}). With the sharp (F1)
$|G_j-g-\beta g^2|\le e_jg^2+c\,g^3$, $\sum_{j<N}e_j\le E$ (\texttt{SharpHyp}), $1/g$ decreases by at most
$\beta+e_j+c/(1/\epsilon+\beta(N-j))$ per step along the worst case, and
$\sum_{k=1}^N(1/\epsilon+\beta k)^{-1}\le\beta^{-1}\log(1+\beta N\epsilon)$. Hence with
\[
  \frac1{b_{\max}^S}=\frac1\epsilon+\beta N+E+\frac c\beta\log(1+\beta N\epsilon)
\]
every clamped trajectory from $b\in(0,b^S_{\max}]$ stays in $D_\epsilon$, with $1/\check g_j\ge1/\epsilon+\beta(N-j)$
(\texttt{trajc\_InDhS}, \texttt{keepD\_of\_bmaxS}). For the RG map (\texttt{flow\_range\_hyp}: the proof of
\texttt{thm\_flow\_sharp} exporting \texttt{HHyp} and the sharp (F1) with $e_0\le\beta/2$,
$e_j\le K_e(L^{-2j}+m^2L^{2(j+1)})$) one has $E=\beta/2+K_e(1+2\vartheta)$, and \texttt{thm\_flow\_range} gives, for
$N_h$ as in \texttt{thm\_flow\_sharp}:
\begin{itemize}
\item existence and uniqueness of the admissible $\mu_0$ and continuity of $\mu_0(b)$, $\check g_{N_h}(b)$ on
  $(0,b^S_{\max}]$, where $1/b^S_{\max}=\beta N_h+O(\log(1+\beta N_h\epsilon))+O_L(1)+1/\epsilon$;
\item $1/\check g_{N_h}(b^S_{\max})\le1/\epsilon+(3+2\beta)E+\frac{c_g+C}{\beta}\log(1+\beta N_h\epsilon)$, and every
  $\beta_R$ above this bound is a final coupling: $\check g_{N_h}(b)=1/\beta_R$ for some $b\in(0,b^S_{\max}]$;
\item (i) in logarithmic form, along every admissible trajectory and $n\le N_h$ (\texttt{ginv\_log}, using
  $c_g\epsilon\le\beta/4$, so that $1/\check g$ decreases by at least $\beta/5$ per step):
  \[
    \Bigl|\frac1{\check g_n}-\frac1b+\beta n\Bigr|\le2(1+\beta)E+\frac{5C}{\beta}\log\frac{\check g_n}{b},
    \qquad C=2(\beta^2+(1+\beta)c_g).
  \]
\end{itemize}
The final coupling at the edge of the range is therefore of order $(1/\epsilon+O(\log N_h))^{-1}$, not of
order $\epsilon$ uniformly in $N_h$. The $\log N_h$ is genuine for this argument: the $c_g\check g^3$ error of
(F1) has no sign and depends on the history, so along the worst case it shifts $1/\check g$ by
$\sim(c_g/\beta)\log(\beta N_h\epsilon)$, and different $\mu_0$ (needed for the intermediate value argument)
may realise different signs.

\emph{Fixed $\beta_R$} (\texttt{H02UV/Flow2/FixedR}, \texttt{H02UV/Flow2/FixedR2}, \texttt{H02UV/Flow2/FixedR3}).
The intermediate value argument in $\mu_0$ over
\emph{all} clamped trajectories is replaced by a continuation in $b$, in which only the admissible trajectory and
its neighbours have to stay in $D_\epsilon$. Since the coupling is nondecreasing ($c_g\epsilon\le(1-\delta)\beta$),
an admissible trajectory with $\check g_N<y\le\epsilon$ has $\check g_j<y$ for all $j\le N$, and
$|\mu_j|<c_\mu\check g_j$ strictly (the faces $\mu=\pm c_\mu g$ are forward invariant and $\mu_N=0$). These are open
conditions on $(b,\mu_0)$ for the unclamped flow, and they imply that the trajectory is a regular history in
$D_\epsilon$. The cone argument of the uniqueness proof, with signs, gives
$\mu_0-\mu_0'\le\mu_N-\mu_N'$ for two such trajectories from the same $b$ with $\mu_0'\le\mu_0$
(\texttt{cone\_signed}). Hence $\mu_N(b_1,\mu_1\pm\eta)\gtrless0$, and by continuity in $b$ and the intermediate
value theorem on $[\mu_1-\eta,\mu_1+\eta]$ every $b$ near $b_1$ has an admissible $\mu_0$ near $\mu_1$ with
$\check g_N<y$ (\texttt{local\_adm}). The set of $b\in[b_0,b_2]$ with an admissible trajectory and
$\check g_N\le y$ is closed (projection of a compact set). Continuous induction from
$b_0=(1/y+N\beta_+)^{-1}\le b_{\max}$, where $\check g_N(b_0)\le y$, up to $b=y$, where $\check g_N\ge b=y$,
shows: \emph{for every $y\in(0,\epsilon]$ some $b\in(0,y]$ has an admissible trajectory with $\check g_N=y$}
(\texttt{exists\_traj\_eq}). Only \texttt{HHyp}, $\epsilon\le\epsilon_0^u$, $a\le a_0$ and
$c_g\epsilon\le(1-\delta)\beta$ are used; no a priori range of $b$. With \texttt{MonoHyp}, the set
$U=\{b:\text{admissible trajectory},\ \check g_N<\epsilon\}$ is an open interval, $\mu_0(b)$ and
$\check g_N(b)$ are continuous on it, $\check g_N$ is strictly increasing on all $b$ with an admissible trajectory,
and $\check g_N(U)=(0,\epsilon)$ (\texttt{fh\_flow\_fixedR}). For the RG map, the proof
of \texttt{flow\_range\_hyp} with the fading-memory bounds gives \texttt{HHyp},
\texttt{MonoHyp}, \texttt{MonoCoef} and the sharp (F1) with the same constants, and hence:
for every $\beta_R\ge1/\epsilon$ (independent of $N$, $N_h$) there is a unique $b$ with an admissible
trajectory and $\check g_{N_h}(b)=1/\beta_R$, $b\le1/\beta_R$; the admissible $\mu_0$ is unique for every $b$;
continuity and monotonicity as above; and (i) in the logarithmic form of \texttt{ginv\_log}.
\item \emph{The consolidated theorem} (\texttt{H02UV/Step2/Flow/FlowFinalZ.lean}, \texttt{thm\_flow\_final}). For $L\ge L_0$
there is one set of constants $\ell,A,\rho,\kappa,c_g,c_m,c_T,\epsilon,K_e$, independent of $N$, $m^2$
and $N_h$, such that for every $N$, every horizon $N_h+2\le N$ and every mass $0<m^2\le L^{-2(N-1)}$
(no condition $m^2L^{2(N_h+1)}\le\vartheta$) the flow of \texttt{rgData} has the following properties
(\texttt{FixedConcl}, \texttt{H02UV/Flow2/FixedRDiff2.lean}). Write $U=\{b:\text{admissible trajectory},\ \check g_{N_h}<\epsilon\}$.
  \begin{itemize}
  \item (i) in the logarithmic form, with $E=\beta/2+3K_e$.
  \item (ii) For every $\beta_R\ge1/\epsilon$ there is a unique $b(\beta_R)\in(0,1/\beta_R]$ with an admissible
    trajectory and $\check g_{N_h}=1/\beta_R$; the admissible $\mu_0$ is unique for every $b$. The set $U$ is an
    open interval, $\check g_{N_h}$ is strictly increasing and maps $U$ onto $(0,\epsilon)$.
  \item (iii) $\mu_0(b)$, $\check g_{N_h}(b)$ and all coordinates $(\check g_j,\mu_j,x_j)$, $j\le N_h$, are continuous
    and differentiable in $b$ on $U$. On $U$ one has $\check g_j\le\check g_{N_h}<\epsilon$, so the trajectory stays in
    the open part of the domain, and the expansion $|\mu_0-\mu_0'|\le|\mu_{N_h}-\mu_{N_h}'|$ holds near it
    (\texttt{H02UV/Flow2/FixedRDiff.lean}). Jointly in $(m^2,b)$: the set of $(m^2,b)$ with $b\in U(m^2)$ is relatively
    open, and $\mu_0$, $\check g_{N_h}$ and the trajectory are jointly continuous on it. The shooting solution
    $b(m^2,\beta_R)$ is jointly continuous for $\beta_R>1/\epsilon$ (\texttt{MsqConcl},
    \texttt{H02UV/Step2/FlowCont/FlowFixedRMsq.lean}). The proofs are a parametric version of \texttt{local\_adm}
    (\texttt{H02UV/Flow2/FixedRParam.lean}) and, for $b(m^2,\beta_R)$, strict monotonicity together with joint
    continuity of $\check g_{N_h}$.
  \item The domain bounds hold along every admissible trajectory: $0<b\le\check g_j\le\check g_{N_h}\le\epsilon$,
    $|\mu_j|\le c_\mu\check g_j$ and $\|K_j\|\le c_x\check g_j^3$, with strict inequalities $\check g_j<\epsilon$,
    $|\mu_j|<c_\mu\check g_j$ on $U$.
  \end{itemize}
  As in \texttt{thm\_flow\_hor2}, \texttt{HHyp} is used with large constants (existence, band, continuity), and
  everything that needs small Lipschitz constants goes through a separate channel for the hidden coupling
  $\zeta_j$ (\texttt{ZHyp}, \texttt{MonoZHyp}): the signed cone of the continuation argument
  (\texttt{cone\_signed\_z}), the expansion used for differentiability, and the strict monotonicity of
  $\check g_{N_h}$ (\texttt{traj\_strict\_mono\_z}; the $\mu$-cone carries $\|\delta\zeta_k\|\le s|\delta\mu_k|$ with
  $s$ of order one, the $\check g$-cone $\|\delta\zeta_k\|\le z_2\delta\check g_k$). The fixed-$\beta_R$ arguments are
  stated generically with the signed cone and the monotonicity as hypotheses (\texttt{H02UV/Flow2/FixedRParamG*.lean}).
  The (F1) mass term uses $|\zeta_j|\le c_\zeta\check g_j$, $c_\zeta=O(L^2)$, so $K_e=O(L^2)$.
  The horizon enters only through the hypothesis bundles (\texttt{FinConstZ}, \texttt{FinHypZ},
  \texttt{flow\_finalZ\_hyp}); \texttt{flow\_finalZ\_of} derives all conclusions on any set of masses where the
  bundles hold. Remaining deviations (see the status item below): $m^2=0$ is not covered;
  differentiability is in $b$ only; at $\beta_R=1/\epsilon$ (the right end of $U$) only existence and uniqueness are
  proved; the $\zeta_j$ bounds are not part of the statement; the constants include the extra conditions of
  differentiability ($4\ell^2(L^2c_\mu+c_m)\epsilon\le\frac18$), of $m^2$-continuity (the radius $\rho$) and of the
  $\zeta$-channel.
\item \emph{Status of Theorem~\ref{thm:flow}.} With the constants of \texttt{thm\_flow\_final} and
  $\beta_1\log L=\beta=\frac{3\log L}{2\pi}$, all three parts are proved for $N_h=N-2$, and
  \texttt{thm\_flow\_complete} carries the admissible trajectory through the last
  two steps (bounds, sharp band at $j=N-2$ and weak band for $\check g_N$, continuity of all last-step outputs,
  existence, uniqueness and $m^2$-continuity for the shooting on $\check g_N$; \texttt{thm\_flow\_complete\_band},
  \texttt{thm\_flow\_complete2}). Remaining deviations:
  \begin{itemize}
  \item \emph{Horizon and the last two steps.} The flow (with the final condition $\mu_{N_h}=0$, shooting
    and regularity) runs to $N_h=N-2$ (\texttt{thm\_flow\_hor2}, \texttt{thm\_flow\_final}). The steps $j=N-2$ and
    $j=N-1$ (the latter integrating only $C'^{\ne0}_N$) are appended by \texttt{thm\_flow\_complete}
    (\texttt{H02UV/Step2/Flow/FlowComplete.lean}) and its extension \texttt{thm\_flow\_complete\_band}
    (\texttt{H02UV/Step2/Flow/FlowCompleteBand.lean}), with one set of constants for all of the following:
    the bounds \texttt{GlueOut} ($|\check g_{N-1}|\le3\check g_{N-2}$, all final couplings
    $O(\check g_{N-2})$, $\|K_N(\Lambda)\|=O(\check g_{N-2}^3)$); the \emph{sharp band at $j=N-2$}
    (\texttt{BandSharp}, from \texttt{F1\_rg\_sharpL}: the (F1) chain under \texttt{ValidAtL}; the one-loop inputs
    of prop:d2 hold for every $j+1<N$): $|\operatorname{Re}\check g_{N-1}-\check g_{N-2}-\beta\check g_{N-2}^2|\le
    e\,\check g_{N-2}^2+C_s\check g_{N-2}^3$ with $e\le\beta/2$ and, for $N\ge3$,
    $e\le K_e(L^{-2(N-2)}+m^2L^{2(N-1)})$ --- at this step the mass term $m^2L^{2(N-1)}\le1$ does not vanish,
    it is $O(1)$ with an $L$-independent constant ($c_\zeta m^2L^{2(N-2)}=O(1)$) and fits into the band for $L$
    large; the weak band $|\check g_N-\operatorname{Re}\check g_{N-1}|\le C_b\check g_{N-2}^2$ for the final
    coupling $\check g_N=\kappa_{N,g}/Z_N^2$ of $V_N=V_{\kappa_N}$. This is the natural statement at the last
    step: there is no localisation at $j=N-1$, so $\check g_N=\check g_{N-1}/(1+n_{N-1})$ exactly
    (\texttt{lastG\_eq}) and the second order ($\beta\check g^2$) of that step sits in the unlocalised
    $\mathcal W_N=S$; a band for an ``effective'' $\check g_N$ would need a localisation at scale $N$ and the
    one-loop value of the bubble of $C'^{\ne0}_N$, neither of which is available. So (i) holds with the band up to
    $j=N-1$ and at $N$ with an $O(\check g^2)$ error. Continuity: \emph{all} outputs of the last step
    ($Z_N$, $\kappa_N$ hence $\nu_N$ and $\check g_N$, $u_N(\Lambda)$, $V_N$, $\mathcal W_N$, $K_N(\Lambda)$
    coefficient-wise) are jointly continuous in $(m^2,b)$ on the open set where $\check g_{N_h}<\epsilon$
    (\texttt{lastB\_contOn}, via \texttt{lastCov\_contM}, \texttt{rawStepN\_contM}, \texttt{KN\_holC}); and
    the shooting on the final coupling, \emph{existence}: for every $\beta_R\ge2/\epsilon$ there is $b$ with
    $\operatorname{Re}\check g_N(b)=1/\beta_R$, and every such $b$ has $|\check g_{N-2}(b)-1/\beta_R|\le
    4C_b\beta_R^{-2}$. \emph{Uniqueness and continuity in $m^2$} (\texttt{thm\_flow\_complete2},
    \texttt{H02UV/Step2/Flow/FlowComplete2.lean}): $b\mapsto\operatorname{Re}\check g_N(b)$ is strictly increasing on the
    admissible $b$, so the solution of $\operatorname{Re}\check g_N(b)=1/\beta_R$ is unique, and
    $m^2\mapsto b(m^2,1/\beta_R)$ is continuous. The proof adds one step to the monotonicity cone: the Lipschitz
    layer is redone at $j=N-2$ under the weaker validity condition (\texttt{H02UV/Step2/Last/FlowLastLip*}, same $C_L$),
    giving along the cone $\delta\check g_{N-1}\ge\frac34\delta\check g_{N-2}$ and
    $|\delta\mu_{N-1}|,\,L^{-2(N-1)}|\delta(Z_Nw^{(1)}_{N-1})|=O(\delta\check g_{N-2})$; the explicit last map
    $\check g_N=\check g_{N-1}/(1+n_{N-1})$ with $n_{N-1},\check g_{N-1}=O(\epsilon)$ then gives
    $\delta\operatorname{Re}\check g_N\ge\frac14\delta\check g_{N-2}$. Not proved: differentiability of the last two
    steps in $b$. The zero mode $tQ$ is left to Section~\ref{sec:corr}.
  \item \emph{The condition on $b_0$.} In (i)--(ii) the hypothesis ``$\check g_N\le2/\beta_0$ for the truncated
    flow'' is replaced by $\check g_{N_h}\le\epsilon$ for the flow itself; $\beta_0(L)=1/\epsilon(L)$.
  \item \emph{(i).} Proved in the form $|1/\check g_n-1/b+\beta n|\le C+C'\log(\check g_n/b)$, $1\le n\le N_h$,
    with $L^{2j}|\nu_j|\lesssim|\mu_j|\le c_\mu\check g_j$ and $\|K_j\|\le c_x\check g_j^3$. The bound on $\zeta_j$
    is $|\zeta_j|\le c_\zeta\check g_j$ with $c_\zeta=O(L^2)$ (\texttt{zeta\_hist}), not $O(\check g_j\log L)$, and it
    is not part of the statement. The step $j=0$ enters with the band constant only (a sharp constant needs
    $\lim_T(\mathfrak B(w_T)-\frac{\log T}{2\pi})$).
  \item \emph{(iii).} Continuity in $m^2$ and differentiability in $b$ are proved for $m^2>0$, on the open set
    where $\check g_{N_h}<\epsilon$; $m^2=0$ is not covered, differentiability is in $b$ only, and at
    $\beta_R=1/\epsilon$ only existence and uniqueness hold. The last two steps are continuous but not shown to be
    differentiable in $b$, and the shooting solution on $\check g_N$ is continuous in $m^2$ at fixed $\beta_R$
    (not jointly in $(m^2,\beta_R)$).
  \item Only the direction $i_0=0$.
  \end{itemize}
  \emph{Superseded intermediate versions} (still in Lean, used as building blocks): \texttt{thm\_flow}
  ($b\le b_{\max}$, $m^2L^{2(N_h+1)}\le\vartheta$), \texttt{thm\_flow\_sharp},
  \texttt{thm\_flow\_range}, \texttt{thm\_flow\_horizon}.
  The ingredients of the final versions:
  \begin{itemize}
  \item \emph{No mass condition.} $\zeta_{j+1}=-\check g_rZ_{j+1}C_j(e,0)$ up to remainders
    (\texttt{step\_zeta\_sharp}) and $|Z_{j+1}\Gamma^{\rm int}_j(x,y)|\le C_4L^2+KM_z$ (\texttt{Zcov\_hist}) give
    $c_\zeta=O(L^2)$ independent of $\ell$, so for $j\le N-3$ the mass term of (F1) is $O(L^{-2})\check g^2$; in (Lm)
    the mass term is kept as $4|\delta\zeta_j|$ and the uniqueness, continuation and monotonicity arguments carry
    the hidden differences $\delta\zeta$ in a separate channel (\texttt{ZHyp}, \texttt{MonoZHyp},
    \texttt{admissible\_unique\_z}, \texttt{cone\_signed\_z}, \texttt{traj\_strict\_mono\_z}).
  \item \emph{The step $j=N-2$} is the unchanged map: small $k$-sets are regions for $\Loc$ at every $k<N$ once
    $L>2^{d+4}+4$ (\texttt{region\_small\_L}), giving \texttt{Last.rg\_step\_bounds} under \texttt{ValidAtL}.
    The step $j=N-1$ is \texttt{LastN.rg\_last\_bounds} (Section~\ref{sec:corr}). The bound on $\check g_{N-1}$
    comes from the $g$-component of \texttt{Last.rawStep\_coup} (\texttt{g\_step\_bounds}).
  \item \emph{Uniqueness of $b_0$} uses, beyond Proposition~\ref{prop:C6}, the difference form of (LG) and fading
    memory ($w^{j-1-i}\mathrm{dist}(t_i,t_i')\le M$, $w=\frac34$) of the Lipschitz bounds (\texttt{MonoHyp},
    \texttt{lip\_rg\_z}): the past enters only through errors generated with geometric decay.
  \item \emph{Differentiability}: the step is holomorphic except for clamps, which are strictly inactive on the
    open part of the domain (\texttt{clamps\_of\_InDhs}); $\mu_0$ is differentiable by a one-dimensional implicit
    function argument in which the expansion of the uniqueness proof replaces invertibility
    (\texttt{hasDerivAt\_implicit}).
  \end{itemize}
\end{itemize}

%======================================================================
\section{The flow}
\label{sec:flow}

\subsection{Perturbative flow equations}
\label{sec:pt-flow}
By Lemmas~\ref{lem:first} and~\ref{lem:second} and Proposition~\ref{prop:d2}, the map of Section~\ref{sec:onestep} is,
for $j+1<N$,
\begin{align}
  \zeta_{j+1}&=-g^\flat\,C'_{j+1}(0,e_1)+\mathfrak z^{(2)}_j+r^\zeta_j,\label{e:zflow}\\
  \nu_{j+1}&=\nu^\flat+g^\flat\,\Delta C'_{j+1}(0)+\mathfrak v^{(2)}_j+r^\nu_j,\label{e:nuflow}\\
  g_{j+1}&=g^\flat+(g^\flat)^2\bigl(\mathfrak B(w'_{j+1})-\mathfrak B(w'_j)\bigr)-4g^\flat\mathfrak n_j+r^g_j.\label{e:gflow}
\end{align}
Here:
\begin{itemize}
\item $\mathfrak n_j=\nu' w'^{(1)}_{j+1}-\nu^\flat w'^{(1)}_j$, with $\nu'=\nu^\flat+g^\flat\Delta C'_{j+1}(0)$
  and $w^{(1)}=\hat w(0)$.
\item $\mathfrak z^{(2)}_j$ and $\mathfrak v^{(2)}_j$ are the explicit second-order coefficients
  $-\frac12\Loc[F_{w'_{j+1}}(V',V')-\E\theta F_{w'_j}(V^\flat,V^\flat)]$ in the gradient and mass
  monomials. They are $O(\bar g^2)$.
\item The remainders are $r_j=O(L^2\|K_j\|_j)+O(\bar g_j^3)$, using Theorem~\ref{thm:step}
  and Expected result~\ref{thm:step-rest}.
\end{itemize}
By \eqref{e:univ}, $C'_{j+1}(0,e_1)$ and $\mathfrak B(w'_{j+1})-\mathfrak B(w'_j)$ both equal
$\frac{\log L}{2\pi}+O(L^{-j}+m^2L^{2j}+\bar g)$.

\begin{corollary}[$\beta_1$]\label{cor:beta1}
Let $\check g_j=g^\flat_j(1+4\nu^\flat_jw'^{(1)}_j)$. Then
\[
  \check g_{j+1}=\check g_j+\frac{3\log L}{2\pi}\check g_j^2+O(\check g_j^3+\check g_j^2L^{-j}+\|K_j\|_j).
\]
So $\beta_j=1/\check g_j$ satisfies $\beta_{j+1}=\beta_j-\beta_1\log L+O(\beta_j^{-1})$ with
$\beta_1=\frac3{2\pi}$.
\end{corollary}

\begin{proof}
By \eqref{e:gflow}, $g_{j+1}=g^\flat+\frac{\log L}{2\pi}(g^\flat)^2-4g^\flat\mathfrak n_j+\dots$
Then $g^\flat_{j+1}=g_{j+1}(1+\zeta_{j+1})^{-2}$, and
$(1+\zeta_{j+1})^{-2}=1+2\frac{\log L}{2\pi}g^\flat+O(g^2)$ by \eqref{e:zflow}. The
$\mathfrak n$-term telescopes into the change of variables. This is Corollary~\ref{cor:oneloop}.
\end{proof}

\begin{remark}[How $\beta_1$ was confirmed]
\begin{itemize}
\item The field strength coefficient $-g\,C'(0,e_1)$ is an exact first-order Wick computation
  (Lemma~\ref{lem:first}, as in the proof of \texttt{thm:flow}).
\item The bubble coefficient is the double contraction of two copies of $M$. Computer algebra
  confirms that its local part is exactly proportional to $M$, with coefficient density
  $-f(s^3f')'$ for a radial propagator $f(|p|^2)$
  (\texttt{scripts/pt-checks/oneloop.py}, \texttt{general.py}).
\item Frullani's formula then gives the universal value $\frac{\log L}{2\pi}$.
\item The mass cross term $-4g\mathfrak n$ is checked in \texttt{masscross.py}.
\end{itemize}
\end{remark}

\subsection{An abstract finite-horizon flow}
\label{sec:fh-flow}
The flow of Theorem~\ref{thm:flow} has one relevant coordinate $\mu$ (the rescaled
mass, $\mu_j=L^{2j}\nu_j$), one marginal coordinate $g$, and irrelevant coordinates $x$ (here
$(W,K)$). The stable manifold theorem of \texttt{Flow} does not apply, because $g$ is not
contracted. It is replaced by the following finite-horizon statement, formalised in the package
\texttt{Flow2}.

\begin{proposition}[Finite-horizon flow; proved in Lean, \texttt{Flow2.fh\_flow}]\label{prop:fh-flow}
Let $\Lambda>1$, $\beta>0$, $0\le\kappa<1$, $c>0$, $A,a\ge0$ and $N\ge0$. Let $X_0,\dots,X_N$ be real
normed spaces. For $j<N$ let $G_j,m_j:\R\times\R\times X_j\to\R$ and
$T_j:\R\times\R\times X_j\to X_{j+1}$ be continuous. Consider the flow
\[
  g_{j+1}=G_j(s_j),\qquad\mu_{j+1}=\Lambda\mu_j+m_j(s_j),\qquad x_{j+1}=T_j(s_j),\qquad
  s_j=(g_j,\mu_j,x_j),
\]
started at $s_0=(b,\mu_0,0)$. Put $c_\mu=2c/(\Lambda-1)$ and $c_x=2c/(1-\kappa)$, and let
\[
  D_\epsilon=\{(g,\mu,x):0<g\le\epsilon,\ |\mu|\le c_\mu g,\ \|x\|\le c_xg^3\}.
\]
Assume, for $s=(g,\mu,x)$ and $s'=(g',\mu',x')$ in $D_\epsilon$:
\begin{enumerate}
\item[(F1)] $|G_j(s)-g+\beta g^2|\le cg^3$;
\item[(F2)] $|m_j(s)|\le cg$;
\item[(F3)] $\|T_j(s)\|\le\kappa\|x\|+cg^3$;
\item[(L)] $|G_j(s)-G_j(s')|\le(1+a)|g-g'|+a|\mu-\mu'|+A\|x-x'\|$,\\
  $|m_j(s)-m_j(s')|\le A|g-g'|+a|\mu-\mu'|+A\|x-x'\|$,\\
  $\|T_j(s)-T_j(s')\|\le\kappa\|x-x'\|+a(|g-g'|+|\mu-\mu'|)$.
\end{enumerate}
Put
\[
  \epsilon_0=\min\Bigl\{1,\frac\beta c,\frac{1-\kappa}{6(\beta+c)}\Bigr\},\qquad
  K_1=\frac{4(1+2A/(1-\kappa))}{\Lambda-1},\qquad
  a_0=\min\Bigl\{\frac1{K_1},\frac{\Lambda-1}{2(1+AK_1+2A/(1-\kappa))}\Bigr\},
\]
so that $\epsilon_0,a_0>0$ depend only on $(\Lambda,\beta,\kappa,c,A)$. If $\epsilon\le\epsilon_0$
and $a\le a_0$, then the following hold.
(Existence and (ii) need only $\epsilon\le\epsilon_0$; (L) and $a\le a_0$ are used only for
uniqueness and continuity.)
\begin{enumerate}
\item[(i)] (Existence and uniqueness.) For every $b\in(0,\epsilon]$ there is a unique
  $\mu_0=\mu_0^c(b)\in\R$ such that $s_j\in D_\epsilon$ for all $j\le N$ and $\mu_N=0$. It lies in
  $[-c_\mu b,c_\mu b]$.
\item[(ii)] (The marginal coordinate.) Along this flow, for $j<N$,
  $|g_{j+1}^{-1}-g_j^{-1}-\beta|\le C\,g_j$ with $C=2(\beta^2+\beta c+c)$. Hence
  $|g_j^{-1}-b^{-1}-\beta j|\le C\sum_{k<j}g_k$ for $j\le N$.
\item[(iii)] (Continuity and shooting.) $b\mapsto\mu_0^c(b)$ and $b\mapsto g_N(b)$ are continuous
  on $(0,\epsilon]$. Hence every value between $g_N(b_1)$ and $g_N(b_2)$ is $g_N(b)$ for some $b$
  between $b_1$ and $b_2$.
\end{enumerate}
\end{proposition}

\emph{Proof strategy.} The proof has three steps.
\begin{itemize}
\item \emph{Existence: clamping.} Replace $\mu_{j+1}$ by its clamp to $[-c_\mu g_{j+1},c_\mu g_{j+1}]$.
  The clamped flow stays in $D_\epsilon$ by (F1)--(F3), and it is continuous in $\mu_0$. The
  boundary is absorbing, because the $\mu$-direction expands: $(\Lambda-1)c_\mu\ge c$. Starting at
  $\mu_0=\pm c_\mu b$ therefore gives $\mu_N=\pm c_\mu g_N$. By the intermediate value theorem
  some $\mu_0$ has clamped $\mu_N=0$. Its clamped trajectory never touches the boundary, so it
  is the true trajectory.
\item \emph{Uniqueness.} Uniqueness follows from (L). For two trajectories in $D_\epsilon$, the
  difference of the $\mu$'s grows by a factor $\ge\frac{1+\Lambda}2$ per step and dominates the
  differences of $g$ and $x$.
\item \emph{Continuity.} Continuity follows from uniqueness and the continuity of the flow.
\end{itemize}
No fixed point theorem in an infinite-dimensional space is needed.

\emph{Lean.} \texttt{H02UV/Flow2/FiniteHorizon.lean}. The data are \texttt{FHData} ($G_j,m_j,T_j$ as
functions of the triple $s_j$), and the hypotheses are \texttt{FHHyp}. The pieces are:
\begin{itemize}
\item \texttt{exists\_admissible} (clamped flow \texttt{trajc}, absorbing boundary
  \texttt{trajc\_top}/\texttt{trajc\_bot}, IVT);
\item \texttt{admissible\_unique} (the $\mu$-difference grows by $\frac{1+\Lambda}2$ and dominates
  $|\delta g|\le K_1a|\delta\mu|$ and $\|\delta x\|\le\frac{2a}{1-\kappa}|\delta\mu|$);
\item \texttt{ginv\_admissible}, \texttt{ginv\_sum}, \texttt{continuousOn\_muc},
  \texttt{continuousOn\_gN}, \texttt{shooting};
\item positivity of the constants: \texttt{eps0\_pos}, \texttt{a0\_pos}.
\end{itemize}
The packaged statement is \texttt{fh\_flow}. Audit: \texttt{H02UV/Flow2/Audit.lean} (axioms
\texttt{propext}, \texttt{Classical.choice}, \texttt{Quot.sound}).

The assumption is $c>0$ rather than $c\ge0$. With $c>0$, a clamped trajectory that touches the
boundary ends at $\mu_N=\pm c_\mu g_N\ne0$. (For $c=0$ one has $m\equiv0$ on $D_\epsilon$ and the
statement is trivial.) In the application $c>0$.

\subsection{The flow theorem}
\begin{theorem}[Flow; proved in Lean, \texttt{Step2.thm\_flow\_final}, \texttt{thm\_flow\_complete\_band},
\texttt{thm\_flow\_complete2}]\label{thm:flow}
Let $d=2$, $i_0=0$ and $\beta=\frac{3\log L}{2\pi}$. There is $L_0$ such that for every $L\ge L_0$ there are
$\ell,A,\epsilon>0$ and constants $c_\mu,c_x,c_\zeta,C,C',C_b,C_s,K_e$ such that the following holds for every
$N\ge2$ and $m^2\in(0,L^{-2(N-1)}]$, with horizon $N_h=N-2$. Call $b>0$ \emph{admissible} if there is $\mu_0$
such that the flow with $\check g_0=b$, $\mu_0$ and the final condition $\mu_{N_h}=0$ stays in the domain
$\check g_j\le\epsilon$, $|\mu_j|\le c_\mu\check g_j$, $\|K_j\|\le c_x\check g_j^3$ for $j\le N_h$, and let
$U(m^2)$ be the set of admissible $b$ with $\check g_{N_h}<\epsilon$.
\begin{enumerate}[label=(\roman*)]
\item For admissible $b$, $\mu_0=a_0^c(b,m^2;N)$ is unique, the flow exists for all $j\le N$, and
  \[
    0<b\le\check g_j\le\check g_{N_h}\le\epsilon,\quad L^{2j}|\nu_j|\lesssim|\mu_j|\le c_\mu\check g_j,\quad
    \|K_j\|_j\le c_x\check g_j^3,\quad |\zeta_j|\le c_\zeta\check g_j\quad(j\le N_h),
  \]
  \[
    \Bigl|\frac1{\check g_n}-\frac1b+\beta n\Bigr|\le C+C'\log\frac{\check g_n}{b}\quad(1\le n\le N_h),
  \]
  with $c_\zeta=O(L^2)$. The last two steps satisfy $|\check g_{N-1}|\le3\check g_{N-2}$, all final couplings
  are $O(\check g_{N-2})$, $\|K_N(\Lambda)\|=O(\check g_{N-2}^3)$, the band
  $|\operatorname{Re}\check g_{N-1}-\check g_{N-2}-\beta\check g_{N-2}^2|\le e\,\check g_{N-2}^2+C_s\check g_{N-2}^3$
  with $e\le\beta/2$, and for the final coupling $\check g_N=\kappa_{N,g}/Z_N^2$ only
  $|\check g_N-\operatorname{Re}\check g_{N-1}|\le C_b\check g_{N-2}^2$.
\item The map $b\mapsto\operatorname{Re}\check g_N$ is continuous and strictly increasing on $U(m^2)$. For every
  $\beta_R\ge2/\epsilon$ there is a unique $b=1/\beta_\epsilon\in U(m^2)$ with $\operatorname{Re}\check g_N=1/\beta_R$; it
  satisfies $|\check g_{N-2}-1/\beta_R|\le4C_b\beta_R^{-2}$ and is continuous in $m^2$. The same holds for the
  shooting on $\check g_{N_h}$, for every value in $(0,\epsilon]$.
\item On $\{(m^2,b):b\in U(m^2)\}$, $a_0^c$ and all couplings $(\check g_j,\mu_j,x_j)$, $j\le N_h$, are jointly
  continuous in $(m^2,b)$ and differentiable in $b$; all outputs of the last two steps ($Z_N$, $\nu_N$,
  $\check g_N$, $u_N(\Lambda)$, $V_N$, $\mathcal W_N$, $K_N(\Lambda)$) are jointly continuous.
\end{enumerate}
\end{theorem}
Compared with the original target: $m^2=0$ and $i_0=1$ are not covered; the condition on $b_0$ is
$\check g_{N_h}\le\epsilon$ for the flow itself (so $\beta_0=1/\epsilon$); (i) is in log form, with a non-sharp
$\zeta$ bound and only a weak band at the last step (there is no localisation at scale $N$); the last two steps
are not shown to be differentiable in $b$. See the status item above for details.

\emph{Proof strategy.} The stable manifold theorem (\texttt{H02RG/Flow/StableManifold}) is replaced
by a finite-horizon fixed point:
\begin{itemize}
\item The relevant coordinate $\nu$ is solved backwards from $\nu_N=0$, using the expansion
  $L^2$ per step.
\item The marginal coordinate $\check g$ is solved forwards; its increments are explicit up to
  $O(\check g^3)$.
\item The irrelevant coordinates $(W,K)$ are solved forwards with the bounds of
  Theorem~\ref{thm:step}, Proposition~\ref{prop:W-contract} and Corollary~\ref{cor:W-uniform}
  (abstractly: Proposition~\ref{prop:fh-flow}, with $x=(W,K)$).
\end{itemize}
The sum $\sum_j\check g_j^3<\infty$ controls the remainders. The divergence of
$\sum_j\check g_j\approx\beta_1^{-1}\log(\beta_\epsilon/\beta_R)$ is the source of the $\log\log$
terms. Shooting in $b_0$ uses the intermediate value theorem, as in \texttt{H02RG/Ward/IVT}. The
increments $\zeta_j$ that enter the covariances are determined by the flow itself. By
\ref{D5} this is a triangular system: $\Gamma^{\rm int}_{j+1}$ depends only on
$\zeta_0,\dots,\zeta_j$, and these are determined by $V_0,\dots,V_j$.

%======================================================================
\section{The observable flow}
\label{sec:obs-flow}

\subsection{Localisation of observables in $d=2$}
\label{sec:loc-obs2}
\begin{lemma}[Observable localisation in $d=2$]\label{lem:loc-obs2}
There are $C,L_0$ such that for $L\ge L_0$, $j<N$, $\ell'\le2\ell$, $X\in\cS_j$ and $F\in\cN(X^\square)$:
\begin{enumerate}[label=(\roman*)]
\item if $a\in X$ and $F$ has the charge of the $a$-component, then
  $\|F-\Loc^a_XF\|_{T_{j+1}(\ell')}\le CL^{-1}\|F\|_{T_j(\ell)}$, where $\Loc^a_XF=(\sum_x F_{\bar x})\bar\psi_a$
  in Case~(1) and $\Loc^a_XF=F_\emptyset+(\sum_{x,y}F_{\bar xy})\bar\psi_a\psi_a$ in Case~(2); the same for $b$;
\item if $a\in X$ and $F$ is $U(1)$-invariant, then with $\hat x=r(x-x_X)$, $e=\hat b-\hat a$,
  $t(x)=\langle\hat x-\hat a,e\rangle/|e|^2$ and $\delta=\psi_b-\psi_a$, $\bar\delta=\bar\psi_b-\bar\psi_a$,
  \[\Loc^{ab}_XF=F_\emptyset+M_{00}\bar\psi_a\psi_a+M_{10}\bar\delta\psi_a+M_{01}\bar\psi_a\delta+M_{11}\bar\delta\delta
     +M_4\bar\psi_a\psi_a\bar\psi_b\psi_b,\]
  $M_{kl}=\sum_{x,y}F_{\bar xy}t(x)^kt(y)^l$, $M_4=\sum F_{\bar x_0x_1\bar x_2x_3}t(x_0)t(x_1)$
  (the substitution $\psi_x\mapsto\psi_a+t(x)\delta$), satisfies
  $\|\Loc^{ab}_XF\|_{T_j(\ell)}\le C\|F\|_{T_j(\ell)}$ and
  $\|F-\Loc^{ab}_XF\|_{T_{j+1}(\ell')}\le CL^{-1}\|F\|_{T_j(\ell)}$, uniformly in $|a-b|$.
  It reproduces $1$ and $\bar\psi_u\psi_v$, $u,v\in\{a,b\}$ (and $\bar\psi_a\psi_a\bar\psi_b\psi_b$).
\end{enumerate}
\end{lemma}
\begin{proof}[Sketch]
(i) $G=F-\Loc F$ has vanishing coefficients at length $0$ and (by $U(1)$) at the non-marginal short
lengths, and vanishing masses at the marginal length $m\in\{1,2\}$; the Taylor expansion of order $0$
at the base point gains $L^{-1}$ there. At lengths $n\ge3$, $G_n=F_n$, and the single-block
decomposition of Lemma~\ref{lem:pauli} freezes at most one variable of each spin, so $n-2\ge1$
variables carry a remainder $O(L^{-1})$; the bound $(4C)^2(4CA/L)^{n-2}\le C'L^{-1}$ is uniform in $n$.
The cost $(\ell'/\ell)^n\le2^n$ is absorbed per frozen variable as in Lemma~\ref{lem:loc2}.
(ii) $F-\Loc^{ab}F=(F-\Loc^aF)-D$ with $D$ the $\delta$-terms. $M_{kl},M_4$ are pairings with test
functions of weights $t$, which is affine with gradient $|e|^{-1}$; with
$\|\delta\|_{T_h(\ell')}\le\ell'|e|_1/h$ and $|e|_1^2\le d|e|^2$ each product is $O((L^j/h)(\ell'/\ell)^k)$,
i.e.\ $O(1)$ for $h=L^j$ and $O(L^{-1})$ for $h=L^{j+1}$.
\end{proof}
\begin{remark}
The sharp-region choice ($\psi_x\mapsto\psi_a$ for $x\in R_a$, else $\psi_b$) is not bounded: for
$F=\bar\psi_x\psi_y-\bar\psi_{x'}\psi_y$ across $\partial R_a$, $\|F\|_{T_j}\sim L^{-j}$ but
$\|\Loc F\|_{T_j}\sim|a-b|L^{-j}$. The barycentric weight is the fix. On scales where $a,b$ lie in a
common small set, $\bar\psi_a\psi_b-\bar\psi_a\psi_a$ and $\bar\psi_a\psi_a\bar\psi_b\psi_b$ are irrelevant;
keeping $\eta,\chi$ in $V^{\rm obs}$ is a bookkeeping choice, but with $\Loc^{ab}=\Loc_{\{a,b\}}$ the
non-mass part of $\mu_a\mu_bC'(a,b)\bar\psi_a\psi_b$ would land in $K^{\rm obs}$ without a small factor.
\end{remark}
\paragraph{Lean.} \texttt{H02UV/Obs2/Loc/LocContract.lean} (\texttt{exists\_odd\_contract2},
\texttt{exists\_even\_contract2}, \texttt{exists\_locA\_contract2}, \texttt{exists\_locB\_contract2}),
\texttt{H02UV/Obs2/Loc/LocContractLarge.lean} (the Pauli bound at lengths $\ge m$), \texttt{H02UV/Obs2/Loc/LocABt.lean}
(\texttt{locABt}, \texttt{exists\_locABt\_bound}, \texttt{exists\_locABt\_contract}),
\texttt{H02UV/Obs2/Loc/LocT.lean} (\texttt{locABt\_m4}: the reproduction of $\bar\psi_a\psi_a\bar\psi_b\psi_b$),
\texttt{H02UV/Obs2/Loc/LocRepro.lean} (\texttt{LocO2\_evalO}: $\Loc_{X,X}$ reproduces all observable couplings
for $a,b\in X$, $a\ne b$). The observable map (\texttt{Obs2.LocO2}, used by \texttt{H02UV/Obs2/Map/Map.lean})
now uses the barycentric $\Loc^{ab}$, based at $a$ if $a\in X$ and at $b$ otherwise
(\texttt{locABtX}, \texttt{locABtXY}); the sharp-region operator is kept but unused.
$\Loc\circ\Loc=\Loc$ is not formalised.

\subsection{Observable couplings}
The sources are those of \texttt{sec:obsfields}:
\begin{itemize}
\item Case~(1): anticommuting $\sigma_a,\bar\sigma_b$;
\item Case~(2): commuting $\sigma_x=\bar\vartheta_x\vartheta_x$.
\end{itemize}
In $d=2$, every source monomial times a derivative-free polynomial in the fields at $a,b$ has
dimension $0$, so it is marginal. By nilpotency and $U(1)$ invariance there are finitely many.
In normalised fields, the observable spaces are as in \texttt{def:cVobs}, plus a coupling $\chi$:
\begin{align*}
  \text{Case (1):}\quad V^{\rm obs}&=-\mu^{(1)}_a\sigma_a\bar\psi_a-\mu^{(1)}_b\psi_b\bar\sigma_b
  +\sigma_a\bar\sigma_b\bigl(\tfrac r2(\bar\psi_a\psi_a+\bar\psi_b\psi_b)+\chi\,\bar\psi_a\psi_a\bar\psi_b\psi_b\bigr),\\
  \text{Case (2):}\quad V^{\rm obs}&=-\mu^{(2)}_a\sigma_a\bar\psi_a\psi_a-\mu^{(2)}_b\sigma_b\bar\psi_b\psi_b
  -\sigma_a\sigma_b\tfrac\eta2(\bar\psi_a\psi_b+\bar\psi_b\psi_a)\\
  &\qquad+\sigma_a\sigma_b\bigl(\tfrac r2(\bar\psi_a\psi_a+\bar\psi_b\psi_b)+\chi\,\bar\psi_a\psi_a\bar\psi_b\psi_b\bigr),
\end{align*}
with $u^{\rm obs}$ as in \hr{}. The bare vertices are $\lambda^{(1)}_j=Z_j^{1/2}\mu^{(1)}_j$ and
$\lambda^{(2)}_j=Z_j\mu^{(2)}_j$.

The observable components of (A)--(D) are defined as in \texttt{sec:rgmap-ext}. The second-order
observable term $P^{\rm obs}$ of \texttt{e:Pobs-def} is replaced by the second-order map of
Section~\ref{sec:pt-map} applied to $V^\flat+V^{\rm obs,\flat}$, with remainder $W^{\rm obs}$, and with the
second-order term split over the blocks so that its observable part stays with the block carrying the observable:
\[S^\sharp(B)=\tfrac12F_C(V'_B,V'(\Lambda))+\tfrac12\bigl(F_C(V'_B,V'(\Lambda)^\emptyset)-F_C(V'^{\,\emptyset}_B,V'(\Lambda))\bigr)
  +\E\theta W(B),\]
i.e.\ $S^\sharp(B)^\emptyset=S(B)^\emptyset$ and $S^\sharp(B)_a=F_C(V'_{B,a},V'(\Lambda)_\emptyset)+(\E\theta W(B))_a$ (the same
for $b$). The correction is $\frac12\sum_{B'}G(B,B')$ with $G(B,B')=F_C(V'^{\,\rm obs}_B,V'_{B'})-F_C(V'_B,V'^{\,\rm obs}_{B'})$
antisymmetric, so $\sum_BS^\sharp(B)=\sum_BS(B)$; the single-block part of $K^{(2)}$ is changed by $-(S^\sharp(B)-S(B))$
and the split $J(B_+,U)$ of its two-block part carries $\frac12G$, which keeps $\sum_UJ(B_+,U)=0$. In step
(B), $\mu^{(1)}$ is multiplied by $(1+\zeta_j)^{-1/2}$; $\mu^{(2)},\eta,r$ by $(1+\zeta_j)^{-1}$;
and $\chi$ by $(1+\zeta_j)^{-2}$.

\subsection{Flow equations}
By Lemma~\ref{lem:second}, writing $C'=C'_{j+1}$ and $\mathfrak l_j=\mathfrak l(w'_{j+1})-\mathfrak l(w'_j)$:
\begin{align*}
  \mu^{(1)}_{x,j+1}&=\mu^{(1),\flat}_x(1-\mathfrak n_j)+O(\bar g^2|\mu|),\\
  \mu^{(2)}_{x,j+1}&=\mu^{(2),\flat}_x\bigl(1+g^\flat\mathfrak l_j-2\mathfrak n_j\bigr)+O(\bar g^2|\mu|),\\
  \gamma_{x,j+1}&=\gamma_{x,j}+\mu^{(2),\flat}_xC'(0,0),\qquad
  \eta_{j+1}=\eta^\flat(1-2\mathfrak n_j)-2\mu^{(2),\flat}_a\mu^{(2),\flat}_bC'(a,b),\\
  q_{j+1}&=q_j+\mu^{(1),\flat}_a\mu^{(1),\flat}_bC'(a,b)+r^\flat C'(0,0)\quad\text{(Case 1)},\\
  q_{j+1}&=q_j+\eta^\flat C'(a,b)-\mu^{(2),\flat}_a\mu^{(2),\flat}_bC'(a,b)^2+r^\flat C'(0,0)\quad\text{(Case 2)}.
\end{align*}
Up to the stated order, all errors are $O(L^2\|K^{\rm obs}_j\|)$. The couplings $r$ and $\chi$
are generated at order $\bar g$ and $\bar g^2$, and enter $q$ only as written.

There is no one-loop correction to $\mu^{(1)}$ other than the telescoping $\mathfrak n$. This is
because $\sigma_a\bar\psi_a$ has a single field (Lemma~\ref{lem:second}(iii)).

\begin{lemma}[Coalescence]\label{lem:coal}
If $|a-b|_\infty\ge\frac12L^{j+1}$, then $C'_{j+1}(a,b)=0$, by \ref{D2}. Hence
$q_j=\eta_j=r_j=\chi_j=0$ for $j<j_{ab}=\lfloor\log_L(2|a-b|_\infty)\rfloor$, exactly as in \hr{}.
\end{lemma}
\paragraph{Lean status of Section~\ref{sec:obs-flow}.}
\begin{itemize}
\item \emph{Lemma~\ref{lem:coal}} is proved: \texttt{Obs2.obsHist\_coal} (\texttt{H02UV/Obs2/Pair/CoalHist.lean}),
  and \texttt{obsHist\_coal\_valid} with the range of $C'_{k+1}$ from \texttt{ValidAt}. For $L\ge 2^5$
  (with $2^5+16<L$), along any history, if the observable couplings start with $q=\eta=r=\chi=0$
  (\texttt{coalInv\_init}), then the $ab$ components of $V^{\rm obs}_j(B)$ and of the accumulated
  $u^{\rm obs}_j$ vanish for all $j\le j_{ab}$ (one scale more than stated). The proof is a support
  argument (\texttt{H02UV/Obs2/Map/Supp.lean}, \texttt{H02UV/Obs2/Pair/Coal.lean}): $\sigma_a$-free elements form a
  subalgebra; $V(B)$ has no $\sigma_a$ if $a\notin B$, $\mathcal W(B)$ and $K(X)$ have none at distance
  $\ge L^j$ from $a$, and $V(B)_a\in\mathcal N(\{a\})$. The only new $\sigma_a\sigma_b$ terms are
  $F_C(V'_{B,a},V'_{B',b})$, which vanish since $C(a,b)=0$. Products $(\sigma_a\cdots)(\sigma_b\cdots)$
  occur only in objects close to both points. Below $j_{ab}$ no block, $(j+1)$-block or small set is
  close to both. Note that $\mathcal W^{\rm obs}_{j+1}(B_+)_{ab}$ and $K^{\rm obs}_{j+1}(U)_{ab}$
  ($U$ large) need not vanish below $j_{ab}$; only small sets enter $\Loc$.
\item \emph{Flow equations.} \texttt{H02UV/Obs2/Flow/FlowMu.lean}: the exact recursion at the block $B\ni a$,
  $\beta_{j+1}=\beta_j-\sum_{X\in\cS_j,X\supseteq B}m(K_j(X)_a)-m(S(B_+)_a)$, with $\beta=-\mu^{(1)}$,
  $m=\sum_xF_{\bar x}$ (Case~1, \texttt{VpO\_a\_one}) resp.\ $\beta=-\mu^{(2)}$, $m=\sum_{x,y}F_{\bar xy}$
  (Case~2, \texttt{mass2\_VpO\_a}, \texttt{VpO\_a\_form\_two}). $S(B_+)$ is the second-order term
  ($\frac12F_C(V'_B,V'(\Lambda))+\E_C\theta\mathcal W(B)$ summed over $B\subseteq B_+$). The $K$ error is
  $\le c\,C\ell^{-p}\sup_X\|K(X)_a\|_{T_j(\ell)}$ (\texttt{K\_term\_bound}). 
\item \emph{Wick evaluation} (\texttt{H02UV/Obs2/Flow/Wick1}, \texttt{WickPair}, \texttt{WickV}, \texttt{WickVcoup},
  \texttt{Wick2}, \texttt{WickTotal}), exact for every complex $C$: $F_C(\iota v,B)=-X_C(\iota v,B)$, so
  $m_1F_C(\bar\psi_a,B)=0$ for $B$ homogeneous of degree $\ne2$ (Lemma~\ref{lem:second}(iii), first part);
  $F_C(\iota u\iota w,B)=X_C+\frac12X_2$; with $V_\kappa(\Lambda)$:
  $m_1F_C(\bar\psi_a,V_\kappa(\Lambda))=\nu\sum_xC(a,x)$ and
  $m_2F_C(\bar\psi_a\psi_a,V_\kappa(\Lambda))=\nu\sum_x(C(a,x)+C(x,a))-g\,\mathfrak l_C(a)$,
  $\mathfrak l_C(a)=\sum_i\sum_z\frac12\sum_{e=\pm e_i}\nabla_eC(a,\cdot)(z)\,\nabla_eC(\cdot,a)(z)$
  (\texttt{mass1\_FC\_psibar\_Vcoup}, \texttt{mass2\_FC\_pbp\_Vcoup}, lifted to sources in \texttt{Wick2}).
  The per-block term $m(S(B_+)_a)$ is \emph{not} symmetric: $\frac12F_C(V'_{\rm bulk}(\Lambda\setminus B_+),V'_a)$
  stays in $\mathcal W^{\rm obs}$ of neighbouring blocks and is re-localised later. Summed over all blocks the
  localisation cancels (\texttt{total\_step}): $\sum_B[m(V_{j+1}(B)_a)-m(\mathcal W_{j+1}(B)_a)]$ equals
  $\sum_Bm(V'_{B,a})-\sum_Bm((\E\theta Q(B))_a)-m(F_C(V'(\Lambda)_a,V'(\Lambda)_\varnothing))-\sum_Bm((\E\theta\mathcal W(B))_a)$,
  with the symmetric one-loop term evaluated in \texttt{total\_FC\_one}, \texttt{total\_FC\_two}. So the
  flow equations above hold for $\mu+w$, $w$ the total $\mathcal W^{\rm obs}$ mass coefficient. For a
  translation-invariant even covariance $C(x,y)=c(y-x)$, $\mathfrak l_C(a)=\sum_xc(x)((-\Delta)c)(x)=\mathfrak l(C)$
  (summation by parts). Not done:
  the dressed $\mathfrak l_j$ and the tadpole part $-(\nu'-\nu)w^{(1)}$ of $\mathfrak n_j$ inside
  $m((\E\theta\mathcal W)_a)$.
\item \emph{Forward iteration.} \texttt{H02UV/Obs2/Flow/LinIter.lean}: the Lyapunov bound
  $\Phi_{j+1}\le(1+a_j+ABg_j^2/(1-\theta))\Phi_j$, $\Phi=m+\frac A{1-\theta}k$, for the linear system
  $m_{j+1}\le(1+a_j)m_j+Ak_j$, $k_{j+1}\le\theta k_j+Bg_j^2m_j$ (\texttt{lin\_iter}, \texttt{lin\_iter\_K}),
  and $\mu_n=P_n(\mu_0+\sum_je_j/P_{j+1})$ (\texttt{renorm\_eq}, \texttt{renorm\_bound}).
\item \emph{Block bounds} (single sources, weights $w_{ab}=0$; \texttt{H02UV/Obs2/Est/EstNorm}, \texttt{EstBlock},
  \texttt{EstMu}). In the observable part of the norm, with bulk radius $r$ of $V^\flat$ and observable
  sizes $m_V,m_{\mathcal W},k$: $\|S(B)^{\rm obs}\|\le 2c\,Q_1^2rm_V+Q_1m_{\mathcal W}$ (\texttt{Sb\_o}),
  $\|V^{\rm obs}_{j+1}(B)\|\le 2(Q_1m_V+Ck+CL^d(rm_V+m_{\mathcal W}))$ (\texttt{Vp\_o}),
  $\|\mathcal W^{\rm obs}_{j+1}(B)\|\le C(L)(rm_V+m_{\mathcal W})$ (\texttt{Wp\_o}). Together with the exact
  recursion this gives $|\beta_{j+1}-\beta_j|\le C\ell^{-p}k+C(L)\ell'^{-p}(r|\beta|\ell^p+m_{\mathcal W})$
  (\texttt{mu\_step\_one}, \texttt{mu\_step\_two}): the $m$-half of \texttt{lin\_iter}, with $a_j=O(\bar g_j)$.
\item \emph{Pointed localisation.} With $\Loc^a_X=0$ for $a\notin X$, $(1-\Loc_X)$ would be the identity on
  $\sigma_a$-components of $K(X)$ for $X\not\ni a$; such components are generated by $\mathcal W^{\rm obs}(B)$,
  $B\not\ni a$, and would not be contracted. The map therefore localises at $a$ for every $X$:
  the pointed operator is $\operatorname{loc}^a_{X,Y}F$ for $a\in X$ and
  $(|Y|/|X|)\operatorname{loc}^a_{\{a\},\{a\}}F$ for $a\notin X$
  (\texttt{locAp}, \texttt{locBp}, \texttt{H02UV/Obs2/Loc/Loc.lean}; for $a\notin X^\square$ the components vanish, so only
  $a\in X^\square\setminus X$ matters). Then $\Loc_{X,X}$ extracts the full mass of $F_a$ for every polymer, and
  the contraction $\|F-m(F)\bar\psi_a\|_{T_{j+1}(\ell')}\le CL^{-1}\|F\|_{T_j(\ell)}$ for $a\in X^\square$ holds
  (\texttt{exists\_odd\_contract\_pt}, \texttt{exists\_even\_contract\_pt}). The coalescence, support and
  reproduction statements are adapted (\texttt{H02UV/Obs2/Map/Supp}, \texttt{Coal}, \texttt{CoalHist}, \texttt{EstLoc}:
  far-from-$a$ conditions are now distances $\ge L^j$, and $\operatorname{loc}^a$ is supported on $Y^\square\cup\{a\}$).
\item \emph{Total coupling} (\texttt{H02UV/Obs2/Flow/TotalMu.lean}). With the pointed operator the coupling is the total
  $\beta=\sum_Bm(V(B)_a)$, the coefficient of $\sigma_a\bar\psi_a$ in $V^{\rm obs}(\Lambda)$ (a per-block flow is no
  longer meaningful). $\sum_Bm(\mathcal W_{j+1}(B)_a)=0$ (\texttt{total\_W\_mass}), and
  $\beta'=\beta(1-\nu'\sum_xC(a,x))-\sum_{X\in\mathcal S_j}m_1(K(X)_a)-\sum_Bm_1((\E\theta\mathcal W(B))_a)$, resp.
  $\beta'=\beta(1-\nu'\sum_x(C(a,x)+C(x,a))+g'\mathfrak l_C(a))-\sum_Xm_2(K(X)_a)-\sum_Bm_2((\E\theta\mathcal W(B))_a)$
  (\texttt{total\_mu\_one}, \texttt{total\_mu\_two}). The $\mathcal W$ term is the tadpole
  $\sum_B[m(e^{\mathcal L_C}\mathcal W_j(B)_a)-m(\mathcal W_j(B)_a)]$ of the cubic part of $\mathcal W^{\rm obs}$
  (order $\bar g|\mu|$, containing the $(\nu'-\nu)w^{(1)}$ part of $\mathfrak n_j$) plus the product part of
  $\prod(1+\mathcal W)-1$ (order $\bar g^2|\mu|$).
\item \emph{The $\mathcal W^{\rm obs}$ tadpole and the $\mu^{(1)}$ flow} (\texttt{H02UV/Obs2/Flow/MuTad0--3.lean}). With
  $G_j(\Gamma)=\sum_Bm(e^{\mathcal L_\Gamma}\mathcal W_j(B)_a)$, the $\mathcal W$ term is $G_j(C_j)$, and exactly
  $G_{j+1}(\Gamma)=T_j(\Gamma)+G_j(\Gamma+C_j)-G_j(C_j)+P_{j+1}(\Gamma)$, with $T_j$ the tadpole of the one-loop term
  $F_{C_j}(V'_a,V'_\varnothing)$ and $P_{j+1}$ the product part (\texttt{wTad\_step}). The tadpole is evaluated exactly:
  $m_1(e^{\mathcal L_\Gamma}F_C(\bar\psi_p,V_\kappa(\Lambda)))-m_1(F_C(\bar\psi_p,V_\kappa(\Lambda)))=g\,\tau(C,\Gamma)$ with $\tau$ bilinear,
  and $\tau(C,\Gamma)=\Delta\Gamma(0)\sum_zC(p,z)$ for translation-invariant $\Gamma$
  (\texttt{mass1\_expLC\_FC\_psibar\_Vcoup}, \texttt{tauC\_TI}). Along a sequence of steps (hypothesis bundle
  \texttt{MuChain}), \texttt{mu\_flow\_exact} gives the exact identity
  $\beta_{j+1}=\beta_j(1-\mathfrak n_j)-\sum_{X\in\cS_j}m_1(K_j(X)_a)-R_j$, $\mathfrak n_j=\nu'_jC_j^{(1)}(a)+g'_j\tau(w_j,C_j)$,
  which equals $\nu'w_{j+1}^{(1)}-\nu^\flat w_j^{(1)}$ for translation-invariant $C_j$ (\texttt{frakN\_TI}). The remainder
  $R_j=\sum_{k<j}(\beta_kg'_k-\beta_jg'_j)\tau(C_k,C_j)+(\text{product parts})$ is explicit, and
  \texttt{mu\_flow\_bound} gives $|\beta_{j+1}-\beta_j(1-\mathfrak n_j)|\le E_K+ABq/(1-q)^2+E_P$ under the bounds
  $|\sum_Xm_1(K_j(X)_a)|\le E_K$, $|\beta_kg'_k-\beta_jg'_j|\le A(j-k)$, $|\tau(C_k,C_j)|\le Bq^{j-k}$ and product
  parts $\le E_P$ (hypothesis bundle \texttt{MuErr}). Case~(1) only.
\item \emph{$K^{\rm obs}$ step, reduction} (\texttt{H02UV/Obs2/Est/KObs1--2.lean}). The locality hypothesis of the block
  bounds now has supports $Y^\square\cup\{a,b\}$ and is discharged for the pointed $\Loc$ (\texttt{bho\_hlocl}).
  With sources rescaled by $\lambda=\bar g/m$, $\|F\|_\lambda=\|F_\varnothing\|+\lambda\|F^{\rm obs}\|$ (\texttt{oTj\_scale}), and a
  bound $\|K'\|_\lambda\le\theta_0\|K\|_\lambda+C_1((2\bar g)^3+2\bar g\|K\|_\lambda+\|K\|_\lambda^2)$ implies
  $k'\le(\theta_0+O(\bar g))k+B\bar g^2m$ (\texttt{lam\_trick}, \texttt{kobs\_step}).
\item \emph{Exact one-step identity with sources} (\texttt{H02UV/Obs2/Map/ExactStep.lean}). The map with sources is an instance of the
  generic map of \texttt{Obs.GenericDefs} ($\theta$, $\iota$, $\E_C$ componentwise; \texttt{Knext2O\_eq\_KnextG}), so
  $\E_C\theta Z_j=e^{-u_{j+1}(\Lambda)}\sum_U\tilde I_{j+1}^{\Lambda\setminus U}K_{j+1}(U)$ holds exactly (\texttt{stepO\_ZVK}).
  The factorisation $K_{j+1}(U_1\cup U_2)=K_{j+1}(U_1)K_{j+1}(U_2)$ (\texttt{Knext2O\_eq\_KextO}) uses the supports
  $\mathcal N(U^{\square\square})$: the block factors of scale $j+1$ (with the nonlocal $\mathcal W_{j+1}$) enter only through
  $\iota$ (\texttt{SuppHypW}), and the pointed $\operatorname{loc}^a_{Y,B}$ is local in $B^{\square\square}$ provided the
  $\sigma_a$ ($\sigma_b$, $\sigma_a\sigma_b$) components of $K(Y)$, $Y$ small, vanish unless $a\in Y^\square$ ($b\in Y^\square$)
  (hypothesis \texttt{OLocHyp}, together with the locality of all components of $V$, $\mathcal W$, $K$; \texttt{J2O\_mem}).
  This needs $L\ge2^{d+3}$. With the change of variables (\texttt{reblock\_ZVKO}, \texttt{KoutO\_union}):
  $\E_C\theta Z_j=e^{-u_{j+1}(\Lambda)}Z_{j+1}$ with sources (\texttt{stepK\_ZVKO}).
\item \emph{Locality along the flow} (\texttt{H02UV/Obs2/Loc/LocFlow.lean}, \texttt{H02UV/PtCorr2/Obs/ObsLoc.lean}). Let $\mathcal N^{\rm obs}(S)$
  (\texttt{osub}) be the even elements whose components lie in $\mathcal N(S)$ and whose $\sigma_a$ ($\sigma_b$) and
  $\sigma_a\sigma_b$ components vanish unless $a\in S$ ($b\in S$). If $V(B),\mathcal W(B)\in\mathcal N^{\rm obs}(B^\square)$ and
  $K(Y)\in\mathcal N^{\rm obs}(Y^\square)$ for connected $Y$ (\texttt{OLocS}), the same holds at scale $j+1$ for
  $V_{j+1}$, $\mathcal W_{j+1}$, $K_{j+1}$ (\texttt{OLocS.step}; $L\ge2^{d+3}$, covariance of range $L^{j+1}/2$). The
  pointed $\operatorname{loc}^a_{Y,B}$ is local in $B^{\square\square}\subseteq(B_+)^\square$ (\texttt{nb2\_subset\_succ}).
  \texttt{OLocS} implies \texttt{OLocHyp}; the initial state \texttt{obsInit} is local for all couplings, so
  \texttt{OLocHyp} holds at every step of every trajectory (\texttt{oLocH\_traj}).
\item \emph{The $\mu^{(1)}$ flow along the history} (\texttt{H02UV/Obs2/Flow/MuHist.lean}). Since both branches of the pointed
  $\operatorname{loc}^a$ are multiples of $\bar\psi_a$, $V_{j+1}(B)_a\in\mathbb C\bar\psi_a$ for every block $B$
  (\texttt{VpO\_a\_one\_span}). The bulk part of $\E_{C_k}\theta V_k(\Lambda)$ is $V_{\kappa'_k}(\Lambda)+{\rm const}$, where
  $\kappa'_k$ is the shift of the couplings of $V^\flat_k$ (\texttt{histVprLO\_e}). Hence \texttt{MuChain} holds along every
  history (\texttt{muChain\_obsHist}; initial state with $\mathcal W_0(B)_a=0$, $V_0(B)_a\in\mathbb C\bar\psi_a$), and the exact
  identity and the conditional bound hold along the flow (\texttt{mu\_flow\_exact\_hist}, \texttt{mu\_flow\_bound\_hist}).
\item \emph{The decay of the tadpole} (\texttt{H02UV/Obs2/Flow/MuTau.lean}). For $k\le j$,
  $|\tau(C_k,C_j)(a)|\le18\ell^4L^{-2(j-k)}$ (\texttt{norm\_tauC\_le}): $\tau=2d\delta_j\sum_zC_k(a,z)$ with
  $|\delta_j|\le\ell^2/(2L^{2(j+1)})$ and $|\sum_zC_k(a,z)|\le9L^{2(k+1)}\ell^2$. Along every history this is the field
  \texttt{htau} of \texttt{MuErr} with $B=18\ell_M^4$, $q=L^{-2}$ (\texttt{histCovO\_htau}). The drift field reduces to a
  per-step bound on $\beta_kg'_k$ (\texttt{drift\_of\_step}, \texttt{muErr\_hist}).
\item \emph{The crude bound with sources} (\texttt{H02UV/Obs2/Est/CrudeG}, \texttt{OSNorm}, \texttt{BlockO1--5}). The crude bound of the
  polymer map holds for the generic map over any submultiplicative norm (\texttt{KHypG.KnextG\_le}). The norm with sources
  $\|F\|=\|F_\varnothing\|+w_a\|F_a\|+w_b\|F_b\|$ ($w_{ab}=0$) is such a norm, also on the doubled algebra with two field
  scales (\texttt{oSN}). The block bounds hold in the full norm: the bulk parts are the bulk map, and the observable
  parts are bounded by the observable block bounds (\texttt{BHF}, \texttt{gdIO\_full}, \texttt{thetaKO\_full},
  \texttt{thetaJO\_full}, \texttt{prefO\_full}, \texttt{BHF.khyp}). Hence (\texttt{Knext2O\_crude\_ts}): if the bulk inputs
  satisfy the hypotheses of the bulk crude bound at radius $r_0$, all components are local, and the observable parts of
  $V(B)$, $\mathcal W(B)$ and $A^{(|\mathcal B_j(X)|-2^d)_+}K(X)$ have norm $\le r_0$, then
  $A^{(|\mathcal B_{j+1}(U)|-2^d)_+}\|K_{j+1}(U)\|_{T_{j+1}(\ell')}\le1$ in the full norm. This is the bound on the polydisc
  for the Cauchy estimates with sources.
\item \emph{The second-order step with sources} (\texttt{H02UV/Obs2/Map/HolW}, \texttt{HolO}, \texttt{OrdW}, \texttt{SecondO},
  \texttt{LinO}, \texttt{LinBO}, \texttt{CauchyW}, \texttt{LinComp}, \texttt{LinBoundO}, \texttt{Step2O}, \texttt{KObs3}). The map
  with sources is holomorphic in $(V,\mathcal W,K)$ in the sense that all linear functionals are entire (\texttt{HolW},
  \texttt{Knext2O\_holW}, \texttt{KplusO\_holW}); $K_{j+1}(tV,t^2\mathcal W,0)=t^2K^{(2)}+O(t^3)$ (\texttt{Knext2O\_expansion});
  $K_{j+1}(0,0,sK)=s\mathcal L(K)+O(s^2)$ with $\mathcal L(K)(U)=\sum_{X\in\cS_j,\bar X=U}\E\theta(1-\Loc_X)K(X)
  +\sum_{X\notin\cS_j,\bar X=U}\E\theta K(X)$ (\texttt{Knext2O\_lin}, \texttt{Llin2O\_eq}). Since the $\sigma_a$ components of $K(X)$
  vanish unless $a\in X^\square$, only $O(1)$ small sets contribute to $\mathcal L(K)_a$, and the pointed contraction gives
  $A^{(|\mathcal B_{j+1}(U)|-2^d)_+}\|\mathcal L(K)(U)^{\rm obs}\|_{T_{j+1}(\ell')}\le C(L^{-1}+A^{-\eta})k$ (\texttt{Llin2O\_obs\_bound};
  it uses the $U(1)$ invariance \texttt{IsSymS} of $K(X)$). With the crude bound on the polydisc and the Cauchy estimates for the
  weighted pairings of $F_\varnothing,F_a,F_b$ (\texttt{cauchy\_assemble}):
  \[A^{(|\mathcal B_{j+1}(U)|-2^d)_+}\|K_{j+1}(U)-K^{(2)}(U)\|_{T_{j+1}(\ell')}\le C_0(L^{d-3}+A^{-\eta})k+C_1(\bar g^3+\bar gk+k^2)\]
  in the full norm with sources ($w_{ab}=0$; \texttt{step2O\_ts}). In the $\lambda$-norm this is the bound assumed by
  \texttt{kobs\_step}, which therefore holds for $K_{j+1}-K^{(2)}$ (\texttt{kobs2\_step}):
  $A^{(\cdots)_+}\|(K_{j+1}-K^{(2)})^{\rm obs}\|\le(\theta_0+O(\bar g))k+B\bar g^2m$, $\theta_0=C_0(L^{d-3}+A^{-\eta})$.
\item \emph{The activity after the change of variables} (\texttt{H02UV/Obs2/Est/SplitO}, \texttt{ReblockO}, \texttt{LinCompL},
  \texttt{LinBoundOL}, \texttt{BlockO6}, \texttt{Step2OL}, \texttt{Step3O}, \texttt{Step3OL}, \texttt{StepKO}, \texttt{KObs4W}).
  Under the strong locality of the inputs, $\sum_{U\supseteq B_+}J(B_+,U)=0$ with sources (\texttt{JhO\_sum\_U}), so
  $K_{j+1}(U)=K^\sharp_{j+1}(U)+\sum_{\mathcal R}R_g(U)$ (\texttt{KoutO\_split}); the correction is bounded in any
  submultiplicative norm (\texttt{KoutO\_sub\_le}). With the output regulator $\lambda A$ in the second-order step
  (\texttt{step2O\_tsL}, \texttt{step3OL}) and the full-norm bounds $\|\tilde I_{j+1}(B_+)\|\le2$, $\|\bar J\|\le C r^2$
  ($r=\max(\bar g,k)$; \texttt{ItO\_full}, \texttt{JhO\_full}):
  \[A^{(|\mathcal B_{j+1}(U)|-2^d)_+}\|K_{j+1}(U)\|_{T_{j+1}(\ell')}\le C_0(L^{d-3}+A^{-\eta})k+C_1(\bar g^3+\bar gk+k^2)\]
  in the full norm with sources (\texttt{stepKO\_ts}), and in the $\lambda$-norm
  $A^{(\cdots)_+}\|K_{j+1}^{\rm obs}\|\le(\theta_0+O(\bar g))k+B\bar g^2m$ (\texttt{kobsK\_stepW} with $\mu=1$): the $k$-half of the
  linear recursion of \texttt{lin\_iter}; also with input regulator $A$ and output regulator $\mu A$ (\texttt{stepKO\_tsW}, \texttt{kobsK\_stepW}), the form used along the flow.
\item \emph{$U(1)$ invariance with sources} (\texttt{H02UV/Obs2/Map/SymO}, \texttt{H02UV/PtCorr2/Obs/ObsSym}). $\Loc_{X,Y}F$ is $U(1)$ invariant for
  every $F$ (\texttt{isSymS\_LocO2}), and all other operations of the map preserve invariance; hence invariance of
  $V,\mathcal W,K$ is preserved by a step (\texttt{OSymS.step}), holds for the initial state, and holds for the full inputs of
  every step along the history (\texttt{obsSymInv\_hist}, \texttt{osymS\_hist}).
\item \emph{The activity term of the $\mu$ flow} (\texttt{H02UV/Obs2/Flow/MuK}, \texttt{H02UV/PtCorr2/Obs/KObsEK}):
  $|\sum_{X\in\cS_j}m_1(K_j(X)_a)|\le c\,(k/w_a)\ell^{-1}$ from the norm of $K^{\rm obs}_j$ (\texttt{sum\_mass1\_le},
  \texttt{histKO\_hK}).
\item \emph{The $K^{\rm obs}$ recursion along the flow} (\texttt{H02UV/Obs2/Est/KFlatO}, \texttt{KObsInd}, \texttt{H02UV/Step2/Flow/FlowRGIn},
  \texttt{H02UV/PtCorr2/Obs/KObsStep}, \texttt{KObsHist}, \texttt{KObsEK}). At the raw level, with the clamp conditions and the radius
  $r$ of the bulk step, observable sizes $\|V^{\rm obs}_j\|\le m$, $\|\mathcal W^{\rm obs}_j\|\le rm$ and
  $A^{(|\mathcal B_j(X)|-4)_+}\|K^{\rm obs}_j(X)\|_{T_j(s_j)}\le k$ with $8k\le c_kr^2m$,
  \[A^{(|\mathcal B_{j+1}(U)|-4)_+}\|K^{\rm obs}_{j+1}(U)\|_{T_{j+1}(s_{j+1})}\le(\theta_0+C_1(2+2c_b+c_k)r)\,8k+B\,r^2m,\]
  $\theta_0=C_0(L^{-1}+(A/2)^{-\eta})$, $B=\theta_0c_b+C_1(8+2c_b+c_b^2)$ (\texttt{kobs\_raw\_step}; the factor $8\ge e^{4\delta}$ and
  the regulator $A/2$ come from $K^\flat=\prod e^{\Delta(B)}K$, \texttt{oTo\_KflatO\_le}). Along a regular history, with
  $r_j=c_Rg_j$, $c_b=8c_x$, $c_k=64B$ and $\bar\theta=\theta_0+C_1(2+2c_b+c_k)c_R\varepsilon$:
  if $\theta_0\le1/72$, $C_1(2+2c_b+c_k)c_R\varepsilon\le1/72$, $\varepsilon$ is small, $m_j\le2m_{j+1}$ and $K^{\rm obs}_0=0$, then
  \[A^{(|\mathcal B_j(X)|-4)_+}\|K^{\rm obs}_j(X)\|_{T_j(s_j)}\le k_j:=\sum_{i<j}(8\bar\theta)^{j-1-i}B\,r_i^2m_i,\qquad 8k_j\le c_kr_j^2m_j\]
  (\texttt{kobs\_hist}, \texttt{kseq\_eq\_sum}), and the activity term of the $\mu$ flow is bounded by
  $c\,(8k_j/w_a)s_j^{-1}$ (\texttt{kobs\_EK}). The observable sizes $m_j$ (with $\|\mathcal W^{\rm obs}_j\|\le r_jm_j$) are inputs.
\item \emph{The contraction of $\mathcal W^{\rm obs}$} (\texttt{H02UV/Obs2/Est/WObs1--3}, \texttt{H02UV/PtCorr2/Obs/WObsStep}). In
  $\mathcal W_{j+1}(B)=S(B)-\Loc_{B_+,B}S(B_+)$ the localisation acts after the expectation. For a component $G$ of
  $\mathcal W(B)$ with $p\in B^\square$ write $G=A+P$, $P=\operatorname{loc}^pG$: $A$ gains $L^{-1}$ at scale $j+1$ with field
  scale $2\ell$, and $\E\theta$ maps $T_{j+1}(2\ell)$ to $T_{j+1}(\ell')$; $P$ has degree $\le2$, so
  $\|\E\theta P\|_{T_j(\ell)}\le16\|P\|_{T_j(\ell)}$, and the pointed contraction applies to $\E\theta P$. Hence
  $\|\E\theta G-\operatorname{loc}^p\E\theta G\|_{T_{j+1}(\ell')}\le KL^{-1}\|G\|_{T_j(\ell)}$ (\texttt{ecc\_odd\_contract},
  \texttt{ecc\_even\_contract}, \texttt{EthO\_subLoc\_contract}). Since $\operatorname{loc}^a_{X,X}=\operatorname{loc}^a_{\{a\},\{a\}}$
  for every $X$, $\sum_{B\subseteq B_+}\mathcal W_{j+1}(B)=(1-\Loc_{B_+,B_+})S(B_+)$ (\texttt{sum\_WpO\_eq}), and only $O(1)$
  blocks have $a\in B^\square$ (\texttt{card\_blocks\_near}), the $\mathcal W$ part of $S(B_+)$ contracts
  (\texttt{sum\_EthW\_subLoc}); the one-loop part is $O(rm_V)$; the products $\prod(1+\mathcal W_{j+1}(B))-1-\sum$ are linear
  in the observable sizes (\texttt{oTo\_mul\_le0}, \texttt{oTo\_prod\_sub\_le}). Result (\texttt{Wnext2O\_obs\_le},
  \texttt{wobs\_raw\_step}): $\|\mathcal W^{\rm obs}_{j+1}\|_{T_{j+1}(s')}\le(KL^{-1}+C_1(L)r)\,m_{\mathcal W}+C_1(L)\,r\,m_V$ for
  $r\le\varepsilon_1(L)$, with $K$ independent of $L$.
\item \emph{The $m$-half} (\texttt{H02UV/Obs2/Flow/VObs1--2}, \texttt{H02UV/PtCorr2/Obs/VObsStep}; Case~(1), the $\sigma_a$ part, weights $(1,0,0)$).
  With the pointed $\operatorname{loc}^a$ the coupling is spread over the blocks, and the quantity with leading coefficient $1$
  is the $\ell^1$ sum $\beta^{\rm abs}=\sum_B|m_1(V(B)_a)|$: $\sum_{B\subseteq X}|m(\operatorname{loc}^a_{X,B}F)|=|m(F)|$
  (\texttt{sum\_norm\_locAp\_blocks}), so $\beta^{\rm abs}_{j+1}\le\beta^{\rm abs}_j+\sum_{X\in\cS_j}|m(K(X)_a)|+\sum_B|m(S(B)_a)|$
  (\texttt{abs\_mass\_step}), and with pair bounds with block-dependent sizes and the symmetry of the near relation,
  $\beta^{\rm abs}_{j+1}\le\beta^{\rm abs}_j+cC_m\ell^{-1}k+C_m\ell'^{-1}(2c_{\rm near}Q_1^2r\ell\beta^{\rm abs}_j+c'Q_1m_{\mathcal W})$
  (\texttt{beta\_step}, \texttt{beta\_raw\_step}). Since $\|V(B)_a\|_{T_j(\ell)}=|m_1(V(B)_a)|\ell$, $\ell_j\beta^{\rm abs}_j$ bounds
  all blocks.
\item \emph{The coupled recursion} (\texttt{H02UV/PtCorr2/Obs/ObsHist}, \texttt{ObsHist2}; Case~(1), $\sigma_a$). With
  $q_j=\max(1,s_{j+1}/s_j)$, $c_W=2C_1+1$ and $M_{j+1}=M_jq_j(1+C_Mr_j)$, a joint induction along a regular history gives
  $s_j\beta^{\rm abs}_j\le M_j$, $\|\mathcal W^{\rm obs}_j\|\le c_Wr_jM_j$ and the $K^{\rm obs}$ bound
  $k_j=\sum_{i<j}(8\bar\theta)^{j-1-i}Br_i^2c_WM_i$ (\texttt{obs\_hist}), under the smallness conditions of \texttt{kobs\_hist},
  $K/L\le1/6$ and $C_1c_R\varepsilon\le1/6$. This discharges the observable sizes of \texttt{kobs\_hist}
  ($m_j=c_WM_j$, nondecreasing).
\item \emph{Along the bulk flow} (\texttt{H02UV/Step2/Flow/FlowComplete2X}, \texttt{H02UV/PtCorr2/Obs/ObsFlow}). The constants of the bulk flow
  theorem can be chosen such that the coupled recursion holds unconditionally along every admissible trajectory
  (\texttt{thm\_obs\_flow}): the bulk construction accepts requirements $A\ge A_0(L)$ and $\varepsilon\le\varepsilon_0(L,A,c_R,c_x)$
  (\texttt{thm\_flow\_complete2X}, with the same conclusions as \texttt{thm\_flow\_complete2}), and $\theta_0\le1/72$,
  $K/L\le1/6$ hold for $L\ge L_o$ and these $A$. The field-scale ratio satisfies $q_j\le1+|\zeta_j|$
  (\texttt{histS\_ratio\_le}), so $M_j\le M_0\exp\sum_{i<j}(|\zeta_i|+C_Mr_i)$ (\texttt{Mseq\_le\_exp}). Since
  $|\zeta_i|,r_i=O(\check g_i)$ and $\sum_i\check g_i$ grows like $\log j$, $M_j$ grows polynomially in $j$
  (with an $L$-dependent exponent); this upper bound is not sharp enough for a relative error of the $\mu$ flow.
\item \emph{The shares} (\texttt{H02UV/Obs2/Flow/BlockMass}, \texttt{H02UV/Obs2/Flow/BlockFlow}; Case~(1)). Write $V_j(B)_a=c_j(B)\bar\psi_a$,
  $\beta_j=\sum_Bc_j(B)$, and $\beta_j(X)=\sum_{B\subseteq X}c_j(B)$. The map \texttt{Obs2.Map.Map} uses $S^\sharp$
  (\texttt{SbO}, \texttt{SFO}, \texttt{SCorO}, \texttt{GO}). Since $\operatorname{Loc}_{B_+,B}S(B_+)$ moves mass only inside
  $B_+$, $c_{j+1}(B_+)=\sum_{B\subseteq B_+}[c_j(B)-m_1(S^\sharp(B)_a)-m_1((\mathbb E\theta Q(B))_a)]$ (\texttt{share\_step}),
  and $m_1(S^\sharp(B)_a)=c_j(B)\bar\nu+m_1((\mathbb E\theta\mathcal W(B))_a)$, $\bar\nu=m_1F_C(\bar\psi_a,V'(\Lambda)_\emptyset)$
  (\texttt{mass\_SbO\_a\_one}, \texttt{share\_step\_one}). The tadpole recursion of \texttt{mu\_flow\_exact} holds in every
  region $X$ that is a union of blocks of the coarsest scale (\texttt{wTadX\_step}, \texttt{wTadX\_unroll}), so for every
  $(j+1)$-block $c_{j+1}(B_+)=\beta_j(B_+)(1-\mathfrak n_j)-{\rm err}_j(B_+)$ with the same $\mathfrak n_j$ as the total flow;
  ${\rm err}_j(B_+)$ is the activity term $\sum_{B\subseteq B_+}m_1((\mathbb E\theta Q(B))_a)$, the drift of $\beta_k(B_+)g'_k$, and
  the product parts of $\mathcal W$ in $B_+$ (\texttt{share\_flow\_exact}). With $D_j=\sum_B|c_j(B)-\beta_j1_{a\in B}|$ and
  $E_j=\sum_{B_+}|{\rm err}_j(B_+)|$: $D_{j+1}\le|1-\mathfrak n_j|D_j+2E_j$ (\texttt{shareDev\_succ\_le}), and if $D_0=0$,
  $E_j\le\delta_j|1-\mathfrak n_j||\beta_j|$, $\delta_j\le\frac12$, $\sum_j\delta_j\le\frac16$, then
  $D_j\le6(\sum_{k<j}\delta_k)|\beta_j|$ and $\sum_B|c_j(B)|\le(1+6\sum_{k<j}\delta_k)|\beta_j|$ (\texttt{phase\_align},
  \texttt{sum\_norm\_share\_le}): the shares stay at the block of $a$. $E_j$ is bounded by
  $\sum_{X\in\mathcal S_j}|m_1(K_j(X)_a)|$, the drift and the product parts (\texttt{shareErrL1\_le}); the drift of the regions
  telescopes into per-step drifts $\le|g'_{i+1}-g'_i|A_{i+1}+|g'_i|(|\mathfrak n_i|A_i+E_i)$, $A_i=\sum_B|c_i(B)|$
  (\texttt{drift\_step\_l1}, \texttt{drift\_l1\_le}).
\item \emph{The relative coupled recursion} (\texttt{H02UV/Obs2/Flow/RelRec}, \texttt{H02UV/Obs2/Flow/BlockFlow2}, \texttt{H02UV/PtCorr2/Obs/ObsHistRel};
  Case~(1)). The bounds of $\mathcal W^{\rm obs}_j$, $K^{\rm obs}_j$ are run with sizes proportional to the coupling,
  $M_j=2s_j|\beta_j|$, in one induction with the concentration of the shares: if $r_j$ is nondecreasing,
  $s_j\le\frac65s_{j+1}$, $\frac67\le|1-\mathfrak n_j|\le2$, $|\tau(C_k,C_j)|\le B_\tau q^{j-k}$ with $q\le\frac14$,
  $|g'_{i+1}-g'_i|\le c_gr_i^2$, $|g'_i|,|\mathfrak n_i|\le c_gr_i$, the product parts are bounded in $\ell^1$ by
  $c_P\sum_{i\le j}q^{j-i}r_i^2m_i/s_i$ for observable sizes $m_i$, and $c_\delta\sum_{k<n}r_k^2\le\frac18$ with $c_\delta$
  large (a linear condition in \texttt{RelIn}), then for $j\le n$: $|\beta_j|>0$, $D_j\le6c_\delta(\sum_{k<j}r_k^2)|\beta_j|$,
  $\|\mathcal W^{\rm obs}_j\|\le c_Wr_jM_j$, the weighted $K^{\rm obs}_j$ bound is $\le16c_WBr_j^2s_j|\beta_j|$, and
  $|\beta_{j+1}-\beta_j(1-\mathfrak n_j)|\le c_\delta r_j^2|\beta_j(1-\mathfrak n_j)|$ (\texttt{RelRec.close},
  \texttt{MuChain.rel\_flow}); hence $|\beta_j-\beta_0\prod_{k<j}(1-\mathfrak n_k)|\le(e^{c_\delta\sum_{k<j}r_k^2}-1)
  |\beta_0\prod_{k<j}(1-\mathfrak n_k)|$ (\texttt{rel\_prod\_le}). Along the history (\texttt{obs\_hist\_rel}) the
  $\mathcal W$ and $K$ steps, the activity term in $\ell^1$ (\texttt{kobs\_EK1}), the tadpole decay, the field-scale ratio
  ($\rho\le\frac18$) and $\sum_kr_k^2\le c_R^2\varepsilon/\beta'$ (from $g_{k+1}\ge g_k+\beta'g_k^2$) are discharged; the
  remaining inputs are the bulk bounds on $g'_i$, $\mathfrak n_i$, the $\ell^1$ product bound (\texttt{ProdBd}), and
  $D_0=0$, $\beta_0\neq0$ for the initial state.
\item \emph{Uniform constants and the master construction} (\texttt{H02UV/Obs2/Flow/BlockFlow3}, \texttt{H02UV/PtCorr2/Obs/ObsHistRelW},
  \texttt{H02UV/Spin2/Multiscale/ObsRel}; Case~(1)). Since $C_k$ lives at the scale $|Z_{k+1}|^{-1/2}\ell$ and $|Z_j|\to0$, the sizes
  $|\tau(C_k,C_j)|$ and $|g'_i|$ are not separately uniform (and \texttt{ellM} is a sum over scales); the recursion is
  run with weights $w_i=|Z_{i+1}|^2$ (\texttt{rel\_flow\_w}): $|\tau(C_k,C_j)|\,|Z_{i+1}|^2\le\frac{9\ell^4}{128}(2L^{-2})^{j-k}$
  (\texttt{hist\_tau\_w}), $|g'_i|/|Z_{i+1}|^2\le\frac87\check g_i$, $|g'_{i+1}-g'_i|/|Z_{i+1}|^2=O(\check g_i^2)$,
  $|\mathfrak n_i|=O(\check g_i)$ (\texttt{hist\_G\_w}, \texttt{hist\_dg\_w}, \texttt{hist\_frakN\_le}, with the explicit couplings
  \texttt{kapH}); the band increment comes from the regularity of the history (\texttt{obs\_hist\_rel\_w}). Under the
  constants of the master construction, for point initial shares $c_0(B)=\beta_01_{a\in B}$, $\beta_0\ne0$, and conditional only on
  \texttt{ProdBd}: $\beta_j=\beta_0\prod_{k<j}(1-\mathfrak n_k)(1+O(\sum_{k<j}\check g_k^2))$ uniformly, with phase alignment
  (\texttt{thm\_obs\_rel}). \emph{\texttt{ProdBd} is proved} (\texttt{H02UV/Obs2/Loc/PauliMass}, \texttt{GramDiff}, \texttt{LocGram},
  \texttt{PauliLip}, \texttt{ProdObs}, \texttt{H02UV/PtCorr2/Obs/ProdStep}, \texttt{ProdBdHist}): with $J$ the all-ones matrix,
  $\bar D=\sum_x\partial_{\bar\psi_x}$, $D=\sum_y\partial_{\psi_y}$ one has $\mathcal L_J=D\bar D$, $\bar D^2=D^2=0$,
  hence $e^{\mathcal L_{\Gamma+cJ}}=e^{\mathcal L_\Gamma}(1+cD\bar D)$ and $m(e^{\mathcal L_{\Gamma+cJ}}F)=m(e^{\mathcal L_\Gamma}F)$
  for $m=m_1=\varepsilon\circ\bar D$ and also for $m=m_2=\varepsilon\circ D\bar D$ (\texttt{mass1\_expLC\_add\_J},
  \texttt{mass2\_expLC\_add\_J}); $e^{\mathcal L_\Gamma}F$ only depends on $\Gamma|_{Y\times Y}$ for $F\in\mathcal N(Y)$
  (\texttt{expLC\_congr\_NX}); on a box of radius $R\le h$ around $x_0$, $C-C(x_0,x_0)$ agrees with a covariance with Gram
  constant $\gamma'^2=16dK^2R\gamma^2/h$ at all scales $\le R$ (smooth cut-off, \texttt{gram\_loc}), and
  $\|e^{\mathcal L_R}F-F\|_{T(\ell)}\le(\gamma/\ell)^2\|F\|_{T(2\ell)}$ (\texttt{tNorm\_expLC\_sub\_le}); so
  $|m(e^{\mathcal L_{S+C}}F)-m(e^{\mathcal L_S}F)|\le M_\sigma(\gamma'_C/2\sigma)^2\|F\|_{T(4\sigma)}$ (\texttt{pauli\_lip\_loc}).
  For the product parts ($\|\cdot\|\le C_P(L)r_k(r_km_V+m_W)$, \texttt{prod\_raw\_step}) with $R=10L^{k+1}$ this gives the
  decay $(8/(7L))^{j-k}$ (\texttt{delta\_hist\_le}) and \texttt{ProdBd} with ratio $q=2/L$ (\texttt{prodBd\_hist});
  hence the relative $\sigma_a$ flow of \texttt{Spin2.thm\_obs\_rel} holds under the master constants with \texttt{ProdBd} discharged, \emph{unconditionally}.
  The Case~(2) increment bound \texttt{delta\_hist\_le2} ($m_2$) is proved in the same (generic) form.
\item \emph{Case (2), $\mu^{(2)}$} (\texttt{H02UV/Obs2/Flow/MuChain2.lean}, \texttt{H02UV/Obs2/Flow/MuHist2.lean}): the $\sigma_a$ component of
  every block is $\alpha_k(B)+c_k(B)\bar\psi_a\psi_a$ along the history; with $\tau_2(C,\Gamma)=m_2(e^{\mathcal L_\Gamma}
  F_C(\bar\psi_a\psi_a,V'(\Lambda)_\emptyset))-m_2(F_C(\bar\psi_a\psi_a,V'(\Lambda)_\emptyset))$ (kept as a functional) and
  $\mathfrak n^{(2)}_j=m_2F_{C_j}(\bar\psi_a\psi_a,V'_j(\Lambda)_\emptyset)+\sum_{k<j}[\tau_2(C_k,C_j+S_{k,j})-\tau_2(C_k,S_{k,j})]$,
  the exact per-block and total flows $c_{j+1}(B_+)=\beta_j(B_+)(1-\mathfrak n^{(2)}_j)-\mathrm{err}_j(B_+)$
  (\texttt{MuChain2.share2\_flow\_exact}, \texttt{betaTot2\_succ\_eq}, \texttt{muChain2\_obsHist}), $\beta_0=-\mu_a$
  (\texttt{betaTot2\_obsInit}). The tadpole explicitly (\texttt{H02UV/Obs2/Flow/Tau2.lean}, \texttt{Tau2Hist.lean}):
  $e^{\mathcal L_\Gamma}F_C(A,B)=F_{C+\Gamma}(A_\Gamma,B_\Gamma)-F_\Gamma(A_\Gamma,B_\Gamma)$ (\texttt{expLC\_FC\_eq}), hence
  for $\Gamma(x,x)=c_0$, $\Gamma(x\pm e,x)=c_0+\delta$: $\tau_2(C,\Gamma)=g\,(2d\delta\sum_x(C(a,x)+C(x,a))-\mathfrak l_{C,\Gamma}(a))$
  with $\mathfrak l_{C,\Gamma}=\mathfrak l_{C+\Gamma}-\mathfrak l_C-\mathfrak l_\Gamma$ the polarised bubble
  (\texttt{tau2\_Vcoup}); along the history ($\mathtt{ValidAt}$) $\mathfrak n^{(2)}_j=\nu'_jS_{j+1}-\nu^\flat_jS_j
  -g_j(\mathfrak l_{w_{j+1}}-\mathfrak l_{w_j})-4\delta_j\sum_{k<j}(g_j-g_k)S^{(1)}_{C_k}+\sum_{k<j}(g_j-g_k)\mathfrak l_{C_k,C_j}$
  exactly (\texttt{frakN2\_hist}). The bulk estimates of \texttt{Spin2.Case2Bulk} are proved under the master
  constants (\texttt{H02UV/Spin2/Case2/Case2Thm.lean}, \texttt{LoopConcl}). The observable remainders
  (\texttt{Case2Rel}) are bounded for every $\gamma_a$ (exact decoupling of the scale-$0$ constant, \texttt{RelConcl2G};
  for $\gamma_a=0$: \texttt{Spin2.thm\_multiscale\_case2D}: $\ell^1$ share
  machinery \texttt{H02UV/Obs2/Flow/BlockFlowE.lean}, \texttt{PtCorr2.obs\_hist\_rel2}, \texttt{PtCorr2.prodBd2\_hist},
  $|\mathrm{tadInc2}_{kj}|\le c(2/L)^{j-k}\check g_k$).
\item \emph{Not done:}
  the $\sigma_b$ part of the $m$-half, the bulk bounds on $g'_i,\mathfrak n_i$ of the previous item,
  the bounds of Case~(2), the assembly along the flow, the $ab$ parts, the second-order coefficients,
  the flows of $\gamma,\eta,q,r,\chi$, and the hypotheses of the conditional part of
  Proposition~\ref{prop:gamma} (see its Lean status).
\end{itemize}

\subsection{Spin renormalisation and the anomalous dimension}
\label{sec:spin}

\begin{proposition}[$\gamma$]\label{prop:gamma}
Along the flow of Theorem~\ref{thm:flow}, with $\beta_j=1/\check g_j$, there are two
expressions for the spin renormalisation, and they agree:
\[
  Z_j=\Bigl(\frac{\beta_j}{\beta_\epsilon}\Bigr)^{1/3}(1+O(\beta_j^{-1})),\qquad
  \frac{\lambda^{(1)}_j}{\lambda^{(1)}_0}=1+O(\beta_j^{-1}),\qquad
  \frac{\lambda^{(2)}_j}{\lambda^{(2)}_0}=\Bigl(\frac{\beta_\epsilon}{\beta_j}\Bigr)^{1/3}(1+O(\beta_j^{-1})).
\]
\begin{enumerate}[label=(\roman*)]
\item \emph{Case 1 (from $Z_j$ and $\beta_j$):} the spin renormalisation is
  $\vartheta^{(1)}_j=(\beta_j/(\beta_\epsilon Z_j))^{1/2}\,\lambda^{(1)}_j/\lambda^{(1)}_0$.
\item \emph{Case 2 (from the $\bar\psi\psi$ vertex):} it is the magnetisation $M_j$ of
  \eqref{e:PM}.
\end{enumerate}
Both equal $\vartheta_j=(\beta_j/\beta_\epsilon)^{1/3}(1+O(\beta_j^{-1}))$. So the spin field has
anomalous dimension $\gamma_j=\frac1{2\pi\beta_j}$ per unit of $\log$ length:
$\vartheta_{j+1}/\vartheta_j=1-\gamma_j\log L+O(\beta_j^{-2})$.
\end{proposition}

\begin{proof}
\begin{itemize}
\item \emph{$Z_j$.} By \eqref{e:zflow} and \eqref{e:univ},
  $\log Z_{j+1}-\log Z_j=-\frac{\log L}{2\pi\beta_j}+O(\beta_j^{-2})$. By
  Corollary~\ref{cor:beta1}, $\log L=-(\beta_{j+1}-\beta_j)/\beta_1+O(\beta_j^{-1})$. Hence
  $d\log Z=\frac{d\beta}{2\pi\beta_1\beta}=\frac13\,d\log\beta$.
\item \emph{$\lambda^{(1)}$.} $\lambda^{(1)}_{j+1}/\lambda^{(1)}_j=1-\mathfrak n_j+O(\bar g^2)$.
  The factors $Z^{1/2}$ cancel against step~(B). The $\mathfrak n_j$ telescope, and $\nu w'^{(1)}=O(\check g)$.
\item \emph{$\lambda^{(2)}$.} Similarly,
  $\lambda^{(2)}_{j+1}/\lambda^{(2)}_j=1+\frac{\log L}{2\pi\beta_j}-2\mathfrak n_j+O(\beta_j^{-2})$
  by \eqref{e:univ}. Hence $d\log\lambda^{(2)}=-\frac13\,d\log\beta$.
\item \emph{(i).} $\vartheta^{(1)}_j=(\beta_j/\beta_\epsilon)^{1/2-1/6}=(\beta_j/\beta_\epsilon)^{1/3}$.
\item \emph{(ii).} By \eqref{e:PM}, $M_j=1-b_0\sum_{k<j}\hat\lambda_k\Gamma^{\rm int}_{k+1}(0,0)$,
  with $\Gamma^{\rm int}_{k+1}(0,0)=C'_{k+1}(0,0)/Z_{k+1}=\frac{\log L}{2\pi}Z_{k+1}^{-1}(1+O(\cdot))$
  and $b_0=1/\beta_\epsilon$. So
  \[
    M_j=1-\frac1{\beta_\epsilon}\sum_{k<j}\Bigl(\frac{\beta_\epsilon}{\beta_k}\Bigr)^{2/3}\frac{\log L}{2\pi}
    =1-\frac1{2\pi\beta_1}\beta_\epsilon^{-1/3}\int_{\beta_j}^{\beta_\epsilon}\beta^{-2/3}d\beta
    =\Bigl(\frac{\beta_j}{\beta_\epsilon}\Bigr)^{1/3},
  \]
  up to relative errors $O(\beta_j^{-1})$, since $\frac{3}{2\pi\beta_1}=1$.
\end{itemize}
\end{proof}

\begin{remark}
The agreement of (i) and (ii) is a non-trivial check. Route (i) uses $\beta_1=\frac3{2\pi}$ and
the field strength. Route (ii) uses the independently computed $\bar\psi\psi$ vertex
$\frac{\log L}{2\pi}$ (Lemma~\ref{lem:second}(ii), \texttt{scripts/pt-checks/obs.py}). They agree as
required by the $\mathfrak{osp}(1|2)$ symmetry, which relates $\xi,\eta$ and $z=1-\xi\eta$.
\end{remark}

\paragraph{Lean status of Proposition~\ref{prop:gamma} (\texttt{H02UV/Step2/Flow/ZetaSharp*.lean}, \texttt{H02UV/Obs2/Flow/Gamma1--2.lean}).}
\begin{itemize}
\item \emph{\eqref{e:zflow} with the explicit coefficient.} \texttt{zeta\_sharp\_hist}: along regular histories
  $|\zeta_{j+1}+\check g_{r,j}\,Z_{j+1}C_j(e_1,0)|\le 2c_E\check g_j^2$ ($\mathfrak z^{(2)}_j$, $r^\zeta_j$ are in the
  $O(\check g^2)$); $|\check g_r-\check g|\le\frac{32}7|\mu|\ell^2\check g$ (\texttt{stGr\_sub\_le}); and
  $|Z_{k+1}C_k(e_1,0)-\frac{\log L}{2\pi}|\le(2\pi)^{-2}(E_{C1}L^{-2k}+E_Bm^2(L^{2k}+L^{2(k+1)}))+KM_z(k+1)$ for
  $k\ge1$ (\texttt{X\_value\_hist}, from \texttt{Cprime\_values}). Hence (\texttt{zs\_step})
  $|\log|Z_{k+2}|-\log|Z_{k+1}|+\frac{\log L}{2\pi}\check g_k|\le A_Z(\check g_k^2+\check g_ka_k)$, $a_0=1$,
  $a_k=L^{-2k}+m^2L^{2(k+1)}$ (in Lean $Z_{k+1}=Z_k(1+\zeta_k)$ is the normalisation of step $k$).
\item \emph{$Z_j$, unconditional.} \texttt{Step2.thm\_flow\_final\_Z}: the conclusions of
  \texttt{thm\_flow\_final} (with its own constants, same construction) and, along every admissible
  trajectory, $Z_{n+1}\ne0$ and $\bigl|\log|Z_{n+1}|-\frac13\log\frac{\beta_n}{\beta_\epsilon}\bigr|\le K_Z/\beta_n$ for
  $n\le N_h$, with $K_Z$ depending on $L$ only (the $j=0$ step and the summable $L^{-2j}$, $m^2L^{2(j+1)}$
  errors contribute $O(\check g_0)\le O(\check g_n)$). Only $|Z|$ is controlled. The summation is the pure
  lemma \texttt{log\_sum\_bound} ($\sum_k\check g_k^2\le(\check g_n-\check g_0)/\beta'$, $\log(\check g_{k+1}/\check g_k)
  =\beta\check g_k+O(\check g_k^2+e_k\check g_k)$ from the band) and its flow form \texttt{LogSum}
  (\texttt{logSum\_of}).
\item \emph{$\lambda^{(1)},\lambda^{(2)}$, $\vartheta^{(1)}$, conditional.} \texttt{Obs2.prop\_gamma}: for any vertex sequence
  $\lambda_{k+1}=\lambda_k(1+x_k)$ with $|x_k|\le\min(\frac14,B'\check g_k)$, $q_0=0$, $|q_k|\le c_q\check g_k$ and
  $|\mathrm{Re}\,x_k-s\frac{\log L}{2\pi}\check g_k+(q_{k+1}-q_k)|\le B(\check g_k^2+\check g_ka_k)$ (the flow equations
  above with $s=0$, $q=n/4$ resp.\ $s=1$, $q=n/2$, $n_k=4\nu_kw^{(1)}_k$, and an $O(\bar g^2)$ remainder):
  $|\log\frac{|\lambda_n|}{|\lambda_0|}-\frac s3\log\frac{\beta_\epsilon}{\beta_n}|\le K_\lambda/\beta_n$; and in Case~1
  $\vartheta^{(1)}_n=(\beta_n/(\beta_\epsilon|Z_{n+1}|))^{1/2}|\lambda_n/\lambda_0|$ satisfies
  $|\log\vartheta^{(1)}_n-\frac13\log\frac{\beta_n}{\beta_\epsilon}|\le(K_Z/2+K_\lambda)/\beta_n$.
\item \emph{(ii), conditional} (\texttt{H02UV/Spin2/Multiscale/Sum,Flow,Thm.lean}). In Lean $\Gamma^{\rm int}_{k+1}(0,0)$ is
  \texttt{Spin2.GamInt} (the diagonal of the dressed $\Gamma^{\rm int}$ of step $k$ along the history), and
  \texttt{gamInt\_value\_hist} gives $|Z_{k+1}\Gamma^{\rm int}_{k+1}(0,0)-\frac{\log L}{2\pi}|\le(2\pi)^{-2}(E_CL^{-2k}+E_Bm^2(L^{2k}+L^{2(k+1)}))+KM_z(k+1)$
  for $k\ge1$. \texttt{Spin2.prop\_gamma\_M}: with the constants of \texttt{thm\_flow\_final\_spin} (= \texttt{thm\_flow\_final\_Z}
  plus the summation principles below), along every admissible trajectory, if
  $\hat\lambda_k=(\beta_\epsilon/\beta_k)^{1/3}(1+O(\check g_k+a_k))$ (complex),
  $Z_{k+1}\Gamma^{\rm int}_{k+1}(0,0)=\frac{\log L}{2\pi}+O(\check g_k+a_k)$ and $Z_{k+1}/|Z_{k+1}|=1+O(\check g_k)$, then
  $|M_j-(\beta_j/\beta_\epsilon)^{1/3}|\le K_M(\beta_j/\beta_\epsilon)^{1/3}\check g_j$ for $j\le N_h$; $|Z_{k+1}|$ is used
  unconditionally. The summation is the pure lemma \texttt{spin\_sum} (an algebraic Riemann sum in $w=\check g^{1/3}$,
  using $\frac{\beta}{3}=\frac{\log L}{2\pi}$, i.e.\ $\frac3{2\pi\beta_1}=1$), along the flow \texttt{SpinSum}.
\item \emph{Running coupling} (\texttt{Spin2.BetaAsym}, part of \texttt{thm\_flow\_final\_spin}): along every admissible
  trajectory $|\beta_n-\beta_\epsilon+n\beta_1\log L|\le C+C_{gu}(\frac1{\beta_1\log L}\log\frac{\beta_\epsilon}{\beta_n}+K'\check g_n)$,
  the summed form of $\beta_{j+1}=\beta_j-\beta_1\log L+O(\beta_j^{-1})+e_j$ of Theorem~\ref{thm:beta}.
\item \emph{(ii) with only the vertex hypothesis} (\texttt{H02UV/Spin2/Multiscale/Gamma.lean}, \texttt{H02UV/Step2/Flow/ZetaSharpHypX.lean}).
  \texttt{flow\_finalZ\_hyp0\_zetaX} exports $M_z(k+1)\le2c_z\check g_k$, $|C'_1(0,0)|\le K_z$ and the value of
  $C'_{k+1}(0,0)$; hence (\texttt{gam\_hist\_bound}) $|Z_{k+1}\Gamma^{\rm int}_{k+1}(0,0)-\frac{\log L}{2\pi}|\le K_G(\check g_k+a_k)$
  for $k<N_h$. The raw flow is real (\texttt{PtCorr2.traj\_real}), so $Z_j=|Z_j|>0$ (\texttt{Z\_eq\_norm\_hist}): the phase
  of $Z$ is trivial. \texttt{Spin2.prop\_gamma\_M2}: along every admissible trajectory, if
  $\hat\lambda_k=(\beta_\epsilon/\beta_k)^{1/3}(1+O(\check g_k+a_k))$ (complex) for $k<N_h$, then
  $|M_j-(\beta_j/\beta_\epsilon)^{1/3}|\le K_M(\beta_j/\beta_\epsilon)^{1/3}\check g_j$, $j\le N_h$. By \texttt{lamHyp\_of\_log} the
  hypothesis follows from the ($\lambda$) output of \texttt{Obs2.prop\_gamma} ($s=1$) and $\hat\lambda_k/|\hat\lambda_k|=1+O(\check g_k)$.
\item \emph{One construction} (\texttt{H02UV/Spin2/Master/Master.lean}): \texttt{Spin2.thm\_flow\_master} gives the conclusions of
  \texttt{thm\_flow\_complete2X} (with the requirements \texttt{Aext}, \texttt{Eext} and the bulk data used by the observable flow)
  and of \texttt{thm\_flow\_final\_spinX} with one set of constants; \texttt{magConcl\_of} gives \texttt{prop\_gamma\_M2} under
  these constants (\texttt{MagConcl}, in \texttt{thm\_flow\_master\_all}).
\item \emph{$\beta_\epsilon$ vs.\ $\log\frac1\epsilon$} (\texttt{Spin2.BetaLog}, \texttt{betaLog\_of}): along every admissible trajectory
  with $\beta_{N_h}\le T$, $|\beta_\epsilon-N_h\beta_1\log L|\le D+a\log(2(N_h\beta_1\log L+D+a|\log2a|))$ ($D,a$ depending on
  $L,T$); with $N_h=N-2$, $\epsilon=L^{-N}$ this is $\beta_\epsilon=\frac3{2\pi}\log\frac1\epsilon+O(\log\log\frac1\epsilon)$.
\item \emph{Not done:} the discharge of the hypotheses on $x_k$ (observable remainders, the dressed
  $\mathfrak l$ increment) and the phase of $\hat\lambda$, i.e.\ the vertex hypothesis of \texttt{prop\_gamma\_M2}
  (\texttt{MagConcl}).
\end{itemize}

%======================================================================
\section{Correlation functions and the zero mode}
\label{sec:corr}

\subsection{The last scale}
Since $\Gamma_N=0$, the last step $j=N-1$ integrates all of
$\Gamma^{\rm int}_N=\Gamma_{N-1}(1+z_{N-1}A\Gamma_{N-1})^{-1}$. We split it in Fourier space:
\[
  C'_N=Z_N\Gamma^{\rm int}_N=C'^{\ne0}_N+t\,Q,\qquad Q_{xy}=|\Lambda|^{-1},\qquad
  t=Z_N\hat\Gamma^{\rm int}_N(0),
\]
where $C'^{\ne0}_N$ has symbol $\hat C'_N\mathbf 1_{p\ne0}$, so $\sum_yC'^{\ne0}_N(x,y)=0$.
\begin{proposition}[Last scale]\label{prop:last}
There are $c_S,\bar\zeta,K$ (depending on $L$) such that for $N\ge1$, $0<m^2\le L^{-2(N-1)}$
and $|\zeta_i|\le\bar\zeta$:
\begin{enumerate}[label=(\roman*)]
\item $m^2t=\dfrac{1-\epsilon_{N-1}(0)}{1-\zeta'_{N-1}\epsilon_{N-1}(0)}$ exactly; for $\zeta=0$,
  $t=t_N$.
\item $|t-t_N|\le K|\Lambda|L^{-(N-1)(d-2)}|\zeta|_{<N}$ and
  $|t-m^{-2}|\le KL^{2N}+K|\Lambda|L^{-(N-1)(d-2)}|\zeta|_{<N}$. In $d=2$:
  $|t-t_N|\le KL^{2N}|\zeta|_{<N}$ and $|m^2t-1|\le Km^2|\Lambda|$, so
  $t=m^{-2}(1+O(m^2|\Lambda|))$.
\item $|Z_N|\,\|\Gamma^{{\rm int},\ne0}_N\|_{(N-1)}\le c_S+K|\zeta|_{<N}$. Hence
  $C'^{\ne0}_N$ satisfies \texttt{GoodCov2} at scale $N-1$ with
  $\ell=(c_S+K\bar\zeta)^{1/2}L^{-(N-1)(d-2)/2}L^{p_\Phi}$, uniformly in $m^2$. It is
  invariant under the lattice symmetries and symmetric.
\item $t$ and the entries of $C'^{\ne0}_N$ are holomorphic in $(\zeta_0,\dots,\zeta_{N-1})$ on
  $\{|\zeta_i|<\bar\zeta\}$, real for real $\zeta$, and continuous in $m^2\in(0,L^{-2(N-1)}]$.
\end{enumerate}
\end{proposition}
\begin{proof}[Sketch]
(i) $Z_jA\Gamma_j=1-\epsilon_j$ and $1+\zeta_j(1-\epsilon_j)=(1+\zeta_j)(1-\zeta_j'\epsilon_j)$, so
$Z_{j+1}A\Gamma_j(1+z_jA\Gamma_j)^{-1}=(1-\epsilon_j)/(1-\zeta'_j\epsilon_j)$; at $p=0$, $A=m^2$.
(ii) Theorem~F-3.1(vi) gives $\|Z_N\Gamma^{\rm int}_N-C_N\|_{(N-1)}\le2K|\zeta|_{<N}$; its $\alpha=0$
moment contains $|\Lambda|^{-1}L^{(N-1)(d-2)}|t-\hat C_N(0)|$, and $\hat C_N(0)=t_N$. Then use
$0\le m^{-2}-t_N\le KL^{2N}$.
(iii) Removing the zero mode decreases every Fourier moment; by (S5) and the lattice moment bound,
$\|\hat C_N\mathbf 1_{p\ne0}\|_{(N-1)}\le c_SL^d(1+2\cdot5^d(\pi/2)^M)$ uniformly in $m^2$. Then (ii)
on $p\ne0$.
(iv) Finite rational expressions with non-vanishing denominators.
\end{proof}
The zero mode cannot stay in the Gram bound: the $\alpha=0$ moment of $C'_N$ contains
$|\Lambda|^{-1}|t|\approx(m^2|\Lambda|)^{-1}$. For $N_h\le j<N-1$ the covariances satisfy
\texttt{GoodCov2} by Proposition~\ref{prop:C5}, uniformly in $m^2\le L^{-2(N-1)}$. By
Lemma~\ref{lem:onestep} the last step is
$\E_{\Gamma_{N-1}}[e^{-z_{N-1}(\psi,A\bar\psi)}F]
 =\det(1+z_{N-1}A\Gamma_{N-1})\,\E_{tQ/Z_N}\E_{\Gamma^{{\rm int},\ne0}_N}\theta F$.
The integral over $C'^{\ne0}_N$ is an ordinary RG step at scale $N-1$; the zero mode $tQ$ is
integrated as in \texttt{prop:ZNN}, with $t$ in place of $t_N$ and coupling $\check g_N=1/\beta_R$.
After the zero-mode integration, as in \hr{}: restricted to constant fields,
$\tilde\cZ_{N,N}=1+\tilde u_{N,N}-|\Lambda|\tilde a_{N,N}\psi\bar\psi$, and the Ward identity determines
$\tilde a_{N,N}$; all errors are $O(\beta_R^{-1})$ relative.

\paragraph{Lean.} \texttt{DressedFRD/ZeroMode.lean} (\texttt{kmat\_split}, \texttt{snorm\_nzP\_CN\_le},
\texttt{Data.absorbed\_closed}, \texttt{Data.last\_zero\_formula}, \texttt{Concl.last\_nz},
\texttt{Concl.last\_zero}); \texttt{H02UV/Step2/Last/ZeroMode.lean} (\texttt{prop\_last}, \texttt{prop\_last\_d2},
\texttt{last\_continuousOn}, \texttt{validCov\_concrete}, \texttt{validLast\_concrete}).
The RG step at $j=N-2$ is the unchanged second-order map: for $L>2^{d+4}+4$ small $(N-1)$-sets
are regions for $\Loc$ (\texttt{region\_small\_L}), and \texttt{Last.rg\_step\_bounds},
\texttt{RawInv.stepL}, \texttt{validAtL\_concrete} (\texttt{H02UV/Step2/Last/Last*.lean}) give its one-step bounds
under \texttt{ValidAtL} ($j+1<N$).

\emph{The step $j=N-1$} (\texttt{H02UV/Step2/Last/LastN*.lean}) integrates $C'^{\ne0}_N$ only. The map of
Theorem~\ref{thm:step} cannot be used unchanged: its $V_{j+1}$ contains $\Loc_{B_+,B}S(B_+)$ with
$B_+=\Lambda$, which is not a small set at any scale $<N$, so Lemma~\texttt{lem:loc2} does not apply.
As in \hr{} ($\Loc=0$ at the last step) the last map has no localisation and $Q=J=0$:
\[
  u_N(B)+V_N(B)=V'_B,\qquad \cW_N(B)=S(B)=\tfrac12F(V'_B,V'(\Lambda))+\E\theta\cW(B),
  \qquad K_N(\Lambda)=e^{u_N(\Lambda)}\textstyle\sum_{\mathcal G(\Lambda)}\cdots
\]
(the polymer formula with these block factors). Then, exactly (\texttt{rawStepN\_ZVK}),
\[
  \E_{C'^{\ne0}_N}\theta Z_{N-1}=e^{-u_N(\Lambda)}\Bigl(\prod_{B\in\cB_{N-1}}e^{-V_N(B)}(1+\cW_N(B))
  +K_N(\Lambda)\Bigr),
\]
with $V_N(B)=V_{\kappa_N}(B)$ a coupling polynomial, $\kappa_N$ the first-order shift of the couplings
of $V^\flat$ by $c_0=C(0,0)$, $\delta=C(e_1,0)-c_0$ (\texttt{rawStepN\_V}). No range of $C$ is used:
$F(V'_B,V'(\Lambda))$ is summed over all $L^d$ blocks. The second-order part $K^{(2)}(\Lambda)$ vanishes
at scale $N$ since $\sum_UK^{(2)}(U)=0$ and $\cP_N=\{\emptyset,\Lambda\}$ (\texttt{K2\_last}), and
the part linear in $K$ is bounded by the crude bound (Schwarz lemma), so
(\texttt{lastK\_ts}) $\|K_N(\Lambda)\|_{T_N}\le C_1(L,A)(\bar g^3+\|K\|_{N-1})$, without contraction.
\texttt{LastN.rg\_last\_bounds} is the analogue of \texttt{rg\_step\_bounds} under \texttt{ValidLast}
(the covariance part of \texttt{ValidAt} for $C'^{\ne0}_N$, supplied by \texttt{validLastN\_concrete}):
$|z_N/Z_N|,|y_N/Z_N|\le\ell^2\check g/14$, $|L^{2N}\nu_N/Z_N-L^2\mu|\le\frac{16}{7}L^2\eta+\ell^2\check g/7$,
$|g_N/Z_N^2|\le\frac87\check g$, $|u_N(\Lambda)|\le L^2Q_1v$, $\|V_N(B)\|\le2Q_1v$,
$\|\cW_N(B)\|\le L^2Q_1^2v^2+Q_1w$, $\|K_N(\Lambda)\|\le C_K(r^3+8k)$, all at the field scale $s_N\ell$.
\emph{Gluing.} \texttt{thm\_flow\_complete} (with \texttt{H02UV/Step2/Last/FlowGlue*.lean}) appends both steps to the
flow with horizon $N_h=N-2$: along every admissible trajectory, the step $j=N-2$ is the unchanged map
under \texttt{ValidAtL} (\texttt{glue\_step}, with $\check g_{N-1}\le3\check g_{N-2}$ and
$\operatorname{Re}\check g_{N-1}\ge\check g_{N-2}/12$ from \texttt{g\_step\_bounds}), and its output, with the flow
coordinates $\operatorname{Re}\check g_{N-1}$, $\mu_{N-1}=\operatorname{Re}\nu_{N-1}$ and the raw $D_{N-1}$, $K_{N-1}$, is fed to
\texttt{LastN.rawStepN} (\texttt{lastOut}, \texttt{glue\_last}). With $\check g=\check g_{N-2}$ and
$C$ depending on $L$ only (\texttt{GlueOut}): $|z_N/Z_N|,|y_N/Z_N|,|L^{2N}\nu_N/Z_N|,|g_N/Z_N^2|,
|u_N(\Lambda)|\le C\check g$, $V_N(B)=V_{\kappa_N}(B)$, $\|V_N(B)\|\le C\check g$, $\|\cW_N(B)\|\le C\check g^2$,
$\|K_N(\Lambda)\|\le C\check g^3$ at the field scale $s_N\ell$. All outputs of the last step are jointly
continuous in $(m^2,b)$ and the step $N-2$ satisfies the sharp band (F1) (\texttt{thm\_flow\_complete\_band}); the
shooting on $\check g_N$ is unique and continuous in $m^2$ (\texttt{thm\_flow\_complete2}); see
Theorem~\ref{thm:flow}.

\emph{Zero mode} (\texttt{H02UV/PtCorr2/Zero/*.lean}). In raw fields the last step integrates $C'^{\ne0}_N=\operatorname{kmat}(\mathrm{nzP}\,\Gamma^{\rm int}_N)$,
and $\Gamma^{\rm int}_N=C'^{\ne0}_N+t_{\rm raw}Q$ with $t_{\rm raw}=\hat\Gamma^{\rm int}_N(0)$ complex, $t=Z_Nt_{\rm raw}$
(\texttt{fullCovN\_eq}). With the final-scale object
$F_N=\prod_{B}e^{-V_N(B)}(1+\cW_N(B))+K_N(\Lambda)$ and $\tilde\cZ_{N,N}=\E_{t_{\rm raw}Q}\theta F_N$,
$\E_{\Gamma^{\rm int}_N}\theta Z_{N-1}=e^{-u_N(\Lambda)}\tilde\cZ_{N,N}$ (\texttt{lastN\_full\_ZVK}). The Wick formula for a
constant covariance holds for complex $c$: $\E_{cQ}\theta F|_{\rm const}=(\epsilon(F)-c\,F^2)+F^2\psi\bar\psi$ for even $F$
(\texttt{Rc\_ECC\_cstC\_even}). $\cW_N$ is not localised: restricted to constant fields each block factor is
$\alpha_B+\beta_B\psi\bar\psi$ with $\alpha_B=1+\epsilon(\cW_N(B))$, $\beta_B=\cW_N(B)^2-|B|\nu_N\alpha_B$ (since
$V_{\kappa}(Y)|_{\rm const}=|Y|\nu\,\psi\bar\psi$, \texttt{c2\_Vcoup}), and the product over the $\le L^2$ blocks is exact
because $(\psi\bar\psi)^2=0$. Hence (\texttt{prop\_ZNN2})
\[
  \tilde\cZ_{N,N}|_{\rm const}=1+\tilde u_{N,N}-|\Lambda|\tilde a_{N,N}\psi\bar\psi,\qquad
  \tilde a_{N,N}=\nu_NP_0-\frac{k^{2,\cW}+K_N(\Lambda)^2}{|\Lambda|},\qquad
  \tilde u_{N,N}=k^0+\tilde a_{N,N}t_{\rm raw},
\]
with $P_0=\prod_B\alpha_B$, $k^0=P_0-1+\epsilon(K_N(\Lambda))$, $k^{2,\cW}=\sum_B\cW_N(B)^2\prod_{B'\ne B}\alpha_{B'}$.
Along every admissible trajectory of Theorem~\ref{thm:flow} (\texttt{thm\_zero\_mode\_flow}, which restates
\texttt{thm\_flow\_complete2} with $\rho\le1/8$ added, and \texttt{ZeroModeOut}), with $\check g=\check g_{N-2}$ and
$E=L^2C\check g^2e^{L^2C\check g^2}$: $|k^0|\le E+C\check g^3$, $|L^{2N}\nu_N/Z_N|\le C\check g$ and
$\bigl||\Lambda|\tilde a_{N,N}/Z_N-L^{2N}\nu_N/Z_N\bigr|\le C\check gE+\ell^{-2}(E+C\check g^3)$. In normalised fields the
$\psi\bar\psi$ coefficient is $-|\Lambda|\tilde a_{N,N}/Z_N$ and $\tilde a_{N,N}t_{\rm raw}=(|\Lambda|\tilde a_{N,N}/Z_N)(t/|\Lambda|)$.
At the shooting solution $\check g_{N-2}=\beta_R^{-1}+O(\beta_R^{-2})$, so the corrections are $O(\beta_R^{-1})$ relative.
Along every history $Z_N=\prod_{i<N}(1+\zeta_i)$, and with the constants of \texttt{prop\_last\_d2} and $\rho\le\bar\zeta$,
$t=Z_Nt_{\rm raw}$ satisfies $|t-t_N|\le KL^{2N}\rho$ and $|m^2t-1|\le Km^2L^{2N}$ (\texttt{zero\_mode\_t}).

\emph{Partition function.} The dressed flow acts on Gaussian integrals $\avg{F}_\Gamma$ (Lemma~\ref{lem:onestep}), not on
$\E_C\theta Z_0$ as a function of an external field. One raw step gives
$\avg{Z_j}_{\Gamma_j}=\det(1+a_jA\Gamma_j)e^{-u_{j+1}(\Lambda)}\avg{Z_{j+1}}_{\Gamma_{j+1}}$ (\texttt{rawStep\_omega}), and the last
step, with $\Gamma_N=0$, gives $\avg{Z_{N-1}}_{\Gamma_{N-1}}=\det(1+aA\Gamma_{N-1})e^{-u_N(\Lambda)}(1+\tilde u_{N,N})$
(\texttt{lastStep\_omega}). Hence (\texttt{hist\_omega\_traj})
\[
  \avg{Z_0}_{\Gamma_0}=\prod_{j\le N-2}\bigl[\det(1+a_jA\Gamma_j)e^{-u_{j+1}(\Lambda)}\bigr]\,
  \det(1+a_{N-1}A\Gamma_{N-1})\,e^{-u_N(\Lambda)}\,(1+\tilde u_{N,N}),
\]
along every trajectory of Theorem~\ref{thm:flow}. Here $\Gamma_0=C=(-\Delta+m^2)^{-1}$. The output of each step is exactly the input of
the next one: the raw flow is real for real data, so the real flow coordinates $(\operatorname{Re}\check g_{j+1},\operatorname{Re}\mu_{j+1},x_{j+1})$
reproduce the raw output (\texttt{rawStep\_real}, \texttt{exactStep\_of}, \texttt{traj\_exact}). The splitting of Theorem~F-3.1(i) holds for
$|\zeta_i|\le\rho$ (\texttt{SplitSmall}, exported by \texttt{thm\_flow\_complete2}). In particular $1+\tilde u_{N,N}\ne0$ iff
$\avg{Z_0}_C\ne0$.

\emph{Susceptibility.} Since $A\mathbb 1=m^2\mathbb 1$, a shift $C\to C+s\mathbb 1$ passes through every absorption step as a zero-mode shift,
$(\Gamma_j+x_j/|\Lambda|)(1+a_jA(\Gamma_j+x_j/|\Lambda|))^{-1}=\Gamma^{\rm int}_{j+1}+\Gamma_{j+1}+x_{j+1}/|\Lambda|$, so the RG maps do not see it
and it ends in the last zero-mode integral (\texttt{split\_shift}, \texttt{hist\_omega\_shift}). Here $x_0=s|\Lambda|$, the next shift is
$x_{j+1}=\frac{\gamma_j+x_j}{1+c_j(\gamma_j+x_j)}-\frac{\gamma_j}{1+c_j\gamma_j}$, and $\gamma_j=\hat\Gamma_j(0)$, $c_j=a_jm^2$.
Since $\avg{Z_0}_{C+s\mathbb 1}=\avg{Z_0}_C-s\,(\E_C\theta Z_0)^2$ exactly, differentiating at $s=0$ gives (\texttt{c2\_formula}) the
$\psi\bar\psi$ coefficient of $\E_C\theta Z_0$. Gaussian integration by parts then gives the dressed form of \texttt{prop:suscept}
(\texttt{suscept2}):
\[
  \frac{1}{|\Lambda|}\sum_{x,y}\frac{\avg{Z_0\bar\psi_x\psi_y}_C}{\avg{Z_0}_C}
  =\frac1{m^2}-\frac1{m^4}\Bigl[\frac{\tilde a_{N,N}\pi_N}{1+\tilde u_{N,N}}+\sum_{j<N}\frac{c_j\pi_j}{\rho_j}\Bigr],\qquad
  \rho_j=1+c_j\gamma_j,\ \pi_j=\prod_{i<j}\rho_i^{-2}.
\]
Without dressing this is \texttt{prop:suscept}.

\emph{The initial object and $\avg{Z_0}_C\ne0$} (\texttt{H02UV/PtCorr2/Zero/ZeroInit.lean}). Along a trajectory started at $(b,\mu_0,0)$ all clamps are
inactive ($Z=1$, $\zeta=0$), so $Z_0=e^{-V_0(\Lambda)}$ with $V_0=V_{(0,0,\mu_0,b)}$. For $z=y=0$ the symmetric quartic
$\sum_i\hat M^{ii}_x=\frac12\sum_e(\nabla_e\psi)(\nabla_e\bar\psi)\psi\bar\psi$ is $\psi_x\bar\psi_x(\nabla\psi\nabla\bar\psi)_x$
(\texttt{vcoup\_zy\_eq\_eval}), and $\Gamma_0=(-\Delta+m^2)^{-1}$ as a matrix (\texttt{$\Gamma$h\_zero\_eq\_Cov}). Hence $V_0$ is the coupling \texttt{Ward.Vh} of the bridge with $s_0=0$, $\beta=1/b$, $h=(\mu_0+m^2)/b-1$ (\texttt{Zh\_zero\_eq\_Z0},
\texttt{Vh\_eq\_V0}), and $Z_{\beta,h}=\beta^{|\Lambda|}\det(-\Delta+m^2)\avg{Z_0}_C$. By the forest representation $Z_{\beta,h}>0$ for $h\ge0$,
so (\texttt{omega\_Zh\_zero\_ne}, \texttt{one\_add\_uFin\_ne\_zero})
\[
  b\le\mu_0+m^2\ \Longrightarrow\ \avg{Z_0}_C\ne0\ \text{ and }\ 1+\tilde u_{N,N}\ne0.
\]
The condition $h\ge0$ is a sign condition on the initial mass coupling; it is a hypothesis of all statements below.

\emph{Ward identity} (\texttt{H02UV/PtCorr2/Zero/Ward2.lean}). With $c^2=b$, $\avg{z_0}=1-b\avg{\bar\psi_0\psi_0}$ and
$\avg{\xi_0\eta_x}=b\avg{\bar\psi_0\psi_x}$, so $\avg{z_0}=h\sum_x\avg{\xi_0\eta_x}$ becomes (\texttt{ward\_V0}, \texttt{ward2})
\[
  (\mu_0+m^2-b)\sum_x\avg{\bar\psi_0\psi_x}=1-b\avg{\bar\psi_0\psi_0},\qquad
  \sum_x\avg{\bar\psi_0\psi_x}=\frac1{m^2}-\frac1{m^4}\Bigl[\frac{\tilde a_{N,N}\pi_N}{1+\tilde u_{N,N}}+\sum_{j<N}\frac{c_j\pi_j}{\rho_j}\Bigr],
\]
the second by translation invariance. For $b<\mu_0+m^2$ this determines $\tilde a_{N,N}\pi_N/(1+\tilde u_{N,N})$ through
$\avg{\bar\psi_0\psi_0}$ (\texttt{ward2\_atilde}). The initial coupling $\mu_0$ plays the role of $h$ in \texttt{cor:ward}: $hb=\mu_0+m^2-b$.

\emph{Dressing terms} (\texttt{H02UV/PtCorr2/Zero/Dressing.lean}). Let $e_j=1-Z_jm^2\gamma_j=\epsilon_j(0)$ (\texttt{eh\_eq\_eps}). Then $c_j=m^2Z_j\zeta_j$ and
$\rho_j=1+\zeta_j(1-e_j)$ exactly. With $q_j=(1+\zeta_j)/\rho_j=1+\zeta_je_j/\rho_j$ and $Q_j=\prod_{i<j}q_i$,
\[
  \pi_j=\frac{Q_j^2}{Z_j^2},\qquad \frac{c_j\pi_j}{\rho_j}=m^2\Bigl(\frac1{Z_j}-\frac1{Z_{j+1}}\Bigr)Q_j^2q_j,\qquad
  \frac1{Z_j}-\frac1{Z_{j+1}}=\frac{\zeta_j}{Z_j(1+\zeta_j)}
\]
(\texttt{dress\_pi}, \texttt{dress\_term}). So each term is $O(|\zeta_j|)$ relative to $Z_j^{-1}$
(\texttt{dress\_term\_norm\_le}). Since $Z_0=1$ the sum telescopes, and the dressing terms reproduce the field strength
(\texttt{dress\_suscept}, \texttt{suscept2\_dressed}):
\[
  \Bigl|\sum_x\avg{\bar\psi_0\psi_x}-\frac1{Z_N}\Bigl[\frac1{m^2}-\frac1{m^4}\frac{\tilde a_{N,N}/Z_N}{1+\tilde u_{N,N}}\Bigr]\Bigr|
  \le\Bigl(\frac1{m^2}\sum_{j<N}\frac{|\zeta_j|}{|Z_j||1+\zeta_j|}+\frac1{m^4}\frac{|\tilde a_{N,N}/Z_N|}{|Z_N||1+\tilde u_{N,N}|}\Bigr)(e^{2S}-1),
\]
$S=\sum_{i<N}|\zeta_ie_i/\rho_i|$.

\emph{Zero modes} (\texttt{DressedFRD/EpsZero.lean}, \texttt{H02UV/PtCorr2/Zero/EpsFlow.lean}). For $|\zeta_i|\le\bar\zeta$ and $j<N$,
$|\zeta_j'\epsilon_j(0)|\le\frac12$ and $|\epsilon_j(0)|\le Bm^2L^{2j}$ (\texttt{eps\_zero\_bound}): $\epsilon_j(0)=m^2\hat C_{\le j}(0)-Z_jm^2\sum_{r>j}\hat R_{j,r}(0)$,
$\hat C_k(0)\le c_dL^{2k}$ by (S4)--(S5) (\texttt{std\_d}), and $|\hat R_{j,r}(0)|\le(2\operatorname{rng}+1)^2\|R_{j,r}\|_{(r-1)}\le4L^{2r}\|R_{j,r}\|_{(r-1)}$ with
$L^{2(r-j)}\le6\cdot128^3W_{r-j}$, so the tails are bounded by $\mathcal S_j\le1$. Hence $|\rho_j|\ge\frac7{16}$, $|q_j-1|\le\frac{16}7|\zeta_j||e_j|$ and
$S\le\frac{16}7\rho B\sum_{i<N}m^2L^{2i}\le5B\rho$ (\texttt{sum\_dQ\_sub\_one\_le}). The bound is carried by \texttt{SplitSmall} to the clamp radius $\rho$ of
Theorem~\ref{thm:flow}. Along the admissible trajectories (\texttt{thm\_ward\_flow\_rho}), if $b\le\mu_0+m^2$,
\[
  \Bigl|\sum_x\avg{\bar\psi_0\psi_x}-\frac1{Z_N}\Bigl[\frac1{m^2}-\frac1{m^4}\frac{\tilde a_{N,N}/Z_N}{1+\tilde u_{N,N}}\Bigr]\Bigr|
  \le\Bigl(\frac1{m^2}\sum_{j<N}\frac{|\zeta_j|}{|Z_j||1+\zeta_j|}+\frac1{m^4}\frac{|\tilde a_{N,N}/Z_N|}{|Z_N||1+\tilde u_{N,N}|}\Bigr)\,10Be^{5B/4}\rho.
\]

\emph{The one-point function} (\texttt{H02UV/PtCorr2/Zero/OnePt.lean}, \texttt{H02UV/}\allowbreak\texttt{PtCorr2/Zero/OnePtFlow.lean}). For $h\ge0$ the ghost-forest
representation gives $\avg{z_0}=\Pb_{\beta,h}[0\leftrightarrow g]$ and $\avg{\xi_0\eta_x}=\Pb_{\beta,h}[0\leftrightarrow x,\ 0\not\leftrightarrow g]$, so
(\texttt{onept\_avg0\_mem}, \texttt{onept\_avg0\_two\_mem})
\[
  0\le1-b\avg{\bar\psi_0\psi_0}\le1,\qquad 0\le b\avg{\bar\psi_0\psi_x}\le1,
\]
and at $h=0$ ($b=\mu_0+m^2$) the Ward identity gives $\avg{\bar\psi_0\psi_0}=1/b$ exactly (\texttt{onept\_avg0\_h0}).
With $r=1-b\avg{\bar\psi_0\psi_0}$ and $X=\tilde a_{N,N}\pi_N/(1+\tilde u_{N,N})+\sum_{j<N}c_j\pi_j/\rho_j$, for $h>0$
\[
  X=m^2-\frac{m^4r}{\mu_0+m^2-b},\qquad\text{so $X$ is real and}\quad m^2-\frac{m^4}{\mu_0+m^2-b}\le X\le m^2
\]
(\texttt{ward2\_onept}). Along the admissible trajectories this is added to the conclusions of \texttt{thm\_ward\_flow\_rho}
(\texttt{thm\_ward\_flow\_onept}, same construction of the constants). So the size of $\avg{\bar\psi_0\psi_0}$ is controlled by
the model, $b\avg{\bar\psi_0\psi_0}\in[0,1]$; its value for $h>0$ (equivalently $\avg{z_0}$) needs the observable flow.

\emph{Sources at the last scale} (\texttt{H02UV/PtCorr2/Obs/ObsZero.lean}). The zero-mode integration acts componentwise on
$F=F_\varnothing+\sigma_aF_a+\sigma_bF_b+\sigma_a\sigma_bF_{ab}$. For the even components ($F_\varnothing$, $F_{ab}$ in both cases, and
$F_a,F_b$ in Case~2) $\epsilon((\E_{cQ}\theta F)_\bullet)=\epsilon(F_\bullet)-c\,F_\bullet^2$ and the restriction to constant fields is
as in \texttt{Rc\_ECC\_cstC\_even} (\texttt{eps\_ECCO\_cstC}, \texttt{eps\_ECCO\_cstC\_two}, \texttt{Rc\_ECCO\_cstC\_e},
\texttt{Rc\_ECCO\_cstC\_ab}, \texttt{Rc\_ECCO\_cstC\_two}). This is the analogue of $1+\tilde u_{N,N}=\epsilon(F_\varnothing)-t_{\rm raw}F_\varnothing^2$
for the observable components.

\emph{The observable chain} (\texttt{H02UV/PtCorr2/Obs/ObsStep.lean}, \texttt{ObsChain.lean}, \texttt{ObsInit.lean}). On the algebra with
sources, $\omega^{\rm obs}_\Gamma(F)=\epsilon(e^{\mathcal L_\Gamma}F)$ componentwise (\texttt{omegaCO}); constants such as
$e^{-u_{j+1}(\Lambda)}=e^{-u^\varnothing}e^{-u^{\rm obs}}$ factor out. One raw step (\texttt{rawStepO\_omega}, $j+1<N$) is
$\omega^{\rm obs}_{\Gamma_j}(Z_j)=\det(1+a_jA\Gamma_j)\,e^{-u_{j+1}(\Lambda)}\,\omega^{\rm obs}_{\Gamma_{j+1}}(Z_{j+1})$ (absorption componentwise,
\texttt{omegaCO\_absorb}, and the exact step \texttt{stepK\_ZVKO} of Section~\ref{sec:obs-flow}); the last step
(\texttt{lastStepO\_omega}) needs no factorisation and ends with the zero mode $t_{\rm raw}Q$, integrated componentwise
($\epsilon F_\bullet-(t_{\rm raw}/|\Lambda|)F_\bullet^2$). Exactness with sources (\texttt{zoutO\_eq}) identifies the output of a step with the
input of the next one built from \texttt{rawStep} and the observable flow \texttt{obsStep}. Along every trajectory with the
hypotheses of \texttt{hist\_omega\_traj} (\texttt{hist\_omegaO\_traj}),
\[
  \omega^{\rm obs}_{\Gamma_0}(Z_0^{\rm obs})=\prod_{j\le N_h}\bigl[\det(1+a_jA\Gamma_j)e^{-u_{j+1}(\Lambda)}\bigr]\cdot\det(1+a_{N-1}A\Gamma_{N-1})\,e^{-u_N(\Lambda)}\,
  \omega^{\rm obs}_{t_{\rm raw}Q}(F_N^{\rm obs}),
\]
where the locality hypothesis \texttt{OLocHyp} of each step is discharged along the flow (\texttt{oLocH\_traj},
Section~\ref{sec:obs-flow}; \texttt{hist\_omegaO\_traj\_loc}). For the initial observable \texttt{obsInit}, $Z_0^{\rm obs}=Z_0e^{-\sum_BV^{\rm obs}_0(B)}$
(\texttt{ZhO\_zero\_eq}), and the components of the left side are $\avg{Z_0}_C$, in Case~2 $\avg{Z_0(\gamma_a+\mu_a\bar\psi_a\psi_a)}_C$,
in Case~1 (no $ab$ couplings) $\mu_a\mu_b\avg{Z_0\bar\psi_a\psi_b}_C$ (\texttt{thm\_obs\_chain\_onept},
\texttt{thm\_obs\_chain\_twopt}; without the locality hypothesis: \texttt{thm\_obs\_chain\_onept\_loc},
\texttt{thm\_obs\_chain\_twopt\_loc}). This is the exact renormalisation-group representation of $\avg{\bar\psi_0\psi_0}$ and
$\avg{\bar\psi_0\psi_x}$. The observable last step uses the map of Section~\ref{sec:obs-flow} at $j=N-1$ ($Q=J=0$) with the
nonzero-mode covariance; correlations are ratios of components of this single identity.

Open: the estimates of the factors
(\texttt{hK'}, \texttt{MuErr}, the $\mathcal W^{\rm obs}$ contraction; \texttt{MuChain} holds along the history, \texttt{muChain\_obsHist}). Without them the value of $\avg{z_0}$ for $h>0$ is represented exactly but not
evaluated; at $h=0$ the two-point function and Expected result~\ref{thm:rginput} are evaluated on route (a1)
(Theorem~\ref{thm:multiscale}).

\subsection{Two-point function}
Combining the observable flow (Section~\ref{sec:obs-flow}) with the zero mode, as in
\texttt{sec:obs}, and using the Ward identity as in Lemma~\ref{lem:square}, should give
Expected result~\ref{thm:rginput}. Proposition~\ref{prop:gamma} then gives
\[
  P_\epsilon(x)=M_{j_x}^2(1+O(\beta_{j_x}^{-1}))=\Bigl(\frac{\beta_{j_x}}{\beta_\epsilon}\Bigr)^{2/3}(1+O(\beta_{j_x}^{-1})),
\]
which is Theorem~\ref{thm:multiscale}.

The factor $2$ is worth noting. The nonzero-mode part of $\avg{\xi_0\eta_x}$ in Case~1 contributes
increments $\frac12(\vartheta_k^2-\vartheta_{k+1}^2)$. The full increments of $P$ arise only
after the zero mode, through the perfect square of Lemma~\ref{lem:square} and its $d=2$ analogue. This is the same
mechanism that produces $2\theta$ in \texttt{cor:conn4pt}.

\subsection{Macroscopic form}\label{sec:macro}
For $|x|\ge\delta$, $j_x\ge N-\log_L(1/\delta)-1$. So at most $D(\delta)$ scales lie above the
coalescence scale, and they are integrated with coupling $\check g_j=\beta_R^{-1}(1+O_\delta(\beta_R^{-1}))$.
At order zero this gives Theorem~\ref{thm:macro0}. The first-order term of Expected result~\ref{thm:macro} comes from
the following sources.
\begin{itemize}
\item The $x$-dependence of $P_\epsilon(x)=b\avg{\bar\psi_0\psi_x}$ lies in the $ab$-component of the pair observable. In the
top split, $P=AB-E$, and the one-point parts $A$, $B$ are expected to coincide by translation invariance. To first
order, the $ab$-coefficient is $\mu_a\mu_b\sum_{k\ge j_x}\Gamma'_k(0,x)$ plus the zero-mode term, which does not depend
on $x$. For $k<j_x$, finite range gives $\Gamma_k(0,x)=0$, so the full sum is $G_{\rm lat}(x)$ up to a constant. The
dressing changes this by $O(\check g)$.
\item The prefactor is $\mu_a\mu_b\approx\vartheta^2\check g_{\rm top}$. The factor $2$ comes from the perfect square, as in
\texttt{cor:conn4pt}.
\item The logarithmic part of $2G_\T$ is already visible in
$\vartheta^2_{j(|x|)}/\vartheta^2_{j(|y|)}=1+\frac23(\beta_{j_x}-\beta_{j_y})/\beta_R+\dots$ with $\frac23\beta_1=\frac1\pi$.
The coalescence scale $j_x$ is discontinuous in $x$, and its jumps must cancel against the pair term. So the
$O(\check g)$ part of the pair layer has to be identified, not only bounded.
\item Finally, $G_{\rm lat}(\lfloor L^Nx\rceil)\to G_\T(x)$ (\texttt{latGreen\_latPtN\_tendstoUniformlyOn}).
\end{itemize}
\emph{Missing estimates.} The proved pair layer bounds the order-$\check g$ terms without identifying them:
$|P-M_{j}^2|\le K_PM_j^2\check g_j$ in \texttt{TwoPtSq}, and $|E|\le K_EM^2_{N-2}\check g_{N-2}$ in \texttt{pairTop\_E\_le}. The
first-order form needs three estimates.
(M1) For $j_x\le j\le N-2$, the $ab$-coefficient of the observable flow is
$\mu_a\mu_b\sum_{j_x\le k<j}\Gamma'_k(0,x)+O(M^2\check g^2)$, with $K^{ab}_j=O(M^2\check g^2)$. This is the second-order
observable map applied in the pair sector.
(M2) In the top steps, $E=\mu_a\mu_b(\mathfrak w^{ab}_N(x)+\Gamma_{N-1}(0,x)+\text{zero mode})+O(M^2\check g^2)$.
(M3) $\check g_N-\check g_{N-2}=O(\check g^2)$, together with the normalisation $Z_\epsilon$.

%======================================================================
\section{Ultraviolet consistency and the continuum limit}
\label{sec:uv}

Compare $\T_\epsilon$ and $\T_{\epsilon/L}$, with $\beta_\epsilon$ and $\beta_{\epsilon/L}$ chosen
for the same $\beta_R$. The finer lattice has one more scale at the bottom.

\begin{conjecture}[Consistency]\label{thm:uv}
At a fixed physical scale $L^{-k}$, i.e.\ at scales $N-k$ and $N+1-k$ respectively, the
normalised coordinates $(\check g,\nu,W,K)$ and the covariances $C'$ of the two flows differ by
$O(L^{-\kappa(N-k)})+O(\epsilon^\kappa)$, after matching $\check g$ exactly.
\end{conjecture}

The memory of the initial data is lost geometrically in the irrelevant directions. The marginal
direction is matched by the choice of $\beta_\epsilon$, and the relevant one by the tuning.
With physical-units convergence of the covariances (Section~\ref{sec:frd}), this gives a Cauchy
sequence in $\epsilon$ at each fixed physical scale. Hence Expected results~\ref{thm:limit}
and~\ref{thm:short} follow, the latter from the explicit $\beta_j$ and $\vartheta_j$.

%======================================================================
\section{What is routine and what is new; formalisation}
\label{sec:new}

\paragraph{Reused from \hr{} as it is}
\begin{itemize}
\item Grassmann calculus and norms, except \texttt{e:normcontract} (Section~\ref{sec:normcontract}).
\item Polymers and the factorisation property.
\item $\Loc$, in the abstract form of \texttt{prop:Loc}.
\item The proof of \texttt{prop:ZVK}.
\item The zero-mode integration.
\item The Ward identity argument, and the generic layers (\texttt{H02RG/Obs/Generic*}, \texttt{Step.Reg},
  \texttt{H02RG/Ward/IVT}).
\end{itemize}

\paragraph{New}
\begin{enumerate}
\item The dressed decomposition (\frd{}, Theorem~F-3.1), and its use in step (A).
\item The exact second-order map with remainder $W$ (Section~\ref{sec:pt}), and the third-order
  estimate (Theorem~\ref{thm:step}).
\item The one-loop coefficients and their universal values in $d=2$ (Proposition~\ref{prop:d2}):
  $\beta_1=\frac3{2\pi}$ and $\gamma=\frac1{2\pi\beta}$. Computing $\beta_2$ needs the two-loop
  coefficients $\mathfrak z^{(2)},\mathfrak v^{(2)}$ and the third-order terms.
\item The finite-horizon flow and shooting (Theorem~\ref{thm:flow}).
\item Ultraviolet consistency (Expected result~\ref{thm:uv}).
\end{enumerate}

\paragraph{Lean plan}
\begin{itemize}
\item \texttt{FRD}:
  \begin{itemize}
  \item set \texttt{stdPt := PtF.prePt.toPtFamily};
  \item the dressed decomposition (Theorem~F-3.1, Corollary~F-3.2(a),(b)): done, Lean
    package \texttt{DressedFRD};
  \item the bounds (S1)--(S5) in $d=2$ and Theorem~F-3.1 for this decomposition: done
    (\texttt{DressedFRD.frd\_stdFull}, \texttt{DressedFRD.thm\_main\_frd});
  \item add the $d=2$ asymptotics: the scaling profile and Frullani.
  \end{itemize}
\item \texttt{Grassmann}:
  \begin{itemize}
  \item complex covariances in $\cL_C$ and $\E_C$: done (\texttt{Grassmann/GaussianC},
    \texttt{Grassmann/IBPC}: \texttt{LopC}, \texttt{expLC}, \texttt{ECC}, \texttt{omegaC}). The
    real operators are the special case \texttt{cplx C} (\texttt{EC\_eq\_ECC},
    \texttt{omega\_eq\_omegaC}); their definitions are unchanged;
  \item the absorption identity \eqref{e:absorb} with the determinant, and
    Lemma~\ref{lem:onestep}: done (\texttt{Grassmann/Absorb}: \texttt{omegaC\_exp\_mul\_det},
    \texttt{ECC\_exp\_mul\_det}, \texttt{omegaC\_onestep}), for every complex $\Gamma$ with
    $\det(1+zA\Gamma)\ne0$;
  \item $F_\Gamma$ and Lemma~\ref{lem:F}: done (\texttt{Grassmann/Truncated});
  \item Gram--Hadamard over $\R$ or $\C$ for Gram matrices $\ip{F_x}{G_y}$ with possibly
    different $F,G$: done (\texttt{Grassmann/HadamardC}: \texttt{norm\_det\_inner\_le});
  \item \texttt{tNorm\_EC\_le} and \texttt{tNorm\_EC\_theta\_le} for covariances with a Gram
    representation (\texttt{IsGram}, \texttt{GramBound}): done (\texttt{Grassmann/GramC}:
    \texttt{tNorm\_ECC\_le}, \texttt{tNorm\_ECC\_theta\_le}; the psd case is
    \texttt{gram\_of\_psd}). For dressed covariances, \texttt{GramBound} follows from the
    Fourier moment norm via Lemma~F-1.1(ii) (\texttt{DressedFRD.kmat\_gramBound},
    \texttt{DressedFRD.Tj\_ECC\_theta\_le\_of\_snorm});
  \item Lemma~\ref{lem:pauli}, replacing \texttt{Tj\_succ\_le} and
    \texttt{Tj\_EC\_theta\_le} (which assume $d\ge3$): done (\texttt{Loc/Pauli*.lean}:
    \texttt{Loc.pauli\_abstract} for abstract partitions of unity, \texttt{Loc.pauli} on the
    torus with the B-spline partition, \texttt{Loc.card\_nearB\_le},
    \texttt{Loc.pauli\_EC} for the Gaussian expectation).
  \end{itemize}
\item \texttt{Loc}:
  \begin{itemize}
  \item \texttt{prop:Loc}(i) in every dimension and for every field scale: done
    (\texttt{Loc.Loc\_bound\_gen});
  \item the contraction for general field scales $\ell'\le\ell$,
    $\|(1-\Loc_X)F\|_{T_{j+1}(\ell')}\le C(r^2L^{-3}+r^4L^{-2}+r^8)\|F\|_{T_j(\ell)}$ with
    $r=\ell'/\ell$, for every $d\ge2$: done (\texttt{Loc.exists\_contract\_bound\_gen}). In
    $d=2$ with $\ell'=2\ell$ this does not apply ($r=2$);
  \item the quartic marginal localisation $\Loc^{(4)}$ and \texttt{prop:Loc}(i) for
    $\Loc^{(2)}+\Loc^{(4)}$: done (\texttt{Loc/Loc4Def.lean}, \texttt{Loc/Loc4Est.lean});
  \item Lemma~\ref{lem:loc2} for $\ell'\le2\ell$ and every $d\ge2$: done
    (\texttt{Loc.loc2\_contract}; the pieces are \texttt{exists\_two\_bound},
    \texttt{exists\_four\_bound} and \texttt{exists\_large\_bound}).
  \end{itemize}
\item \texttt{Step}:
  \begin{itemize}
  \item coordinates with $W$, step (D) and Proposition~\ref{prop:ZVK} for general block
    factors, and Theorem~\ref{thm:step} (third-order bound and closure, with explicit
    $K^{(2)}$, in $d=2$ and $j+1<N$): done (package \texttt{Step2}: \texttt{Step2.gZVK},
    \texttt{Step2.step2\_ZVK}, \texttt{Step2.step2\_V}, \texttt{Step2.step2},
    \texttt{Step2.stepK}, \texttt{Step2.stepK\_d2}, \texttt{Step2.InK2\_Kplus},
    \texttt{Step2.coords\_bound}). The old $d\ge3$ step (\texttt{Step.step}) is unchanged. Still
    open: a uniform bound on $W_{j+1}$ along the flow, and the last step;
  \item steps (A)--(C);
  \item \texttt{GoodCov} with a Gram field instead of \texttt{psd}; all uses of
    \texttt{hC.psd} in \texttt{H02RG/Step/Est*}, \texttt{H02RG/Obs/StepObs*} and \texttt{Coupling} go
    through the two \texttt{Grassmann} lemmas above;
  \item the local factor of step (A);
  \item Proposition~\ref{prop:ZVK};
  \item the first-order coefficient lemma (Lemma~\ref{lem:first}), which is already in the proof of
    \texttt{thm:flow};
  \item the second-order coefficient lemma (Lemma~\ref{lem:second}).
  \end{itemize}
\item \texttt{Flow}: a finite-horizon fixed point with one relevant, one marginal and irrelevant
  coordinates, replacing \texttt{StableManifold}.
\item \texttt{Obs}, \texttt{PtCorr}, \texttt{Ward}: the observable spaces with $\chi$, the
  flow of $\mu^{(2)}$ and $r$, and Lemma~\ref{lem:square} with its $d=2$ analogue.
\end{itemize}

\appendix
\section{Table of one-loop coefficients}
\label{app:table}
All coefficients are in normalised fields, at step $j\to j+1$, with $C'=C'_{j+1}$ and
$w=w'_j$. Each universal value is $\frac{\log L}{2\pi}$ up to $O(L^{-j}+m^2L^{2j}+\bar g)$.

\medskip
\begin{center}
\renewcommand{\arraystretch}{1.3}
\begin{tabular}{l|l|l|l}
quantity & formula & universal value & check\\\hline
gradient, $\zeta_{j+1}/(-g)$ & $C'(0,e_1)$ & $\frac{\log L}{2\pi}$ & Lemma~\ref{lem:first}\\
mass, $(\nu'-\nu)/g$ & $\Delta C'(0)$ & not universal & Lemma~\ref{lem:first}\\
bubble, $\delta g/g^2$ & $\mathfrak B(w+C')-\mathfrak B(w)$ & $\frac{\log L}{2\pi}$ & \texttt{oneloop.py}, \texttt{general.py}\\
$\bar\psi\psi$ vertex, $\delta\mu^{(2)}/(g\mu^{(2)})$ & $\mathfrak l(w+C')-\mathfrak l(w)$ & $\frac{\log L}{2\pi}$ & \texttt{obs.py}, \texttt{obs\_gen.py}\\
$\bar\psi$ vertex, $\delta\mu^{(1)}/(g\mu^{(1)})$ & $0$ & $0$ & Lemma~\ref{lem:second}(iii)\\
mass cross term, $\delta g/(g\,\mathfrak n)$ & $-4$ & --- & \texttt{masscross.py}
\end{tabular}
\end{center}
\medskip

The scripts are in \texttt{scripts/pt-checks/}. They perform the fermionic Wick contractions
symbolically, in momentum space with a general radial propagator, and extract the local
coefficient by Taylor expansion in the external momenta and angular averaging.

\begin{thebibliography}{9}
\bibitem{MR3332939} D.~C.~Brydges and G.~Slade, \emph{A renormalisation group method. II.
Approximation by local polynomials}, J.\ Stat.\ Phys.\ 159 (2015).
\end{thebibliography}

\end{document}
